\documentclass[14pt,a4paper,oneside]{extbook} % Dùng extbook để nhận cỡ chữ 14pt chuẩn
\usepackage[utf8]{vietnam}
\usepackage{amsmath,amssymb,mathrsfs,fancyhdr,enumerate,multirow,makecell,venndiagram}
\usepackage{tkz-euclide,tikz-3dplot,tikz,tkz-tab}
\usetikzlibrary{arrows,calc,intersections,angles,quotes,backgrounds,shapes.geometric,patterns,through}
\usepackage{yhmath}
\usepackage{pgfplots}
\pgfplotsset{compat=1.9}
\usepgfplotslibrary{fillbetween}
\usepackage[tikz]{bclogo}
\usepackage[top=1.8cm, bottom=1.8cm, left=1.8cm, right=1.8cm, headheight=16pt]{geometry}
\usepackage[hidelinks,unicode]{hyperref}
\usepackage{qrcode}

\usepackage[many]{tcolorbox}

% Gói ex_test (Chọn [dethi] để in đề hoặc [loigiai] để in cả lời giải)
%\usepackage[loigiai]{ex_test}
\usepackage[dethi]{ex_test}

\def\colorEX{\color{blue}}

\usepackage{fontawesome}
\usepackage{framed}
\usepackage{polynom}

\colorlet{tcbcol@back}{tcbcolback}
\colorlet{tcbcol@frame}{tcbcolframe}
\newtcolorbox[auto counter]{danggg}[1]{
	fonttitle=\bfseries\upshape,
	colframe=blue,colback=blue!10,
	sharp corners, breakable, 
	title={\color{white}\faFolderO \ DẠNG~\thetcbcounter.\ #1}
	\addcontentsline{toc}{subsection}{Dạng~\thetcbcounter. #1}
}
\newenvironment{dang}[1]{
	\begin{danggg}{#1}
		\textit{Phương pháp giải:~}
	}{\end{danggg}}
\newenvironment{tomtat}{
	\begin{tcolorbox}[boxrule=1pt,arc=0mm,breakable]
	}{\end{tcolorbox}}

\newtheorem{bt}{Bài}
\newtheorem{vd}{Ví dụ}
\theorembodyfont{\it}
\newtheorem{dl}{ Định lí}
\theoremseparator{.}
\theorembodyfont{\rm}
\newtheorem{hq}{ Hệ quả}
\newtheorem{tc}{ Tính chất}
\newtheorem{dn}{ Định nghĩa}
\theoremstyle{nonumberplain}
\newtheorem{kn}{ Khái niệm}
\theoremseparator{.}
\theorembodyfont{\it}
\usepackage{stackengine}
\usepackage{scalerel}

\newcommand\dangersign[1][2ex]{%
	\renewcommand\stacktype{L}%
	\scaleto{\stackon[1.3pt]{\color{red}$\triangle$}{\fontsize{4}{0}\selectfont !}}{#1}%
}
\renewcommand{\parallel}{\mathbin{\!\raisebox{0.25pt}/\mkern-2mu\raisebox{0.25pt}/\!}}
\renewcommand{\backsim}{\ \mathord{\raisebox{0.025\baselineskip}{\rotatebox{90}{\fontfamily{qag}\selectfont S}}}\ }
\theoremseparator{}
\newtheorem{note}{{\dangersign[3ex]}}
\newtheorem{nx}{Nhận xét.}
\newtcolorbox{luuy}[1][]{colback=yellow!20,title={\bfseries Chú ý},#1}

% Tiêu đề chương / bài
\setcounter{secnumdepth}{4}
\setcounter{tocdepth}{2}
\usepackage{titlesec}
\usepackage{titletoc}

\makeatletter
\renewcommand\thesection{\@arabic\c@section}
\makeatother

\titleformat{\section}{\centering\Large\bfseries\color{blue!80!black}}{Chủ đề~\thesection.}{1ex}{
	\setcounter{vd}{0}
	\setcounter{ex}{0}
	\setcounter{bt}{0}
	\setcounter{dl}{0}
	\setcounter{dn}{0}
	\setcounter{hq}{0}
	\setcounter{tc}{0}
	\MakeUppercase
}
\titleformat{\subsection}{\large\bfseries\color{blue!70!black}}{}{0ex}{}
\usepackage[final]{pdfpages}

% Cài đặt riêng
\def\hideboxEX{}
\renewcommand{\boxEX}[2][15pt]{%
	\setbox1=\hbox{#2\strut}
	\def\sizeboxEX{#1}%
	\ \EXbox{\hideboxEX#2\strut}\ %
}

\renewcommand\arraystretch{1.4} 
\renewcommand{\tabcolsep}{4mm}

% Dòng kẻ làm bài TikZ chuẩn 8mm
\renewcommand{\dotlineEX}[1]{%
  \par\vspace{2mm}%
  \noindent%
  \begin{tikzpicture}[x=1mm, y=1mm]%
    \foreach \i in {1,...,#1} {%
      \draw[gray!60, loosely dotted, line width=0.6pt] (0, -\i*8) -- (\linewidth, -\i*8);%
    }%
  \end{tikzpicture}%
}

% Nút Casio
\renewcommand\boxed[1]{%
	\tikz[baseline=(key.base)]
	\node[%
	draw,
	fill=black!40,
	drop shadow={shadow xshift=0.25ex,shadow yshift=-0.25ex,fill=black!90,opacity=0.75},
	rectangle,minimum width=0.6cm,minimum height=0.6cm,
	rounded corners=2pt,
	inner sep=1pt,
	line width=0.5pt,
	font=\sffamily
	](key) {\color{white}#1\strut}
	;
}

\newcommand{\hoac}[1]{\left[\begin{aligned}#1\end{aligned}\right.}
\newcommand{\heva}[1]{\left\{\begin{aligned}#1\end{aligned}\right.}

\pagestyle{fancy}
\fancyhf{} % Xóa sạch toàn bộ header/footer cũ
\fancyhead[L]{\small\itshape\color{gray} Tuyển tập Đề kiểm tra Toán 2 theo từng bài}
\fancyhead[R]{\small\bfseries\color{blue!80!black} Lớp Toán Thầy Lĩnh -- loptoanthaylinh.vn}
\fancyfoot[C]{\bfseries\thepage}
\renewcommand{\headrulewidth}{0.4pt}

\newcommand{\caidat}{2}
\newcommand{\indapan}[2]{
	\begin{center}
		\color{blue!50!green} \fbox{\fontfamily{qag}\bfseries\selectfont BẢNG ĐÁP ÁN}
	\end{center}
	{\fontfamily{lmtt}\bfseries\selectfont \inputansbox[\caidat]{#1}{#2}}
}

\everymath{\color{blue}}

% -------------------------------------------------------------------------
% KHUNG TIÊU ĐỀ PHIẾU KIỂM TRA KÈM MÃ QR CODE CHO TỪNG BÀI
% Cú pháp: \TaoPhieuKiemTra{Chủ đề}{Số bài}{Tên bài}{Link QR thi online}
% -------------------------------------------------------------------------
\newcommand{\TaoPhieuKiemTra}[4]{%
  \refstepcounter{subsection}%
  \addcontentsline{toc}{subsection}{Bài #2. #3}%
  \noindent
  \begin{tcolorbox}[
    colback=white,
    colframe=blue!85!black,
    arc=2.5mm,
    boxrule=1pt,
    left=2.5mm, right=2.5mm, top=1.8mm, bottom=1.8mm
  ]
    % LOGO TRUNG TÂM & HOTLINE | TIÊU ĐỀ BÀI KIỂM TRA | MÃ QR CODE
    \begin{minipage}[c]{0.27\textwidth}
      \centering
      \begin{tcolorbox}[colback=blue!8, colframe=blue!80!black, boxrule=0.6pt, arc=1.5mm, halign=center, top=1.2mm, bottom=1.2mm, left=1mm, right=1mm]
        {\bfseries\color{blue!95!black}\small LỚP TOÁN THẦY LĨNH}\\[1.5pt]
        {\scriptsize\color{blue!70!black}\faGlobe~loptoanthaylinh.vn}\\[1.5pt]
        {\scriptsize\bfseries\color{blue!85!black}\faPhone~Hotline: 0866976878}
      \end{tcolorbox}
    \end{minipage}%
    \hfill
    \begin{minipage}[c]{0.56\textwidth}
      \centering
      {\footnotesize\bfseries\color{blue!80!black}\MakeUppercase{#1}}\\[2pt]
      {\large\bfseries\color{red!85!black}\MakeUppercase{BÀI #2: #3}}\\[2pt]
      {\scriptsize\itshape Thời gian làm bài: 35 phút -- Cấu trúc 3 phần chuẩn}
    \end{minipage}%
    \hfill
    \begin{minipage}[c]{0.14\textwidth}
      \centering
      \qrcode[height=1.35cm]{#4}\\[2pt]
      {\tiny\bfseries\color{blue!90!black}\faQrcode~Quét QR thi}
    \end{minipage}
  \end{tcolorbox}
  \vspace{1.5mm}
}

\begin{document}

\tableofcontents
\newpage

% =========================================================================
% CHỦ ĐỀ 1 - ÔN TẬP VÀ BỔ SUNG
% =========================================================================
\section{Ôn tập và bổ sung}

\TaoPhieuKiemTra{Chủ đề 1 -- Ôn tập và bổ sung}{01}{Ôn tập các số đến 100}{https://loptoanthaylinh.vn/thi-online/toan-2/bai-01}
\textbf{Phần 1. Trắc nghiệm lựa chọn}
\setcounterpageref{ex}{0}

\begin{ex} %[2D1N01]
Số gồm $5$ chục và $7$ đơn vị viết là
\choice
{$75$}
{$507$}
{$50$}
{\True $57$}
\loigiai{
Số gồm $5$ chục và $7$ đơn vị là $57$.\\
\textbf{Đáp án của câu hỏi là D}
}
\end{ex}

\begin{ex} %[2D1N01]
Số lớn nhất có hai chữ số là
\choice
{\True $99$}
{$90$}
{$100$}
{$98$}
\loigiai{
Số lớn nhất có hai chữ số là $99$.\\
\textbf{Đáp án của câu hỏi là A}
}
\end{ex}

\begin{ex} %[2D1H01]
Số bé nhất có hai chữ số giống nhau là
\choice
{$10$}
{\True $11$}
{$22$}
{$00$}
\loigiai{
Số bé nhất có hai chữ số giống nhau là $11$.\\
\textbf{Đáp án của câu hỏi là B}
}
\end{ex}

\begin{ex} %[2D1H01]
Số $84$ đọc là
\choice
{\True Tám mươi tư}
{Tám mươi bốn}
{Tám bốn}
{Bốn mươi tám}
\loigiai{
Số $84$ đọc là Tám mươi tư (hoặc tám mươi bốn).\\
\textbf{Đáp án của câu hỏi là A}
}
\end{ex}

\begin{ex} %[2D1H01]
Trong số $68$, chữ số 6 có giá trị là
\choice
{$6$}
{$68$}
{\True $60$}
{$8$}
\loigiai{
Chữ số 6 ở hàng chục nên có giá trị là $60$.\\
\textbf{Đáp án của câu hỏi là C}
}
\end{ex}

\begin{ex} %[2D1H01]
Điền dấu thích hợp: $45 \dots 54$.
\choice
{$>$}
{$=$}
{Không so sánh được}
{\True $<$}
\loigiai{
Hàng chục: $4 < 5$ nên $45 < 54$.\\
\textbf{Đáp án của câu hỏi là D}
}
\end{ex}

\begin{ex} %[2D1V01]
Sắp xếp các số $35, 53, 30, 50$ theo thứ tự từ bé đến lớn:
\choice
{$53, 50, 35, 30$}
{\True $30, 35, 50, 53$}
{$30, 50, 35, 53$}
{$35, 30, 50, 53$}
\loigiai{
Thứ tự từ bé đến lớn: $30 < 35 < 50 < 53$.\\
\textbf{Đáp án của câu hỏi là B}
}
\end{ex}

\begin{ex} %[2D1V01]
Từ 2 chữ số $3$ và $8$ lập được tất cả bao nhiêu số có hai chữ số khác nhau?
\choice
{$1$ số}
{$3$ số}
{\True $2$ số ($38$ và $83$)}
{$4$ số}
\loigiai{
Các số là $38$ và $83$ (có 2 số).\\
\textbf{Đáp án của câu hỏi là C}
}
\end{ex}

\begin{ex} %[2D1V01]
Viết số $76$ thành tổng của chục và đơn vị:
\choice
{\True $70 + 6$}
{$7 + 6$}
{$70 + 60$}
{$76 + 0$}
\loigiai{
$76 = 70 + 6$.\\
\textbf{Đáp án của câu hỏi là A}
}
\end{ex}

\begin{ex} %[2D1V01]
Số tròn chục liền sau số $40$ là
\choice
{$41$}
{\True $50$}
{$30$}
{$49$}
\loigiai{
Số tròn chục liền sau của $40$ là $50$.\\
\textbf{Đáp án của câu hỏi là B}
}
\end{ex}

\begin{ex} %[2D1V01]
Có bao nhiêu số tròn chục có hai chữ số?
\choice
{$8$ số}
{$10$ số}
{\True $9$ số (từ $10$ đến $90$)}
{$11$ số}
\loigiai{
Các số là $10, 20, 30, 40, 50, 60, 70, 80, 90$ (có 9 số).\\
\textbf{Đáp án của câu hỏi là C}
}
\end{ex}

\begin{ex} %[2D1C01]
Tìm một số có hai chữ số biết chữ số hàng chục là số lớn nhất có 1 chữ số, chữ số hàng đơn vị kém chữ số hàng chục 5 đơn vị.
\choice
{$95$}
{$83$}
{$93$}
{\True $94$}
\loigiai{
Hàng chục: 9. Hàng đơn vị: $9 - 5 = 4$. Số đó là $94$.\\
\textbf{Đáp án của câu hỏi là D}
}
\end{ex}

\textbf{Phần 2. Trắc nghiệm đúng sai}
\setcounterpageref{ex}{0}

\begin{ex} %[2D1H01]
Xét tính đúng sai của các nhận định về số có hai chữ số:
\choiceTF
{\True $10$ là số bé nhất có hai chữ số.}
{\True $99$ là số lớn nhất có hai chữ số.}
{Số $100$ là số có hai chữ số.}
{\True Số $100$ là số có ba chữ số.}
\loigiai{
\begin{itemize}
    \item $10$ là số bé nhất có hai chữ số: Đúng.
    \item $99$ là số lớn nhất có hai chữ số: Đúng.
    \item Số $100$ có 3 chữ số. Ý 3 Sai, Ý 4 Đúng.
\end{itemize}
\begin{itemize}
    \item a) \textbf{Đúng}.
    \item b) \textbf{Đúng}.
    \item c) \textbf{Sai}.
    \item d) \textbf{Đúng}.
\end{itemize}
}
\end{ex}

\begin{ex} %[2D1V01]
Cho các số: $48, 84, 40, 80$. Xét tính đúng sai:
\choiceTF
{\True Số lớn nhất trong các số trên là $84$.}
{\True Số bé nhất trong các số trên là $40$.}
{Có 3 số tròn chục trong 4 số trên.}
{\True Chỉ có 2 số tròn chục là $40$ và $80$.}
\loigiai{
\begin{itemize}
    \item Số lớn nhất là $84$: Đúng.
    \item Số bé nhất là $40$: Đúng.
    \item Có 2 số tròn chục ($40, 80$). Ý 3 Sai, Ý 4 Đúng.
\end{itemize}
\begin{itemize}
    \item a) \textbf{Đúng}.
    \item b) \textbf{Đúng}.
    \item c) \textbf{Sai}.
    \item d) \textbf{Đúng}.
\end{itemize}
}
\end{ex}

\textbf{Phần 3. Trả lời ngắn}
\setcounterpageref{ex}{0}

\begin{ex} %[2D1H01]
Viết số gồm $8$ chục và $3$ đơn vị.
\shortans{83}
\loigiai{
\textbf{Đáp số:} $83$\\
83.
}
\end{ex}

\begin{ex} %[2D1H01]
Tìm số bé nhất có hai chữ số.
\shortans{10}
\loigiai{
\textbf{Đáp số:} $10$\\
10.
}
\end{ex}

\begin{ex} %[2D1V01]
Tìm số lớn nhất có hai chữ số khác nhau.
\shortans{98}
\loigiai{
\textbf{Đáp số:} $98$\\
98.
}
\end{ex}

\begin{ex} %[2D1V01]
Tìm số bé nhất có hai chữ số khác nhau.
\shortans{10}
\loigiai{
\textbf{Đáp số:} $10$\\
10.
}
\end{ex}

\begin{ex} %[2D1V01]
Tính tổng chữ số hàng chục và hàng đơn vị của số $45$.
\shortans{9}
\loigiai{
\textbf{Đáp số:} $9$\\
$4 + 5 = 9$.
}
\end{ex}

\begin{ex} %[2D1C01]
Tìm số tròn chục lớn nhất có hai chữ số mà tổng các chữ số của nó bằng 8.
\shortans{80}
\loigiai{
\textbf{Đáp số:} $80$\\
Số tròn chục có dạng $\overline{a0} \Rightarrow a + 0 = 8 \Rightarrow a = 8$. Số đó là $80$.
}
\end{ex}

\newpage

\TaoPhieuKiemTra{Chủ đề 1 -- Ôn tập và bổ sung}{02}{Tia số. Số liền trước, số liền sau}{https://loptoanthaylinh.vn/thi-online/toan-2/bai-02}
\textbf{Phần 1. Trắc nghiệm lựa chọn}
\setcounterpageref{ex}{0}

\begin{ex} %[2D1N02]
Trên tia số, gốc của tia số là số nào?
\choice
{\True $0$}
{$1$}
{$10$}
{$100$}
\loigiai{
Gốc của tia số là số 0.\\
\textbf{Đáp án của câu hỏi là A}
}
\end{ex}

\begin{ex} %[2D1N02]
Số liền sau của số $39$ là
\choice
{$38$}
{\True $40$}
{$41$}
{$30$}
\loigiai{
$39 + 1 = 40$.\\
\textbf{Đáp án của câu hỏi là B}
}
\end{ex}

\begin{ex} %[2D1H02]
Số liền trước của số $70$ là
\choice
{$71$}
{$60$}
{\True $69$}
{$80$}
\loigiai{
$70 - 1 = 69$.\\
\textbf{Đáp án của câu hỏi là C}
}
\end{ex}

\begin{ex} %[2D1H02]
Số liền trước của số $1$ là
\choice
{\True $0$}
{$2$}
{$10$}
{Không có}
\loigiai{
$1 - 1 = 0$.\\
\textbf{Đáp án của câu hỏi là A}
}
\end{ex}

\begin{ex} %[2D1H02]
Trên tia số, số đứng liền bên phải số 45 là
\choice
{$44$}
{$47$}
{$50$}
{\True $46$}
\loigiai{
Đứng liền bên phải là số liền sau: $45 + 1 = 46$.\\
\textbf{Đáp án của câu hỏi là D}
}
\end{ex}

\begin{ex} %[2D1H02]
Trên tia số, số đứng liền bên trái số 80 là
\choice
{$81$}
{$70$}
{\True $79$}
{$90$}
\loigiai{
Đứng liền bên trái là số liền trước: $80 - 1 = 79$.\\
\textbf{Đáp án của câu hỏi là C}
}
\end{ex}

\begin{ex} %[2D1V02]
Số liền sau của số lớn nhất có hai chữ số là
\choice
{\True $100$}
{$98$}
{$99$}
{$101$}
\loigiai{
Số lớn nhất có 2 chữ số là 99. Số liền sau là $99 + 1 = 100$.\\
\textbf{Đáp án của câu hỏi là A}
}
\end{ex}

\begin{ex} %[2D1V02]
Số liền trước của số bé nhất có hai chữ số là
\choice
{$10$}
{$11$}
{\True $9$}
{$8$}
\loigiai{
Số bé nhất có 2 chữ số là 10. Số liền trước là $10 - 1 = 9$.\\
\textbf{Đáp án của câu hỏi là C}
}
\end{ex}

\begin{ex} %[2D1V02]
Ba số liên tiếp tăng dần là
\choice
{$18, 20, 19$}
{\True $28, 29, 30$}
{$40, 39, 38$}
{$50, 52, 54$}
\loigiai{
Ba số liên tiếp: $28, 29, 30$.\\
\textbf{Đáp án của câu hỏi là B}
}
\end{ex}

\begin{ex} %[2D1V02]
Tìm số ở giữa số $68$ và số $70$:
\choice
{$67$}
{$71$}
{$72$}
{\True $69$}
\loigiai{
Số ở giữa 68 và 70 là 69.\\
\textbf{Đáp án của câu hỏi là D}
}
\end{ex}

\begin{ex} %[2D1V02]
Có bao nhiêu số tự nhiên nằm giữa hai số $25$ và $29$?
\choice
{$2$ số}
{\True $3$ số ($26, 27, 28$)}
{$4$ số}
{$5$ số}
\loigiai{
Các số là $26, 27, 28$ (có 3 số).\\
\textbf{Đáp án của câu hỏi là B}
}
\end{ex}

\begin{ex} %[2D1C02]
Tìm một số biết rằng số liền trước của nó là số lớn nhất có 1 chữ số.
\choice
{$9$}
{$11$}
{$8$}
{\True $10$}
\loigiai{
Số lớn nhất có 1 chữ số là 9. Số cần tìm là $9 + 1 = 10$.\\
\textbf{Đáp án của câu hỏi là D}
}
\end{ex}

\textbf{Phần 2. Trắc nghiệm đúng sai}
\setcounterpageref{ex}{0}

\begin{ex} %[2D1H02]
Xét tính đúng sai về số liền trước và số liền sau:
\choiceTF
{\True Số liền sau hơn số liền trước 1 đơn vị.}
{\True Số liền trước của 50 là 49.}
{Số liền sau của 89 là 88.}
{\True Số liền sau của 89 là 90.}
\loigiai{
\begin{itemize}
    \item Liền sau hơn liền trước 1 đơn vị: Đúng.
    \item $50 - 1 = 49$: Đúng.
    \item $89 + 1 = 90 \neq 88$. Ý 3 Sai, Ý 4 Đúng.
\end{itemize}
\begin{itemize}
    \item a) \textbf{Đúng}.
    \item b) \textbf{Đúng}.
    \item c) \textbf{Sai}.
    \item d) \textbf{Đúng}.
\end{itemize}
}
\end{ex}

\begin{ex} %[2D1V02]
Cho số 35. Xét tính đúng sai:
\choiceTF
{\True Số liền trước của 35 là 34.}
{\True Số liền sau của 35 là 36.}
{Tổng của số liền trước và số liền sau của 35 là 60.}
{\True Tổng của số liền trước và số liền sau của 35 là 70 ($34 + 36 = 70$).}
\loigiai{
\begin{itemize}
    \item Liền trước là 34: Đúng.
    \item Liền sau là 36: Đúng.
    \item $34 + 36 = 70 \neq 60$. Ý 3 Sai, Ý 4 Đúng.
\end{itemize}
\begin{itemize}
    \item a) \textbf{Đúng}.
    \item b) \textbf{Đúng}.
    \item c) \textbf{Sai}.
    \item d) \textbf{Đúng}.
\end{itemize}
}
\end{ex}

\textbf{Phần 3. Trả lời ngắn}
\setcounterpageref{ex}{0}

\begin{ex} %[2D1H02]
Tìm số liền trước của số $85$.
\shortans{84}
\loigiai{
\textbf{Đáp số:} $84$\\
$85 - 1 = 84$.
}
\end{ex}

\begin{ex} %[2D1H02]
Tìm số liền sau của số $79$.
\shortans{80}
\loigiai{
\textbf{Đáp số:} $80$\\
$79 + 1 = 80$.
}
\end{ex}

\begin{ex} %[2D1V02]
Tìm số liền sau của số liền sau số $48$.
\shortans{50}
\loigiai{
\textbf{Đáp số:} $50$\\
$48 + 1 + 1 = 50$.
}
\end{ex}

\begin{ex} %[2D1V02]
Tìm số liền trước của số liền trước số $62$.
\shortans{60}
\loigiai{
\textbf{Đáp số:} $60$\\
$62 - 1 - 1 = 60$.
}
\end{ex}

\begin{ex} %[2D1V02]
Có bao nhiêu số tự nhiên nằm giữa số $40$ và số $45$?
\shortans{4}
\loigiai{
\textbf{Đáp số:} $4$\\
Các số là $41, 42, 43, 44$ (có 4 số).
}
\end{ex}

\begin{ex} %[2D1C02]
Tìm một số biết rằng số liền sau của nó là số bé nhất có hai chữ số giống nhau.
\shortans{10}
\loigiai{
\textbf{Đáp số:} $10$\\
Số bé nhất có 2 chữ số giống nhau là 11. Số cần tìm là $11 - 1 = 10$.
}
\end{ex}

\newpage

\TaoPhieuKiemTra{Chủ đề 1 -- Ôn tập và bổ sung}{03}{Các thành phần của phép cộng, phép trừ}{https://loptoanthaylinh.vn/thi-online/toan-2/bai-03}
\textbf{Phần 1. Trắc nghiệm lựa chọn}
\setcounterpageref{ex}{0}

\begin{ex} %[2D1N03]
Trong phép tính $32 + 15 = 47$, số 47 được gọi là
\choice
{Số hạng}
{\True Tổng}
{Số bị trừ}
{Hiệu}
\loigiai{
Kết quả của phép cộng gọi là Tổng.\\
\textbf{Đáp án của câu hỏi là B}
}
\end{ex}

\begin{ex} %[2D1N03]
Trong phép tính $58 - 23 = 35$, số 58 được gọi là
\choice
{Số trừ}
{Hiệu}
{\True Số bị trừ}
{Tổng}
\loigiai{
Số đứng trước dấu trừ gọi là Số bị trừ.\\
\textbf{Đáp án của câu hỏi là C}
}
\end{ex}

\begin{ex} %[2D1H03]
Trong phép tính $58 - 23 = 35$, số 23 được gọi là
\choice
{\True Số trừ}
{Số bị trừ}
{Hiệu}
{Số hạng}
\loigiai{
Số đứng sau dấu trừ gọi là Số trừ.\\
\textbf{Đáp án của câu hỏi là A}
}
\end{ex}

\begin{ex} %[2D1H03]
Trong phép tính $58 - 23 = 35$, số 35 được gọi là
\choice
{Tổng}
{Số trừ}
{Số bị trừ}
{\True Hiệu}
\loigiai{
Kết quả phép trừ gọi là Hiệu.\\
\textbf{Đáp án của câu hỏi là D}
}
\end{ex}

\begin{ex} %[2D1H03]
Cho hai số hạng lần lượt là $24$ và $15$. Tổng của hai số đó là
\choice
{$38$}
{\True $39$}
{$40$}
{$9$}
\loigiai{
$24 + 15 = 39$.\\
\textbf{Đáp án của câu hỏi là B}
}
\end{ex}

\begin{ex} %[2D1H03]
Số bị trừ là $68$, số trừ là $24$. Hiệu là
\choice
{$42$}
{$46$}
{$92$}
{\True $44$}
\loigiai{
$68 - 24 = 44$.\\
\textbf{Đáp án của câu hỏi là D}
}
\end{ex}

\begin{ex} %[2D1V03]
Tìm tổng của số bé nhất có hai chữ số và số lớn nhất có một chữ số.
\choice
{$18$}
{$20$}
{\True $19$}
{$10$}
\loigiai{
Số bé nhất có 2 chữ số: 10. Số lớn nhất có 1 chữ số: 9. Tổng: $10 + 9 = 19$.\\
\textbf{Đáp án của câu hỏi là C}
}
\end{ex}

\begin{ex} %[2D1V03]
Tìm hiệu của số tròn chục lớn nhất có hai chữ số và $40$.
\choice
{\True $50$}
{$40$}
{$60$}
{$70$}
\loigiai{
Số tròn chục lớn nhất: 90. Hiệu: $90 - 40 = 50$.\\
\textbf{Đáp án của câu hỏi là A}
}
\end{ex}

\begin{ex} %[2D1V03]
Trong một phép trừ có số bị trừ bằng $45$ và số trừ bằng $0$. Hiệu là
\choice
{$0$}
{$54$}
{$1$}
{\True $45$}
\loigiai{
$45 - 0 = 45$.\\
\textbf{Đáp án của câu hỏi là D}
}
\end{ex}

\begin{ex} %[2D1V03]
Trong một phép trừ có số bị trừ bằng số trừ. Hiệu là
\choice
{\True $0$}
{$1$}
{Bằng số bị trừ}
{Không tính được}
\loigiai{
Hai số bằng nhau trừ cho nhau có hiệu bằng 0.\\
\textbf{Đáp án của câu hỏi là A}
}
\end{ex}

\begin{ex} %[2D1V03]
Nếu một số hạng bằng 0 thì tổng sẽ
\choice
{Bằng 0}
{\True Bằng số hạng còn lại}
{Bằng 1}
{Tăng gấp đôi}
\loigiai{
$a + 0 = a$ (bằng số hạng còn lại).\\
\textbf{Đáp án của câu hỏi là B}
}
\end{ex}

\begin{ex} %[2D1C03]
Trong một phép trừ, nếu tăng số bị trừ thêm 5 đơn vị và giữ nguyên số trừ thì hiệu sẽ
\choice
{Giảm đi 5 đơn vị}
{Không thay đổi}
{\True Tăng thêm 5 đơn vị}
{Tăng gấp đôi}
\loigiai{
Tăng số bị trừ bao nhiêu đơn vị thì hiệu tăng bấy nhiêu đơn vị.\\
\textbf{Đáp án của câu hỏi là C}
}
\end{ex}

\textbf{Phần 2. Trắc nghiệm đúng sai}
\setcounterpageref{ex}{0}

\begin{ex} %[2D1H03]
Xét tính đúng sai về tên gọi thành phần phép tính:
\choiceTF
{\True Trong phép cộng $20 + 30 = 50$, $20$ và $30$ là các số hạng.}
{\True Trong phép trừ $50 - 20 = 30$, $50$ là số bị trừ.}
{Trong phép trừ $50 - 20 = 30$, $30$ là số trừ.}
{\True Trong phép trừ $50 - 20 = 30$, $20$ là số trừ và $30$ là hiệu.}
\loigiai{
\begin{itemize}
    \item $20$ và $30$ là số hạng: Đúng.
    \item $50$ là số bị trừ: Đúng.
    \item $30$ là hiệu (không phải số trừ). Ý 3 Sai, Ý 4 Đúng.
\end{itemize}
\begin{itemize}
    \item a) \textbf{Đúng}.
    \item b) \textbf{Đúng}.
    \item c) \textbf{Sai}.
    \item d) \textbf{Đúng}.
\end{itemize}
}
\end{ex}

\begin{ex} %[2D1V03]
Cho phép tính $45 + 12 = 57$. Xét tính đúng sai:
\choiceTF
{\True Tổng là $57$.}
{\True Các số hạng là $45$ và $12$.}
{Nếu đổi chỗ hai số hạng thì tổng thay đổi.}
{\True Khi đổi chỗ các số hạng thì tổng không thay đổi ($12 + 45 = 57$).}
\loigiai{
\begin{itemize}
    \item Tổng là 57: Đúng.
    \item Số hạng 45 và 12: Đúng.
    \item Phép cộng có tính chất giao hoán nên không đổi. Ý 3 Sai, Ý 4 Đúng.
\end{itemize}
\begin{itemize}
    \item a) \textbf{Đúng}.
    \item b) \textbf{Đúng}.
    \item c) \textbf{Sai}.
    \item d) \textbf{Đúng}.
\end{itemize}
}
\end{ex}

\textbf{Phần 3. Trả lời ngắn}
\setcounterpageref{ex}{0}

\begin{ex} %[2D1H03]
Tính tổng của hai số hạng $35$ và $24$.
\shortans{59}
\loigiai{
\textbf{Đáp số:} $59$\\
$35 + 24 = 59$.
}
\end{ex}

\begin{ex} %[2D1H03]
Tính hiệu biết số bị trừ là $76$, số trừ là $32$.
\shortans{44}
\loigiai{
\textbf{Đáp số:} $44$\\
$76 - 32 = 44$.
}
\end{ex}

\begin{ex} %[2D1V03]
Tính tổng của số tròn chục bé nhất có hai chữ số và số lớn nhất có hai chữ số.
\shortans{109}
\loigiai{
\textbf{Đáp số:} $109$\\
$10 + 99 = 109$.
}
\end{ex}

\begin{ex} %[2D1V03]
Tính hiệu của số lớn nhất có hai chữ số và số bé nhất có hai chữ số.
\shortans{89}
\loigiai{
\textbf{Đáp số:} $89$\\
$99 - 10 = 89$.
}
\end{ex}

\begin{ex} %[2D1V03]
Tìm tổng của ba số hạng: $10, 20, 30$.
\shortans{60}
\loigiai{
\textbf{Đáp số:} $60$\\
$10 + 20 + 30 = 60$.
}
\end{ex}

\begin{ex} %[2D1C03]
Trong một phép trừ, số bị trừ là $50$, hiệu là $20$. Hỏi số trừ là bao nhiêu?
\shortans{30}
\loigiai{
\textbf{Đáp số:} $30$\\
Số trừ: $50 - 20 = 30$.
}
\end{ex}

\newpage

\TaoPhieuKiemTra{Chủ đề 1 -- Ôn tập và bổ sung}{04}{Hơn, kém nhau bao nhiêu đơn vị}{https://loptoanthaylinh.vn/thi-online/toan-2/bai-04}
\textbf{Phần 1. Trắc nghiệm lựa chọn}
\setcounterpageref{ex}{0}

\begin{ex} %[2D1N04]
Muốn biết số thứ nhất hơn số thứ hai bao nhiêu đơn vị, ta làm phép tính gì?
\choice
{Phép cộng}
{Phép nhân}
{\True Phép trừ (lấy số lớn trừ số bé)}
{Phép chia}
\loigiai{
Ta lấy số lớn trừ số bé.\\
\textbf{Đáp án của câu hỏi là C}
}
\end{ex}

\begin{ex} %[2D1N04]
Số 15 hơn số 10 bao nhiêu đơn vị?
\choice
{$3$}
{\True $5$}
{$25$}
{$4$}
\loigiai{
$15 - 10 = 5$.\\
\textbf{Đáp án của câu hỏi là B}
}
\end{ex}

\begin{ex} %[2D1H04]
Số 20 kém số 28 bao nhiêu đơn vị?
\choice
{\True $8$}
{$6$}
{$10$}
{$48$}
\loigiai{
$28 - 20 = 8$.\\
\textbf{Đáp án của câu hỏi là A}
}
\end{ex}

\begin{ex} %[2D1H04]
Bạn Nam có 18 viên bi, bạn Hải có 12 viên bi. Nam có nhiều hơn Hải số viên bi là
\choice
{$5$ viên}
{\True $6$ viên}
{$7$ viên}
{$30$ viên}
\loigiai{
$18 - 12 = 6$ viên bi.\\
\textbf{Đáp án của câu hỏi là B}
}
\end{ex}

\begin{ex} %[2D1H04]
Cây cam có 35 quả, cây bưởi có 20 quả. Cây bưởi có ít hơn cây cam số quả là
\choice
{\True $15$ quả}
{$10$ quả}
{$25$ quả}
{$55$ quả}
\loigiai{
$35 - 20 = 15$ quả.\\
\textbf{Đáp án của câu hỏi là A}
}
\end{ex}

\begin{ex} %[2D1H04]
Một đoạn dây dài $48\text{ cm}$, một đoạn dây khác dài $30\text{ cm}$. Đoạn thứ nhất dài hơn đoạn thứ hai là
\choice
{$16\text{ cm}$}
{$20\text{ cm}$}
{\True $18\text{ cm}$}
{$78\text{ cm}$}
\loigiai{
$48 - 30 = 18\text{ cm}$.\\
\textbf{Đáp án của câu hỏi là C}
}
\end{ex}

\begin{ex} %[2D1V04]
Tổ 1 trồng được 25 cây, tổ 2 trồng được 15 cây. Tổ 1 trồng nhiều hơn tổ 2 số cây là
\choice
{\True $10$ cây}
{$8$ cây}
{$12$ cây}
{$40$ cây}
\loigiai{
$25 - 15 = 10$ cây.\\
\textbf{Đáp án của câu hỏi là A}
}
\end{ex}

\begin{ex} %[2D1V04]
Bình có 24 cái kẹo, Lan có 36 cái kẹo. Bình có ít hơn Lan bao nhiêu cái kẹo?
\choice
{$10$ cái}
{$14$ cái}
{$60$ cái}
{\True $12$ cái}
\loigiai{
$36 - 24 = 12$ cái kẹo.\\
\textbf{Đáp án của câu hỏi là D}
}
\end{ex}

\begin{ex} %[2D1V04]
Số lớn nhất có hai chữ số hơn số bé nhất có hai chữ số bao nhiêu đơn vị?
\choice
{$80$}
{\True $89$}
{$90$}
{$88$}
\loigiai{
$99 - 10 = 89$.\\
\textbf{Đáp án của câu hỏi là B}
}
\end{ex}

\begin{ex} %[2D1V04]
Một lớp học có 20 bạn nam và 18 bạn nữ. Số bạn nam nhiều hơn số bạn nữ là
\choice
{$1$ bạn}
{$3$ bạn}
{$4$ bạn}
{\True $2$ bạn}
\loigiai{
$20 - 18 = 2$ bạn.\\
\textbf{Đáp án của câu hỏi là D}
}
\end{ex}

\begin{ex} %[2D1V04]
Bao gạo to nặng $50\text{ kg}$, bao gạo nhỏ nặng $20\text{ kg}$. Bao to nặng hơn bao nhỏ bao nhiêu ki-lô-gam?
\choice
{$25\text{ kg}$}
{$35\text{ kg}$}
{\True $30\text{ kg}$}
{$70\text{ kg}$}
\loigiai{
$50 - 20 = 30\text{ kg}$.\\
\textbf{Đáp án của câu hỏi là C}
}
\end{ex}

\begin{ex} %[2D1C04]
An có 15 viên bi, Bình có 25 viên bi. Hỏi Bình phải cho An bao nhiêu viên bi để số bi của hai bạn bằng nhau?
\choice
{$10$ viên}
{$6$ viên}
{$4$ viên}
{\True $5$ viên}
\loigiai{
Bình hơn An: $25 - 15 = 10$ viên bi.\\
Bình cần cho An: $10 : 2 = 5$ viên bi.\\
\textbf{Đáp án của câu hỏi là D}
}
\end{ex}

\textbf{Phần 2. Trắc nghiệm đúng sai}
\setcounterpageref{ex}{0}

\begin{ex} %[2D1H04]
Xét tính đúng sai về so sánh hơn kém:
\choiceTF
{\True Số 45 hơn số 30 là 15 đơn vị.}
{\True Số 20 kém số 50 là 30 đơn vị.}
{Số 60 hơn số 40 là 30 đơn vị.}
{\True Số 60 hơn số 40 là 20 đơn vị ($60 - 40 = 20$).}
\loigiai{
\begin{itemize}
    \item $45 - 30 = 15$: Đúng.
    \item $50 - 20 = 30$: Đúng.
    \item $60 - 40 = 20 \neq 30$. Ý 3 Sai, Ý 4 Đúng.
\end{itemize}
\begin{itemize}
    \item a) \textbf{Đúng}.
    \item b) \textbf{Đúng}.
    \item c) \textbf{Sai}.
    \item d) \textbf{Đúng}.
\end{itemize}
}
\end{ex}

\begin{ex} %[2D1V04]
Mai có 28 bông hoa, Cúc có 16 bông hoa. Xét tính đúng sai:
\choiceTF
{\True Mai có nhiều hơn Cúc 12 bông hoa ($28 - 16 = 12$).}
{\True Cúc có ít hơn Mai 12 bông hoa.}
{Cả hai bạn có tất cả 40 bông hoa.}
{\True Cả hai bạn có tất cả 44 bông hoa ($28 + 16 = 44$).}
\loigiai{
\begin{itemize}
    \item Mai hơn Cúc: $28 - 16 = 12$: Đúng.
    \item Cúc ít hơn Mai 12: Đúng.
    \item Tổng: $28 + 16 = 44 \neq 40$. Ý 3 Sai, Ý 4 Đúng.
\end{itemize}
\begin{itemize}
    \item a) \textbf{Đúng}.
    \item b) \textbf{Đúng}.
    \item c) \textbf{Sai}.
    \item d) \textbf{Đúng}.
\end{itemize}
}
\end{ex}

\textbf{Phần 3. Trả lời ngắn}
\setcounterpageref{ex}{0}

\begin{ex} %[2D1H04]
Số 58 hơn số 25 bao nhiêu đơn vị?
\shortans{33}
\loigiai{
\textbf{Đáp số:} $33$\\
$58 - 25 = 33$.
}
\end{ex}

\begin{ex} %[2D1H04]
Số 32 kém số 67 bao nhiêu đơn vị?
\shortans{35}
\loigiai{
\textbf{Đáp số:} $35$\\
$67 - 32 = 35$.
}
\end{ex}

\begin{ex} %[2D1V04]
Bác Ba nuôi 48 con gà và 25 con vịt. Hỏi số gà nhiều hơn số vịt bao nhiêu con?
\shortans{23}
\loigiai{
\textbf{Đáp số:} $23$\\
$48 - 25 = 23$ con.
}
\end{ex}

\begin{ex} %[2D1V04]
Hùng có 35 viên bi, Dũng có 14 viên bi. Hỏi Dũng có ít hơn Hùng bao nhiêu viên bi?
\shortans{21}
\loigiai{
\textbf{Đáp số:} $21$\\
$35 - 14 = 21$ viên.
}
\end{ex}

\begin{ex} %[2D1V04]
Số tròn chục lớn nhất có hai chữ số hơn số tròn chục bé nhất có hai chữ số bao nhiêu đơn vị?
\shortans{80}
\loigiai{
\textbf{Đáp số:} $80$\\
$90 - 10 = 80$.
}
\end{ex}

\begin{ex} %[2D1C04]
Mẹ 36 tuổi, con 12 tuổi. Mẹ hơn con bao nhiêu tuổi?
\shortans{24}
\loigiai{
\textbf{Đáp số:} $24$\\
$36 - 12 = 24$ tuổi.
}
\end{ex}

\newpage

\TaoPhieuKiemTra{Chủ đề 1 -- Ôn tập và bổ sung}{05}{Bảng cộng và bảng trừ (có nhớ)}{https://loptoanthaylinh.vn/thi-online/toan-2/bai-05}
\textbf{Phần 1. Trắc nghiệm lựa chọn}
\setcounterpageref{ex}{0}

\begin{ex} %[2D1N05]
Kết quả của phép tính $40 + 30$ là
\choice
{$60$}
{\True $70$}
{$80$}
{$75$}
\loigiai{
$40 + 30 = 70$.\\
\textbf{Đáp án của câu hỏi là B}
}
\end{ex}

\begin{ex} %[2D1N05]
Kết quả của phép tính $80 - 50$ là
\choice
{$20$}
{$40$}
{\True $30$}
{$35$}
\loigiai{
$80 - 50 = 30$.\\
\textbf{Đáp án của câu hỏi là C}
}
\end{ex}

\begin{ex} %[2D1H05]
Kết quả của phép tính $34 + 23$ là
\choice
{\True $57$}
{$56$}
{$58$}
{$67$}
\loigiai{
$34 + 23 = 57$.\\
\textbf{Đáp án của câu hỏi là A}
}
\end{ex}

\begin{ex} %[2D1H05]
Kết quả của phép tính $68 - 25$ là
\choice
{$42$}
{$44$}
{$33$}
{\True $43$}
\loigiai{
$68 - 25 = 43$.\\
\textbf{Đáp án của câu hỏi là D}
}
\end{ex}

\begin{ex} %[2D1H05]
Tìm $x$ biết: $x - 20 = 45$.
\choice
{$25$}
{$55$}
{\True $65$}
{$60$}
\loigiai{
$x = 45 + 20 = 65$.\\
\textbf{Đáp án của câu hỏi là C}
}
\end{ex}

\begin{ex} %[2D1H05]
Tính giá trị biểu thức: $45 + 20 - 15$.
\choice
{$40$}
{\True $50$}
{$60$}
{$55$}
\loigiai{
$65 - 15 = 50$.\\
\textbf{Đáp án của câu hỏi là B}
}
\end{ex}

\begin{ex} %[2D1V05]
Một rổ có 25 quả cam, mẹ bỏ thêm vào 14 quả cam nữa. Trong rổ có tất cả số quả cam là
\choice
{$38$ quả}
{$40$ quả}
{\True $39$ quả}
{$49$ quả}
\loigiai{
$25 + 14 = 39$ quả.\\
\textbf{Đáp án của câu hỏi là C}
}
\end{ex}

\begin{ex} %[2D1V05]
Đàn vịt có 58 con, bác nông dân bán đi 24 con. Đàn vịt còn lại số con là
\choice
{\True $34$ con}
{$32$ con}
{$36$ con}
{$82$ con}
\loigiai{
$58 - 24 = 34$ con.\\
\textbf{Đáp án của câu hỏi là A}
}
\end{ex}

\begin{ex} %[2D1V05]
Tính nhanh: $23 + 40 + 15$.
\choice
{$76$}
{$80$}
{$68$}
{\True $78$}
\loigiai{
$63 + 15 = 78$.\\
\textbf{Đáp án của câu hỏi là D}
}
\end{ex}

\begin{ex} %[2D1V05]
Tìm $y$ biết: $87 - y = 35$.
\choice
{$51$}
{\True $52$}
{$53$}
{$54$}
\loigiai{
$y = 87 - 35 = 52$.\\
\textbf{Đáp án của câu hỏi là B}
}
\end{ex}

\begin{ex} %[2D1V05]
Đoạn thẳng $AB$ dài $32\text{ cm}$, đoạn thẳng $BC$ dài $15\text{ cm}$. Đoạn thẳng $AC$ gồm $AB$ và $BC$ dài là
\choice
{\True $47\text{ cm}$}
{$46\text{ cm}$}
{$48\text{ cm}$}
{$45\text{ cm}$}
\loigiai{
$32 + 15 = 47\text{ cm}$.\\
\textbf{Đáp án của câu hỏi là A}
}
\end{ex}

\begin{ex} %[2D1C05]
Tổng của hai số là 68. Nếu giữ nguyên một số hạng và bớt số hạng kia đi 8 đơn vị thì tổng mới là
\choice
{$76$}
{$68$}
{$58$}
{\True $60$}
\loigiai{
$68 - 8 = 60$.\\
\textbf{Đáp án của câu hỏi là D}
}
\end{ex}

\textbf{Phần 2. Trắc nghiệm đúng sai}
\setcounterpageref{ex}{0}

\begin{ex} %[2D1H05]
Xét tính đúng sai của các phép tính:
\choiceTF
{\True $25 + 34 = 59$.}
{\True $78 - 45 = 33$.}
{$60 + 30 = 80$.}
{\True $60 + 30 = 90$.}
\loigiai{
\begin{itemize}
    \item $25 + 34 = 59$: Đúng.
    \item $78 - 45 = 33$: Đúng.
    \item $60 + 30 = 90 \neq 80$. Ý 3 Sai, Ý 4 Đúng.
\end{itemize}
\begin{itemize}
    \item a) \textbf{Đúng}.
    \item b) \textbf{Đúng}.
    \item c) \textbf{Sai}.
    \item d) \textbf{Đúng}.
\end{itemize}
}
\end{ex}

\begin{ex} %[2D1V05]
Một cửa hàng buổi sáng bán được 35 gói kẹo, buổi chiều bán được 24 gói kẹo. Xét tính đúng sai:
\choiceTF
{\True Cả ngày bán được 59 gói kẹo ($35 + 24 = 59$).}
{\True Buổi sáng bán nhiều hơn buổi chiều 11 gói kẹo ($35 - 24 = 11$).}
{Nếu buổi chiều bán thêm 5 gói kẹo thì tổng số kẹo bán trong ngày là 65 gói.}
{\True Khi đó cả ngày bán được $59 + 5 = 64$ gói kẹo.}
\loigiai{
\begin{itemize}
    \item Tổng: $35 + 24 = 59$: Đúng.
    \item Hiệu: $35 - 24 = 11$: Đúng.
    \item $59 + 5 = 64 \neq 65$. Ý 3 Sai, Ý 4 Đúng.
\end{itemize}
\begin{itemize}
    \item a) \textbf{Đúng}.
    \item b) \textbf{Đúng}.
    \item c) \textbf{Sai}.
    \item d) \textbf{Đúng}.
\end{itemize}
}
\end{ex}

\textbf{Phần 3. Trả lời ngắn}
\setcounterpageref{ex}{0}

\begin{ex} %[2D1H05]
Tính tổng: $42 + 35$.
\shortans{77}
\loigiai{
\textbf{Đáp số:} $77$\\
$42 + 35 = 77$.
}
\end{ex}

\begin{ex} %[2D1H05]
Tính hiệu: $89 - 46$.
\shortans{43}
\loigiai{
\textbf{Đáp số:} $43$\\
$89 - 46 = 43$.
}
\end{ex}

\begin{ex} %[2D1V05]
Tính giá trị biểu thức: $30 + 50 - 20$.
\shortans{60}
\loigiai{
\textbf{Đáp số:} $60$\\
$80 - 20 = 60$.
}
\end{ex}

\begin{ex} %[2D1V05]
Tìm $x$ biết: $x + 23 = 58$.
\shortans{35}
\loigiai{
\textbf{Đáp số:} $35$\\
$x = 58 - 23 = 35$.
}
\end{ex}

\begin{ex} %[2D1V05]
Một vườn cây có 45 cây chuối và 32 cây ổi. Tính tổng số cây chuối và ổi trong vườn.
\shortans{77}
\loigiai{
\textbf{Đáp số:} $77$\\
$45 + 32 = 77$ cây.
}
\end{ex}

\begin{ex} %[2D1C05]
Tìm một số biết rằng lấy số đó cộng với 25 rồi trừ đi 10 thì được kết quả là 65.
\shortans{50}
\loigiai{
\textbf{Đáp số:} $50$\\
Số đó là: $65 + 10 - 25 = 50$.
}
\end{ex}

\newpage

\TaoPhieuKiemTra{Chủ đề 1 -- Ôn tập và bổ sung}{06}{Luyện tập chung}{https://loptoanthaylinh.vn/thi-online/toan-2/bai-06}
\textbf{Phần 1. Trắc nghiệm lựa chọn}
\setcounterpageref{ex}{0}

\begin{ex} %[2D1N06]
Số gồm $7$ chục và $0$ đơn vị là
\choice
{$7$}
{$77$}
{$17$}
{\True $70$}
\loigiai{
Số đó là $70$.\\
\textbf{Đáp án của câu hỏi là D}
}
\end{ex}

\begin{ex} %[2D1N06]
Số liền sau của số $89$ là
\choice
{$88$}
{$91$}
{\True $90$}
{$99$}
\loigiai{
$89 + 1 = 90$.\\
\textbf{Đáp án của câu hỏi là C}
}
\end{ex}

\begin{ex} %[2D1H06]
Trong phép trừ $75 - 32 = 43$, số trừ là
\choice
{$75$}
{$43$}
{$107$}
{\True $32$}
\loigiai{
Số trừ là 32.\\
\textbf{Đáp án của câu hỏi là D}
}
\end{ex}

\begin{ex} %[2D1H06]
Kết quả của phép tính $43 + 36$ là
\choice
{\True $79$}
{$78$}
{$80$}
{$69$}
\loigiai{
$43 + 36 = 79$.\\
\textbf{Đáp án của câu hỏi là A}
}
\end{ex}

\begin{ex} %[2D1H06]
Kết quả của phép tính $96 - 54$ là
\choice
{$41$}
{$43$}
{\True $42$}
{$44$}
\loigiai{
$96 - 54 = 42$.\\
\textbf{Đáp án của câu hỏi là C}
}
\end{ex}

\begin{ex} %[2D1H06]
Số 58 hơn số 20 bao nhiêu đơn vị?
\choice
{$28$}
{\True $38$}
{$48$}
{$78$}
\loigiai{
$58 - 20 = 38$.\\
\textbf{Đáp án của câu hỏi là B}
}
\end{ex}

\begin{ex} %[2D1V06]
Sắp xếp các số $42, 24, 65, 56$ theo thứ tự từ lớn đến bé:
\choice
{\True $65, 56, 42, 24$}
{$24, 42, 56, 65$}
{$65, 42, 56, 24$}
{$56, 65, 42, 24$}
\loigiai{
Thứ tự giảm dần: $65 > 56 > 42 > 24$.\\
\textbf{Đáp án của câu hỏi là A}
}
\end{ex}

\begin{ex} %[2D1V06]
Một sợi dây dài $65\text{ cm}$, người ta cắt bớt đi $23\text{ cm}$. Sợi dây còn lại dài là
\choice
{$41\text{ cm}$}
{\True $42\text{ cm}$}
{$43\text{ cm}$}
{$44\text{ cm}$}
\loigiai{
$65 - 23 = 42\text{ cm}$.\\
\textbf{Đáp án của câu hỏi là B}
}
\end{ex}

\begin{ex} %[2D1V06]
Tính giá trị biểu thức: $32 + 14 + 20$.
\choice
{\True $66$}
{$64$}
{$68$}
{$70$}
\loigiai{
$46 + 20 = 66$.\\
\textbf{Đáp án của câu hỏi là A}
}
\end{ex}

\begin{ex} %[2D1V06]
Tìm $x$ biết: $x - 30 = 25 + 14$.
\choice
{$68$}
{$70$}
{$59$}
{\True $69$}
\loigiai{
$x - 30 = 39 \Rightarrow x = 39 + 30 = 69$.\\
\textbf{Đáp án của câu hỏi là D}
}
\end{ex}

\begin{ex} %[2D1V06]
Lớp 2A có 20 bạn nam và 15 bạn nữ. Hỏi lớp 2A có tất cả bao nhiêu học sinh?
\choice
{$34$ bạn}
{\True $35$ bạn}
{$36$ bạn}
{$37$ bạn}
\loigiai{
$20 + 15 = 35$ bạn.\\
\textbf{Đáp án của câu hỏi là B}
}
\end{ex}

\begin{ex} %[2D1C06]
Tìm số có hai chữ số biết chữ số hàng chục gấp đôi chữ số hàng đơn vị và tổng hai chữ số bằng 9.
\choice
{$36$}
{$42$}
{\True $63$}
{$84$}
\loigiai{
$6 = 3 \times 2$ và $6 + 3 = 9$. Số đó là $63$.\\
\textbf{Đáp án của câu hỏi là C}
}
\end{ex}

\textbf{Phần 2. Trắc nghiệm đúng sai}
\setcounterpageref{ex}{0}

\begin{ex} %[2D1H06]
Xét tính đúng sai về số học và phép tính:
\choiceTF
{\True Số liền trước của số 100 là 99.}
{\True $40 + 50 = 90$.}
{Số 75 gồm 5 chục và 7 đơn vị.}
{\True Số 75 gồm 7 chục và 5 đơn vị.}
\loigiai{
\begin{itemize}
    \item $100 - 1 = 99$: Đúng.
    \item $40 + 50 = 90$: Đúng.
    \item 75 gồm 7 chục và 5 đơn vị (không phải 5 chục). Ý 3 Sai, Ý 4 Đúng.
\end{itemize}
\begin{itemize}
    \item a) \textbf{Đúng}.
    \item b) \textbf{Đúng}.
    \item c) \textbf{Sai}.
    \item d) \textbf{Đúng}.
\end{itemize}
}
\end{ex}

\begin{ex} %[2D1V06]
Một trang trại có 46 con thỏ trắng và 32 con thỏ nâu. Xét tính đúng sai:
\choiceTF
{\True Trang trại có tất cả 78 con thỏ ($46 + 32 = 78$).}
{\True Số thỏ trắng nhiều hơn thỏ nâu là 14 con ($46 - 32 = 14$).}
{Số thỏ nâu nhiều hơn thỏ trắng 14 con.}
{\True Thỏ trắng nhiều hơn thỏ nâu nên thỏ nâu ít hơn thỏ trắng 14 con.}
\loigiai{
\begin{itemize}
    \item Tổng: $46 + 32 = 78$: Đúng.
    \item Trắng hơn nâu: $46 - 32 = 14$: Đúng.
    \item Thỏ nâu ít hơn (không phải nhiều hơn). Ý 3 Sai, Ý 4 Đúng.
\end{itemize}
\begin{itemize}
    \item a) \textbf{Đúng}.
    \item b) \textbf{Đúng}.
    \item c) \textbf{Sai}.
    \item d) \textbf{Đúng}.
\end{itemize}
}
\end{ex}

\textbf{Phần 3. Trả lời ngắn}
\setcounterpageref{ex}{0}

\begin{ex} %[2D1H06]
Viết số liền trước của số $60$.
\shortans{59}
\loigiai{
\textbf{Đáp số:} $59$\\
$60 - 1 = 59$.
}
\end{ex}

\begin{ex} %[2D1H06]
Tính tổng: $53 + 24$.
\shortans{77}
\loigiai{
\textbf{Đáp số:} $77$\\
$53 + 24 = 77$.
}
\end{ex}

\begin{ex} %[2D1V06]
Tính hiệu: $78 - 35$.
\shortans{43}
\loigiai{
\textbf{Đáp số:} $43$\\
$78 - 35 = 43$.
}
\end{ex}

\begin{ex} %[2D1V06]
Tìm $x$ biết: $x - 24 = 43$.
\shortans{67}
\loigiai{
\textbf{Đáp số:} $67$\\
$x = 43 + 24 = 67$.
}
\end{ex}

\begin{ex} %[2D1V06]
Số lớn nhất có hai chữ số hơn số tròn chục lớn nhất có hai chữ số bao nhiêu đơn vị?
\shortans{9}
\loigiai{
\textbf{Đáp số:} $9$\\
$99 - 90 = 9$.
}
\end{ex}

\begin{ex} %[2D1C06]
Tìm số tròn chục nhỏ nhất có hai chữ số mà tổng hai chữ số bằng 5.
\shortans{50}
\loigiai{
\textbf{Đáp số:} $50$\\
Số tròn chục có dạng $\overline{a0} \Rightarrow a = 5 \Rightarrow 50$.
}
\end{ex}

\newpage

% =========================================================================
% CHỦ ĐỀ 2 - PHÉP CỘNG, PHÉP TRỪ TRONG PHẠM VI 20
% =========================================================================
\section{Phép cộng, phép trừ trong phạm vi 20}

\TaoPhieuKiemTra{Chủ đề 2 -- Phép cộng, phép trừ trong phạm vi 20}{07}{Phép cộng có tổng là số tròn chục}{https://loptoanthaylinh.vn/thi-online/toan-2/bai-07}
\textbf{Phần 1. Trắc nghiệm lựa chọn}
\setcounterpageref{ex}{0}

\begin{ex} %[2D2N07]
Kết quả của phép tính $9 + 2$ là
\choice
{$10$}
{$12$}
{$13$}
{\True $11$}
\loigiai{
$9 + 2 = 11$.\\
\textbf{Đáp án của câu hỏi là D}
}
\end{ex}

\begin{ex} %[2D2N07]
Kết quả của phép tính $8 + 3$ là
\choice
{\True $11$}
{$10$}
{$12$}
{$13$}
\loigiai{
$8 + 3 = 11$.\\
\textbf{Đáp án của câu hỏi là A}
}
\end{ex}

\begin{ex} %[2D2H07]
Kết quả của phép tính $9 + 5$ là
\choice
{$13$}
{$15$}
{\True $14$}
{$16$}
\loigiai{
$9 + 5 = 14$.\\
\textbf{Đáp án của câu hỏi là C}
}
\end{ex}

\begin{ex} %[2D2H07]
Kết quả của phép tính $8 + 6$ là
\choice
{$13$}
{\True $14$}
{$15$}
{$12$}
\loigiai{
$8 + 6 = 14$.\\
\textbf{Đáp án của câu hỏi là B}
}
\end{ex}

\begin{ex} %[2D2H07]
Kết quả của phép tính $7 + 5$ là
\choice
{$11$}
{$13$}
{$14$}
{\True $12$}
\loigiai{
$7 + 5 = 12$.\\
\textbf{Đáp án của câu hỏi là D}
}
\end{ex}

\begin{ex} %[2D2H07]
Kết quả của phép tính $6 + 7$ là
\choice
{$12$}
{$14$}
{\True $13$}
{$15$}
\loigiai{
$6 + 7 = 13$.\\
\textbf{Đáp án của câu hỏi là C}
}
\end{ex}

\begin{ex} %[2D2V07]
Điền dấu thích hợp: $9 + 4 \dots 8 + 6$.
\choice
{\True $<$}
{$>$}
{$=$}
{Không so sánh được}
\loigiai{
$9 + 4 = 13$ và $8 + 6 = 14$. Vì $13 < 14$ nên điền dấu $<$.\\
\textbf{Đáp án của câu hỏi là A}
}
\end{ex}

\begin{ex} %[2D2V07]
Tìm $x$ biết: $x - 6 = 8$.
\choice
{$13$}
{\True $14$}
{$15$}
{$12$}
\loigiai{
$x = 8 + 6 = 14$.\\
\textbf{Đáp án của câu hỏi là B}
}
\end{ex}

\begin{ex} %[2D2V07]
Nam có 9 viên bi, An có 7 viên bi. Cả hai bạn có tất cả số viên bi là
\choice
{$15$ viên}
{$17$ viên}
{\True $16$ viên}
{$18$ viên}
\loigiai{
$9 + 7 = 16$ viên bi.\\
\textbf{Đáp án của câu hỏi là C}
}
\end{ex}

\begin{ex} %[2D2V07]
Mẹ mua 8 quả táo và 7 quả cam. Mẹ mua tất cả số quả là
\choice
{\True $15$ quả}
{$14$ quả}
{$16$ quả}
{$17$ quả}
\loigiai{
$8 + 7 = 15$ quả.\\
\textbf{Đáp án của câu hỏi là A}
}
\end{ex}

\begin{ex} %[2D2V07]
Tính giá trị biểu thức: $9 + 3 + 4$.
\choice
{$15$}
{\True $16$}
{$17$}
{$18$}
\loigiai{
$12 + 4 = 16$.\\
\textbf{Đáp án của câu hỏi là B}
}
\end{ex}

\begin{ex} %[2D2C07]
Tìm số thích hợp điền vào ô trống: $8 + \dots = 9 + 6$.
\choice
{$6$}
{$8$}
{$9$}
{\True $7$}
\loigiai{
$9 + 6 = 15 \Rightarrow 8 + 7 = 15$.\\
\textbf{Đáp án của câu hỏi là D}
}
\end{ex}

\textbf{Phần 2. Trắc nghiệm đúng sai}
\setcounterpageref{ex}{0}

\begin{ex} %[2D2H07]
Xét tính đúng sai của các phép cộng:
\choiceTF
{\True $9 + 6 = 15$.}
{\True $8 + 7 = 15$.}
{$7 + 6 = 14$.}
{\True $7 + 6 = 13$.}
\loigiai{
\begin{itemize}
    \item $9 + 6 = 15$: Đúng.
    \item $8 + 7 = 15$: Đúng.
    \item $7 + 6 = 13 \neq 14$. Ý 3 Sai, Ý 4 Đúng.
\end{itemize}
\begin{itemize}
    \item a) \textbf{Đúng}.
    \item b) \textbf{Đúng}.
    \item c) \textbf{Sai}.
    \item d) \textbf{Đúng}.
\end{itemize}
}
\end{ex}

\begin{ex} %[2D2V07]
Cho biểu thức $A = 9 + 7$ và $B = 8 + 8$. Xét tính đúng sai:
\choiceTF
{\True Giá trị của biểu thức $A$ là 16.}
{\True Giá trị của biểu thức $B$ là 16.}
{Biểu thức $A$ có giá trị lớn hơn biểu thức $B$.}
{\True Hai biểu thức $A$ và $B$ có giá trị bằng nhau ($A = B = 16$).}
\loigiai{
\begin{itemize}
    \item $A = 9 + 7 = 16$: Đúng.
    \item $B = 8 + 8 = 16$: Đúng.
    \item $A = B$ (không lớn hơn). Ý 3 Sai, Ý 4 Đúng.
\end{itemize}
\begin{itemize}
    \item a) \textbf{Đúng}.
    \item b) \textbf{Đúng}.
    \item c) \textbf{Sai}.
    \item d) \textbf{Đúng}.
\end{itemize}
}
\end{ex}

\textbf{Phần 3. Trả lời ngắn}
\setcounterpageref{ex}{0}

\begin{ex} %[2D2H07]
Tính tổng: $9 + 8$.
\shortans{17}
\loigiai{
\textbf{Đáp số:} $17$\\
$9 + 8 = 17$.
}
\end{ex}

\begin{ex} %[2D2H07]
Tính tổng: $8 + 5$.
\shortans{13}
\loigiai{
\textbf{Đáp số:} $13$\\
$8 + 5 = 13$.
}
\end{ex}

\begin{ex} %[2D2V07]
Tính tổng: $7 + 7$.
\shortans{14}
\loigiai{
\textbf{Đáp số:} $14$\\
$7 + 7 = 14$.
}
\end{ex}

\begin{ex} %[2D2V07]
Tính giá trị biểu thức: $6 + 8 + 2$.
\shortans{16}
\loigiai{
\textbf{Đáp số:} $16$\\
$6 + 10 = 16$.
}
\end{ex}

\begin{ex} %[2D2V07]
Tổ 1 có 8 bạn, tổ 2 có 9 bạn. Cả hai tổ có tất cả bao nhiêu bạn?
\shortans{17}
\loigiai{
\textbf{Đáp số:} $17$\\
$8 + 9 = 17$ bạn.
}
\end{ex}

\begin{ex} %[2D2C07]
Tìm số thích hợp điền vào dấu ?: $9 + ? = 18$.
\shortans{9}
\loigiai{
\textbf{Đáp số:} $9$\\
$18 - 9 = 9$.
}
\end{ex}

\newpage

\TaoPhieuKiemTra{Chủ đề 2 -- Phép cộng, phép trừ trong phạm vi 20}{08}{Phép cộng (có nhớ) trong phạm vi 20}{https://loptoanthaylinh.vn/thi-online/toan-2/bai-08}
\textbf{Phần 1. Trắc nghiệm lựa chọn}
\setcounterpageref{ex}{0}

\begin{ex} %[2D2N08]
Trong bảng cộng (qua 10), phép tính $9 + 9$ có kết quả là
\choice
{$17$}
{\True $18$}
{$19$}
{$20$}
\loigiai{
$9 + 9 = 18$.\\
\textbf{Đáp án của câu hỏi là B}
}
\end{ex}

\begin{ex} %[2D2N08]
Phép tính nào sau đây có kết quả bằng 11?
\choice
{$9 + 3$}
{$8 + 4$}
{$7 + 5$}
{\True $9 + 2$}
\loigiai{
$9 + 2 = 11$.\\
\textbf{Đáp án của câu hỏi là D}
}
\end{ex}

\begin{ex} %[2D2H08]
Phép tính nào sau đây có kết quả bằng 13?
\choice
{$8 + 4$}
{\True $8 + 5$}
{$9 + 5$}
{$7 + 7$}
\loigiai{
$8 + 5 = 13$.\\
\textbf{Đáp án của câu hỏi là B}
}
\end{ex}

\begin{ex} %[2D2H08]
Phép tính nào sau đây có kết quả lớn nhất?
\choice
{$8 + 7$}
{$9 + 7$}
{$8 + 8$}
{\True $9 + 8$}
\loigiai{
$9 + 8 = 17$ (lớn nhất).\\
\textbf{Đáp án của câu hỏi là D}
}
\end{ex}

\begin{ex} %[2D2H08]
Phép tính nào sau đây có kết quả bé nhất?
\choice
{$8 + 4$}
{$7 + 5$}
{\True $6 + 5$}
{$9 + 3$}
\loigiai{
$6 + 5 = 11$ (bé nhất).\\
\textbf{Đáp án của câu hỏi là C}
}
\end{ex}

\begin{ex} %[2D2H08]
Điền số thích hợp: $7 + \dots = 16$.
\choice
{\True $9$}
{$8$}
{$7$}
{$6$}
\loigiai{
$7 + 9 = 16$.\\
\textbf{Đáp án của câu hỏi là A}
}
\end{ex}

\begin{ex} %[2D2V08]
Có bao nhiêu phép cộng hai số có 1 chữ số có kết quả bằng 12? (Ví dụ: $9+3, 8+4, 7+5, 6+6, 5+7, 4+8, 3+9$)
\choice
{$5$ phép tính}
{$6$ phép tính}
{\True $7$ phép tính}
{$8$ phép tính}
\loigiai{
Các phép tính: $9+3, 8+4, 7+5, 6+6, 5+7, 4+8, 3+9$ (có 7 phép tính).\\
\textbf{Đáp án của câu hỏi là C}
}
\end{ex}

\begin{ex} %[2D2V08]
Tính: $8 + 7 + 3 = \dots$
\choice
{$17$}
{$19$}
{$20$}
{\True $18$}
\loigiai{
$8 + (7 + 3) = 8 + 10 = 18$.\\
\textbf{Đáp án của câu hỏi là D}
}
\end{ex}

\begin{ex} %[2D2V08]
Điền dấu thích hợp: $7 + 8 \dots 9 + 6$.
\choice
{$>$}
{\True $=$}
{$<$}
{Không so sánh được}
\loigiai{
$7 + 8 = 15$ và $9 + 6 = 15$. Vậy $7 + 8 = 9 + 6$.\\
\textbf{Đáp án của câu hỏi là B}
}
\end{ex}

\begin{ex} %[2D2V08]
Một luống hoa có 8 cây hoa hồng và 6 cây hoa cúc. Số cây hoa có tất cả là
\choice
{\True $14$ cây}
{$13$ cây}
{$15$ cây}
{$16$ cây}
\loigiai{
$8 + 6 = 14$ cây.\\
\textbf{Đáp án của câu hỏi là A}
}
\end{ex}

\begin{ex} %[2D2V08]
Tìm $x$ biết: $x - 7 = 9$.
\choice
{$15$}
{$17$}
{\True $16$}
{$14$}
\loigiai{
$x = 9 + 7 = 16$.\\
\textbf{Đáp án của câu hỏi là C}
}
\end{ex}

\begin{ex} %[2D2C08]
Tổng của hai số có 1 chữ số là 17. Hai số đó là
\choice
{\True $9$ và $8$}
{$8$ và $8$}
{$9$ và $9$}
{$7$ và $9$}
\loigiai{
Vì các số có 1 chữ số nên $9 + 8 = 17$.\\
\textbf{Đáp án của câu hỏi là A}
}
\end{ex}

\textbf{Phần 2. Trắc nghiệm đúng sai}
\setcounterpageref{ex}{0}

\begin{ex} %[2D2H08]
Xét tính đúng sai của các phép tính trong bảng cộng:
\choiceTF
{\True $9 + 4 = 13$.}
{\True $8 + 6 = 14$.}
{$7 + 8 = 16$.}
{\True $7 + 8 = 15$.}
\loigiai{
\begin{itemize}
    \item $9 + 4 = 13$: Đúng.
    \item $8 + 6 = 14$: Đúng.
    \item $7 + 8 = 15 \neq 16$. Ý 3 Sai, Ý 4 Đúng.
\end{itemize}
\begin{itemize}
    \item a) \textbf{Đúng}.
    \item b) \textbf{Đúng}.
    \item c) \textbf{Sai}.
    \item d) \textbf{Đúng}.
\end{itemize}
}
\end{ex}

\begin{ex} %[2D2V08]
Cho các phép tính: (1) $9 + 5$, (2) $8 + 6$, (3) $7 + 7$. Xét tính đúng sai:
\choiceTF
{\True Phép tính (1) có kết quả bằng 14.}
{\True Phép tính (2) có kết quả bằng 14.}
{Phép tính (3) có kết quả lớn hơn phép tính (1).}
{\True Cả 3 phép tính đều có kết quả bằng 14.}
\loigiai{
\begin{itemize}
    \item $9 + 5 = 14$: Đúng.
    \item $8 + 6 = 14$: Đúng.
    \item Cả 3 đều bằng 14. Ý 3 Sai, Ý 4 Đúng.
\end{itemize}
\begin{itemize}
    \item a) \textbf{Đúng}.
    \item b) \textbf{Đúng}.
    \item c) \textbf{Sai}.
    \item d) \textbf{Đúng}.
\end{itemize}
}
\end{ex}

\textbf{Phần 3. Trả lời ngắn}
\setcounterpageref{ex}{0}

\begin{ex} %[2D2H08]
Tính kết quả: $6 + 9$.
\shortans{15}
\loigiai{
\textbf{Đáp số:} $15$\\
$6 + 9 = 15$.
}
\end{ex}

\begin{ex} %[2D2H08]
Tính kết quả: $8 + 9$.
\shortans{17}
\loigiai{
\textbf{Đáp số:} $17$\\
$8 + 9 = 17$.
}
\end{ex}

\begin{ex} %[2D2V08]
Tính giá trị biểu thức: $9 + 5 + 1$.
\shortans{15}
\loigiai{
\textbf{Đáp số:} $15$\\
$(9 + 1) + 5 = 15$.
}
\end{ex}

\begin{ex} %[2D2V08]
Điền số thích hợp: $8 + ? = 17$.
\shortans{9}
\loigiai{
\textbf{Đáp số:} $9$\\
$17 - 8 = 9$.
}
\end{ex}

\begin{ex} %[2D2V08]
Bình có 7 cái nhãn vở, mẹ cho Bình thêm 8 cái nhãn vở. Hỏi Bình có tất cả bao nhiêu cái nhãn vở?
\shortans{15}
\loigiai{
\textbf{Đáp số:} $15$\\
$7 + 8 = 15$ cái.
}
\end{ex}

\begin{ex} %[2D2C08]
Tìm tổng lớn nhất có thể của hai số có 1 chữ số.
\shortans{18}
\loigiai{
\textbf{Đáp số:} $18$\\
$9 + 9 = 18$.
}
\end{ex}

\newpage

\TaoPhieuKiemTra{Chủ đề 2 -- Phép cộng, phép trừ trong phạm vi 20}{09}{Phép trừ (có nhớ) trong phạm vi 20}{https://loptoanthaylinh.vn/thi-online/toan-2/bai-09}
\textbf{Phần 1. Trắc nghiệm lựa chọn}
\setcounterpageref{ex}{0}

\begin{ex} %[2D2N09]
Trên cành cây có 8 con chim, có thêm 4 con chim bay đến. Hỏi trên cành cây có tất cả bao nhiêu con chim?
\choice
{$11$ con}
{$13$ con}
{$14$ con}
{\True $12$ con}
\loigiai{
$8 + 4 = 12$ con chim.\\
\textbf{Đáp án của câu hỏi là D}
}
\end{ex}

\begin{ex} %[2D2N09]
Trong bể có 15 con cá, vớt ra 6 con cá. Trong bể còn lại số con cá là
\choice
{\True $9$ con}
{$8$ con}
{$10$ con}
{$11$ con}
\loigiai{
$15 - 6 = 9$ con cá.\\
\textbf{Đáp án của câu hỏi là A}
}
\end{ex}

\begin{ex} %[2D2H09]
Mai có 9 cái kẹo, mẹ cho Mai thêm 5 cái kẹo nữa. Mai có tất cả số cái kẹo là
\choice
{$13$ cái}
{$15$ cái}
{\True $14$ cái}
{$16$ cái}
\loigiai{
$9 + 5 = 14$ cái kẹo.\\
\textbf{Đáp án của câu hỏi là C}
}
\end{ex}

\begin{ex} %[2D2H09]
Lan có 14 quả bóng bay, Lan làm vỡ 5 quả. Lan còn lại số quả bóng bay là
\choice
{$8$ quả}
{\True $9$ quả}
{$10$ quả}
{$11$ quả}
\loigiai{
$14 - 5 = 9$ quả.\\
\textbf{Đáp án của câu hỏi là B}
}
\end{ex}

\begin{ex} %[2D2H09]
Một xe buýt có 12 hành khách, đến bến đỗ có thêm 6 hành khách lên xe. Trên xe có tất cả số hành khách là
\choice
{$17$ người}
{$19$ người}
{\True $18$ người}
{$20$ người}
\loigiai{
$12 + 6 = 18$ người.\\
\textbf{Đáp án của câu hỏi là C}
}
\end{ex}

\begin{ex} %[2D2H09]
Một bãi đỗ xe có 16 chiếc ô tô, sau đó có 7 chiếc ô tô rời bãi. Trong bãi còn lại số ô tô là
\choice
{\True $9$ chiếc}
{$8$ chiếc}
{$10$ chiếc}
{$11$ chiếc}
\loigiai{
$16 - 7 = 9$ chiếc.\\
\textbf{Đáp án của câu hỏi là A}
}
\end{ex}

\begin{ex} %[2D2V09]
Hà có 8 quyển vở, Hà được thưởng thêm 7 quyển vở nữa. Hỏi Hà có tất cả bao nhiêu quyển vở?
\choice
{$14$ quyển}
{\True $15$ quyển}
{$16$ quyển}
{$17$ quyển}
\loigiai{
$8 + 7 = 15$ quyển vở.\\
\textbf{Đáp án của câu hỏi là B}
}
\end{ex}

\begin{ex} %[2D2V09]
Đàn gà có 17 con, đã bán đi 8 con. Hỏi đàn gà còn lại bao nhiêu con?
\choice
{$8$ con}
{$10$ con}
{$11$ con}
{\True $9$ con}
\loigiai{
$17 - 8 = 9$ con.\\
\textbf{Đáp án của câu hỏi là D}
}
\end{ex}

\begin{ex} %[2D2V09]
Hộp thứ nhất có 9 cái bút, hộp thứ hai có nhiều hơn hộp thứ nhất 4 cái bút. Hỏi hộp thứ hai có bao nhiêu cái bút?
\choice
{\True $13$ cái}
{$12$ cái}
{$14$ cái}
{$15$ cái}
\loigiai{
$9 + 4 = 13$ cái bút.\\
\textbf{Đáp án của câu hỏi là A}
}
\end{ex}

\begin{ex} %[2D2V09]
Cửa hàng có 13 thùng sữa, đã bán đi 6 thùng. Cửa hàng còn lại số thùng sữa là
\choice
{$6$ thùng}
{$8$ thùng}
{$9$ thùng}
{\True $7$ thùng}
\loigiai{
$13 - 6 = 7$ thùng.\\
\textbf{Đáp án của câu hỏi là D}
}
\end{ex}

\begin{ex} %[2D2V09]
Trên cành có 9 quả táo chín và 8 quả táo xanh. Trên cành có tất cả số quả táo là
\choice
{$16$ quả}
{$18$ quả}
{\True $17$ quả}
{$19$ quả}
\loigiai{
$9 + 8 = 17$ quả.\\
\textbf{Đáp án của câu hỏi là C}
}
\end{ex}

\begin{ex} %[2D2C09]
Một xe khách có 15 hành khách. Đến một bến đỗ, có 4 người xuống xe và 6 người lên xe. Hỏi trên xe lúc này có bao nhiêu hành khách?
\choice
{$16$ người}
{\True $17$ người}
{$18$ người}
{$19$ người}
\loigiai{
$15 - 4 + 6 = 17$ người.\\
\textbf{Đáp án của câu hỏi là B}
}
\end{ex}

\textbf{Phần 2. Trắc nghiệm đúng sai}
\setcounterpageref{ex}{0}

\begin{ex} %[2D2H09]
Bác An có 9 con vịt, bác mua thêm 6 con vịt nữa. Xét tính đúng sai:
\choiceTF
{\True Bác An mua thêm nghĩa là ta làm phép tính cộng.}
{\True Bác An có tất cả 15 con vịt ($9 + 6 = 15$).}
{Nếu bác An bán bớt 5 con vịt thì còn lại 11 con.}
{\True Khi đó bác An còn lại $15 - 5 = 10$ con vịt.}
\loigiai{
\begin{itemize}
    \item Mua thêm là phép cộng: Đúng.
    \item $9 + 6 = 15$: Đúng.
    \item $15 - 5 = 10 \neq 11$. Ý 3 Sai, Ý 4 Đúng.
\end{itemize}
\begin{itemize}
    \item a) \textbf{Đúng}.
    \item b) \textbf{Đúng}.
    \item c) \textbf{Sai}.
    \item d) \textbf{Đúng}.
\end{itemize}
}
\end{ex}

\begin{ex} %[2D2V09]
Một đĩa có 16 quả dâu tây, bé Nam ăn mất 7 quả. Xét tính đúng sai:
\choiceTF
{\True Số dâu tây còn lại là 9 quả ($16 - 7 = 9$).}
{\True Bé Nam ăn mất là bài toán về bớt một số đơn vị.}
{Nếu mẹ cho thêm vào đĩa 4 quả thì trong đĩa có 12 quả.}
{\True Khi đó trong đĩa có $9 + 4 = 13$ quả dâu tây.}
\loigiai{
\begin{itemize}
    \item Còn lại: $16 - 7 = 9$: Đúng.
    \item Bài toán bớt đi: Đúng.
    \item $9 + 4 = 13 \neq 12$. Ý 3 Sai, Ý 4 Đúng.
\end{itemize}
\begin{itemize}
    \item a) \textbf{Đúng}.
    \item b) \textbf{Đúng}.
    \item c) \textbf{Sai}.
    \item d) \textbf{Đúng}.
\end{itemize}
}
\end{ex}

\textbf{Phần 3. Trả lời ngắn}
\setcounterpageref{ex}{0}

\begin{ex} %[2D2H09]
Có 8 con bướm, thêm 5 con bướm nữa bay tới. Hỏi có tất cả bao nhiêu con bướm?
\shortans{13}
\loigiai{
\textbf{Đáp số:} $13$\\
$8 + 5 = 13$ con.
}
\end{ex}

\begin{ex} %[2D2H09]
Có 15 bông hoa, cắm vào lọ 7 bông hoa. Còn lại bao nhiêu bông hoa chưa cắm?
\shortans{8}
\loigiai{
\textbf{Đáp số:} $8$\\
$15 - 7 = 8$ bông.
}
\end{ex}

\begin{ex} %[2D2V09]
Lan có 9 cái kẹo, mẹ cho thêm 7 cái kẹo. Hỏi Lan có tất cả bao nhiêu cái kẹo?
\shortans{16}
\loigiai{
\textbf{Đáp số:} $16$\\
$9 + 7 = 16$ cái.
}
\end{ex}

\begin{ex} %[2D2V09]
Một tổ có 14 bạn, có 6 bạn đi làm nhiệm vụ khác. Hỏi tổ đó còn lại bao nhiêu bạn?
\shortans{8}
\loigiai{
\textbf{Đáp số:} $8$\\
$14 - 6 = 8$ bạn.
}
\end{ex}

\begin{ex} %[2D2V09]
Trong chuồng có 8 con gà trống và 8 con gà mái. Hỏi trong chuồng có tất cả bao nhiêu con gà?
\shortans{16}
\loigiai{
\textbf{Đáp số:} $16$\\
$8 + 8 = 16$ con.
}
\end{ex}

\begin{ex} %[2D2C09]
Nam có 12 viên bi. Nam cho Bình 3 viên bi và được Dũng cho thêm 5 viên bi. Hỏi Nam có bao nhiêu viên bi?
\shortans{14}
\loigiai{
\textbf{Đáp số:} $14$\\
$12 - 3 + 5 = 14$ viên bi.
}
\end{ex}

\newpage

\TaoPhieuKiemTra{Chủ đề 2 -- Phép cộng, phép trừ trong phạm vi 20}{10}{Bảng cộng (có nhớ) trong phạm vi 20}{https://loptoanthaylinh.vn/thi-online/toan-2/bai-10}
\textbf{Phần 1. Trắc nghiệm lựa chọn}
\setcounterpageref{ex}{0}

\begin{ex} %[2D2N10]
Kết quả của phép tính $9 + 7$ là
\choice
{$15$}
{$17$}
{$18$}
{\True $16$}
\loigiai{
$9 + 7 = 16$.\\
\textbf{Đáp án của câu hỏi là D}
}
\end{ex}

\begin{ex} %[2D2N10]
Kết quả của phép tính $8 + 8$ là
\choice
{\True $16$}
{$15$}
{$17$}
{$14$}
\loigiai{
$8 + 8 = 16$.\\
\textbf{Đáp án của câu hỏi là A}
}
\end{ex}

\begin{ex} %[2D2H10]
Kết quả của phép tính $7 + 6$ là
\choice
{$12$}
{\True $13$}
{$14$}
{$15$}
\loigiai{
$7 + 6 = 13$.\\
\textbf{Đáp án của câu hỏi là B}
}
\end{ex}

\begin{ex} %[2D2H10]
Kết quả của phép tính $6 + 8$ là
\choice
{\True $14$}
{$13$}
{$15$}
{$16$}
\loigiai{
$6 + 8 = 14$.\\
\textbf{Đáp án của câu hỏi là A}
}
\end{ex}

\begin{ex} %[2D2H10]
Điền dấu thích hợp: $8 + 5 \dots 9 + 4$.
\choice
{$>$}
{$<$}
{Không so sánh được}
{\True $=$}
\loigiai{
$8 + 5 = 13$ và $9 + 4 = 13$. Vậy $8 + 5 = 9 + 4$.\\
\textbf{Đáp án của câu hỏi là D}
}
\end{ex}

\begin{ex} %[2D2H10]
Tìm $x$ biết: $x - 8 = 7$.
\choice
{$14$}
{$16$}
{\True $15$}
{$13$}
\loigiai{
$x = 7 + 8 = 15$.\\
\textbf{Đáp án của câu hỏi là C}
}
\end{ex}

\begin{ex} %[2D2V10]
Tính giá trị biểu thức: $9 + 4 + 3$.
\choice
{\True $16$}
{$15$}
{$17$}
{$18$}
\loigiai{
$13 + 3 = 16$.\\
\textbf{Đáp án của câu hỏi là A}
}
\end{ex}

\begin{ex} %[2D2V10]
Mai có 8 bông hoa đỏ và 7 bông hoa vàng. Mai có tất cả số bông hoa là
\choice
{$14$ bông}
{$16$ bông}
{\True $15$ bông}
{$17$ bông}
\loigiai{
$8 + 7 = 15$ bông hoa.\\
\textbf{Đáp án của câu hỏi là C}
}
\end{ex}

\begin{ex} %[2D2V10]
Tổ 1 có 9 bạn, có thêm 5 bạn chuyển vào tổ 1. Hỏi tổ 1 có tất cả bao nhiêu bạn?
\choice
{$13$ bạn}
{\True $14$ bạn}
{$15$ bạn}
{$16$ bạn}
\loigiai{
$9 + 5 = 14$ bạn.\\
\textbf{Đáp án của câu hỏi là B}
}
\end{ex}

\begin{ex} %[2D2V10]
Tính nhanh: $7 + 9 + 3$.
\choice
{$18$}
{$20$}
{$17$}
{\True $19$}
\loigiai{
$(7 + 3) + 9 = 10 + 9 = 19$.\\
\textbf{Đáp án của câu hỏi là D}
}
\end{ex}

\begin{ex} %[2D2V10]
Phép tính nào sau đây có kết quả lớn hơn 15?
\choice
{$8 + 7$}
{\True $9 + 8$}
{$9 + 6$}
{$8 + 6$}
\loigiai{
$9 + 8 = 17 > 15$.\\
\textbf{Đáp án của câu hỏi là B}
}
\end{ex}

\begin{ex} %[2D2C10]
Tìm số thích hợp điền vào ô trống: $8 + \dots > 16$.
\choice
{$8$}
{$7$}
{\True $9$}
{$6$}
\loigiai{
$8 + 9 = 17 > 16$.\\
\textbf{Đáp án của câu hỏi là C}
}
\end{ex}

\textbf{Phần 2. Trắc nghiệm đúng sai}
\setcounterpageref{ex}{0}

\begin{ex} %[2D2H10]
Xét tính đúng sai của các phép tính:
\choiceTF
{\True $9 + 8 = 17$.}
{\True $7 + 5 = 12$.}
{$6 + 7 = 14$.}
{\True $6 + 7 = 13$.}
\loigiai{
\begin{itemize}
    \item $9 + 8 = 17$: Đúng.
    \item $7 + 5 = 12$: Đúng.
    \item $6 + 7 = 13 \neq 14$. Ý 3 Sai, Ý 4 Đúng.
\end{itemize}
\begin{itemize}
    \item a) \textbf{Đúng}.
    \item b) \textbf{Đúng}.
    \item c) \textbf{Sai}.
    \item d) \textbf{Đúng}.
\end{itemize}
}
\end{ex}

\begin{ex} %[2D2V10]
Một đàn gà có 9 con gà mái và 8 con gà trống. Xét tính đúng sai:
\choiceTF
{\True Tổng số gà là 17 con ($9 + 8 = 17$).}
{\True Số gà mái nhiều hơn gà trống là 1 con ($9 - 8 = 1$).}
{Nếu mua thêm 4 con gà nữa thì đàn gà có 20 con.}
{\True Khi đó đàn gà có $17 + 4 = 21$ con gà.}
\loigiai{
\begin{itemize}
    \item Tổng: $9 + 8 = 17$: Đúng.
    \item Hiệu: $9 - 8 = 1$: Đúng.
    \item $17 + 4 = 21 \neq 20$. Ý 3 Sai, Ý 4 Đúng.
\end{itemize}
\begin{itemize}
    \item a) \textbf{Đúng}.
    \item b) \textbf{Đúng}.
    \item c) \textbf{Sai}.
    \item d) \textbf{Đúng}.
\end{itemize}
}
\end{ex}

\textbf{Phần 3. Trả lời ngắn}
\setcounterpageref{ex}{0}

\begin{ex} %[2D2H10]
Tính: $8 + 6$.
\shortans{14}
\loigiai{
\textbf{Đáp số:} $14$\\
$8 + 6 = 14$.
}
\end{ex}

\begin{ex} %[2D2H10]
Tính: $9 + 5$.
\shortans{14}
\loigiai{
\textbf{Đáp số:} $14$\\
$9 + 5 = 14$.
}
\end{ex}

\begin{ex} %[2D2V10]
Tính: $7 + 8$.
\shortans{15}
\loigiai{
\textbf{Đáp số:} $15$\\
$7 + 8 = 15$.
}
\end{ex}

\begin{ex} %[2D2V10]
Tính giá trị biểu thức: $8 + 5 + 2$.
\shortans{15}
\loigiai{
\textbf{Đáp số:} $15$\\
$(8 + 2) + 5 = 15$.
}
\end{ex}

\begin{ex} %[2D2V10]
Hà hái được 9 quả táo, An hái được 8 quả táo. Cả hai bạn hái được bao nhiêu quả táo?
\shortans{17}
\loigiai{
\textbf{Đáp số:} $17$\\
$9 + 8 = 17$ quả.
}
\end{ex}

\begin{ex} %[2D2C10]
Tìm số thích hợp điền vào dấu ?: $7 + ? = 16$.
\shortans{9}
\loigiai{
\textbf{Đáp số:} $9$\\
$16 - 7 = 9$.
}
\end{ex}

\newpage

\TaoPhieuKiemTra{Chủ đề 2 -- Phép cộng, phép trừ trong phạm vi 20}{11}{Bảng trừ (có nhớ) trong phạm vi 20}{https://loptoanthaylinh.vn/thi-online/toan-2/bai-11}
\textbf{Phần 1. Trắc nghiệm lựa chọn}
\setcounterpageref{ex}{0}

\begin{ex} %[2D2N11]
Kết quả của phép tính $11 - 2$ là
\choice
{$8$}
{$10$}
{$7$}
{\True $9$}
\loigiai{
$11 - 2 = 9$.\\
\textbf{Đáp án của câu hỏi là D}
}
\end{ex}

\begin{ex} %[2D2N11]
Kết quả của phép tính $12 - 3$ là
\choice
{$8$}
{\True $9$}
{$10$}
{$7$}
\loigiai{
$12 - 3 = 9$.\\
\textbf{Đáp án của câu hỏi là B}
}
\end{ex}

\begin{ex} %[2D2H11]
Kết quả của phép tính $11 - 5$ là
\choice
{\True $6$}
{$5$}
{$7$}
{$8$}
\loigiai{
$11 - 5 = 6$.\\
\textbf{Đáp án của câu hỏi là A}
}
\end{ex}

\begin{ex} %[2D2H11]
Kết quả của phép tính $13 - 7$ là
\choice
{$5$}
{$7$}
{$8$}
{\True $6$}
\loigiai{
$13 - 7 = 6$.\\
\textbf{Đáp án của câu hỏi là D}
}
\end{ex}

\begin{ex} %[2D2H11]
Kết quả của phép tính $14 - 8$ là
\choice
{\True $6$}
{$5$}
{$7$}
{$8$}
\loigiai{
$14 - 8 = 6$.\\
\textbf{Đáp án của câu hỏi là A}
}
\end{ex}

\begin{ex} %[2D2H11]
Kết quả của phép tính $15 - 9$ là
\choice
{$5$}
{$7$}
{\True $6$}
{$8$}
\loigiai{
$15 - 9 = 6$.\\
\textbf{Đáp án của câu hỏi là C}
}
\end{ex}

\begin{ex} %[2D2V11]
Điền dấu thích hợp: $13 - 6 \dots 14 - 8$.
\choice
{$<$}
{$=$}
{Không so sánh được}
{\True $>$}
\loigiai{
$13 - 6 = 7$ và $14 - 8 = 6$. Vì $7 > 6$ nên điền dấu $>$.\\
\textbf{Đáp án của câu hỏi là D}
}
\end{ex}

\begin{ex} %[2D2V11]
Tìm $x$ biết: $x + 6 = 15$.
\choice
{$8$}
{\True $9$}
{$10$}
{$7$}
\loigiai{
$x = 15 - 6 = 9$.\\
\textbf{Đáp án của câu hỏi là B}
}
\end{ex}

\begin{ex} %[2D2V11]
Có 16 quả táo, ăn hết 9 quả táo. Còn lại số quả táo là
\choice
{$6$ quả}
{$8$ quả}
{\True $7$ quả}
{$9$ quả}
\loigiai{
$16 - 9 = 7$ quả.\\
\textbf{Đáp án của câu hỏi là C}
}
\end{ex}

\begin{ex} %[2D2V11]
Tính giá trị biểu thức: $17 - 8 + 3$.
\choice
{\True $12$}
{$11$}
{$13$}
{$10$}
\loigiai{
$9 + 3 = 12$.\\
\textbf{Đáp án của câu hỏi là A}
}
\end{ex}

\begin{ex} %[2D2V11]
Một chuồng có 14 con thỏ, người ta bán đi 5 con thỏ. Số thỏ còn lại trong chuồng là
\choice
{$8$ con}
{\True $9$ con}
{$10$ con}
{$7$ con}
\loigiai{
$14 - 5 = 9$ con.\\
\textbf{Đáp án của câu hỏi là B}
}
\end{ex}

\begin{ex} %[2D2C11]
Tìm số thích hợp điền vào ô trống: $18 - \dots = 9$.
\choice
{$8$}
{$7$}
{\True $9$}
{$10$}
\loigiai{
$18 - 9 = 9$.\\
\textbf{Đáp án của câu hỏi là C}
}
\end{ex}

\textbf{Phần 2. Trắc nghiệm đúng sai}
\setcounterpageref{ex}{0}

\begin{ex} %[2D2H11]
Xét tính đúng sai của các phép trừ:
\choiceTF
{\True $12 - 7 = 5$.}
{\True $15 - 8 = 7$.}
{$16 - 9 = 8$.}
{\True $16 - 9 = 7$.}
\loigiai{
\begin{itemize}
    \item $12 - 7 = 5$: Đúng.
    \item $15 - 8 = 7$: Đúng.
    \item $16 - 9 = 7 \neq 8$. Ý 3 Sai, Ý 4 Đúng.
\end{itemize}
\begin{itemize}
    \item a) \textbf{Đúng}.
    \item b) \textbf{Đúng}.
    \item c) \textbf{Sai}.
    \item d) \textbf{Đúng}.
\end{itemize}
}
\end{ex}

\begin{ex} %[2D2V11]
Cho biểu thức $M = 13 - 5$ và $N = 16 - 8$. Xét tính đúng sai:
\choiceTF
{\True $M = 8$.}
{\True $N = 8$.}
{$M > N$.}
{\True $M = N = 8$.}
\loigiai{
\begin{itemize}
    \item $M = 13 - 5 = 8$: Đúng.
    \item $N = 16 - 8 = 8$: Đúng.
    \item $M = N$ (không lớn hơn). Ý 3 Sai, Ý 4 Đúng.
\end{itemize}
\begin{itemize}
    \item a) \textbf{Đúng}.
    \item b) \textbf{Đúng}.
    \item c) \textbf{Sai}.
    \item d) \textbf{Đúng}.
\end{itemize}
}
\end{ex}

\textbf{Phần 3. Trả lời ngắn}
\setcounterpageref{ex}{0}

\begin{ex} %[2D2H11]
Tính: $11 - 8$.
\shortans{3}
\loigiai{
\textbf{Đáp số:} $3$\\
$11 - 8 = 3$.
}
\end{ex}

\begin{ex} %[2D2H11]
Tính: $14 - 6$.
\shortans{8}
\loigiai{
\textbf{Đáp số:} $8$\\
$14 - 6 = 8$.
}
\end{ex}

\begin{ex} %[2D2V11]
Tính: $17 - 9$.
\shortans{8}
\loigiai{
\textbf{Đáp số:} $8$\\
$17 - 9 = 8$.
}
\end{ex}

\begin{ex} %[2D2V11]
Tìm $x$ biết: $x + 8 = 14$.
\shortans{6}
\loigiai{
\textbf{Đáp số:} $6$\\
$x = 14 - 8 = 6$.
}
\end{ex}

\begin{ex} %[2D2V11]
Có 15 bông hoa, đem cắm vào bình 6 bông hoa. Hỏi còn lại bao nhiêu bông hoa chưa cắm?
\shortans{9}
\loigiai{
\textbf{Đáp số:} $9$\\
$15 - 6 = 9$ bông hoa.
}
\end{ex}

\begin{ex} %[2D2C11]
Tìm số thích hợp điền vào dấu ?: $16 - ? = 8$.
\shortans{8}
\loigiai{
\textbf{Đáp số:} $8$\\
$16 - 8 = 8$.
}
\end{ex}

\newpage

\TaoPhieuKiemTra{Chủ đề 2 -- Phép cộng, phép trừ trong phạm vi 20}{12}{Bài toán nhiều hơn, ít hơn một số đơn vị}{https://loptoanthaylinh.vn/thi-online/toan-2/bai-12}
\textbf{Phần 1. Trắc nghiệm lựa chọn}
\setcounterpageref{ex}{0}

\begin{ex} %[2D2N12]
Trong bảng trừ (qua 10), kết quả của phép tính $11 - 9$ là
\choice
{\True $2$}
{$1$}
{$3$}
{$4$}
\loigiai{
$11 - 9 = 2$.\\
\textbf{Đáp án của câu hỏi là A}
}
\end{ex}

\begin{ex} %[2D2N12]
Phép tính nào sau đây có kết quả bằng 8?
\choice
{$14 - 7$}
{$13 - 6$}
{$15 - 8$}
{\True $14 - 6$}
\loigiai{
$14 - 6 = 8$.\\
\textbf{Đáp án của câu hỏi là D}
}
\end{ex}

\begin{ex} %[2D2H12]
Phép tính nào sau đây có kết quả bằng 7?
\choice
{\True $15 - 8$}
{$15 - 7$}
{$14 - 6$}
{$13 - 5$}
\loigiai{
$15 - 8 = 7$.\\
\textbf{Đáp án của câu hỏi là A}
}
\end{ex}

\begin{ex} %[2D2H12]
Phép tính nào sau đây có kết quả lớn nhất?
\choice
{$12 - 5$}
{$13 - 6$}
{\True $17 - 8$}
{$14 - 7$}
\loigiai{
$17 - 8 = 9$ (lớn nhất).\\
\textbf{Đáp án của câu hỏi là C}
}
\end{ex}

\begin{ex} %[2D2H12]
Phép tính nào sau đây có kết quả bé nhất?
\choice
{$13 - 8$}
{$12 - 7$}
{$14 - 8$}
{\True $11 - 7$}
\loigiai{
$11 - 7 = 4$ (bé nhất).\\
\textbf{Đáp án của câu hỏi là D}
}
\end{ex}

\begin{ex} %[2D2H12]
Điền số thích hợp: $13 - \dots = 9$.
\choice
{$3$}
{\True $4$}
{$5$}
{$6$}
\loigiai{
$13 - 4 = 9$.\\
\textbf{Đáp án của câu hỏi là B}
}
\end{ex}

\begin{ex} %[2D2V12]
Có bao nhiêu phép trừ dạng $11 - a$ trong bảng trừ? ($a$ từ 2 đến 9)
\choice
{$7$ phép tính}
{$9$ phép tính}
{\True $8$ phép tính}
{$10$ phép tính}
\loigiai{
Từ $11 - 2$ đến $11 - 9$ có $9 - 2 + 1 = 8$ phép tính.\\
\textbf{Đáp án của câu hỏi là C}
}
\end{ex}

\begin{ex} %[2D2V12]
Tính: $16 - 7 - 3 = \dots$
\choice
{\True $6$}
{$5$}
{$7$}
{$8$}
\loigiai{
$9 - 3 = 6$.\\
\textbf{Đáp án của câu hỏi là A}
}
\end{ex}

\begin{ex} %[2D2V12]
Điền dấu thích hợp: $15 - 7 \dots 16 - 8$.
\choice
{$>$}
{$<$}
{Không so sánh được}
{\True $=$}
\loigiai{
$15 - 7 = 8$ và $16 - 8 = 8$. Vậy $15 - 7 = 16 - 8$.\\
\textbf{Đáp án của câu hỏi là D}
}
\end{ex}

\begin{ex} %[2D2V12]
Một hộp có 18 viên bi, lấy ra 9 viên bi. Trong hộp còn lại số viên bi là
\choice
{$8$ viên}
{\True $9$ viên}
{$10$ viên}
{$11$ viên}
\loigiai{
$18 - 9 = 9$ viên.\\
\textbf{Đáp án của câu hỏi là B}
}
\end{ex}

\begin{ex} %[2D2V12]
Tìm $y$ biết: $14 - y = 8$.
\choice
{$5$}
{$7$}
{\True $6$}
{$8$}
\loigiai{
$y = 14 - 8 = 6$.\\
\textbf{Đáp án của câu hỏi là C}
}
\end{ex}

\begin{ex} %[2D2C12]
Hiệu của số lớn nhất có hai chữ số và số bé nhất có hai chữ số là bao nhiêu?
\choice
{$88$}
{\True $89$}
{$90$}
{$80$}
\loigiai{
$99 - 10 = 89$.\\
\textbf{Đáp án của câu hỏi là B}
}
\end{ex}

\textbf{Phần 2. Trắc nghiệm đúng sai}
\setcounterpageref{ex}{0}

\begin{ex} %[2D2H12]
Xét tính đúng sai của các phép trừ trong bảng trừ:
\choiceTF
{\True $13 - 6 = 7$.}
{\True $15 - 9 = 6$.}
{$17 - 9 = 9$.}
{\True $17 - 9 = 8$.}
\loigiai{
\begin{itemize}
    \item $13 - 6 = 7$: Đúng.
    \item $15 - 9 = 6$: Đúng.
    \item $17 - 9 = 8 \neq 9$. Ý 3 Sai, Ý 4 Đúng.
\end{itemize}
\begin{itemize}
    \item a) \textbf{Đúng}.
    \item b) \textbf{Đúng}.
    \item c) \textbf{Sai}.
    \item d) \textbf{Đúng}.
\end{itemize}
}
\end{ex}

\begin{ex} %[2D2V12]
Cho 3 phép tính: $11 - 6, 12 - 7, 13 - 8$. Xét tính đúng sai:
\choiceTF
{\True Cả 3 phép tính đều có kết quả bằng 5.}
{\True $11 - 6 = 5$.}
{Phép tính $13 - 8$ có kết quả lớn nhất.}
{\True Ba phép tính này có kết quả bằng nhau.}
\loigiai{
\begin{itemize}
    \item Cả 3 đều bằng 5: Đúng.
    \item $11 - 6 = 5$: Đúng.
    \item Chúng bằng nhau (không có lớn nhất). Ý 3 Sai, Ý 4 Đúng.
\end{itemize}
\begin{itemize}
    \item a) \textbf{Đúng}.
    \item b) \textbf{Đúng}.
    \item c) \textbf{Sai}.
    \item d) \textbf{Đúng}.
\end{itemize}
}
\end{ex}

\textbf{Phần 3. Trả lời ngắn}
\setcounterpageref{ex}{0}

\begin{ex} %[2D2H12]
Tính: $12 - 9$.
\shortans{3}
\loigiai{
\textbf{Đáp số:} $3$\\
$12 - 9 = 3$.
}
\end{ex}

\begin{ex} %[2D2H12]
Tính: $14 - 7$.
\shortans{7}
\loigiai{
\textbf{Đáp số:} $7$\\
$14 - 7 = 7$.
}
\end{ex}

\begin{ex} %[2D2V12]
Tính giá trị biểu thức: $15 - 8 + 4$.
\shortans{11}
\loigiai{
\textbf{Đáp số:} $11$\\
$7 + 4 = 11$.
}
\end{ex}

\begin{ex} %[2D2V12]
Điền số thích hợp: $18 - ? = 9$.
\shortans{9}
\loigiai{
\textbf{Đáp số:} $9$\\
$18 - 9 = 9$.
}
\end{ex}

\begin{ex} %[2D2V12]
Một cây có 16 quả bưởi, đã hái đi 8 quả. Hỏi trên cây còn lại bao nhiêu quả bưởi?
\shortans{8}
\loigiai{
\textbf{Đáp số:} $8$\\
$16 - 8 = 8$ quả.
}
\end{ex}

\begin{ex} %[2D2C12]
Tìm số thích hợp điền vào dấu ?: $13 - ? = 5$.
\shortans{8}
\loigiai{
\textbf{Đáp số:} $8$\\
$13 - 5 = 8$.
}
\end{ex}

\newpage

\TaoPhieuKiemTra{Chủ đề 2 -- Phép cộng, phép trừ trong phạm vi 20}{13}{Luyện tập chung}{https://loptoanthaylinh.vn/thi-online/toan-2/bai-13}
\textbf{Phần 1. Trắc nghiệm lựa chọn}
\setcounterpageref{ex}{0}

\begin{ex} %[2D2N13]
Lan có 8 cái kẹo, Mai có nhiều hơn Lan 5 cái kẹo. Hỏi Mai có bao nhiêu cái kẹo?
\choice
{\True $13$ cái}
{$12$ cái}
{$14$ cái}
{$15$ cái}
\loigiai{
$8 + 5 = 13$ cái kẹo.\\
\textbf{Đáp án của câu hỏi là A}
}
\end{ex}

\begin{ex} %[2D2N13]
An có 15 viên bi, Bình có ít hơn An 6 viên bi. Hỏi Bình có bao nhiêu viên bi?
\choice
{$8$ viên}
{$10$ viên}
{$11$ viên}
{\True $9$ viên}
\loigiai{
$15 - 6 = 9$ viên bi.\\
\textbf{Đáp án của câu hỏi là D}
}
\end{ex}

\begin{ex} %[2D2H13]
Bể thứ nhất có 9 con cá, bể thứ hai có nhiều hơn bể thứ nhất 7 con cá. Bể thứ hai có số con cá là
\choice
{\True $16$ con}
{$15$ con}
{$17$ con}
{$18$ con}
\loigiai{
$9 + 7 = 16$ con cá.\\
\textbf{Đáp án của câu hỏi là A}
}
\end{ex}

\begin{ex} %[2D2H13]
Tổ 1 có 14 bạn, tổ 2 có ít hơn tổ 1 là 5 bạn. Tổ 2 có số bạn là
\choice
{$8$ bạn}
{$10$ bạn}
{$11$ bạn}
{\True $9$ bạn}
\loigiai{
$14 - 5 = 9$ bạn.\\
\textbf{Đáp án của câu hỏi là D}
}
\end{ex}

\begin{ex} %[2D2H13]
Bác Ba nuôi 8 con bò, số con lợn nhiều hơn số con bò là 9 con. Bác Ba nuôi số con lợn là
\choice
{$16$ con}
{\True $17$ con}
{$18$ con}
{$15$ con}
\loigiai{
$8 + 9 = 17$ con lợn.\\
\textbf{Đáp án của câu hỏi là B}
}
\end{ex}

\begin{ex} %[2D2H13]
Cây ổi có 16 quả, cây xoài có ít hơn cây ổi 7 quả. Cây xoài có số quả là
\choice
{$8$ quả}
{$10$ quả}
{$11$ quả}
{\True $9$ quả}
\loigiai{
$16 - 7 = 9$ quả.\\
\textbf{Đáp án của câu hỏi là D}
}
\end{ex}

\begin{ex} %[2D2V13]
Đoạn thẳng $AB$ dài $9\text{ cm}$, đoạn thẳng $CD$ dài hơn đoạn thẳng $AB$ là $6\text{ cm}$. Độ dài đoạn thẳng $CD$ là
\choice
{$14\text{ cm}$}
{$16\text{ cm}$}
{\True $15\text{ cm}$}
{$17\text{ cm}$}
\loigiai{
$9 + 6 = 15\text{ cm}$.\\
\textbf{Đáp án của câu hỏi là C}
}
\end{ex}

\begin{ex} %[2D2V13]
Sợi dây màu xanh dài $18\text{ cm}$, sợi dây màu đỏ ngắn hơn sợi dây màu xanh $9\text{ cm}$. Độ dài sợi dây màu đỏ là
\choice
{$8\text{ cm}$}
{\True $9\text{ cm}$}
{$10\text{ cm}$}
{$11\text{ cm}$}
\loigiai{
$18 - 9 = 9\text{ cm}$.\\
\textbf{Đáp án của câu hỏi là B}
}
\end{ex}

\begin{ex} %[2D2V13]
Hùng có 9 viên bi, Dũng có nhiều hơn Hùng 8 viên bi. Cả hai bạn có tất cả bao nhiêu viên bi?
\choice
{\True $26$ viên}
{$25$ viên}
{$27$ viên}
{$17$ viên}
\loigiai{
Dũng có: $9 + 8 = 17$ viên.\\
Cả hai bạn có: $9 + 17 = 26$ viên bi.\\
\textbf{Đáp án của câu hỏi là A}
}
\end{ex}

\begin{ex} %[2D2V13]
Hoa có 14 bông hoa, Cúc có ít hơn Hoa 6 bông hoa. Cả hai bạn có tất cả bao nhiêu bông hoa?
\choice
{$21$ bông}
{$23$ bông}
{\True $22$ bông}
{$24$ bông}
\loigiai{
Cúc có: $14 - 6 = 8$ bông.\\
Cả hai bạn có: $14 + 8 = 22$ bông hoa.\\
\textbf{Đáp án của câu hỏi là C}
}
\end{ex}

\begin{ex} %[2D2V13]
Nam có 8 cái bút chì, số bút chì của Minh nhiều hơn Nam là 4 cái. Minh có số bút chì là
\choice
{$11$ cái}
{\True $12$ cái}
{$13$ cái}
{$14$ cái}
\loigiai{
$8 + 4 = 12$ cái.\\
\textbf{Đáp án của câu hỏi là B}
}
\end{ex}

\begin{ex} %[2D2C13]
Có hai ngăn sách. Ngăn thứ nhất có 9 quyển sách. Nếu chuyển 3 quyển từ ngăn thứ hai sang ngăn thứ nhất thì hai ngăn bằng nhau. Hỏi ban đầu ngăn thứ hai có bao nhiêu quyển sách?
\choice
{$14$ quyển}
{$16$ quyển}
{\True $15$ quyển}
{$12$ quyển}
\loigiai{
Sau khi chuyển, ngăn 1 có: $9 + 3 = 12$ quyển.\\
Vậy lúc sau ngăn 2 cũng có 12 quyển $\Rightarrow$ Ban đầu ngăn 2 có: $12 + 3 = 15$ quyển.\\
\textbf{Đáp án của câu hỏi là C}
}
\end{ex}

\textbf{Phần 2. Trắc nghiệm đúng sai}
\setcounterpageref{ex}{0}

\begin{ex} %[2D2H13]
Tổ 1 có 8 bạn, tổ 2 có nhiều hơn tổ 1 là 5 bạn. Xét tính đúng sai:
\choiceTF
{\True Tìm số bạn tổ 2 ta làm phép tính cộng ($8 + 5$).}
{\True Tổ 2 có 13 bạn.}
{Cả hai tổ có tất cả 20 bạn.}
{\True Cả hai tổ có tất cả 21 bạn ($8 + 13 = 21$).}
\loigiai{
\begin{itemize}
    \item Phép tính cộng $8 + 5$: Đúng.
    \item Tổ 2 có 13 bạn: Đúng.
    \item Tổng: $8 + 13 = 21 \neq 20$. Ý 3 Sai, Ý 4 Đúng.
\end{itemize}
\begin{itemize}
    \item a) \textbf{Đúng}.
    \item b) \textbf{Đúng}.
    \item c) \textbf{Sai}.
    \item d) \textbf{Đúng}.
\end{itemize}
}
\end{ex}

\begin{ex} %[2D2V13]
Thùng to có 15 lít dầu, thùng nhỏ có ít hơn thùng to 7 lít dầu. Xét tính đúng sai:
\choiceTF
{\True Thùng nhỏ có 8 lít dầu ($15 - 7 = 8$).}
{\True Cả hai thùng có $23\text{ lít}$ dầu ($15 + 8 = 23$).}
{Thùng nhỏ có nhiều hơn 10 lít dầu.}
{\True Thùng nhỏ có ít hơn 10 lít dầu (chỉ có $8\text{ lít}$).}
\loigiai{
\begin{itemize}
    \item Thùng nhỏ: $15 - 7 = 8\text{ l}$: Đúng.
    \item Tổng: $15 + 8 = 23\text{ l}$: Đúng.
    \item $8 < 10$ (không nhiều hơn). Ý 3 Sai, Ý 4 Đúng.
\end{itemize}
\begin{itemize}
    \item a) \textbf{Đúng}.
    \item b) \textbf{Đúng}.
    \item c) \textbf{Sai}.
    \item d) \textbf{Đúng}.
\end{itemize}
}
\end{ex}

\textbf{Phần 3. Trả lời ngắn}
\setcounterpageref{ex}{0}

\begin{ex} %[2D2H13]
Lan có 9 cái kẹo, Huệ có nhiều hơn Lan 6 cái kẹo. Hỏi Huệ có bao nhiêu cái kẹo?
\shortans{15}
\loigiai{
\textbf{Đáp số:} $15$\\
$9 + 6 = 15$ cái.
}
\end{ex}

\begin{ex} %[2D2H13]
Đàn vịt có 17 con, đàn gà có ít hơn đàn vịt 9 con. Hỏi đàn gà có bao nhiêu con?
\shortans{8}
\loigiai{
\textbf{Đáp số:} $8$\\
$17 - 9 = 8$ con.
}
\end{ex}

\begin{ex} %[2D2V13]
Đoạn dây thứ nhất dài $8\text{ dm}$, đoạn dây thứ hai dài hơn đoạn thứ nhất $7\text{ dm}$. Tính độ dài đoạn dây thứ hai theo đề-xi-mét.
\shortans{15}
\loigiai{
\textbf{Đáp số:} $15$\\
$8 + 7 = 15\text{ dm}$.
}
\end{ex}

\begin{ex} %[2D2V13]
Hộp màu xanh có 14 cái bút, hộp màu đỏ có ít hơn hộp màu xanh 8 cái bút. Hỏi hộp màu đỏ có bao nhiêu cái bút?
\shortans{6}
\loigiai{
\textbf{Đáp số:} $6$\\
$14 - 8 = 6$ cái.
}
\end{ex}

\begin{ex} %[2D2V13]
Bác An nuôi 9 con thỏ trắng, số thỏ xám nhiều hơn thỏ trắng 8 con. Hỏi bác An nuôi bao nhiêu con thỏ xám?
\shortans{17}
\loigiai{
\textbf{Đáp số:} $17$\\
$9 + 8 = 17$ con.
}
\end{ex}

\begin{ex} %[2D2C13]
Mai có 8 bông hoa. Nếu Mai hái thêm 4 bông hoa nữa thì Mai vẫn ít hơn Lan 3 bông hoa. Hỏi Lan có bao nhiêu bông hoa?
\shortans{15}
\loigiai{
\textbf{Đáp số:} $15$\\
$8 + 4 + 3 = 15$ bông hoa.
}
\end{ex}

\newpage

\TaoPhieuKiemTra{Chủ đề 2 -- Phép cộng, phép trừ trong phạm vi 20}{14}{Đề kiểm tra chủ đề}{https://loptoanthaylinh.vn/thi-online/toan-2/bai-14}
\textbf{Phần 1. Trắc nghiệm lựa chọn}
\setcounterpageref{ex}{0}

\begin{ex} %[2D2N14]
Kết quả của phép tính $9 + 6$ là
\choice
{\True $15$}
{$14$}
{$16$}
{$17$}
\loigiai{
$9 + 6 = 15$.\\
\textbf{Đáp án của câu hỏi là A}
}
\end{ex}

\begin{ex} %[2D2N14]
Kết quả của phép tính $15 - 7$ là
\choice
{$7$}
{$9$}
{$10$}
{\True $8$}
\loigiai{
$15 - 7 = 8$.\\
\textbf{Đáp án của câu hỏi là D}
}
\end{ex}

\begin{ex} %[2D2H14]
Kết quả của phép tính $8 + 7$ là
\choice
{\True $15$}
{$14$}
{$16$}
{$17$}
\loigiai{
$8 + 7 = 15$.\\
\textbf{Đáp án của câu hỏi là A}
}
\end{ex}

\begin{ex} %[2D2H14]
Kết quả của phép tính $17 - 9$ là
\choice
{$7$}
{$9$}
{\True $8$}
{$10$}
\loigiai{
$17 - 9 = 8$.\\
\textbf{Đáp án của câu hỏi là C}
}
\end{ex}

\begin{ex} %[2D2H14]
Điền dấu thích hợp: $13 - 5 \dots 14 - 6$.
\choice
{$>$}
{\True $=$}
{$<$}
{Không so sánh được}
\loigiai{
$13 - 5 = 8$ và $14 - 6 = 8$. Vậy $13 - 5 = 14 - 6$.\\
\textbf{Đáp án của câu hỏi là B}
}
\end{ex}

\begin{ex} %[2D2H14]
Tính giá trị biểu thức: $16 - 8 + 5$.
\choice
{\True $13$}
{$12$}
{$14$}
{$15$}
\loigiai{
$8 + 5 = 13$.\\
\textbf{Đáp án của câu hỏi là A}
}
\end{ex}

\begin{ex} %[2D2V14]
Một cửa hàng buổi sáng bán được 8 cái xe đạp, buổi chiều bán nhiều hơn buổi sáng 5 cái. Buổi chiều bán được số xe đạp là
\choice
{$12$ cái}
{$14$ cái}
{$15$ cái}
{\True $13$ cái}
\loigiai{
$8 + 5 = 13$ cái.\\
\textbf{Đáp án của câu hỏi là D}
}
\end{ex}

\begin{ex} %[2D2V14]
Đàn vịt có 16 con, đàn gà có ít hơn đàn vịt 7 con. Đàn gà có số con là
\choice
{$8$ con}
{$10$ con}
{\True $9$ con}
{$11$ con}
\loigiai{
$16 - 7 = 9$ con.\\
\textbf{Đáp án của câu hỏi là C}
}
\end{ex}

\begin{ex} %[2D2V14]
Tìm $x$ biết: $x + 7 = 16$.
\choice
{$8$}
{$10$}
{$11$}
{\True $9$}
\loigiai{
$x = 16 - 7 = 9$.\\
\textbf{Đáp án của câu hỏi là D}
}
\end{ex}

\begin{ex} %[2D2V14]
Tìm $y$ biết: $15 - y = 6$.
\choice
{$8$}
{\True $9$}
{$10$}
{$7$}
\loigiai{
$y = 15 - 6 = 9$.\\
\textbf{Đáp án của câu hỏi là B}
}
\end{ex}

\begin{ex} %[2D2V14]
Tính nhanh: $9 + 6 + 1$.
\choice
{$15$}
{$17$}
{\True $16$}
{$18$}
\loigiai{
$(9 + 1) + 6 = 16$.\\
\textbf{Đáp án của câu hỏi là C}
}
\end{ex}

\begin{ex} %[2D2C14]
Tìm một số biết rằng lấy số đó cộng với 8 rồi trừ đi 5 thì được kết quả là 12.
\choice
{$8$}
{\True $9$}
{$10$}
{$7$}
\loigiai{
$12 + 5 - 8 = 9$.\\
\textbf{Đáp án của câu hỏi là B}
}
\end{ex}

\textbf{Phần 2. Trắc nghiệm đúng sai}
\setcounterpageref{ex}{0}

\begin{ex} %[2D2H14]
Xét tính đúng sai của các phép toán trong phạm vi 20:
\choiceTF
{\True $9 + 9 = 18$.}
{\True $18 - 9 = 9$.}
{$14 - 7 = 8$.}
{\True $14 - 7 = 7$.}
\loigiai{
\begin{itemize}
    \item $9 + 9 = 18$: Đúng.
    \item $18 - 9 = 9$: Đúng.
    \item $14 - 7 = 7 \neq 8$. Ý 3 Sai, Ý 4 Đúng.
\end{itemize}
\begin{itemize}
    \item a) \textbf{Đúng}.
    \item b) \textbf{Đúng}.
    \item c) \textbf{Sai}.
    \item d) \textbf{Đúng}.
\end{itemize}
}
\end{ex}

\begin{ex} %[2D2V14]
Hà có 14 quyển truyện, Nam có ít hơn Hà 5 quyển truyện. Xét tính đúng sai:
\choiceTF
{\True Nam có 9 quyển truyện ($14 - 5 = 9$).}
{\True Cả hai bạn có tất cả 23 quyển truyện ($14 + 9 = 23$).}
{Hà có ít truyện hơn Nam.}
{\True Hà có nhiều truyện hơn Nam (14 quyển $> 9$ quyển).}
\loigiai{
\begin{itemize}
    \item Nam có 9 quyển: Đúng.
    \item Cả hai có 23 quyển: Đúng.
    \item Hà nhiều hơn (không phải ít hơn). Ý 3 Sai, Ý 4 Đúng.
\end{itemize}
\begin{itemize}
    \item a) \textbf{Đúng}.
    \item b) \textbf{Đúng}.
    \item c) \textbf{Sai}.
    \item d) \textbf{Đúng}.
\end{itemize}
}
\end{ex}

\textbf{Phần 3. Trả lời ngắn}
\setcounterpageref{ex}{0}

\begin{ex} %[2D2H14]
Tính: $8 + 9$.
\shortans{17}
\loigiai{
\textbf{Đáp số:} $17$\\
$8 + 9 = 17$.
}
\end{ex}

\begin{ex} %[2D2H14]
Tính: $16 - 8$.
\shortans{8}
\loigiai{
\textbf{Đáp số:} $8$\\
$16 - 8 = 8$.
}
\end{ex}

\begin{ex} %[2D2V14]
Tính giá trị biểu thức: $17 - 9 + 4$.
\shortans{12}
\loigiai{
\textbf{Đáp số:} $12$\\
$8 + 4 = 12$.
}
\end{ex}

\begin{ex} %[2D2V14]
Tìm $x$ biết: $x - 7 = 8$.
\shortans{15}
\loigiai{
\textbf{Đáp số:} $15$\\
$x = 8 + 7 = 15$.
}
\end{ex}

\begin{ex} %[2D2V14]
Một vườn hoa có 9 bông hoa cúc, số hoa hồng nhiều hơn hoa cúc là 7 bông. Hỏi có bao nhiêu bông hoa hồng?
\shortans{16}
\loigiai{
\textbf{Đáp số:} $16$\\
$9 + 7 = 16$ bông hoa.
}
\end{ex}

\begin{ex} %[2D2C14]
Có hai rổ táo, rổ thứ nhất có 8 quả táo. Nếu thêm 6 quả vào rổ thứ nhất thì số táo ở rổ thứ nhất nhiều hơn rổ thứ hai là 2 quả. Hỏi rổ thứ hai có bao nhiêu quả táo?
\shortans{12}
\loigiai{
\textbf{Đáp số:} $12$\\
Rổ 1 lúc sau: $8 + 6 = 14$ quả.\\
Rổ 2 có: $14 - 2 = 12$ quả.
}
\end{ex}

\newpage

% =========================================================================
% CHỦ ĐỀ 3 - LÀM QUEN VỚI KHỐI LƯỢNG, DUNG TÍCH
% =========================================================================
\section{Làm quen với khối lượng, dung tích}

\TaoPhieuKiemTra{Chủ đề 3 -- Làm quen với khối lượng, dung tích}{15}{Ki-lô-gam}{https://loptoanthaylinh.vn/thi-online/toan-2/bai-15}
\textbf{Phần 1. Trắc nghiệm lựa chọn}
\setcounterpageref{ex}{0}

\begin{ex} %[2H3N15]
Ki-lô-gam được viết tắt là
\choice
{$\text{km}$}
{$\text{g}$}
{$\text{l}$}
{\True $\text{kg}$}
\loigiai{
Ki-lô-gam viết tắt là $\text{kg}$.\\
\textbf{Đáp án của câu hỏi là D}
}
\end{ex}

\begin{ex} %[2H3N15]
Để đo khối lượng (cân nặng) của một vật, người ta dùng đơn vị
\choice
{Xăng-ti-mét}
{Lít}
{\True Ki-lô-gam}
{Giờ}
\loigiai{
Đo khối lượng dùng ki-lô-gam.\\
\textbf{Đáp án của câu hỏi là C}
}
\end{ex}

\begin{ex} %[2H3H15]
Tính: $15\text{ kg} + 8\text{ kg} = \dots$
\choice
{$22\text{ kg}$}
{\True $23\text{ kg}$}
{$24\text{ kg}$}
{$25\text{ kg}$}
\loigiai{
$15 + 8 = 23\text{ kg}$.\\
\textbf{Đáp án của câu hỏi là B}
}
\end{ex}

\begin{ex} %[2H3H15]
Tính: $35\text{ kg} - 12\text{ kg} = \dots$
\choice
{\True $23\text{ kg}$}
{$22\text{ kg}$}
{$24\text{ kg}$}
{$20\text{ kg}$}
\loigiai{
$35 - 12 = 23\text{ kg}$.\\
\textbf{Đáp án của câu hỏi là A}
}
\end{ex}

\begin{ex} %[2H3H15]
Bao gạo to nặng $30\text{ kg}$, bao gạo nhỏ nặng $15\text{ kg}$. Cả hai bao gạo nặng số ki-lô-gam là
\choice
{$40\text{ kg}$}
{\True $45\text{ kg}$}
{$50\text{ kg}$}
{$55\text{ kg}$}
\loigiai{
$30 + 15 = 45\text{ kg}$.\\
\textbf{Đáp án của câu hỏi là B}
}
\end{ex}

\begin{ex} %[2H3H15]
Con lợn cân nặng $65\text{ kg}$, con chó cân nặng $15\text{ kg}$. Con lợn nặng hơn con chó số ki-lô-gam là
\choice
{$40\text{ kg}$}
{$60\text{ kg}$}
{\True $50\text{ kg}$}
{$80\text{ kg}$}
\loigiai{
$65 - 15 = 50\text{ kg}$.\\
\textbf{Đáp án của câu hỏi là C}
}
\end{ex}

\begin{ex} %[2H3V15]
Một quả dưa hấu cân nặng $4\text{ kg}$, quả bí đỏ cân nặng $5\text{ kg}$. Cả hai quả cân nặng là
\choice
{$8\text{ kg}$}
{$10\text{ kg}$}
{$11\text{ kg}$}
{\True $9\text{ kg}$}
\loigiai{
$4 + 5 = 9\text{ kg}$.\\
\textbf{Đáp án của câu hỏi là D}
}
\end{ex}

\begin{ex} %[2H3V15]
Có một bao đường nặng $20\text{ kg}$, người ta đã lấy ra $8\text{ kg}$. Số đường còn lại trong bao là
\choice
{$11\text{ kg}$}
{\True $12\text{ kg}$}
{$13\text{ kg}$}
{$14\text{ kg}$}
\loigiai{
$20 - 8 = 12\text{ kg}$.\\
\textbf{Đáp án của câu hỏi là B}
}
\end{ex}

\begin{ex} %[2H3V15]
Điền dấu thích hợp: $18\text{ kg} + 5\text{ kg} \dots 25\text{ kg}$.
\choice
{$>$}
{$=$}
{\True $<$}
{Không so sánh được}
\loigiai{
$18\text{ kg} + 5\text{ kg} = 23\text{ kg} < 25\text{ kg}$.\\
\textbf{Đáp án của câu hỏi là C}
}
\end{ex}

\begin{ex} %[2H3V15]
Tính giá trị biểu thức: $40\text{ kg} + 25\text{ kg} - 15\text{ kg}$.
\choice
{\True $50\text{ kg}$}
{$45\text{ kg}$}
{$55\text{ kg}$}
{$60\text{ kg}$}
\loigiai{
$65 - 15 = 50\text{ kg}$.\\
\textbf{Đáp án của câu hỏi là A}
}
\end{ex}

\begin{ex} %[2H3V15]
Mẹ mua về $15\text{ kg}$ gạo tẻ và $8\text{ kg}$ gạo nếp. Hỏi mẹ mua tất cả bao nhiêu ki-lô-gam gạo?
\choice
{$22\text{ kg}$}
{$24\text{ kg}$}
{$25\text{ kg}$}
{\True $23\text{ kg}$}
\loigiai{
$15 + 8 = 23\text{ kg}$.\\
\textbf{Đáp án của câu hỏi là D}
}
\end{ex}

\begin{ex} %[2H3C15]
Có một cân đĩa và các quả cân loại $1\text{ kg}, 2\text{ kg}, 5\text{ kg}$. Để cân một túi hoa quả nặng $8\text{ kg}$, ta có thể đặt ở đĩa bên kia các quả cân nào?
\choice
{\True $5\text{ kg}, 2\text{ kg}$ và $1\text{ kg}$}
{$5\text{ kg}$ và $2\text{ kg}$}
{2 quả $5\text{ kg}$}
{3 quả $2\text{ kg}$}
\loigiai{
$5 + 2 + 1 = 8\text{ kg}$.\\
\textbf{Đáp án của câu hỏi là A}
}
\end{ex}

\textbf{Phần 2. Trắc nghiệm đúng sai}
\setcounterpageref{ex}{0}

\begin{ex} %[2H3H15]
Xét tính đúng sai của các phép tính với đơn vị đo:
\choiceTF
{\True $12\text{ kg} + 8\text{ kg} = 20\text{ kg}$.}
{\True $30\text{ kg} - 10\text{ kg} = 20\text{ kg}$.}
{$15\text{ kg} + 7 = 22\text{ kg}$.}
{\True Khi tính toán với đại lượng, các số hạng phải cùng kèm đơn vị đo.}
\loigiai{
\begin{itemize}
    \item $12\text{ kg} + 8\text{ kg} = 20\text{ kg}$: Đúng.
    \item $30\text{ kg} - 10\text{ kg} = 20\text{ kg}$: Đúng.
    \item Ý 3 thiếu đơn vị đo ở số thứ hai. Sai.
    \item Cần cùng kèm đơn vị đo: Đúng.
\end{itemize}
\begin{itemize}
    \item a) \textbf{Đúng}.
    \item b) \textbf{Đúng}.
    \item c) \textbf{Sai}.
    \item d) \textbf{Đúng}.
\end{itemize}
}
\end{ex}

\begin{ex} %[2H3V15]
Bao ngô nặng $45\text{ kg}$, bao thóc nặng hơn bao ngô $10\text{ kg}$. Xét tính đúng sai:
\choiceTF
{\True Bao thóc cân nặng $55\text{ kg}$ ($45 + 10 = 55$).}
{\True Cả hai bao cân nặng $100\text{ kg}$ ($45 + 55 = 100$).}
{Bao thóc nhẹ hơn bao ngô.}
{\True Bao ngô nhẹ hơn bao thóc 10 kg.}
\loigiai{
\begin{itemize}
    \item Thóc: $45 + 10 = 55\text{ kg}$: Đúng.
    \item Tổng: $45 + 55 = 100\text{ kg}$: Đúng.
    \item Thóc nặng hơn (không phải nhẹ hơn). Ý 3 Sai, Ý 4 Đúng.
\end{itemize}
\begin{itemize}
    \item a) \textbf{Đúng}.
    \item b) \textbf{Đúng}.
    \item c) \textbf{Sai}.
    \item d) \textbf{Đúng}.
\end{itemize}
}
\end{ex}

\textbf{Phần 3. Trả lời ngắn}
\setcounterpageref{ex}{0}

\begin{ex} %[2H3H15]
Tính: $18\text{ kg} + 9\text{ kg}$ theo ki-lô-gam.
\shortans{27}
\loigiai{
\textbf{Đáp số:} $27$\\
$18 + 9 = 27\text{ kg}$.
}
\end{ex}

\begin{ex} %[2H3H15]
Tính: $45\text{ kg} - 20\text{ kg}$ theo ki-lô-gam.
\shortans{25}
\loigiai{
\textbf{Đáp số:} $25$\\
$45 - 20 = 25\text{ kg}$.
}
\end{ex}

\begin{ex} %[2H3V15]
Bao gạo nặng $35\text{ kg}$, mẹ mua thêm $15\text{ kg}$ nữa. Hỏi có tất cả bao nhiêu ki-lô-gam gạo?
\shortans{50}
\loigiai{
\textbf{Đáp số:} $50$\\
$35 + 15 = 50\text{ kg}$.
}
\end{ex}

\begin{ex} %[2H3V15]
Một con dê nặng $28\text{ kg}$, con cừu nặng $22\text{ kg}$. Hỏi cả hai con nặng bao nhiêu ki-lô-gam?
\shortans{50}
\loigiai{
\textbf{Đáp số:} $50$\\
$28 + 22 = 50\text{ kg}$.
}
\end{ex}

\begin{ex} %[2H3V15]
Một bao xi măng nặng $50\text{ kg}$, thợ xây đã dùng hết $35\text{ kg}$. Hỏi còn lại bao nhiêu ki-lô-gam xi măng?
\shortans{15}
\loigiai{
\textbf{Đáp số:} $15$\\
$50 - 35 = 15\text{ kg}$.
}
\end{ex}

\begin{ex} %[2H3C15]
Quả bưởi nặng $2\text{ kg}$, quả dưa hấu nặng hơn quả bưởi $3\text{ kg}$. Hỏi cả hai quả nặng bao nhiêu ki-lô-gam?
\shortans{7}
\loigiai{
\textbf{Đáp số:} $7$\\
Dưa hấu: $2 + 3 = 5\text{ kg}$. Cả hai quả: $2 + 5 = 7\text{ kg}$.
}
\end{ex}

\newpage

\TaoPhieuKiemTra{Chủ đề 3 -- Làm quen với khối lượng, dung tích}{16}{Lít}{https://loptoanthaylinh.vn/thi-online/toan-2/bai-16}
\textbf{Phần 1. Trắc nghiệm lựa chọn}
\setcounterpageref{ex}{0}

\begin{ex} %[2H3N16]
Lít là đơn vị dùng để đo
\choice
{Khối lượng của vật}
{Độ dài đoạn thẳng}
{\True Dung tích (sức chứa) của chất lỏng}
{Thời gian trong ngày}
\loigiai{
Lít dùng để đo dung tích chất lỏng.\\
\textbf{Đáp án của câu hỏi là C}
}
\end{ex}

\begin{ex} %[2H3N16]
Lít được viết tắt là
\choice
{$\text{kg}$}
{$\text{cm}$}
{$\text{m}$}
{\True $\text{l}$}
\loigiai{
Lít viết tắt là $\text{l}$.\\
\textbf{Đáp án của câu hỏi là D}
}
\end{ex}

\begin{ex} %[2H3H16]
Tính: $15\text{ l} + 8\text{ l} = \dots$
\choice
{$22\text{ l}$}
{$24\text{ l}$}
{\True $23\text{ l}$}
{$25\text{ l}$}
\loigiai{
$15 + 8 = 23\text{ l}$.\\
\textbf{Đáp án của câu hỏi là C}
}
\end{ex}

\begin{ex} %[2H3H16]
Tính: $38\text{ l} - 15\text{ l} = \dots$
\choice
{$22\text{ l}$}
{$24\text{ l}$}
{$25\text{ l}$}
{\True $23\text{ l}$}
\loigiai{
$38 - 15 = 23\text{ l}$.\\
\textbf{Đáp án của câu hỏi là D}
}
\end{ex}

\begin{ex} %[2H3H16]
Thùng thứ nhất chứa $20\text{ l}$ nước mắm, thùng thứ hai chứa $15\text{ l}$ nước mắm. Cả hai thùng chứa số lít nước mắm là
\choice
{\True $35\text{ l}$}
{$30\text{ l}$}
{$40\text{ l}$}
{$45\text{ l}$}
\loigiai{
$20 + 15 = 35\text{ l}$.\\
\textbf{Đáp án của câu hỏi là A}
}
\end{ex}

\begin{ex} %[2H3H16]
Can to đựng $18\text{ l}$ dầu, can nhỏ đựng $9\text{ l}$ dầu. Can to đựng nhiều hơn can nhỏ số lít dầu là
\choice
{$8\text{ l}$}
{\True $9\text{ l}$}
{$10\text{ l}$}
{$27\text{ l}$}
\loigiai{
$18 - 9 = 9\text{ l}$.\\
\textbf{Đáp án của câu hỏi là B}
}
\end{ex}

\begin{ex} %[2H3V16]
Một can chứa $10\text{ l}$ nước, người ta rót ra các chai loại $2\text{ l}$. Rót được bao nhiêu chai?
\choice
{\True $5$ chai}
{$4$ chai}
{$6$ chai}
{$7$ chai}
\loigiai{
$10 : 2 = 5$ chai.\\
\textbf{Đáp án của câu hỏi là A}
}
\end{ex}

\begin{ex} %[2H3V16]
Điền dấu thích hợp: $24\text{ l} + 6\text{ l} \dots 32\text{ l}$.
\choice
{$>$}
{\True $<$}
{$=$}
{Không so sánh được}
\loigiai{
$24 + 6 = 30\text{ l} < 32\text{ l}$.\\
\textbf{Đáp án của câu hỏi là B}
}
\end{ex}

\begin{ex} %[2H3V16]
Tính giá trị biểu thức: $35\text{ l} + 15\text{ l} - 20\text{ l}$.
\choice
{$25\text{ l}$}
{$35\text{ l}$}
{\True $30\text{ l}$}
{$40\text{ l}$}
\loigiai{
$50 - 20 = 30\text{ l}$.\\
\textbf{Đáp án của câu hỏi là C}
}
\end{ex}

\begin{ex} %[2H3V16]
Một bể cá chứa $45\text{ l}$ nước, sau khi múc bớt $15\text{ l}$ nước thì trong bể còn lại bao nhiêu lít nước?
\choice
{$25\text{ l}$}
{\True $30\text{ l}$}
{$35\text{ l}$}
{$40\text{ l}$}
\loigiai{
$45 - 15 = 30\text{ l}$.\\
\textbf{Đáp án của câu hỏi là B}
}
\end{ex}

\begin{ex} %[2H3V16]
Bình đựng $5\text{ l}$ mật ong, can đựng $15\text{ l}$ mật ong. Hỏi can đựng gấp mấy lần bình?
\choice
{$2$ lần}
{$4$ lần}
{$5$ lần}
{\True $3$ lần}
\loigiai{
$15 : 5 = 3$ lần.\\
\textbf{Đáp án của câu hỏi là D}
}
\end{ex}

\begin{ex} %[2H3C16]
Có một can $10\text{ l}$ đựng đầy nước và hai can rỗng loại $3\text{ l}$ và $5\text{ l}$. Muốn lấy đúng $2\text{ l}$ nước, ta có thể
\choice
{\True Rót đầy can $5\text{ l}$ rồi từ can $5\text{ l}$ rót đầy can $3\text{ l}$, lượng nước còn lại trong can $5\text{ l}$ là $2\text{ l}$}
{Rót đầy can $3\text{ l}$}
{Rót nửa can $5\text{ l}$}
{Rót nửa can $10\text{ l}$}
\loigiai{
Lấy $5\text{ l} - 3\text{ l} = 2\text{ l}$.\\
\textbf{Đáp án của câu hỏi là A}
}
\end{ex}

\textbf{Phần 2. Trắc nghiệm đúng sai}
\setcounterpageref{ex}{0}

\begin{ex} %[2H3H16]
Xét tính đúng sai của các phép tính với đơn vị lít:
\choiceTF
{\True $18\text{ l} + 12\text{ l} = 30\text{ l}$.}
{\True $45\text{ l} - 25\text{ l} = 20\text{ l}$.}
{$10\text{ l} + 5 = 15\text{ l}$.}
{\True Khi cộng trừ dung tích phải viết kèm đơn vị $\text{l}$.}
\loigiai{
\begin{itemize}
    \item $18 + 12 = 30\text{ l}$: Đúng.
    \item $45 - 25 = 20\text{ l}$: Đúng.
    \item Ý 3 thiếu đơn vị ở số hạng thứ 2. Sai.
    \item Cần kèm đơn vị đo: Đúng.
\end{itemize}
\begin{itemize}
    \item a) \textbf{Đúng}.
    \item b) \textbf{Đúng}.
    \item c) \textbf{Sai}.
    \item d) \textbf{Đúng}.
\end{itemize}
}
\end{ex}

\begin{ex} %[2H3V16]
Bình nước có $16\text{ l}$, người ta dùng hết $7\text{ l}$. Xét tính đúng sai:
\choiceTF
{\True Lượng nước còn lại là $9\text{ l}$ ($16 - 7 = 9$).}
{\True Lượng nước ban đầu nhiều hơn lượng nước đã dùng là $9\text{ l}$.}
{Nếu đổ thêm $4\text{ l}$ nước vào bình thì trong bình có $15\text{ l}$.}
{\True Khi đó trong bình có $9 + 4 = 13\text{ l}$ nước.}
\loigiai{
\begin{itemize}
    \item Còn lại $9\text{ l}$: Đúng.
    \item $16 - 7 = 9\text{ l}$: Đúng.
    \item $9 + 4 = 13 \neq 15$. Ý 3 Sai, Ý 4 Đúng.
\end{itemize}
\begin{itemize}
    \item a) \textbf{Đúng}.
    \item b) \textbf{Đúng}.
    \item c) \textbf{Sai}.
    \item d) \textbf{Đúng}.
\end{itemize}
}
\end{ex}

\textbf{Phần 3. Trả lời ngắn}
\setcounterpageref{ex}{0}

\begin{ex} %[2H3H16]
Tính: $19\text{ l} + 8\text{ l}$ theo lít.
\shortans{27}
\loigiai{
\textbf{Đáp số:} $27$\\
$19 + 8 = 27\text{ l}$.
}
\end{ex}

\begin{ex} %[2H3H16]
Tính: $40\text{ l} - 15\text{ l}$ theo lít.
\shortans{25}
\loigiai{
\textbf{Đáp số:} $25$\\
$40 - 15 = 25\text{ l}$.
}
\end{ex}

\begin{ex} %[2H3V16]
Can to có $25\text{ l}$ nước mắm, can nhỏ có $15\text{ l}$ nước mắm. Cả hai can có tất cả bao nhiêu lít nước mắm?
\shortans{40}
\loigiai{
\textbf{Đáp số:} $40$\\
$25 + 15 = 40\text{ l}$.
}
\end{ex}

\begin{ex} %[2H3V16]
Một thùng có $50\text{ l}$ dầu, đã bán đi $30\text{ l}$ dầu. Hỏi trong thùng còn lại bao nhiêu lít dầu?
\shortans{20}
\loigiai{
\textbf{Đáp số:} $20$\\
$50 - 30 = 20\text{ l}$.
}
\end{ex}

\begin{ex} %[2H3V16]
Tính giá trị biểu thức: $20\text{ l} + 30\text{ l} - 15\text{ l}$ theo lít.
\shortans{35}
\loigiai{
\textbf{Đáp số:} $35$\\
$50 - 15 = 35\text{ l}$.
}
\end{ex}

\begin{ex} %[2H3C16]
Thùng thứ nhất có $18\text{ l}$ dầu, thùng thứ hai có ít hơn thùng thứ nhất $6\text{ l}$ dầu. Hỏi cả hai thùng có tất cả bao nhiêu lít dầu?
\shortans{30}
\loigiai{
\textbf{Đáp số:} $30$\\
Thùng 2: $18 - 6 = 12\text{ l}$. Cả hai thùng: $18 + 12 = 30\text{ l}$.
}
\end{ex}

\newpage

\TaoPhieuKiemTra{Chủ đề 3 -- Làm quen với khối lượng, dung tích}{17}{Thực hành cân đo}{https://loptoanthaylinh.vn/thi-online/toan-2/bai-17}
\textbf{Phần 1. Trắc nghiệm lựa chọn}
\setcounterpageref{ex}{0}

\begin{ex} %[2H3N17]
Cặp sách đi học của em cân nặng khoảng
\choice
{$30\text{ kg}$}
{$300\text{ kg}$}
{\True $3\text{ kg}$}
{$300\text{ g}$ (quá nhẹ)}
\loigiai{
Cặp sách học sinh nặng khoảng $3\text{ kg}$.\\
\textbf{Đáp án của câu hỏi là C}
}
\end{ex}

\begin{ex} %[2H3N17]
Một xô nước đầy có dung tích khoảng
\choice
{$1\text{ l}$}
{$100\text{ l}$}
{$1000\text{ l}$}
{\True $10\text{ l}$}
\loigiai{
Xô nước chứa khoảng $10\text{ l}$.\\
\textbf{Đáp án của câu hỏi là D}
}
\end{ex}

\begin{ex} %[2H3H17]
Con gà cân nặng $2\text{ kg}$, con ngỗng nặng hơn con gà $3\text{ kg}$. Con ngỗng cân nặng là
\choice
{$4\text{ kg}$}
{\True $5\text{ kg}$}
{$6\text{ kg}$}
{$7\text{ kg}$}
\loigiai{
$2 + 3 = 5\text{ kg}$.\\
\textbf{Đáp án của câu hỏi là B}
}
\end{ex}

\begin{ex} %[2H3H17]
Một chai nước ngọt chứa $2\text{ l}$, rót đầy vào 4 cốc như nhau. Mỗi cốc chứa được
\choice
{\True Nửa lít ($500\text{ ml}$)}
{$1\text{ l}$}
{$2\text{ l}$}
{$4\text{ l}$}
\loigiai{
Rót đều được nửa lít mỗi cốc.\\
\textbf{Đáp án của câu hỏi là A}
}
\end{ex}

\begin{ex} %[2H3H17]
Một túi cam nặng $3\text{ kg}$, 4 túi cam như thế nặng tất cả là
\choice
{$10\text{ kg}$}
{\True $12\text{ kg}$}
{$14\text{ kg}$}
{$15\text{ kg}$}
\loigiai{
$3 \times 4 = 12\text{ kg}$.\\
\textbf{Đáp án của câu hỏi là B}
}
\end{ex}

\begin{ex} %[2H3H17]
Có 3 can nước, mỗi can chứa $5\text{ l}$. Cả 3 can chứa số lít nước là
\choice
{$12\text{ l}$}
{$18\text{ l}$}
{$20\text{ l}$}
{\True $15\text{ l}$}
\loigiai{
$5 \times 3 = 15\text{ l}$.\\
\textbf{Đáp án của câu hỏi là D}
}
\end{ex}

\begin{ex} %[2H3V17]
Mẹ mua một con cá nặng $2\text{ kg}$ và một miếng thịt nặng $1\text{ kg}$. Tổng khối lượng mẹ đã mua là
\choice
{$2\text{ kg}$}
{$4\text{ kg}$}
{\True $3\text{ kg}$}
{$5\text{ kg}$}
\loigiai{
$2 + 1 = 3\text{ kg}$.\\
\textbf{Đáp án của câu hỏi là C}
}
\end{ex}

\begin{ex} %[2H3V17]
Bình chứa $8\text{ l}$ nước ép táo, rót ra mời khách hết $5\text{ l}$. Bình còn lại số lít nước ép là
\choice
{$2\text{ l}$}
{$4\text{ l}$}
{$5\text{ l}$}
{\True $3\text{ l}$}
\loigiai{
$8 - 5 = 3\text{ l}$.\\
\textbf{Đáp án của câu hỏi là D}
}
\end{ex}

\begin{ex} %[2H3V17]
Trên đĩa cân bên trái có 1 quả dưa hấu, đĩa bên phải có 2 quả cân loại $2\text{ kg}$ thì cân thăng bằng. Quả dưa hấu nặng là
\choice
{$2\text{ kg}$}
{$6\text{ kg}$}
{\True $4\text{ kg}$}
{$8\text{ kg}$}
\loigiai{
$2 + 2 = 4\text{ kg}$.\\
\textbf{Đáp án của câu hỏi là C}
}
\end{ex}

\begin{ex} %[2H3V17]
Rót hết nước từ can $10\text{ l}$ sang hai bình: bình một được $6\text{ l}$, bình hai được
\choice
{\True $4\text{ l}$}
{$3\text{ l}$}
{$5\text{ l}$}
{$6\text{ l}$}
\loigiai{
$10 - 6 = 4\text{ l}$.\\
\textbf{Đáp án của câu hỏi là A}
}
\end{ex}

\begin{ex} %[2H3V17]
Một bao gạo nặng $50\text{ kg}$, người ta chia đều vào 5 túi nhỏ. Mỗi túi nặng là
\choice
{$8\text{ kg}$}
{\True $10\text{ kg}$}
{$12\text{ kg}$}
{$15\text{ kg}$}
\loigiai{
$50 : 5 = 10\text{ kg}$.\\
\textbf{Đáp án của câu hỏi là B}
}
\end{ex}

\begin{ex} %[2H3C17]
Một con mèo và một con gà cùng đứng lên cân thì cân chỉ $5\text{ kg}$. Biết con mèo nặng hơn con gà $1\text{ kg}$. Con mèo cân nặng là
\choice
{\True $3\text{ kg}$}
{$2\text{ kg}$}
{$4\text{ kg}$}
{$1\text{ kg}$}
\loigiai{
Con mèo: $(5 + 1) : 2 = 3\text{ kg}$.\\
\textbf{Đáp án của câu hỏi là A}
}
\end{ex}

\textbf{Phần 2. Trắc nghiệm đúng sai}
\setcounterpageref{ex}{0}

\begin{ex} %[2H3H17]
Xét tính hợp lý trong thực tế:
\choiceTF
{\True Một quả trứng gà nặng khoảng vài chục gam (bé hơn $1\text{ kg}$).}
{\True Một con bò cân nặng vài trăm ki-lô-gam.}
{Một người trưởng thành uống khoảng $20\text{ l}$ nước mỗi ngày.}
{\True Mỗi ngày người trưởng thành uống khoảng $1\text{ l}$ đến $2\text{ l}$ nước.}
\loigiai{
\begin{itemize}
    \item Trứng gà bé hơn $1\text{ kg}$: Đúng.
    \item Con bò nặng vài trăm $\text{kg}$: Đúng.
    \item Uống $20\text{ l}$ là phi thực tế. Ý 3 Sai, Ý 4 Đúng.
\end{itemize}
\begin{itemize}
    \item a) \textbf{Đúng}.
    \item b) \textbf{Đúng}.
    \item c) \textbf{Sai}.
    \item d) \textbf{Đúng}.
\end{itemize}
}
\end{ex}

\begin{ex} %[2H3V17]
Có 2 can dầu: Can A chứa $15\text{ l}$, can B chứa $20\text{ l}$. Xét tính đúng sai:
\choiceTF
{\True Cả hai can chứa $35\text{ l}$ dầu.}
{\True Can B chứa nhiều hơn can A là $5\text{ l}$ dầu.}
{Nếu rót $5\text{ l}$ từ can B sang can A thì can B còn $10\text{ l}$.}
{\True Khi đó can B còn $20 - 5 = 15\text{ l}$ và hai can bằng nhau.}
\loigiai{
\begin{itemize}
    \item Tổng: $15 + 20 = 35\text{ l}$: Đúng.
    \item Hiệu: $20 - 15 = 5\text{ l}$: Đúng.
    \item $20 - 5 = 15\text{ l} \neq 10\text{ l}$. Ý 3 Sai, Ý 4 Đúng.
\end{itemize}
\begin{itemize}
    \item a) \textbf{Đúng}.
    \item b) \textbf{Đúng}.
    \item c) \textbf{Sai}.
    \item d) \textbf{Đúng}.
\end{itemize}
}
\end{ex}

\textbf{Phần 3. Trả lời ngắn}
\setcounterpageref{ex}{0}

\begin{ex} %[2H3H17]
Một bao xi măng nặng $50\text{ kg}$, một bao cát nặng $40\text{ kg}$. Tính tổng khối lượng cả hai bao theo ki-lô-gam.
\shortans{90}
\loigiai{
\textbf{Đáp số:} $90$\\
$50 + 40 = 90\text{ kg}$.
}
\end{ex}

\begin{ex} %[2H3H17]
Một thùng nước chứa $60\text{ l}$, đã dùng hết $25\text{ l}$. Còn lại bao nhiêu lít nước?
\shortans{35}
\loigiai{
\textbf{Đáp số:} $35$\\
$60 - 25 = 35\text{ l}$.
}
\end{ex}

\begin{ex} %[2H3V17]
Mẹ mua 3 túi táo, mỗi túi nặng $2\text{ kg}$. Tính tổng cân nặng của 3 túi táo theo ki-lô-gam.
\shortans{6}
\loigiai{
\textbf{Đáp số:} $6$\\
$2 \times 3 = 6\text{ kg}$.
}
\end{ex}

\begin{ex} %[2H3V17]
Có 4 can nước mắm, mỗi can chứa $5\text{ l}$. Tính tổng số lít nước mắm trong 4 can.
\shortans{20}
\loigiai{
\textbf{Đáp số:} $20$\\
$5 \times 4 = 20\text{ l}$.
}
\end{ex}

\begin{ex} %[2H3V17]
Bé An nặng $24\text{ kg}$, bé Bình nhẹ hơn bé An $3\text{ kg}$. Hỏi bé Bình nặng bao nhiêu ki-lô-gam?
\shortans{21}
\loigiai{
\textbf{Đáp số:} $21$\\
$24 - 3 = 21\text{ kg}$.
}
\end{ex}

\begin{ex} %[2H3C17]
Một thùng to đựng $30\text{ l}$ nước, rót đầy sang một can $12\text{ l}$ và một can $8\text{ l}$. Hỏi trong thùng to còn lại bao nhiêu lít nước?
\shortans{10}
\loigiai{
\textbf{Đáp số:} $10$\\
$30 - (12 + 8) = 10\text{ l}$.
}
\end{ex}

\newpage

\TaoPhieuKiemTra{Chủ đề 3 -- Làm quen với khối lượng, dung tích}{18}{Luyện tập chung}{https://loptoanthaylinh.vn/thi-online/toan-2/bai-18}
\textbf{Phần 1. Trắc nghiệm lựa chọn}
\setcounterpageref{ex}{0}

\begin{ex} %[2H3N18]
Đơn vị dùng để đo khối lượng là
\choice
{\True Ki-lô-gam}
{Lít}
{Mét}
{Đề-xi-mét}
\loigiai{
Đơn vị đo khối lượng là ki-lô-gam.\\
\textbf{Đáp án của câu hỏi là A}
}
\end{ex}

\begin{ex} %[2H3N18]
Đơn vị dùng để đo sức chứa chất lỏng là
\choice
{Ki-lô-gam}
{Xăng-ti-mét}
{Giờ}
{\True Lít}
\loigiai{
Đo dung tích dùng đơn vị lít.\\
\textbf{Đáp án của câu hỏi là D}
}
\end{ex}

\begin{ex} %[2H3H18]
Tính: $25\text{ kg} + 15\text{ kg} = \dots$
\choice
{$30\text{ kg}$}
{$50\text{ kg}$}
{\True $40\text{ kg}$}
{$35\text{ kg}$}
\loigiai{
$25 + 15 = 40\text{ kg}$.\\
\textbf{Đáp án của câu hỏi là C}
}
\end{ex}

\begin{ex} %[2H3H18]
Tính: $48\text{ l} - 20\text{ l} = \dots$
\choice
{$26\text{ l}$}
{\True $28\text{ l}$}
{$30\text{ l}$}
{$18\text{ l}$}
\loigiai{
$48 - 20 = 28\text{ l}$.\\
\textbf{Đáp án của câu hỏi là B}
}
\end{ex}

\begin{ex} %[2H3H18]
Bao gạo nặng $45\text{ kg}$, bao ngô nhẹ hơn bao gạo $10\text{ kg}$. Bao ngô nặng số ki-lô-gam là
\choice
{$30\text{ kg}$}
{$40\text{ kg}$}
{$55\text{ kg}$}
{\True $35\text{ kg}$}
\loigiai{
$45 - 10 = 35\text{ kg}$.\\
\textbf{Đáp án của câu hỏi là D}
}
\end{ex}

\begin{ex} %[2H3H18]
Can to có $15\text{ l}$ dầu, can nhỏ có $8\text{ l}$ dầu. Cả hai can có số lít dầu là
\choice
{\True $23\text{ l}$}
{$22\text{ l}$}
{$24\text{ l}$}
{$25\text{ l}$}
\loigiai{
$15 + 8 = 23\text{ l}$.\\
\textbf{Đáp án của câu hỏi là A}
}
\end{ex}

\begin{ex} %[2H3V18]
Tính giá trị biểu thức: $50\text{ kg} - 20\text{ kg} + 15\text{ kg}$.
\choice
{$40\text{ kg}$}
{\True $45\text{ kg}$}
{$50\text{ kg}$}
{$55\text{ kg}$}
\loigiai{
$30 + 15 = 45\text{ kg}$.\\
\textbf{Đáp án của câu hỏi là B}
}
\end{ex}

\begin{ex} %[2H3V18]
Tính giá trị biểu thức: $60\text{ l} - 15\text{ l} - 15\text{ l}$.
\choice
{\True $30\text{ l}$}
{$25\text{ l}$}
{$35\text{ l}$}
{$40\text{ l}$}
\loigiai{
$60 - 30 = 30\text{ l}$.\\
\textbf{Đáp án của câu hỏi là A}
}
\end{ex}

\begin{ex} %[2H3V18]
Mẹ mua $12\text{ kg}$ dưa hấu và $8\text{ kg}$ xoài. Mẹ mua tất cả số ki-lô-gam quả là
\choice
{$18\text{ kg}$}
{$22\text{ kg}$}
{\True $20\text{ kg}$}
{$24\text{ kg}$}
\loigiai{
$12 + 8 = 20\text{ kg}$.\\
\textbf{Đáp án của câu hỏi là C}
}
\end{ex}

\begin{ex} %[2H3V18]
Bình thứ nhất có $14\text{ l}$ nước, bình thứ hai có ít hơn bình thứ nhất $5\text{ l}$ nước. Bình thứ hai có số lít nước là
\choice
{$8\text{ l}$}
{$10\text{ l}$}
{$11\text{ l}$}
{\True $9\text{ l}$}
\loigiai{
$14 - 5 = 9\text{ l}$.\\
\textbf{Đáp án của câu hỏi là D}
}
\end{ex}

\begin{ex} %[2H3V18]
Điền dấu thích hợp: $35\text{ kg} + 10\text{ kg} \dots 50\text{ kg}$.
\choice
{$>$}
{\True $<$}
{$=$}
{Không so sánh được}
\loigiai{
$35 + 10 = 45\text{ kg} < 50\text{ kg}$.\\
\textbf{Đáp án của câu hỏi là B}
}
\end{ex}

\begin{ex} %[2H3C18]
Có hai bao đường, bao thứ nhất nặng $25\text{ kg}$. Nếu chuyển $5\text{ kg}$ từ bao thứ nhất sang bao thứ hai thì hai bao nặng bằng nhau. Hỏi bao thứ hai ban đầu nặng bao nhiêu ki-lô-gam?
\choice
{$10\text{ kg}$}
{$20\text{ kg}$}
{\True $15\text{ kg}$}
{$25\text{ kg}$}
\loigiai{
Sau khi chuyển, bao 1 nặng: $25 - 5 = 20\text{ kg}$.\\
Lúc đó bao 2 cũng nặng $20\text{ kg} \Rightarrow$ Ban đầu bao 2 nặng: $20 - 5 = 15\text{ kg}$.\\
\textbf{Đáp án của câu hỏi là C}
}
\end{ex}

\textbf{Phần 2. Trắc nghiệm đúng sai}
\setcounterpageref{ex}{0}

\begin{ex} %[2H3H18]
Xét tính đúng sai về khối lượng và dung tích:
\choiceTF
{\True Ki-lô-gam viết tắt là $\text{kg}$.}
{\True Lít viết tắt là $\text{l}$.}
{$20\text{ kg} + 10\text{ l} = 30\text{ kg}$.}
{\True Không thể cộng hai đại lượng khác đơn vị đo ($	ext{kg}$ và $	ext{l}$).}
\loigiai{
\begin{itemize}
    \item $\text{kg}$ là ki-lô-gam: Đúng.
    \item $\text{l}$ là lít: Đúng.
    \item Khác đơn vị đo không cộng trực tiếp được. Ý 3 Sai, Ý 4 Đúng.
\end{itemize}
\begin{itemize}
    \item a) \textbf{Đúng}.
    \item b) \textbf{Đúng}.
    \item c) \textbf{Sai}.
    \item d) \textbf{Đúng}.
\end{itemize}
}
\end{ex}

\begin{ex} %[2H3V18]
Thùng A có $35\text{ l}$ dầu, thùng B có $25\text{ l}$ dầu. Xét tính đúng sai:
\choiceTF
{\True Cả hai thùng có tất cả $60\text{ l}$ dầu.}
{\True Thùng A nhiều hơn thùng B là $10\text{ l}$ dầu.}
{Nếu múc $10\text{ l}$ từ thùng A thì hai thùng bằng nhau.}
{\True Khi đó thùng A còn $35 - 10 = 25\text{ l}$ bằng thùng B.}
\loigiai{
\begin{itemize}
    \item Tổng: $35 + 25 = 60\text{ l}$: Đúng.
    \item Hiệu: $35 - 25 = 10\text{ l}$: Đúng.
    \item Múc 10 lít thì còn 25 lít: Đúng. Cả hai ý cuối đều Đúng.
\end{itemize}
\begin{itemize}
    \item a) \textbf{Đúng}.
    \item b) \textbf{Đúng}.
    \item c) \textbf{Đúng}.
    \item d) \textbf{Đúng}.
\end{itemize}
}
\end{ex}

\textbf{Phần 3. Trả lời ngắn}
\setcounterpageref{ex}{0}

\begin{ex} %[2H3H18]
Tính: $28\text{ kg} + 12\text{ kg}$.
\shortans{40}
\loigiai{
\textbf{Đáp số:} $40$\\
$28 + 12 = 40\text{ kg}$.
}
\end{ex}

\begin{ex} %[2H3H18]
Tính: $50\text{ l} - 18\text{ l}$.
\shortans{32}
\loigiai{
\textbf{Đáp số:} $32$\\
$50 - 18 = 32\text{ l}$.
}
\end{ex}

\begin{ex} %[2H3V18]
Con lợn to nặng $75\text{ kg}$, con lợn nhỏ nặng $45\text{ kg}$. Con lợn to nặng hơn con lợn nhỏ bao nhiêu ki-lô-gam?
\shortans{30}
\loigiai{
\textbf{Đáp số:} $30$\\
$75 - 45 = 30\text{ kg}$.
}
\end{ex}

\begin{ex} %[2H3V18]
Một can chứa $20\text{ l}$ dầu, mẹ đã dùng $12\text{ l}$ dầu. Hỏi trong can còn lại bao nhiêu lít dầu?
\shortans{8}
\loigiai{
\textbf{Đáp số:} $8$\\
$20 - 12 = 8\text{ l}$.
}
\end{ex}

\begin{ex} %[2H3V18]
Tính giá trị biểu thức: $25\text{ kg} + 35\text{ kg} - 20\text{ kg}$.
\shortans{40}
\loigiai{
\textbf{Đáp số:} $40$\\
$60 - 20 = 40\text{ kg}$.
}
\end{ex}

\begin{ex} %[2H3C18]
Thùng thứ nhất có $24\text{ l}$ sữa, thùng thứ hai có nhiều hơn thùng thứ nhất $6\text{ l}$ sữa. Hỏi cả hai thùng có bao nhiêu lít sữa?
\shortans{54}
\loigiai{
\textbf{Đáp số:} $54$\\
Thùng 2: $24 + 6 = 30\text{ l}$. Cả hai thùng: $24 + 30 = 54\text{ l}$.
}
\end{ex}

\newpage

% =========================================================================
% CHỦ ĐỀ 4 - PHÉP CỘNG, PHÉP TRỪ (CÓ NHỚ) TRONG PHẠM VI 100
% =========================================================================
\section{Phép cộng, phép trừ (có nhớ) trong phạm vi 100}

\TaoPhieuKiemTra{Chủ đề 4 -- Phép cộng, phép trừ (có nhớ) trong phạm vi 100}{19}{Phép cộng (có nhớ) số có hai chữ số với số có một chữ số}{https://loptoanthaylinh.vn/thi-online/toan-2/bai-19}
\textbf{Phần 1. Trắc nghiệm lựa chọn}
\setcounterpageref{ex}{0}

\begin{ex} %[2D4N19]
Kết quả của phép tính $37 + 5$ là
\choice
{$41$}
{$43$}
{$44$}
{\True $42$}
\loigiai{
$37 + 5 = 42$.\\
\textbf{Đáp án của câu hỏi là D}
}
\end{ex}

\begin{ex} %[2D4N19]
Kết quả của phép tính $48 + 4$ là
\choice
{\True $52$}
{$51$}
{$53$}
{$54$}
\loigiai{
$48 + 4 = 52$.\\
\textbf{Đáp án của câu hỏi là A}
}
\end{ex}

\begin{ex} %[2D4H19]
Kết quả của phép tính $56 + 8$ là
\choice
{$63$}
{\True $64$}
{$65$}
{$66$}
\loigiai{
$56 + 8 = 64$.\\
\textbf{Đáp án của câu hỏi là B}
}
\end{ex}

\begin{ex} %[2D4H19]
Kết quả của phép tính $69 + 7$ là
\choice
{$75$}
{$77$}
{\True $76$}
{$78$}
\loigiai{
$69 + 7 = 76$.\\
\textbf{Đáp án của câu hỏi là C}
}
\end{ex}

\begin{ex} %[2D4H19]
Tìm $x$ biết: $x - 7 = 45$.
\choice
{$51$}
{$53$}
{$54$}
{\True $52$}
\loigiai{
$x = 45 + 7 = 52$.\\
\textbf{Đáp án của câu hỏi là D}
}
\end{ex}

\begin{ex} %[2D4H19]
Điền dấu thích hợp: $58 + 6 \dots 49 + 8$.
\choice
{\True $>$}
{$<$}
{$=$}
{Không so sánh được}
\loigiai{
$58 + 6 = 64$ và $49 + 8 = 57$. Vì $64 > 57$ nên điền dấu $>$.\\
\textbf{Đáp án của câu hỏi là A}
}
\end{ex}

\begin{ex} %[2D4V19]
Một lớp học có 28 bạn nam và 6 bạn nữ chuyển vào thêm. Số học sinh của lớp lúc này là
\choice
{$33$ bạn}
{$35$ bạn}
{\True $34$ bạn}
{$36$ bạn}
\loigiai{
$28 + 6 = 34$ bạn.\\
\textbf{Đáp án của câu hỏi là C}
}
\end{ex}

\begin{ex} %[2D4V19]
Vườn nhà Mai có 37 cây chuối, bố trồng thêm 8 cây chuối nữa. Vườn nhà Mai có tất cả số cây chuối là
\choice
{$44$ cây}
{$46$ cây}
{$47$ cây}
{\True $45$ cây}
\loigiai{
$37 + 8 = 45$ cây.\\
\textbf{Đáp án của câu hỏi là D}
}
\end{ex}

\begin{ex} %[2D4V19]
Tính giá trị biểu thức: $46 + 7 + 3$.
\choice
{$55$}
{\True $56$}
{$57$}
{$58$}
\loigiai{
$46 + (7 + 3) = 46 + 10 = 56$.\\
\textbf{Đáp án của câu hỏi là B}
}
\end{ex}

\begin{ex} %[2D4V19]
Tìm một số biết lấy số đó trừ đi 9 thì được 58.
\choice
{\True $67$}
{$66$}
{$68$}
{$69$}
\loigiai{
$58 + 9 = 67$.\\
\textbf{Đáp án của câu hỏi là A}
}
\end{ex}

\begin{ex} %[2D4V19]
Đoạn dây thứ nhất dài $35\text{ cm}$, đoạn dây thứ hai dài hơn đoạn thứ nhất $9\text{ cm}$. Độ dài đoạn dây thứ hai là
\choice
{$43\text{ cm}$}
{$45\text{ cm}$}
{\True $44\text{ cm}$}
{$46\text{ cm}$}
\loigiai{
$35 + 9 = 44\text{ cm}$.\\
\textbf{Đáp án của câu hỏi là C}
}
\end{ex}

\begin{ex} %[2D4C19]
Tìm số thích hợp điền vào ô trống: $48 + \dots = 55$.
\choice
{$6$}
{\True $7$}
{$8$}
{$9$}
\loigiai{
$55 - 48 = 7$.\\
\textbf{Đáp án của câu hỏi là B}
}
\end{ex}

\textbf{Phần 2. Trắc nghiệm đúng sai}
\setcounterpageref{ex}{0}

\begin{ex} %[2D4H19]
Xét tính đúng sai của các phép cộng có nhớ:
\choiceTF
{\True $27 + 6 = 33$.}
{\True $38 + 5 = 43$.}
{$49 + 4 = 52$.}
{\True $49 + 4 = 53$.}
\loigiai{
\begin{itemize}
    \item $27 + 6 = 33$: Đúng.
    \item $38 + 5 = 43$: Đúng.
    \item $49 + 4 = 53 \neq 52$. Ý 3 Sai, Ý 4 Đúng.
\end{itemize}
\begin{itemize}
    \item a) \textbf{Đúng}.
    \item b) \textbf{Đúng}.
    \item c) \textbf{Sai}.
    \item d) \textbf{Đúng}.
\end{itemize}
}
\end{ex}

\begin{ex} %[2D4V19]
Lan có 38 con tem, Bình cho Lan thêm 7 con tem. Xét tính đúng sai:
\choiceTF
{\True Lan có tất cả 45 con tem ($38 + 7 = 45$).}
{\True Số tem của Lan đã tăng thêm 7 con tem.}
{Nếu Lan cho Mai 5 con tem thì Lan còn lại 39 con tem.}
{\True Khi đó Lan còn lại $45 - 5 = 40$ con tem.}
\loigiai{
\begin{itemize}
    \item Tổng: $38 + 7 = 45$: Đúng.
    \item Tăng thêm 7: Đúng.
    \item $45 - 5 = 40 \neq 39$. Ý 3 Sai, Ý 4 Đúng.
\end{itemize}
\begin{itemize}
    \item a) \textbf{Đúng}.
    \item b) \textbf{Đúng}.
    \item c) \textbf{Sai}.
    \item d) \textbf{Đúng}.
\end{itemize}
}
\end{ex}

\textbf{Phần 3. Trả lời ngắn}
\setcounterpageref{ex}{0}

\begin{ex} %[2D4H19]
Tính tổng: $47 + 6$.
\shortans{53}
\loigiai{
\textbf{Đáp số:} $53$\\
$47 + 6 = 53$.
}
\end{ex}

\begin{ex} %[2D4H19]
Tính tổng: $58 + 7$.
\shortans{65}
\loigiai{
\textbf{Đáp số:} $65$\\
$58 + 7 = 65$.
}
\end{ex}

\begin{ex} %[2D4V19]
Tính tổng: $79 + 8$.
\shortans{87}
\loigiai{
\textbf{Đáp số:} $87$\\
$79 + 8 = 87$.
}
\end{ex}

\begin{ex} %[2D4V19]
Tìm $x$ biết: $x - 8 = 65$.
\shortans{73}
\loigiai{
\textbf{Đáp số:} $73$\\
$x = 65 + 8 = 73$.
}
\end{ex}

\begin{ex} %[2D4V19]
Hà có 26 cái kẹo, mẹ cho thêm 8 cái kẹo nữa. Hỏi Hà có tất cả bao nhiêu cái kẹo?
\shortans{34}
\loigiai{
\textbf{Đáp số:} $34$\\
$26 + 8 = 34$ cái.
}
\end{ex}

\begin{ex} %[2D4C19]
Điền số thích hợp: $68 + ? = 75$.
\shortans{7}
\loigiai{
\textbf{Đáp số:} $7$\\
$75 - 68 = 7$.
}
\end{ex}

\newpage

\TaoPhieuKiemTra{Chủ đề 4 -- Phép cộng, phép trừ (có nhớ) trong phạm vi 100}{20}{Phép cộng (có nhớ) số có hai chữ số với số có hai chữ số}{https://loptoanthaylinh.vn/thi-online/toan-2/bai-20}
\textbf{Phần 1. Trắc nghiệm lựa chọn}
\setcounterpageref{ex}{0}

\begin{ex} %[2D4N20]
Kết quả của phép tính $38 + 25$ là
\choice
{\True $63$}
{$62$}
{$64$}
{$65$}
\loigiai{
$38 + 25 = 63$.\\
\textbf{Đáp án của câu hỏi là A}
}
\end{ex}

\begin{ex} %[2D4N20]
Kết quả của phép tính $47 + 36$ là
\choice
{$82$}
{\True $83$}
{$84$}
{$73$}
\loigiai{
$47 + 36 = 83$.\\
\textbf{Đáp án của câu hỏi là B}
}
\end{ex}

\begin{ex} %[2D4H20]
Kết quả của phép tính $59 + 28$ là
\choice
{$86$}
{$88$}
{$89$}
{\True $87$}
\loigiai{
$59 + 28 = 87$.\\
\textbf{Đáp án của câu hỏi là D}
}
\end{ex}

\begin{ex} %[2D4H20]
Kết quả của phép tính $64 + 18$ là
\choice
{$81$}
{\True $82$}
{$83$}
{$72$}
\loigiai{
$64 + 18 = 82$.\\
\textbf{Đáp án của câu hỏi là B}
}
\end{ex}

\begin{ex} %[2D4H20]
Tìm $x$ biết: $x - 26 = 38$.
\choice
{$63$}
{$65$}
{$66$}
{\True $64$}
\loigiai{
$x = 38 + 26 = 64$.\\
\textbf{Đáp án của câu hỏi là D}
}
\end{ex}

\begin{ex} %[2D4H20]
Điền dấu thích hợp: $36 + 28 \dots 45 + 19$.
\choice
{$>$}
{\True $=$}
{$<$}
{Không so sánh được}
\loigiai{
$36 + 28 = 64$ và $45 + 19 = 64$. Vậy $36 + 28 = 45 + 19$.\\
\textbf{Đáp án của câu hỏi là B}
}
\end{ex}

\begin{ex} %[2D4V20]
Một cửa hàng buổi sáng bán được $38\text{ kg}$ đường, buổi chiều bán được $27\text{ kg}$ đường. Cả ngày bán được số ki-lô-gam đường là
\choice
{$64\text{ kg}$}
{$66\text{ kg}$}
{\True $65\text{ kg}$}
{$67\text{ kg}$}
\loigiai{
$38 + 27 = 65\text{ kg}$.\\
\textbf{Đáp án của câu hỏi là C}
}
\end{ex}

\begin{ex} %[2D4V20]
Lớp 2A có 28 bạn, lớp 2B có 29 bạn. Cả hai lớp có tất cả số học sinh là
\choice
{$56$ bạn}
{$58$ bạn}
{$59$ bạn}
{\True $57$ bạn}
\loigiai{
$28 + 29 = 57$ bạn.\\
\textbf{Đáp án của câu hỏi là D}
}
\end{ex}

\begin{ex} %[2D4V20]
Tính giá trị biểu thức: $27 + 35 + 18$.
\choice
{\True $80$}
{$78$}
{$82$}
{$84$}
\loigiai{
$62 + 18 = 80$.\\
\textbf{Đáp án của câu hỏi là A}
}
\end{ex}

\begin{ex} %[2D4V20]
Một cuộn vải hoa dài $45\text{ m}$, cuộn vải trắng dài $38\text{ m}$. Cả hai cuộn vải dài tất cả là
\choice
{$82\text{ m}$}
{$84\text{ m}$}
{\True $83\text{ m}$}
{$85\text{ m}$}
\loigiai{
$45 + 38 = 83\text{ m}$.\\
\textbf{Đáp án của câu hỏi là C}
}
\end{ex}

\begin{ex} %[2D4V20]
Tính nhanh: $36 + 27 + 14 + 13$.
\choice
{\True $90$}
{$88$}
{$92$}
{$94$}
\loigiai{
$(36 + 14) + (27 + 13) = 50 + 40 = 90$.\\
\textbf{Đáp án của câu hỏi là A}
}
\end{ex}

\begin{ex} %[2D4C20]
Tổng của số lớn nhất có hai chữ số khác nhau và số bé nhất có hai chữ số là bao nhiêu?
\choice
{$107$}
{$109$}
{\True $108$}
{$110$}
\loigiai{
$98 + 10 = 108$.\\
\textbf{Đáp án của câu hỏi là C}
}
\end{ex}

\textbf{Phần 2. Trắc nghiệm đúng sai}
\setcounterpageref{ex}{0}

\begin{ex} %[2D4H20]
Xét tính đúng sai của các phép tính:
\choiceTF
{\True $48 + 27 = 75$.}
{\True $56 + 39 = 95$.}
{$37 + 45 = 81$.}
{\True $37 + 45 = 82$.}
\loigiai{
\begin{itemize}
    \item $48 + 27 = 75$: Đúng.
    \item $56 + 39 = 95$: Đúng.
    \item $37 + 45 = 82 \neq 81$. Ý 3 Sai, Ý 4 Đúng.
\end{itemize}
\begin{itemize}
    \item a) \textbf{Đúng}.
    \item b) \textbf{Đúng}.
    \item c) \textbf{Sai}.
    \item d) \textbf{Đúng}.
\end{itemize}
}
\end{ex}

\begin{ex} %[2D4V20]
Một trang trại nuôi 38 con thỏ trắng và 45 con thỏ xám. Xét tính đúng sai:
\choiceTF
{\True Tổng số thỏ của trang trại là 83 con ($38 + 45 = 83$).}
{\True Thỏ xám nhiều hơn thỏ trắng là 7 con ($45 - 38 = 7$).}
{Nếu mua thêm 10 con thỏ nữa thì tổng số thỏ là 90 con.}
{\True Khi đó tổng số thỏ là $83 + 10 = 93$ con.}
\loigiai{
\begin{itemize}
    \item Tổng: $38 + 45 = 83$: Đúng.
    \item Hiệu: $45 - 38 = 7$: Đúng.
    \item $83 + 10 = 93 \neq 90$. Ý 3 Sai, Ý 4 Đúng.
\end{itemize}
\begin{itemize}
    \item a) \textbf{Đúng}.
    \item b) \textbf{Đúng}.
    \item c) \textbf{Sai}.
    \item d) \textbf{Đúng}.
\end{itemize}
}
\end{ex}

\textbf{Phần 3. Trả lời ngắn}
\setcounterpageref{ex}{0}

\begin{ex} %[2D4H20]
Tính tổng: $46 + 28$.
\shortans{74}
\loigiai{
\textbf{Đáp số:} $74$\\
$46 + 28 = 74$.
}
\end{ex}

\begin{ex} %[2D4H20]
Tính tổng: $57 + 35$.
\shortans{92}
\loigiai{
\textbf{Đáp số:} $92$\\
$57 + 35 = 92$.
}
\end{ex}

\begin{ex} %[2D4V20]
Tính tổng: $39 + 48$.
\shortans{87}
\loigiai{
\textbf{Đáp số:} $87$\\
$39 + 48 = 87$.
}
\end{ex}

\begin{ex} %[2D4V20]
Tìm $x$ biết: $x - 37 = 46$.
\shortans{83}
\loigiai{
\textbf{Đáp số:} $83$\\
$x = 46 + 37 = 83$.
}
\end{ex}

\begin{ex} %[2D4V20]
Đội Một trồng được 48 cây, đội Hai trồng được 37 cây. Cả hai đội trồng được bao nhiêu cây?
\shortans{85}
\loigiai{
\textbf{Đáp số:} $85$\\
$48 + 37 = 85$ cây.
}
\end{ex}

\begin{ex} %[2D4C20]
Tính nhanh tổng: $28 + 35 + 22 + 15$.
\shortans{100}
\loigiai{
\textbf{Đáp số:} $100$\\
$(28 + 22) + (35 + 15) = 50 + 50 = 100$.
}
\end{ex}

\newpage

\TaoPhieuKiemTra{Chủ đề 4 -- Phép cộng, phép trừ (có nhớ) trong phạm vi 100}{21}{Phép trừ (có nhớ) số có hai chữ số cho số có một chữ số}{https://loptoanthaylinh.vn/thi-online/toan-2/bai-21}
\textbf{Phần 1. Trắc nghiệm lựa chọn}
\setcounterpageref{ex}{0}

\begin{ex} %[2D4N21]
Kết quả của phép tính $58 + 7$ là
\choice
{$64$}
{$66$}
{$67$}
{\True $65$}
\loigiai{
$58 + 7 = 65$.\\
\textbf{Đáp án của câu hỏi là D}
}
\end{ex}

\begin{ex} %[2D4N21]
Kết quả của phép tính $46 + 37$ là
\choice
{$82$}
{$84$}
{\True $83$}
{$85$}
\loigiai{
$46 + 37 = 83$.\\
\textbf{Đáp án của câu hỏi là C}
}
\end{ex}

\begin{ex} %[2D4H21]
Kết quả của phép tính $69 + 25$ là
\choice
{\True $94$}
{$93$}
{$95$}
{$96$}
\loigiai{
$69 + 25 = 94$.\\
\textbf{Đáp án của câu hỏi là A}
}
\end{ex}

\begin{ex} %[2D4H21]
Kết quả của phép tính $38 + 48$ là
\choice
{$85$}
{$87$}
{$88$}
{\True $86$}
\loigiai{
$38 + 48 = 86$.\\
\textbf{Đáp án của câu hỏi là D}
}
\end{ex}

\begin{ex} %[2D4H21]
Tìm $x$ biết: $x - 39 = 45$.
\choice
{$83$}
{$85$}
{\True $84$}
{$74$}
\loigiai{
$x = 45 + 39 = 84$.\\
\textbf{Đáp án của câu hỏi là C}
}
\end{ex}

\begin{ex} %[2D4H21]
Điền dấu thích hợp: $47 + 28 \dots 56 + 19$.
\choice
{$>$}
{$<$}
{Không so sánh được}
{\True $=$}
\loigiai{
$47 + 28 = 75$ và $56 + 19 = 75$. Vậy $47 + 28 = 56 + 19$.\\
\textbf{Đáp án của câu hỏi là D}
}
\end{ex}

\begin{ex} %[2D4V21]
Một cửa hàng có $48\text{ kg}$ gạo tẻ và $35\text{ kg}$ gạo nếp. Tổng số gạo cửa hàng có là
\choice
{$82\text{ kg}$}
{\True $83\text{ kg}$}
{$84\text{ kg}$}
{$85\text{ kg}$}
\loigiai{
$48 + 35 = 83\text{ kg}$.\\
\textbf{Đáp án của câu hỏi là B}
}
\end{ex}

\begin{ex} %[2D4V21]
Khối Hai có 48 bạn nam và 47 bạn nữ. Khối Hai có tất cả số học sinh là
\choice
{$94$ bạn}
{$96$ bạn}
{\True $95$ bạn}
{$97$ bạn}
\loigiai{
$48 + 47 = 95$ bạn.\\
\textbf{Đáp án của câu hỏi là C}
}
\end{ex}

\begin{ex} %[2D4V21]
Tính giá trị biểu thức: $36 + 28 + 16$.
\choice
{$78$}
{\True $80$}
{$82$}
{$84$}
\loigiai{
$64 + 16 = 80$.\\
\textbf{Đáp án của câu hỏi là B}
}
\end{ex}

\begin{ex} %[2D4V21]
Một sợi dây cắt làm 2 đoạn: đoạn 1 dài $38\text{ cm}$, đoạn 2 dài $47\text{ cm}$. Sợi dây ban đầu dài là
\choice
{\True $85\text{ cm}$}
{$84\text{ cm}$}
{$86\text{ cm}$}
{$87\text{ cm}$}
\loigiai{
$38 + 47 = 85\text{ cm}$.\\
\textbf{Đáp án của câu hỏi là A}
}
\end{ex}

\begin{ex} %[2D4V21]
Tính nhanh: $29 + 36 + 11 + 14$.
\choice
{$88$}
{\True $90$}
{$92$}
{$94$}
\loigiai{
$(29 + 11) + (36 + 14) = 40 + 50 = 90$.\\
\textbf{Đáp án của câu hỏi là B}
}
\end{ex}

\begin{ex} %[2D4C21]
Tổng của hai số là 78. Nếu tăng số thứ nhất thêm 6 đơn vị và tăng số thứ hai thêm 7 đơn vị thì tổng mới là
\choice
{\True $91$}
{$90$}
{$92$}
{$93$}
\loigiai{
$78 + 6 + 7 = 91$.\\
\textbf{Đáp án của câu hỏi là A}
}
\end{ex}

\textbf{Phần 2. Trắc nghiệm đúng sai}
\setcounterpageref{ex}{0}

\begin{ex} %[2D4H21]
Xét tính đúng sai của các phép cộng có nhớ:
\choiceTF
{\True $39 + 45 = 84$.}
{\True $58 + 26 = 84$.}
{$47 + 38 = 86$.}
{\True $47 + 38 = 85$.}
\loigiai{
\begin{itemize}
    \item $39 + 45 = 84$: Đúng.
    \item $58 + 26 = 84$: Đúng.
    \item $47 + 38 = 85 \neq 86$. Ý 3 Sai, Ý 4 Đúng.
\end{itemize}
\begin{itemize}
    \item a) \textbf{Đúng}.
    \item b) \textbf{Đúng}.
    \item c) \textbf{Sai}.
    \item d) \textbf{Đúng}.
\end{itemize}
}
\end{ex}

\begin{ex} %[2D4V21]
Vườn có 46 cây xoài và 38 cây bưởi. Xét tính đúng sai:
\choiceTF
{\True Vườn có tất cả 84 cây ăn quả ($46 + 38 = 84$).}
{\True Số xoài nhiều hơn số bưởi là 8 cây ($46 - 38 = 8$).}
{Nếu trồng thêm 10 cây xoài thì vườn có 90 cây.}
{\True Khi đó vườn có $84 + 10 = 94$ cây.}
\loigiai{
\begin{itemize}
    \item Tổng: $46 + 38 = 84$: Đúng.
    \item Hiệu: $46 - 38 = 8$: Đúng.
    \item $84 + 10 = 94 \neq 90$. Ý 3 Sai, Ý 4 Đúng.
\end{itemize}
\begin{itemize}
    \item a) \textbf{Đúng}.
    \item b) \textbf{Đúng}.
    \item c) \textbf{Sai}.
    \item d) \textbf{Đúng}.
\end{itemize}
}
\end{ex}

\textbf{Phần 3. Trả lời ngắn}
\setcounterpageref{ex}{0}

\begin{ex} %[2D4H21]
Tính tổng: $48 + 36$.
\shortans{84}
\loigiai{
\textbf{Đáp số:} $84$\\
$48 + 36 = 84$.
}
\end{ex}

\begin{ex} %[2D4H21]
Tính tổng: $59 + 27$.
\shortans{86}
\loigiai{
\textbf{Đáp số:} $86$\\
$59 + 27 = 86$.
}
\end{ex}

\begin{ex} %[2D4V21]
Tính tổng: $67 + 25$.
\shortans{92}
\loigiai{
\textbf{Đáp số:} $92$\\
$67 + 25 = 92$.
}
\end{ex}

\begin{ex} %[2D4V21]
Tìm $x$ biết: $x - 48 = 37$.
\shortans{85}
\loigiai{
\textbf{Đáp số:} $85$\\
$x = 37 + 48 = 85$.
}
\end{ex}

\begin{ex} %[2D4V21]
Thùng to có $45\text{ l}$ dầu, thùng nhỏ có $28\text{ l}$ dầu. Hỏi cả hai thùng có bao nhiêu lít dầu?
\shortans{73}
\loigiai{
\textbf{Đáp số:} $73$\\
$45 + 28 = 73\text{ l}$.
}
\end{ex}

\begin{ex} %[2D4C21]
Tính nhanh tổng: $19 + 28 + 31 + 22$.
\shortans{100}
\loigiai{
\textbf{Đáp số:} $100$\\
$(19 + 31) + (28 + 22) = 50 + 50 = 100$.
}
\end{ex}

\newpage

\TaoPhieuKiemTra{Chủ đề 4 -- Phép cộng, phép trừ (có nhớ) trong phạm vi 100}{22}{Phép trừ (có nhớ) số có hai chữ số cho số có hai chữ số}{https://loptoanthaylinh.vn/thi-online/toan-2/bai-22}
\textbf{Phần 1. Trắc nghiệm lựa chọn}
\setcounterpageref{ex}{0}

\begin{ex} %[2D4N22]
Kết quả của phép tính $42 - 5$ là
\choice
{\True $37$}
{$36$}
{$38$}
{$39$}
\loigiai{
$42 - 5 = 37$.\\
\textbf{Đáp án của câu hỏi là A}
}
\end{ex}

\begin{ex} %[2D4N22]
Kết quả của phép tính $53 - 6$ là
\choice
{$46$}
{\True $47$}
{$48$}
{$49$}
\loigiai{
$53 - 6 = 47$.\\
\textbf{Đáp án của câu hỏi là B}
}
\end{ex}

\begin{ex} %[2D4H22]
Kết quả của phép tính $64 - 8$ là
\choice
{\True $56$}
{$55$}
{$57$}
{$58$}
\loigiai{
$64 - 8 = 56$.\\
\textbf{Đáp án của câu hỏi là A}
}
\end{ex}

\begin{ex} %[2D4H22]
Kết quả của phép tính $71 - 4$ là
\choice
{$66$}
{$68$}
{$69$}
{\True $67$}
\loigiai{
$71 - 4 = 67$.\\
\textbf{Đáp án của câu hỏi là D}
}
\end{ex}

\begin{ex} %[2D4H22]
Tìm $x$ biết: $x + 8 = 52$.
\choice
{$43$}
{$45$}
{\True $44$}
{$46$}
\loigiai{
$x = 52 - 8 = 44$.\\
\textbf{Đáp án của câu hỏi là C}
}
\end{ex}

\begin{ex} %[2D4H22]
Điền dấu thích hợp: $62 - 7 \dots 73 - 9$.
\choice
{$>$}
{$=$}
{Không so sánh được}
{\True $<$}
\loigiai{
$62 - 7 = 55$ và $73 - 9 = 64$. Vì $55 < 64$ nên điền dấu $<$.\\
\textbf{Đáp án của câu hỏi là D}
}
\end{ex}

\begin{ex} %[2D4V22]
Một cây bưởi có 42 quả, đã hái đi 7 quả. Trên cây còn lại số quả bưởi là
\choice
{$34$ quả}
{$36$ quả}
{\True $35$ quả}
{$37$ quả}
\loigiai{
$42 - 7 = 35$ quả.\\
\textbf{Đáp án của câu hỏi là C}
}
\end{ex}

\begin{ex} %[2D4V22]
Lan có 51 que tính, Lan cho bạn 6 que tính. Lan còn lại số que tính là
\choice
{$44$ que}
{\True $45$ que}
{$46$ que}
{$47$ que}
\loigiai{
$51 - 6 = 45$ que tính.\\
\textbf{Đáp án của câu hỏi là B}
}
\end{ex}

\begin{ex} %[2D4V22]
Tính giá trị biểu thức: $65 - 8 + 4$.
\choice
{$60$}
{$62$}
{\True $61$}
{$63$}
\loigiai{
$57 + 4 = 61$.\\
\textbf{Đáp án của câu hỏi là C}
}
\end{ex}

\begin{ex} %[2D4V22]
Tìm $y$ biết: $83 - y = 9$.
\choice
{\True $74$}
{$73$}
{$75$}
{$76$}
\loigiai{
$y = 83 - 9 = 74$.\\
\textbf{Đáp án của câu hỏi là A}
}
\end{ex}

\begin{ex} %[2D4V22]
Một sợi dây dài $70\text{ cm}$, người ta cắt bớt đi $8\text{ cm}$. Sợi dây còn lại dài là
\choice
{$61\text{ cm}$}
{$63\text{ cm}$}
{$64\text{ cm}$}
{\True $62\text{ cm}$}
\loigiai{
$70 - 8 = 62\text{ cm}$.\\
\textbf{Đáp án của câu hỏi là D}
}
\end{ex}

\begin{ex} %[2D4C22]
Tìm số thích hợp điền vào ô trống: $54 - \dots = 48$.
\choice
{$5$}
{\True $6$}
{$7$}
{$8$}
\loigiai{
$54 - 48 = 6$.\\
\textbf{Đáp án của câu hỏi là B}
}
\end{ex}

\textbf{Phần 2. Trắc nghiệm đúng sai}
\setcounterpageref{ex}{0}

\begin{ex} %[2D4H22]
Xét tính đúng sai của các phép trừ có nhớ:
\choiceTF
{\True $43 - 7 = 36$.}
{\True $52 - 8 = 44$.}
{$61 - 5 = 55$.}
{\True $61 - 5 = 56$.}
\loigiai{
\begin{itemize}
    \item $43 - 7 = 36$: Đúng.
    \item $52 - 8 = 44$: Đúng.
    \item $61 - 5 = 56 \neq 55$. Ý 3 Sai, Ý 4 Đúng.
\end{itemize}
\begin{itemize}
    \item a) \textbf{Đúng}.
    \item b) \textbf{Đúng}.
    \item c) \textbf{Sai}.
    \item d) \textbf{Đúng}.
\end{itemize}
}
\end{ex}

\begin{ex} %[2D4V22]
Bác nông dân thu hoạch được $62\text{ kg}$ khoai lang, đã bán đi $8\text{ kg}$. Xét tính đúng sai:
\choiceTF
{\True Khối lượng khoai lang còn lại là $54\text{ kg}$ ($62 - 8 = 54$).}
{\True Số khoai bán đi ít hơn số khoai còn lại.}
{Nếu bán thêm $4\text{ kg}$ nữa thì còn lại $51\text{ kg}$.}
{\True Khi đó còn lại $54 - 4 = 50\text{ kg}$.}
\loigiai{
\begin{itemize}
    \item Còn lại: $62 - 8 = 54\text{ kg}$: Đúng.
    \item Đã bán ($8\text{ kg}$) ít hơn còn lại ($54\text{ kg}$): Đúng.
    \item $54 - 4 = 50 \neq 51$. Ý 3 Sai, Ý 4 Đúng.
\end{itemize}
\begin{itemize}
    \item a) \textbf{Đúng}.
    \item b) \textbf{Đúng}.
    \item c) \textbf{Sai}.
    \item d) \textbf{Đúng}.
\end{itemize}
}
\end{ex}

\textbf{Phần 3. Trả lời ngắn}
\setcounterpageref{ex}{0}

\begin{ex} %[2D4H22]
Tính hiệu: $53 - 7$.
\shortans{46}
\loigiai{
\textbf{Đáp số:} $46$\\
$53 - 7 = 46$.
}
\end{ex}

\begin{ex} %[2D4H22]
Tính hiệu: $62 - 9$.
\shortans{53}
\loigiai{
\textbf{Đáp số:} $53$\\
$62 - 9 = 53$.
}
\end{ex}

\begin{ex} %[2D4V22]
Tính hiệu: $81 - 6$.
\shortans{75}
\loigiai{
\textbf{Đáp số:} $75$\\
$81 - 6 = 75$.
}
\end{ex}

\begin{ex} %[2D4V22]
Tìm $x$ biết: $x + 9 = 44$.
\shortans{35}
\loigiai{
\textbf{Đáp số:} $35$\\
$x = 44 - 9 = 35$.
}
\end{ex}

\begin{ex} %[2D4V22]
Một thùng có $50\text{ l}$ mật ong, đã lấy ra $8\text{ l}$. Hỏi trong thùng còn lại bao nhiêu lít mật ong?
\shortans{42}
\loigiai{
\textbf{Đáp số:} $42$\\
$50 - 8 = 42\text{ l}$.
}
\end{ex}

\begin{ex} %[2D4C22]
Điền số thích hợp: $72 - ? = 65$.
\shortans{7}
\loigiai{
\textbf{Đáp số:} $7$\\
$72 - 65 = 7$.
}
\end{ex}

\newpage

\TaoPhieuKiemTra{Chủ đề 4 -- Phép cộng, phép trừ (có nhớ) trong phạm vi 100}{23}{Luyện tập phép cộng, phép trừ có nhớ}{https://loptoanthaylinh.vn/thi-online/toan-2/bai-23}
\textbf{Phần 1. Trắc nghiệm lựa chọn}
\setcounterpageref{ex}{0}

\begin{ex} %[2D4N23]
Kết quả của phép tính $52 - 28$ là
\choice
{\True $24$}
{$23$}
{$25$}
{$26$}
\loigiai{
$52 - 28 = 24$.\\
\textbf{Đáp án của câu hỏi là A}
}
\end{ex}

\begin{ex} %[2D4N23]
Kết quả của phép tính $64 - 37$ là
\choice
{$26$}
{$28$}
{\True $27$}
{$29$}
\loigiai{
$64 - 37 = 27$.\\
\textbf{Đáp án của câu hỏi là C}
}
\end{ex}

\begin{ex} %[2D4H23]
Kết quả của phép tính $71 - 45$ là
\choice
{$25$}
{\True $26$}
{$27$}
{$28$}
\loigiai{
$71 - 45 = 26$.\\
\textbf{Đáp án của câu hỏi là B}
}
\end{ex}

\begin{ex} %[2D4H23]
Kết quả của phép tính $83 - 39$ là
\choice
{$43$}
{$45$}
{$46$}
{\True $44$}
\loigiai{
$83 - 39 = 44$.\\
\textbf{Đáp án của câu hỏi là D}
}
\end{ex}

\begin{ex} %[2D4H23]
Tìm $x$ biết: $x + 28 = 65$.
\choice
{$36$}
{\True $37$}
{$38$}
{$39$}
\loigiai{
$x = 65 - 28 = 37$.\\
\textbf{Đáp án của câu hỏi là B}
}
\end{ex}

\begin{ex} %[2D4H23]
Điền dấu thích hợp: $72 - 36 \dots 81 - 45$.
\choice
{$>$}
{$<$}
{Không so sánh được}
{\True $=$}
\loigiai{
$72 - 36 = 36$ và $81 - 45 = 36$. Vậy $72 - 36 = 81 - 45$.\\
\textbf{Đáp án của câu hỏi là D}
}
\end{ex}

\begin{ex} %[2D4V23]
Một cuộn dây dài $80\text{ cm}$, người ta cắt đi $35\text{ cm}$. Sợi dây còn lại dài là
\choice
{$44\text{ cm}$}
{\True $45\text{ cm}$}
{$46\text{ cm}$}
{$47\text{ cm}$}
\loigiai{
$80 - 35 = 45\text{ cm}$.\\
\textbf{Đáp án của câu hỏi là B}
}
\end{ex}

\begin{ex} %[2D4V23]
Một cửa hàng có $92\text{ kg}$ gạo, đã bán đi $48\text{ kg}$. Cửa hàng còn lại số ki-lô-gam gạo là
\choice
{$43\text{ kg}$}
{$45\text{ kg}$}
{$46\text{ kg}$}
{\True $44\text{ kg}$}
\loigiai{
$92 - 48 = 44\text{ kg}$.\\
\textbf{Đáp án của câu hỏi là D}
}
\end{ex}

\begin{ex} %[2D4V23]
Tính giá trị biểu thức: $75 - 28 - 17$.
\choice
{$28$}
{$32$}
{\True $30$}
{$34$}
\loigiai{
$47 - 17 = 30$.\\
\textbf{Đáp án của câu hỏi là C}
}
\end{ex}

\begin{ex} %[2D4V23]
Tìm $y$ biết: $90 - y = 46$.
\choice
{\True $44$}
{$43$}
{$45$}
{$46$}
\loigiai{
$y = 90 - 46 = 44$.\\
\textbf{Đáp án của câu hỏi là A}
}
\end{ex}

\begin{ex} %[2D4V23]
Bác An nuôi 63 con gà và vịt, trong đó có 27 con vịt. Bác An nuôi số con gà là
\choice
{$35$ con}
{$37$ con}
{\True $36$ con}
{$38$ con}
\loigiai{
$63 - 27 = 36$ con gà.\\
\textbf{Đáp án của câu hỏi là C}
}
\end{ex}

\begin{ex} %[2D4C23]
Hiệu của hai số là 38. Nếu bớt ở số bị trừ 5 đơn vị và giữ nguyên số trừ thì hiệu mới là
\choice
{\True $33$}
{$43$}
{$38$}
{$28$}
\loigiai{
$38 - 5 = 33$.\\
\textbf{Đáp án của câu hỏi là A}
}
\end{ex}

\textbf{Phần 2. Trắc nghiệm đúng sai}
\setcounterpageref{ex}{0}

\begin{ex} %[2D4H23]
Xét tính đúng sai của các phép trừ có nhớ:
\choiceTF
{\True $62 - 37 = 25$.}
{\True $85 - 49 = 36$.}
{$74 - 28 = 45$.}
{\True $74 - 28 = 46$.}
\loigiai{
\begin{itemize}
    \item $62 - 37 = 25$: Đúng.
    \item $85 - 49 = 36$: Đúng.
    \item $74 - 28 = 46 \neq 45$. Ý 3 Sai, Ý 4 Đúng.
\end{itemize}
\begin{itemize}
    \item a) \textbf{Đúng}.
    \item b) \textbf{Đúng}.
    \item c) \textbf{Sai}.
    \item d) \textbf{Đúng}.
\end{itemize}
}
\end{ex}

\begin{ex} %[2D4V23]
Đàn vịt có 72 con, đàn gà có 38 con. Xét tính đúng sai:
\choiceTF
{\True Số vịt nhiều hơn số gà là 34 con ($72 - 38 = 34$).}
{\True Số gà ít hơn số vịt là 34 con.}
{Tổng số gà và vịt là 100 con.}
{\True Tổng số gà và vịt là 110 con ($72 + 38 = 110$).}
\loigiai{
\begin{itemize}
    \item Hiệu: $72 - 38 = 34$: Đúng.
    \item Gà ít hơn vịt 34: Đúng.
    \item Tổng: $72 + 38 = 110 \neq 100$. Ý 3 Sai, Ý 4 Đúng.
\end{itemize}
\begin{itemize}
    \item a) \textbf{Đúng}.
    \item b) \textbf{Đúng}.
    \item c) \textbf{Sai}.
    \item d) \textbf{Đúng}.
\end{itemize}
}
\end{ex}

\textbf{Phần 3. Trả lời ngắn}
\setcounterpageref{ex}{0}

\begin{ex} %[2D4H23]
Tính hiệu: $63 - 28$.
\shortans{35}
\loigiai{
\textbf{Đáp số:} $35$\\
$63 - 28 = 35$.
}
\end{ex}

\begin{ex} %[2D4H23]
Tính hiệu: $75 - 39$.
\shortans{36}
\loigiai{
\textbf{Đáp số:} $36$\\
$75 - 39 = 36$.
}
\end{ex}

\begin{ex} %[2D4V23]
Tính hiệu: $91 - 46$.
\shortans{45}
\loigiai{
\textbf{Đáp số:} $45$\\
$91 - 46 = 45$.
}
\end{ex}

\begin{ex} %[2D4V23]
Tìm $x$ biết: $x + 36 = 82$.
\shortans{46}
\loigiai{
\textbf{Đáp số:} $46$\\
$x = 82 - 36 = 46$.
}
\end{ex}

\begin{ex} %[2D4V23]
Bao gạo nặng $50\text{ kg}$, người ta đã bán đi $28\text{ kg}$. Hỏi còn lại bao nhiêu ki-lô-gam gạo?
\shortans{22}
\loigiai{
\textbf{Đáp số:} $22$\\
$50 - 28 = 22\text{ kg}$.
}
\end{ex}

\begin{ex} %[2D4C23]
Tính nhanh: $85 - 28 - 22$.
\shortans{35}
\loigiai{
\textbf{Đáp số:} $35$\\
$85 - (28 + 22) = 85 - 50 = 35$.
}
\end{ex}

\newpage

\TaoPhieuKiemTra{Chủ đề 4 -- Phép cộng, phép trừ (có nhớ) trong phạm vi 100}{24}{Luyện tập chung}{https://loptoanthaylinh.vn/thi-online/toan-2/bai-24}
\textbf{Phần 1. Trắc nghiệm lựa chọn}
\setcounterpageref{ex}{0}

\begin{ex} %[2D4N24]
Kết quả của phép tính $54 + 28$ là
\choice
{\True $82$}
{$81$}
{$83$}
{$84$}
\loigiai{
$54 + 28 = 82$.\\
\textbf{Đáp án của câu hỏi là A}
}
\end{ex}

\begin{ex} %[2D4N24]
Kết quả của phép tính $82 - 36$ là
\choice
{$45$}
{$47$}
{$48$}
{\True $46$}
\loigiai{
$82 - 36 = 46$.\\
\textbf{Đáp án của câu hỏi là D}
}
\end{ex}

\begin{ex} %[2D4H24]
Kết quả của phép tính $67 + 25$ là
\choice
{$91$}
{\True $92$}
{$93$}
{$94$}
\loigiai{
$67 + 25 = 92$.\\
\textbf{Đáp án của câu hỏi là B}
}
\end{ex}

\begin{ex} %[2D4H24]
Kết quả của phép tính $91 - 48$ là
\choice
{$42$}
{$44$}
{$45$}
{\True $43$}
\loigiai{
$91 - 48 = 43$.\\
\textbf{Đáp án của câu hỏi là D}
}
\end{ex}

\begin{ex} %[2D4H24]
Tìm $x$ biết: $x - 29 = 47$.
\choice
{\True $76$}
{$75$}
{$77$}
{$78$}
\loigiai{
$x = 47 + 29 = 76$.\\
\textbf{Đáp án của câu hỏi là A}
}
\end{ex}

\begin{ex} %[2D4H24]
Tính giá trị biểu thức: $48 + 25 - 36$.
\choice
{$36$}
{\True $37$}
{$38$}
{$39$}
\loigiai{
$73 - 36 = 37$.\\
\textbf{Đáp án của câu hỏi là B}
}
\end{ex}

\begin{ex} %[2D4V24]
Một nông trại có $65$ con cừu, đã bán đi $28$ con cừu. Nông trại còn lại số con cừu là
\choice
{$36$ con}
{$38$ con}
{\True $37$ con}
{$39$ con}
\loigiai{
$65 - 28 = 37$ con cừu.\\
\textbf{Đáp án của câu hỏi là C}
}
\end{ex}

\begin{ex} %[2D4V24]
Đoạn đường thứ nhất dài $38\text{ m}$, đoạn đường thứ hai dài $47\text{ m}$. Cả hai đoạn đường dài tất cả là
\choice
{$84\text{ m}$}
{\True $85\text{ m}$}
{$86\text{ m}$}
{$87\text{ m}$}
\loigiai{
$38 + 47 = 85\text{ m}$.\\
\textbf{Đáp án của câu hỏi là B}
}
\end{ex}

\begin{ex} %[2D4V24]
Tìm $y$ biết: $80 - y = 35 + 17$.
\choice
{$27$}
{$29$}
{\True $28$}
{$30$}
\loigiai{
$80 - y = 52 \Rightarrow y = 80 - 52 = 28$.\\
\textbf{Đáp án của câu hỏi là C}
}
\end{ex}

\begin{ex} %[2D4V24]
Tính nhanh: $47 + 28 + 13 + 12$.
\choice
{$98$}
{$102$}
{$95$}
{\True $100$}
\loigiai{
$(47 + 13) + (28 + 12) = 60 + 40 = 100$.\\
\textbf{Đáp án của câu hỏi là D}
}
\end{ex}

\begin{ex} %[2D4V24]
Một cửa hàng có $72\text{ l}$ dầu, đã bán đi $38\text{ l}$ dầu. Cửa hàng còn lại số lít dầu là
\choice
{\True $34\text{ l}$}
{$33\text{ l}$}
{$35\text{ l}$}
{$36\text{ l}$}
\loigiai{
$72 - 38 = 34\text{ l}$.\\
\textbf{Đáp án của câu hỏi là A}
}
\end{ex}

\begin{ex} %[2D4C24]
Tổng của hai số là 85. Nếu bớt ở mỗi số 5 đơn vị thì tổng mới là
\choice
{$80$}
{$70$}
{\True $75$}
{$85$}
\loigiai{
$85 - 5 - 5 = 75$.\\
\textbf{Đáp án của câu hỏi là C}
}
\end{ex}

\textbf{Phần 2. Trắc nghiệm đúng sai}
\setcounterpageref{ex}{0}

\begin{ex} %[2D4H24]
Xét tính đúng sai của các phép toán:
\choiceTF
{\True $48 + 36 = 84$.}
{\True $84 - 48 = 36$.}
{$72 - 38 = 35$.}
{\True $72 - 38 = 34$.}
\loigiai{
\begin{itemize}
    \item $48 + 36 = 84$: Đúng.
    \item $84 - 48 = 36$: Đúng.
    \item $72 - 38 = 34 \neq 35$. Ý 3 Sai, Ý 4 Đúng.
\end{itemize}
\begin{itemize}
    \item a) \textbf{Đúng}.
    \item b) \textbf{Đúng}.
    \item c) \textbf{Sai}.
    \item d) \textbf{Đúng}.
\end{itemize}
}
\end{ex}

\begin{ex} %[2D4V24]
Mai có 52 cái kẹo, Mai ăn hết 16 cái kẹo. Xét tính đúng sai:
\choiceTF
{\True Mai còn lại 36 cái kẹo ($52 - 16 = 36$).}
{\True Số kẹo ban đầu nhiều hơn số kẹo còn lại 16 cái.}
{Nếu Mai được mẹ cho thêm 10 cái thì Mai có 45 cái.}
{\True Khi đó Mai có $36 + 10 = 46$ cái kẹo.}
\loigiai{
\begin{itemize}
    \item Còn lại: $52 - 16 = 36$: Đúng.
    \item Ban đầu hơn còn lại 16: Đúng.
    \item $36 + 10 = 46 \neq 45$. Ý 3 Sai, Ý 4 Đúng.
\end{itemize}
\begin{itemize}
    \item a) \textbf{Đúng}.
    \item b) \textbf{Đúng}.
    \item c) \textbf{Sai}.
    \item d) \textbf{Đúng}.
\end{itemize}
}
\end{ex}

\textbf{Phần 3. Trả lời ngắn}
\setcounterpageref{ex}{0}

\begin{ex} %[2D4H24]
Tính: $57 + 38$.
\shortans{95}
\loigiai{
\textbf{Đáp số:} $95$\\
$57 + 38 = 95$.
}
\end{ex}

\begin{ex} %[2D4H24]
Tính: $83 - 47$.
\shortans{36}
\loigiai{
\textbf{Đáp số:} $36$\\
$83 - 47 = 36$.
}
\end{ex}

\begin{ex} %[2D4V24]
Tính giá trị biểu thức: $46 + 28 - 35$.
\shortans{39}
\loigiai{
\textbf{Đáp số:} $39$\\
$74 - 35 = 39$.
}
\end{ex}

\begin{ex} %[2D4V24]
Tìm $x$ biết: $x + 27 = 71$.
\shortans{44}
\loigiai{
\textbf{Đáp số:} $44$\\
$x = 71 - 27 = 44$.
}
\end{ex}

\begin{ex} %[2D4V24]
Một kho có $90\text{ kg}$ ngô, đã xuất đi $46\text{ kg}$. Hỏi trong kho còn lại bao nhiêu ki-lô-gam ngô?
\shortans{44}
\loigiai{
\textbf{Đáp số:} $44$\\
$90 - 46 = 44\text{ kg}$.
}
\end{ex}

\begin{ex} %[2D4C24]
Tính nhanh: $92 - 25 - 35$.
\shortans{32}
\loigiai{
\textbf{Đáp số:} $32$\\
$92 - (25 + 35) = 92 - 60 = 32$.
}
\end{ex}

\newpage

% =========================================================================
% CHỦ ĐỀ 5 - LÀM QUEN VỚI HÌNH PHẲNG
% =========================================================================
\section{Làm quen với hình phẳng}

\TaoPhieuKiemTra{Chủ đề 5 -- Làm quen với hình phẳng}{25}{Điểm, đoạn thẳng}{https://loptoanthaylinh.vn/thi-online/toan-2/bai-25}
\textbf{Phần 1. Trắc nghiệm lựa chọn}
\setcounterpageref{ex}{0}

\begin{ex} %[2H5N25]
Người ta thường dùng các chữ cái in hoa nào để đặt tên cho điểm?
\choice
{\True $A, B, C$}
{$a, b, c$}
{$1, 2, 3$}
{$x, y, z$}
\loigiai{
Đặt tên cho điểm dùng chữ in hoa $A, B, C, \dots$\\
\textbf{Đáp án của câu hỏi là A}
}
\end{ex}

\begin{ex} %[2H5N25]
Đoạn thẳng $MN$ được giới hạn bởi mấy điểm?
\choice
{$1$ điểm}
{\True $2$ điểm ($M$ và $N$)}
{$3$ điểm}
{Không có điểm nào}
\loigiai{
Đoạn thẳng $MN$ giới hạn bởi hai đầu mút $M$ và $N$.\\
\textbf{Đáp án của câu hỏi là B}
}
\end{ex}

\begin{ex} %[2H5H25]
Ba điểm cùng nằm trên một đường thẳng được gọi là
\choice
{Ba điểm không thẳng hàng}
{Ba điểm cắt nhau}
{Ba điểm song song}
{\True Ba điểm thẳng hàng}
\loigiai{
Ba điểm cùng nằm trên 1 đường thẳng là ba điểm thẳng hàng.\\
\textbf{Đáp án của câu hỏi là D}
}
\end{ex}

\begin{ex} %[2H5H25]
Hình ảnh sợi dây điện chùng xuống cho ta hình ảnh của
\choice
{Đoạn thẳng}
{Đường thẳng}
{\True Đường cong}
{Tia số}
\loigiai{
Sợi dây chùng cho ta hình ảnh đường cong.\\
\textbf{Đáp án của câu hỏi là C}
}
\end{ex}

\begin{ex} %[2H5H25]
Đoạn thẳng $AB$ dài $8\text{ cm}$, đoạn thẳng $CD$ dài $12\text{ cm}$. Đoạn thẳng $CD$ dài hơn đoạn thẳng $AB$ là
\choice
{\True $4\text{ cm}$}
{$3\text{ cm}$}
{$5\text{ cm}$}
{$6\text{ cm}$}
\loigiai{
$12 - 8 = 4\text{ cm}$.\\
\textbf{Đáp án của câu hỏi là A}
}
\end{ex}

\begin{ex} %[2H5H25]
Cho điểm $B$ nằm giữa hai điểm $A$ và $C$ thẳng hàng. Biết $AB = 5\text{ cm}, BC = 4\text{ cm}$. Độ dài đoạn thẳng $AC$ là
\choice
{$8\text{ cm}$}
{\True $9\text{ cm}$}
{$10\text{ cm}$}
{$1\text{ cm}$}
\loigiai{
$AC = AB + BC = 5 + 4 = 9\text{ cm}$.\\
\textbf{Đáp án của câu hỏi là B}
}
\end{ex}

\begin{ex} %[2H5V25]
Qua hai điểm phân biệt $A$ và $B$, ta vẽ được bao nhiêu đường thẳng?
\choice
{$2$ đường thẳng}
{$3$ đường thẳng}
{\True Chỉ $1$ đường thẳng duy nhất}
{Vô số đường thẳng}
\loigiai{
Qua 2 điểm chỉ vẽ được duy nhất 1 đường thẳng.\\
\textbf{Đáp án của câu hỏi là C}
}
\end{ex}

\begin{ex} %[2H5V25]
Cho 3 điểm thẳng hàng $M, N, P$ theo thứ tự đó. Biết $MP = 15\text{ cm}, MN = 7\text{ cm}$. Độ dài đoạn thẳng $NP$ là
\choice
{\True $8\text{ cm}$}
{$7\text{ cm}$}
{$9\text{ cm}$}
{$22\text{ cm}$}
\loigiai{
$NP = 15 - 7 = 8\text{ cm}$.\\
\textbf{Đáp án của câu hỏi là A}
}
\end{ex}

\begin{ex} %[2H5V25]
Một đoạn thẳng $AB$ dài $14\text{ cm}$. Điểm $M$ là điểm chính giữa (ở giữa) của đoạn thẳng $AB$. Độ dài đoạn thẳng $AM$ là
\choice
{$6\text{ cm}$}
{$8\text{ cm}$}
{$9\text{ cm}$}
{\True $7\text{ cm}$}
\loigiai{
$14 : 2 = 7\text{ cm}$.\\
\textbf{Đáp án của câu hỏi là D}
}
\end{ex}

\begin{ex} %[2H5V25]
Cho 4 điểm $A, B, C, D$ cùng nằm trên một đường thẳng. Có tất cả bao nhiêu đoạn thẳng nối từ các điểm đó?
\choice
{$4$ đoạn thẳng}
{\True $6$ đoạn thẳng ($AB, AC, AD, BC, BD, CD$)}
{$5$ đoạn thẳng}
{$8$ đoạn thẳng}
\loigiai{
Các đoạn: $AB, AC, AD, BC, BD, CD$ (gồm 6 đoạn thẳng).\\
\textbf{Đáp án của câu hỏi là B}
}
\end{ex}

\begin{ex} %[2H5V25]
Đoạn thẳng $MN$ dài $20\text{ cm}$, đoạn thẳng $PQ$ ngắn hơn đoạn thẳng $MN$ là $6\text{ cm}$. Độ dài đoạn thẳng $PQ$ là
\choice
{$13\text{ cm}$}
{$15\text{ cm}$}
{\True $14\text{ cm}$}
{$26\text{ cm}$}
\loigiai{
$20 - 6 = 14\text{ cm}$.\\
\textbf{Đáp án của câu hỏi là C}
}
\end{ex}

\begin{ex} %[2H5C25]
Cho 3 điểm phân biệt không thẳng hàng $A, B, C$. Nối từng cặp điểm lại với nhau ta được hình gì?
\choice
{Đoạn thẳng dài}
{Hình tứ giác}
{Đường thẳng}
{\True Hình tam giác}
\loigiai{
Nối 3 điểm không thẳng hàng tạo thành hình tam giác.\\
\textbf{Đáp án của câu hỏi là D}
}
\end{ex}

\textbf{Phần 2. Trắc nghiệm đúng sai}
\setcounterpageref{ex}{0}

\begin{ex} %[2H5H25]
Xét tính đúng sai của các khẳng định hình học:
\choiceTF
{\True Đoạn thẳng có độ dài đo được cụ thể.}
{\True Đường thẳng kéo dài mãi về hai phía.}
{Ba điểm bất kì luôn luôn thẳng hàng.}
{\True Ba điểm chỉ thẳng hàng khi chúng cùng nằm trên một đường thẳng.}
\loigiai{
\begin{itemize}
    \item Đoạn thẳng đo được độ dài: Đúng.
    \item Đường thẳng kéo dài vô tận 2 phía: Đúng.
    \item 3 điểm bất kì có thể không thẳng hàng. Ý 3 Sai, Ý 4 Đúng.
\end{itemize}
\begin{itemize}
    \item a) \textbf{Đúng}.
    \item b) \textbf{Đúng}.
    \item c) \textbf{Sai}.
    \item d) \textbf{Đúng}.
\end{itemize}
}
\end{ex}

\begin{ex} %[2H5V25]
Cho ba điểm $A, B, C$ thẳng hàng theo thứ tự đó, $AB = 6\text{ cm}, BC = 8\text{ cm}$. Xét tính đúng sai:
\choiceTF
{\True Điểm $B$ nằm giữa hai điểm $A$ và $C$.}
{\True Đoạn thẳng $AC$ dài $14\text{ cm}$ ($6 + 8 = 14$).}
{Đoạn thẳng $AB$ dài hơn đoạn thẳng $BC$.}
{\True Đoạn thẳng $BC$ dài hơn đoạn thẳng $AB$ là $2\text{ cm}$ ($8 - 6 = 2$).}
\loigiai{
\begin{itemize}
    \item $B$ nằm giữa $A$ và $C$: Đúng.
    \item $AC = 6 + 8 = 14\text{ cm}$: Đúng.
    \item $AB (6\text{ cm}) < BC (8\text{ cm})$. Ý 3 Sai, Ý 4 Đúng.
\end{itemize}
\begin{itemize}
    \item a) \textbf{Đúng}.
    \item b) \textbf{Đúng}.
    \item c) \textbf{Sai}.
    \item d) \textbf{Đúng}.
\end{itemize}
}
\end{ex}

\textbf{Phần 3. Trả lời ngắn}
\setcounterpageref{ex}{0}

\begin{ex} %[2H5H25]
Cho đoạn thẳng $AB$ dài $12\text{ cm}$ và đoạn thẳng $CD$ dài $18\text{ cm}$. Tính tổng độ dài hai đoạn thẳng theo xăng-ti-mét.
\shortans{30}
\loigiai{
\textbf{Đáp số:} $30$\\
$12 + 18 = 30\text{ cm}$.
}
\end{ex}

\begin{ex} %[2H5H25]
Cho đoạn thẳng $MN$ dài $25\text{ cm}$, đoạn thẳng $PQ$ dài $14\text{ cm}$. Tính hiệu độ dài hai đoạn thẳng theo xăng-ti-mét.
\shortans{11}
\loigiai{
\textbf{Đáp số:} $11$\\
$25 - 14 = 11\text{ cm}$.
}
\end{ex}

\begin{ex} %[2H5V25]
Cho 3 điểm thẳng hàng $A, B, C$ với $B$ nằm giữa $A$ và $C$. Biết $AB = 8\text{ cm}, AC = 15\text{ cm}$. Tính độ dài $BC$ theo xăng-ti-mét.
\shortans{7}
\loigiai{
\textbf{Đáp số:} $7$\\
$BC = 15 - 8 = 7\text{ cm}$.
}
\end{ex}

\begin{ex} %[2H5V25]
Có 3 điểm không thẳng hàng, hỏi nối 3 điểm đó lại tạo thành bao nhiêu đoạn thẳng?
\shortans{3}
\loigiai{
\textbf{Đáp số:} $3$\\
Có 3 đoạn thẳng nối từng đôi một.
}
\end{ex}

\begin{ex} %[2H5V25]
Đoạn thẳng thứ nhất dài $9\text{ cm}$, đoạn thẳng thứ hai dài gấp đôi đoạn thứ nhất. Tính độ dài đoạn thứ hai theo xăng-ti-mét.
\shortans{18}
\loigiai{
\textbf{Đáp số:} $18$\\
$9 + 9 = 18\text{ cm}$.
}
\end{ex}

\begin{ex} %[2H5C25]
Trên một đường thẳng có 5 điểm phân biệt. Hỏi có tất cả bao nhiêu đoạn thẳng được tạo thành từ 5 điểm đó?
\shortans{10}
\loigiai{
\textbf{Đáp số:} $10$\\
Số đoạn thẳng: $4 + 3 + 2 + 1 = 10$ đoạn thẳng.
}
\end{ex}

\newpage

\TaoPhieuKiemTra{Chủ đề 5 -- Làm quen với hình phẳng}{26}{Đường thẳng, đường cong, đường gấp khúc}{https://loptoanthaylinh.vn/thi-online/toan-2/bai-26}
\textbf{Phần 1. Trắc nghiệm lựa chọn}
\setcounterpageref{ex}{0}

\begin{ex} %[2H5N26]
Đường gấp khúc $ABCD$ gồm mấy đoạn thẳng?
\choice
{$2$ đoạn thẳng}
{\True $3$ đoạn thẳng ($AB, BC, CD$)}
{$4$ đoạn thẳng}
{$5$ đoạn thẳng}
\loigiai{
Gồm 3 đoạn thẳng $AB, BC, CD$.\\
\textbf{Đáp án của câu hỏi là B}
}
\end{ex}

\begin{ex} %[2H5N26]
Hình tứ giác là hình có
\choice
{$3$ cạnh}
{$5$ cạnh}
{$6$ cạnh}
{\True $4$ cạnh}
\loigiai{
Tứ giác có 4 cạnh.\\
\textbf{Đáp án của câu hỏi là D}
}
\end{ex}

\begin{ex} %[2H5H26]
Độ dài đường gấp khúc $ABC$ có $AB = 4\text{ cm}, BC = 5\text{ cm}$ là
\choice
{$8\text{ cm}$}
{$10\text{ cm}$}
{\True $9\text{ cm}$}
{$1\text{ cm}$}
\loigiai{
$4 + 5 = 9\text{ cm}$.\\
\textbf{Đáp án của câu hỏi là C}
}
\end{ex}

\begin{ex} %[2H5H26]
Đường gấp khúc $MNPQ$ gồm các đoạn thẳng có độ dài lần lượt là $3\text{ cm}, 4\text{ cm}, 5\text{ cm}$. Độ dài đường gấp khúc đó là
\choice
{\True $12\text{ cm}$}
{$11\text{ cm}$}
{$13\text{ cm}$}
{$14\text{ cm}$}
\loigiai{
$3 + 4 + 5 = 12\text{ cm}$.\\
\textbf{Đáp án của câu hỏi là A}
}
\end{ex}

\begin{ex} %[2H5H26]
Hình chữ nhật và hình vuông đều là
\choice
{Hình tam giác}
{Đường gấp khúc}
{Hình tròn}
{\True Hình tứ giác}
\loigiai{
Hình chữ nhật, hình vuông đều là tứ giác có 4 cạnh.\\
\textbf{Đáp án của câu hỏi là D}
}
\end{ex}

\begin{ex} %[2H5H26]
Một hình tứ giác có độ dài các cạnh là $3\text{ cm}, 4\text{ cm}, 5\text{ cm}, 6\text{ cm}$. Tổng độ dài các cạnh của tứ giác là
\choice
{$17\text{ cm}$}
{\True $18\text{ cm}$}
{$19\text{ cm}$}
{$20\text{ cm}$}
\loigiai{
$3 + 4 + 5 + 6 = 18\text{ cm}$.\\
\textbf{Đáp án của câu hỏi là B}
}
\end{ex}

\begin{ex} %[2H5V26]
Đường gấp khúc $ABCD$ dài $24\text{ cm}$. Biết tổng độ dài đoạn $AB$ và $BC$ là $15\text{ cm}$. Độ dài đoạn $CD$ là
\choice
{\True $9\text{ cm}$}
{$8\text{ cm}$}
{$10\text{ cm}$}
{$11\text{ cm}$}
\loigiai{
$CD = 24 - 15 = 9\text{ cm}$.\\
\textbf{Đáp án của câu hỏi là A}
}
\end{ex}

\begin{ex} %[2H5V26]
Đường gấp khúc gồm 4 đoạn thẳng có độ dài bằng nhau, mỗi đoạn dài $5\text{ cm}$. Độ dài đường gấp khúc đó là
\choice
{$15\text{ cm}$}
{$25\text{ cm}$}
{$10\text{ cm}$}
{\True $20\text{ cm}$}
\loigiai{
$5 \times 4 = 20\text{ cm}$.\\
\textbf{Đáp án của câu hỏi là D}
}
\end{ex}

\begin{ex} %[2H5V26]
Một hình tứ giác có 4 cạnh bằng nhau, mỗi cạnh dài $6\text{ cm}$. Tổng độ dài các cạnh của hình tứ giác là
\choice
{$20\text{ cm}$}
{$28\text{ cm}$}
{\True $24\text{ cm}$}
{$18\text{ cm}$}
\loigiai{
$6 \times 4 = 24\text{ cm}$.\\
\textbf{Đáp án của câu hỏi là C}
}
\end{ex}

\begin{ex} %[2H5V26]
Một đường gấp khúc gồm 2 đoạn thẳng dài $18\text{ cm}$. Biết đoạn thứ nhất dài $8\text{ cm}$. Đoạn thứ hai dài hơn đoạn thứ nhất là
\choice
{\True $2\text{ cm}$}
{$1\text{ cm}$}
{$3\text{ cm}$}
{$4\text{ cm}$}
\loigiai{
Đoạn thứ hai: $18 - 8 = 10\text{ cm}$.\\
Đoạn 2 hơn đoạn 1: $10 - 8 = 2\text{ cm}$.\\
\textbf{Đáp án của câu hỏi là A}
}
\end{ex}

\begin{ex} %[2H5V26]
Độ dài đường gấp khúc $ABC$ là $16\text{ cm}$, biết đoạn thẳng $AB = BC$. Độ dài mỗi đoạn thẳng là
\choice
{$7\text{ cm}$}
{$9\text{ cm}$}
{\True $8\text{ cm}$}
{$6\text{ cm}$}
\loigiai{
$16 : 2 = 8\text{ cm}$.\\
\textbf{Đáp án của câu hỏi là C}
}
\end{ex}

\begin{ex} %[2H5C26]
Một đường gấp khúc khép kín gồm 4 đoạn thẳng tạo thành hình gì?
\choice
{Hình tam giác}
{\True Hình tứ giác}
{Đoạn thẳng}
{Hình tròn}
\loigiai{
Đường gấp khúc khép kín 4 đoạn tạo thành hình tứ giác.\\
\textbf{Đáp án của câu hỏi là B}
}
\end{ex}

\textbf{Phần 2. Trắc nghiệm đúng sai}
\setcounterpageref{ex}{0}

\begin{ex} %[2H5H26]
Xét tính đúng sai về đường gấp khúc và hình tứ giác:
\choiceTF
{\True Độ dài đường gấp khúc bằng tổng độ dài các đoạn thẳng tạo nên nó.}
{\True Hình tứ giác luôn có 4 cạnh.}
{Hình tam giác có 4 cạnh.}
{\True Hình tam giác có 3 cạnh, hình tứ giác có 4 cạnh.}
\loigiai{
\begin{itemize}
    \item Độ dài bằng tổng các đoạn: Đúng.
    \item Tứ giác có 4 cạnh: Đúng.
    \item Tam giác có 3 cạnh (không phải 4 cạnh). Ý 3 Sai, Ý 4 Đúng.
\end{itemize}
\begin{itemize}
    \item a) \textbf{Đúng}.
    \item b) \textbf{Đúng}.
    \item c) \textbf{Sai}.
    \item d) \textbf{Đúng}.
\end{itemize}
}
\end{ex}

\begin{ex} %[2H5V26]
Đường gấp khúc $MNP$ có $MN = 12\text{ cm}, NP = 15\text{ cm}$. Xét tính đúng sai:
\choiceTF
{\True Độ dài đường gấp khúc $MNP$ là $27\text{ cm}$ ($12 + 15 = 27$).}
{\True Đoạn $NP$ dài hơn đoạn $MN$ là $3\text{ cm}$ ($15 - 12 = 3$).}
{Đường gấp khúc $MNP$ gồm 3 đoạn thẳng.}
{\True Đường gấp khúc $MNP$ chỉ gồm 2 đoạn thẳng là $MN$ và $NP$.}
\loigiai{
\begin{itemize}
    \item Độ dài: $12 + 15 = 27\text{ cm}$: Đúng.
    \item $NP$ hơn $MN$ là $3\text{ cm}$: Đúng.
    \item Gồm 2 đoạn thẳng $MN$ và $NP$. Ý 3 Sai, Ý 4 Đúng.
\end{itemize}
\begin{itemize}
    \item a) \textbf{Đúng}.
    \item b) \textbf{Đúng}.
    \item c) \textbf{Sai}.
    \item d) \textbf{Đúng}.
\end{itemize}
}
\end{ex}

\textbf{Phần 3. Trả lời ngắn}
\setcounterpageref{ex}{0}

\begin{ex} %[2H5H26]
Tính độ dài đường gấp khúc gồm 2 đoạn thẳng có độ dài là $8\text{ cm}$ và $7\text{ cm}$ theo xăng-ti-mét.
\shortans{15}
\loigiai{
\textbf{Đáp số:} $15$\\
$8 + 7 = 15\text{ cm}$.
}
\end{ex}

\begin{ex} %[2H5H26]
Đường gấp khúc $ABC$ có $AB = 14\text{ cm}, BC = 16\text{ cm}$. Tính độ dài đường gấp khúc $ABC$ theo xăng-ti-mét.
\shortans{30}
\loigiai{
\textbf{Đáp số:} $30$\\
$14 + 16 = 30\text{ cm}$.
}
\end{ex}

\begin{ex} %[2H5V26]
Tính tổng độ dài các cạnh của một hình tứ giác có độ dài các cạnh là $5\text{ cm}, 6\text{ cm}, 7\text{ cm}, 8\text{ cm}$ theo xăng-ti-mét.
\shortans{26}
\loigiai{
\textbf{Đáp số:} $26$\\
$5 + 6 + 7 + 8 = 26\text{ cm}$.
}
\end{ex}

\begin{ex} %[2H5V26]
Đường gấp khúc $ABCD$ dài $35\text{ cm}$. Biết tổng độ dài đoạn $AB$ và $BC$ là $20\text{ cm}$. Tính độ dài đoạn $CD$ theo xăng-ti-mét.
\shortans{15}
\loigiai{
\textbf{Đáp số:} $15$\\
$CD = 35 - 20 = 15\text{ cm}$.
}
\end{ex}

\begin{ex} %[2H5V26]
Đường gấp khúc gồm 3 đoạn thẳng bằng nhau, mỗi đoạn dài $6\text{ cm}$. Tính độ dài đường gấp khúc đó theo xăng-ti-mét.
\shortans{18}
\loigiai{
\textbf{Đáp số:} $18$\\
$6 \times 3 = 18\text{ cm}$.
}
\end{ex}

\begin{ex} %[2H5C26]
Một hình tứ giác có tổng độ dài các cạnh là $40\text{ cm}$. Biết 3 cạnh có độ dài lần lượt là $8\text{ cm}, 10\text{ cm}, 12\text{ cm}$. Tính độ dài cạnh còn lại theo xăng-ti-mét.
\shortans{10}
\loigiai{
\textbf{Đáp số:} $10$\\
Cạnh còn lại: $40 - (8 + 10 + 12) = 40 - 30 = 10\text{ cm}$.
}
\end{ex}

\newpage

\TaoPhieuKiemTra{Chủ đề 5 -- Làm quen với hình phẳng}{27}{Ba điểm thẳng hàng}{https://loptoanthaylinh.vn/thi-online/toan-2/bai-27}
\textbf{Phần 1. Trắc nghiệm lựa chọn}
\setcounterpageref{ex}{0}

\begin{ex} %[2H5N27]
Để vẽ một đoạn thẳng dài $7\text{ cm}$, ta cần dùng dụng cụ nào?
\choice
{Com-pa}
{\True Thước kẻ có chia vạch xăng-ti-mét}
{Ê-ke}
{Bút chì màu}
\loigiai{
Dùng thước có chia vạch xăng-ti-mét.\\
\textbf{Đáp án của câu hỏi là B}
}
\end{ex}

\begin{ex} %[2H5N27]
Ghép 2 hình tam giác vuông bằng nhau có thể tạo thành hình nào?
\choice
{\True Hình vuông hoặc hình chữ nhật (hình tứ giác)}
{Hình tròn}
{Hình ngũ giác}
{Đoạn thẳng}
\loigiai{
Ghép 2 tam giác vuông bằng nhau tạo thành hình tứ giác.\\
\textbf{Đáp án của câu hỏi là A}
}
\end{ex}

\begin{ex} %[2H5H27]
Cắt một hình vuông theo đường chéo ta được mấy hình tam giác?
\choice
{$1$ hình}
{\True $2$ hình tam giác}
{$3$ hình tam giác}
{$4$ hình tam giác}
\loigiai{
Cắt theo đường chéo tạo thành 2 hình tam giác.\\
\textbf{Đáp án của câu hỏi là B}
}
\end{ex}

\begin{ex} %[2H5H27]
Muốn vẽ đoạn thẳng $AB$ dài $9\text{ cm}$, ta đặt thước kẻ sao cho vạch số 0 trùng với điểm $A$, điểm $B$ trùng với vạch số
\choice
{$8$}
{$10$}
{\True $9$}
{$1$}
\loigiai{
Điểm $B$ trùng với vạch số 9.\\
\textbf{Đáp án của câu hỏi là C}
}
\end{ex}

\begin{ex} %[2H5H27]
Dùng các que tính dài bằng nhau, để xếp thành 1 hình tam giác cần ít nhất bao nhiêu que tính?
\choice
{$2$ que}
{$4$ que}
{$5$ que}
{\True $3$ que}
\loigiai{
Cần ít nhất 3 que tính.\\
\textbf{Đáp án của câu hỏi là D}
}
\end{ex}

\begin{ex} %[2H5H27]
Dùng các que tính để xếp thành 1 hình vuông cần ít nhất bao nhiêu que tính có độ dài bằng nhau?
\choice
{\True $4$ que}
{$3$ que}
{$5$ que}
{$6$ que}
\loigiai{
Hình vuông có 4 cạnh bằng nhau nên cần 4 que tính.\\
\textbf{Đáp án của câu hỏi là A}
}
\end{ex}

\begin{ex} %[2H5V27]
Để xếp thành 2 hình tam giác rời nhau, cần ít nhất bao nhiêu que tính?
\choice
{$5$ que}
{$7$ que}
{$8$ que}
{\True $6$ que}
\loigiai{
Mỗi tam giác cần 3 que, 2 tam giác rời cần $3 \times 2 = 6$ que tính.\\
\textbf{Đáp án của câu hỏi là D}
}
\end{ex}

\begin{ex} %[2H5V27]
Có 5 que tính bằng nhau, xếp được 2 hình tam giác chung một cạnh. Có tất cả bao nhiêu hình tam giác trong hình vừa xếp?
\choice
{$1$ hình}
{$3$ hình tam giác}
{\True $2$ hình tam giác}
{$4$ hình tam giác}
\loigiai{
Xếp được 2 hình tam giác.\\
\textbf{Đáp án của câu hỏi là C}
}
\end{ex}

\begin{ex} %[2H5V27]
Ghép 4 hình vuông nhỏ cạnh $1\text{ cm}$ lại với nhau thành một hình vuông lớn. Cạnh của hình vuông lớn là
\choice
{$1\text{ cm}$}
{$3\text{ cm}$}
{$4\text{ cm}$}
{\True $2\text{ cm}$}
\loigiai{
Cạnh hình vuông lớn: $1 + 1 = 2\text{ cm}$.\\
\textbf{Đáp án của câu hỏi là D}
}
\end{ex}

\begin{ex} %[2H5V27]
Cắt một tờ giấy hình tứ giác theo một đường thẳng, có thể thu được nhiều nhất mấy hình tam giác?
\choice
{$1$ hình}
{\True $2$ hình tam giác (khi cắt theo đường chéo)}
{$3$ hình}
{$4$ hình}
\loigiai{
Cắt theo đường chéo thu được 2 hình tam giác.\\
\textbf{Đáp án của câu hỏi là B}
}
\end{ex}

\begin{ex} %[2H5V27]
Bạn Nam có 8 que tính bằng nhau. Nam có thể xếp được nhiều nhất bao nhiêu hình vuông rời nhau?
\choice
{\True $2$ hình vuông}
{$1$ hình}
{$3$ hình}
{$4$ hình}
\loigiai{
$8 : 4 = 2$ hình vuông.\\
\textbf{Đáp án của câu hỏi là A}
}
\end{ex}

\begin{ex} %[2H5C27]
Cần ít nhất bao nhiêu que tính để xếp thành 1 hình vuông và 1 hình tam giác chung 1 cạnh?
\choice
{$7$ que}
{$5$ que}
{\True $6$ que ($4 + 3 - 1 = 6$)}
{$8$ que}
\loigiai{
Chung 1 cạnh nên cần: $4 + 2 = 6$ que tính.\\
\textbf{Đáp án của câu hỏi là C}
}
\end{ex}

\textbf{Phần 2. Trắc nghiệm đúng sai}
\setcounterpageref{ex}{0}

\begin{ex} %[2H5H27]
Xét tính đúng sai về xếp ghép hình:
\choiceTF
{\True Ghép hai hình vuông nhỏ bằng nhau có thể tạo thành một hình chữ nhật.}
{\True Một hình vuông có thể cắt thành hai hình chữ nhật bằng nhau.}
{Không thể cắt hình vuông thành các hình tam giác.}
{\True Có thể cắt hình vuông thành 2 hoặc 4 hình tam giác bằng nhau.}
\loigiai{
\begin{itemize}
    \item Ghép 2 hình vuông được hình chữ nhật: Đúng.
    \item Cắt đôi hình vuông được 2 hình chữ nhật: Đúng.
    \item Cắt theo đường chéo ra tam giác được. Ý 3 Sai, Ý 4 Đúng.
\end{itemize}
\begin{itemize}
    \item a) \textbf{Đúng}.
    \item b) \textbf{Đúng}.
    \item c) \textbf{Sai}.
    \item d) \textbf{Đúng}.
\end{itemize}
}
\end{ex}

\begin{ex} %[2H5V27]
Có 7 que tính bằng nhau. Xét tính đúng sai:
\choiceTF
{\True Có thể dùng 7 que tính xếp được 3 hình tam giác chung cạnh.}
{\True Có thể dùng 7 que tính xếp được 1 hình vuông và 1 hình tam giác rời nhau.}
{Không thể xếp được hình tứ giác nào với 7 que tính.}
{\True Có thể xếp được hình tứ giác và thừa que tính.}
\loigiai{
\begin{itemize}
    \item Xếp 3 tam giác chung cạnh cần: $3 + 2 + 2 = 7$ que: Đúng.
    \item 1 hình vuông (4 que) + 1 tam giác (3 que) = 7 que: Đúng.
    \item Dùng 4 que xếp tứ giác vẫn được. Ý 3 Sai, Ý 4 Đúng.
\end{itemize}
\begin{itemize}
    \item a) \textbf{Đúng}.
    \item b) \textbf{Đúng}.
    \item c) \textbf{Sai}.
    \item d) \textbf{Đúng}.
\end{itemize}
}
\end{ex}

\textbf{Phần 3. Trả lời ngắn}
\setcounterpageref{ex}{0}

\begin{ex} %[2H5H27]
Cần bao nhiêu que tính để xếp thành 1 hình tam giác?
\shortans{3}
\loigiai{
\textbf{Đáp số:} $3$\\
3 que tính.
}
\end{ex}

\begin{ex} %[2H5H27]
Cần bao nhiêu que tính bằng nhau để xếp thành 1 hình vuông?
\shortans{4}
\loigiai{
\textbf{Đáp số:} $4$\\
4 que tính.
}
\end{ex}

\begin{ex} %[2H5V27]
Có 9 que tính, xếp được nhiều nhất bao nhiêu hình tam giác rời nhau?
\shortans{3}
\loigiai{
\textbf{Đáp số:} $3$\\
$9 : 3 = 3$ hình.
}
\end{ex}

\begin{ex} %[2H5V27]
Có 12 que tính, xếp được nhiều nhất bao nhiêu hình vuông rời nhau?
\shortans{3}
\loigiai{
\textbf{Đáp số:} $3$\\
$12 : 4 = 3$ hình.
}
\end{ex}

\begin{ex} %[2H5V27]
Cắt một hình chữ nhật bởi 1 nhát kéo thẳng theo đường chéo thì được bao nhiêu hình tam giác?
\shortans{2}
\loigiai{
\textbf{Đáp số:} $2$\\
Được 2 hình tam giác.
}
\end{ex}

\begin{ex} %[2H5C27]
Cần ít nhất bao nhiêu que tính để xếp thành 2 hình vuông có chung 1 cạnh?
\shortans{7}
\loigiai{
\textbf{Đáp số:} $7$\\
$4 + 3 = 7$ que tính.
}
\end{ex}

\newpage

\TaoPhieuKiemTra{Chủ đề 5 -- Làm quen với hình phẳng}{28}{Luyện tập chung}{https://loptoanthaylinh.vn/thi-online/toan-2/bai-28}
\textbf{Phần 1. Trắc nghiệm lựa chọn}
\setcounterpageref{ex}{0}

\begin{ex} %[2H5N28]
Hình nào sau đây có 4 đỉnh và 4 cạnh?
\choice
{Hình tam giác}
{Hình tròn}
{Đường thẳng}
{\True Hình tứ giác}
\loigiai{
Hình tứ giác có 4 đỉnh và 4 cạnh.\\
\textbf{Đáp án của câu hỏi là D}
}
\end{ex}

\begin{ex} %[2H5N28]
Đoạn thẳng $PQ$ dài $15\text{ cm}$, đoạn thẳng $RS$ dài $10\text{ cm}$. Cả hai đoạn thẳng dài là
\choice
{$20\text{ cm}$}
{\True $25\text{ cm}$}
{$30\text{ cm}$}
{$35\text{ cm}$}
\loigiai{
$15 + 10 = 25\text{ cm}$.\\
\textbf{Đáp án của câu hỏi là B}
}
\end{ex}

\begin{ex} %[2H5H28]
Đường gấp khúc $ABCD$ gồm các đoạn thẳng $AB = 5\text{ cm}, BC = 6\text{ cm}, CD = 7\text{ cm}$. Độ dài đường gấp khúc là
\choice
{\True $18\text{ cm}$}
{$17\text{ cm}$}
{$19\text{ cm}$}
{$20\text{ cm}$}
\loigiai{
$5 + 6 + 7 = 18\text{ cm}$.\\
\textbf{Đáp án của câu hỏi là A}
}
\end{ex}

\begin{ex} %[2H5H28]
Hình tứ giác có độ dài các cạnh lần lượt là $4\text{ cm}, 5\text{ cm}, 6\text{ cm}, 7\text{ cm}$. Tổng độ dài các cạnh là
\choice
{$21\text{ cm}$}
{\True $22\text{ cm}$}
{$23\text{ cm}$}
{$24\text{ cm}$}
\loigiai{
$4 + 5 + 6 + 7 = 22\text{ cm}$.\\
\textbf{Đáp án của câu hỏi là B}
}
\end{ex}

\begin{ex} %[2H5H28]
Cho 3 điểm $A, B, C$ thẳng hàng theo thứ tự đó. Biết $AB = 7\text{ cm}, BC = 8\text{ cm}$. Độ dài $AC$ là
\choice
{\True $15\text{ cm}$}
{$14\text{ cm}$}
{$16\text{ cm}$}
{$17\text{ cm}$}
\loigiai{
$AC = 7 + 8 = 15\text{ cm}$.\\
\textbf{Đáp án của câu hỏi là A}
}
\end{ex}

\begin{ex} %[2H5H28]
Đoạn thẳng $CD$ dài $20\text{ cm}$, điểm $I$ là điểm chính giữa của đoạn thẳng $CD$. Độ dài đoạn $CI$ là
\choice
{$8\text{ cm}$}
{$12\text{ cm}$}
{\True $10\text{ cm}$}
{$15\text{ cm}$}
\loigiai{
$20 : 2 = 10\text{ cm}$.\\
\textbf{Đáp án của câu hỏi là C}
}
\end{ex}

\begin{ex} %[2H5V28]
Đường gấp khúc gồm 4 đoạn thẳng có độ dài bằng nhau, mỗi đoạn dài $8\text{ cm}$. Độ dài đường gấp khúc đó là
\choice
{$30\text{ cm}$}
{$34\text{ cm}$}
{$36\text{ cm}$}
{\True $32\text{ cm}$}
\loigiai{
$8 \times 4 = 32\text{ cm}$.\\
\textbf{Đáp án của câu hỏi là D}
}
\end{ex}

\begin{ex} %[2H5V28]
Một hình tứ giác có 4 cạnh dài bằng nhau, tổng độ dài 4 cạnh là $28\text{ cm}$. Độ dài mỗi cạnh của tứ giác là
\choice
{$6\text{ cm}$}
{$8\text{ cm}$}
{\True $7\text{ cm}$}
{$9\text{ cm}$}
\loigiai{
$28 : 4 = 7\text{ cm}$.\\
\textbf{Đáp án của câu hỏi là C}
}
\end{ex}

\begin{ex} %[2H5V28]
Có 10 que tính, dùng để xếp được 1 hình vuông và mấy hình tam giác rời nhau?
\choice
{$1$ hình tam giác}
{$3$ hình tam giác}
{$4$ hình tam giác}
{\True $2$ hình tam giác}
\loigiai{
Hình vuông dùng 4 que, còn $10 - 4 = 6$ que tính.\\
$6 : 3 = 2$ hình tam giác.\\
\textbf{Đáp án của câu hỏi là D}
}
\end{ex}

\begin{ex} %[2H5V28]
Đường gấp khúc $ABC$ dài $25\text{ cm}$. Biết đoạn thẳng $AB$ dài hơn đoạn thẳng $BC$ là $5\text{ cm}$. Độ dài đoạn thẳng $BC$ là
\choice
{$8\text{ cm}$}
{\True $10\text{ cm}$}
{$12\text{ cm}$}
{$15\text{ cm}$}
\loigiai{
Đoạn $BC$: $(25 - 5) : 2 = 10\text{ cm}$.\\
\textbf{Đáp án của câu hỏi là B}
}
\end{ex}

\begin{ex} %[2H5V28]
Cho đoạn thẳng $AB = 18\text{ cm}$, vẽ thêm đoạn thẳng $BC = 7\text{ cm}$ thẳng hàng với $AB$. Đoạn thẳng $AC$ dài là
\choice
{\True $25\text{ cm}$}
{$24\text{ cm}$}
{$26\text{ cm}$}
{$27\text{ cm}$}
\loigiai{
$18 + 7 = 25\text{ cm}$.\\
\textbf{Đáp án của câu hỏi là A}
}
\end{ex}

\begin{ex} %[2H5C28]
Một mảnh giấy hình vuông được cắt thành 4 hình tam giác bằng nhau. Ghép 4 hình tam giác đó lại có thể tạo thành hình nào?
\choice
{Chỉ tạo thành hình tròn}
{Không thể ghép lại}
{\True Hình vuông lớn hoặc hình chữ nhật hoặc hình tam giác lớn}
{Chỉ tạo thành đường thẳng}
\loigiai{
Có thể tạo thành hình vuông, chữ nhật hoặc tam giác lớn.\\
\textbf{Đáp án của câu hỏi là C}
}
\end{ex}

\textbf{Phần 2. Trắc nghiệm đúng sai}
\setcounterpageref{ex}{0}

\begin{ex} %[2H5H28]
Xét tính đúng sai của các nhận định hình học:
\choiceTF
{\True Hình tam giác có 3 cạnh và 3 đỉnh.}
{\True Hình tứ giác có 4 cạnh và 4 đỉnh.}
{Đoạn thẳng không giới hạn ở hai đầu.}
{\True Đoạn thẳng bị giới hạn bởi 2 điểm đầu mút.}
\loigiai{
\begin{itemize}
    \item Tam giác: 3 cạnh, 3 đỉnh: Đúng.
    \item Tứ giác: 4 cạnh, 4 đỉnh: Đúng.
    \item Đoạn thẳng bị giới hạn 2 đầu. Ý 3 Sai, Ý 4 Đúng.
\end{itemize}
\begin{itemize}
    \item a) \textbf{Đúng}.
    \item b) \textbf{Đúng}.
    \item c) \textbf{Sai}.
    \item d) \textbf{Đúng}.
\end{itemize}
}
\end{ex}

\begin{ex} %[2H5V28]
Đường gấp khúc $MNPQ$ có $MN = 8\text{ cm}, NP = 10\text{ cm}, PQ = 12\text{ cm}$. Xét tính đúng sai:
\choiceTF
{\True Độ dài đường gấp khúc $MNPQ$ là $30\text{ cm}$ ($8 + 10 + 12 = 30$).}
{\True Đoạn $PQ$ dài hơn đoạn $MN$ là $4\text{ cm}$ ($12 - 8 = 4$).}
{Đường gấp khúc $MNPQ$ gồm 4 đoạn thẳng.}
{\True Đường gấp khúc $MNPQ$ gồm 3 đoạn thẳng là $MN, NP, PQ$.}
\loigiai{
\begin{itemize}
    \item Độ dài: $8 + 10 + 12 = 30\text{ cm}$: Đúng.
    \item $PQ - MN = 12 - 8 = 4\text{ cm}$: Đúng.
    \item Gồm 3 đoạn thẳng (không phải 4). Ý 3 Sai, Ý 4 Đúng.
\end{itemize}
\begin{itemize}
    \item a) \textbf{Đúng}.
    \item b) \textbf{Đúng}.
    \item c) \textbf{Sai}.
    \item d) \textbf{Đúng}.
\end{itemize}
}
\end{ex}

\textbf{Phần 3. Trả lời ngắn}
\setcounterpageref{ex}{0}

\begin{ex} %[2H5H28]
Tính độ dài đường gấp khúc gồm 3 đoạn thẳng lần lượt dài $6\text{ cm}, 8\text{ cm}, 10\text{ cm}$ theo xăng-ti-mét.
\shortans{24}
\loigiai{
\textbf{Đáp số:} $24$\\
$6 + 8 + 10 = 24\text{ cm}$.
}
\end{ex}

\begin{ex} %[2H5H28]
Một tứ giác có độ dài 4 cạnh là $5\text{ cm}, 7\text{ cm}, 9\text{ cm}, 11\text{ cm}$. Tính tổng độ dài các cạnh theo xăng-ti-mét.
\shortans{32}
\loigiai{
\textbf{Đáp số:} $32$\\
$5 + 7 + 9 + 11 = 32\text{ cm}$.
}
\end{ex}

\begin{ex} %[2H5V28]
Đường gấp khúc gồm 5 đoạn thẳng bằng nhau, mỗi đoạn dài $4\text{ cm}$. Tính độ dài đường gấp khúc theo xăng-ti-mét.
\shortans{20}
\loigiai{
\textbf{Đáp số:} $20$\\
$4 \times 5 = 20\text{ cm}$.
}
\end{ex}

\begin{ex} %[2H5V28]
Cho 3 điểm thẳng hàng $A, B, C$ với $AB = 12\text{ cm}, AC = 20\text{ cm}$ ($B$ nằm giữa $A$ và $C$). Tính độ dài $BC$ theo xăng-ti-mét.
\shortans{8}
\loigiai{
\textbf{Đáp số:} $8$\\
$BC = 20 - 12 = 8\text{ cm}$.
}
\end{ex}

\begin{ex} %[2H5V28]
Có 15 que tính, có thể xếp được nhiều nhất bao nhiêu hình tam giác rời nhau?
\shortans{5}
\loigiai{
\textbf{Đáp số:} $5$\\
$15 : 3 = 5$ hình.
}
\end{ex}

\begin{ex} %[2H5C28]
Một hình tứ giác có tổng độ dài 4 cạnh là $36\text{ cm}$, biết 4 cạnh có độ dài bằng nhau. Tính độ dài mỗi cạnh theo xăng-ti-mét.
\shortans{9}
\loigiai{
\textbf{Đáp số:} $9$\\
$36 : 4 = 9\text{ cm}$.
}
\end{ex}

\newpage

% =========================================================================
% CHỦ ĐỀ 6 - NGÀY – GIỜ, GIỜ – PHÚT, NGÀY – THÁNG
% =========================================================================
\section{Ngày -- giờ, giờ -- phút, ngày -- tháng}

\TaoPhieuKiemTra{Chủ đề 6 -- Ngày -- giờ, giờ -- phút, ngày -- tháng}{29}{Ngày - Giờ}{https://loptoanthaylinh.vn/thi-online/toan-2/bai-29}
\textbf{Phần 1. Trắc nghiệm lựa chọn}
\setcounterpageref{ex}{0}

\begin{ex} %[2H6N29]
Một ngày có bao nhiêu giờ?
\choice
{$12$ giờ}
{$20$ giờ}
{\True $24$ giờ}
{$60$ giờ}
\loigiai{
Một ngày có 24 giờ.\\
\textbf{Đáp án của câu hỏi là C}
}
\end{ex}

\begin{ex} %[2H6N29]
Một giờ có bao nhiêu phút?
\choice
{\True $60$ phút}
{$30$ phút}
{$100$ phút}
{$24$ phút}
\loigiai{
Một giờ có 60 phút.\\
\textbf{Đáp án của câu hỏi là A}
}
\end{ex}

\begin{ex} %[2H6H29]
$1$ ngày được tính từ lúc
\choice
{$6$ giờ sáng đến $6$ giờ tối}
{$7$ giờ sáng đến $12$ giờ trưa}
{$1$ giờ sáng đến $12$ giờ đêm}
{\True $12$ giờ đêm hôm trước đến $12$ giờ đêm hôm sau (từ 0 giờ đến 24 giờ)}
\loigiai{
1 ngày tính từ 12 giờ đêm hôm trước đến 12 giờ đêm hôm sau.\\
\textbf{Đáp án của câu hỏi là D}
}
\end{ex}

\begin{ex} %[2H6H29]
Lúc 8 giờ tối còn gọi là
\choice
{$18$ giờ}
{$21$ giờ}
{\True $20$ giờ}
{$22$ giờ}
\loigiai{
$8 + 12 = 20$ giờ.\\
\textbf{Đáp án của câu hỏi là C}
}
\end{ex}

\begin{ex} %[2H6H29]
Lúc 15 giờ còn gọi là
\choice
{$2$ giờ chiều}
{\True $3$ giờ chiều}
{$4$ giờ chiều}
{$5$ giờ chiều}
\loigiai{
$15 - 12 = 3$ giờ chiều.\\
\textbf{Đáp án của câu hỏi là B}
}
\end{ex}

\begin{ex} %[2H6H29]
Kim ngắn chỉ số 8, kim dài chỉ số 12. Lúc đó là
\choice
{$12$ giờ}
{$8$ giờ $30$ phút}
{\True $8$ giờ đúng}
{$7$ giờ $55$ phút}
\loigiai{
Đồng hồ chỉ 8 giờ đúng.\\
\textbf{Đáp án của câu hỏi là C}
}
\end{ex}

\begin{ex} %[2H6V29]
Kim ngắn chỉ giữa số 3 và số 4, kim dài chỉ số 6. Lúc đó là
\choice
{$3$ giờ $15$ phút}
{\True $3$ giờ $30$ phút (hoặc 3 giờ rưỡi)}
{$4$ giờ $30$ phút}
{$3$ giờ $6$ phút}
\loigiai{
Kim dài chỉ số 6 là 30 phút. Lúc đó là 3 giờ 30 phút.\\
\textbf{Đáp án của câu hỏi là B}
}
\end{ex}

\begin{ex} %[2H6V29]
Lan bắt đầu học bài lúc 7 giờ tối và học xong lúc 8 giờ 30 phút tối. Lan đã học trong thời gian là
\choice
{\True $1$ giờ $30$ phút}
{$1$ giờ}
{$2$ giờ}
{$30$ phút}
\loigiai{
Thời gian học: 8 giờ 30 phút - 7 giờ = 1 giờ 30 phút.\\
\textbf{Đáp án của câu hỏi là A}
}
\end{ex}

\begin{ex} %[2H6V29]
Buổi sáng em vào lớp lúc 7 giờ 30 phút, tan học lúc 10 giờ 30 phút. Em học ở trường trong số giờ là
\choice
{$2$ giờ}
{$4$ giờ}
{$5$ giờ}
{\True $3$ giờ}
\loigiai{
$10 - 7 = 3$ giờ.\\
\textbf{Đáp án của câu hỏi là D}
}
\end{ex}

\begin{ex} %[2H6V29]
Bố đi công tác trong 3 ngày. Bố đi công tác trong số giờ là
\choice
{$48$ giờ}
{\True $72$ giờ}
{$60$ giờ}
{$96$ giờ}
\loigiai{
$24 \times 3 = 72$ giờ.\\
\textbf{Đáp án của câu hỏi là B}
}
\end{ex}

\begin{ex} %[2H6V29]
Một tiết học kéo dài 40 phút. Nếu bắt đầu lúc 8 giờ thì tiết học kết thúc lúc
\choice
{$8$ giờ $30$ phút}
{$8$ giờ $50$ phút}
{$9$ giờ}
{\True $8$ giờ $40$ phút}
\loigiai{
Kết thúc lúc 8 giờ 40 phút.\\
\textbf{Đáp án của câu hỏi là D}
}
\end{ex}

\begin{ex} %[2H6C29]
Bây giờ là 9 giờ sáng. Hỏi sau 15 giờ nữa sẽ là mấy giờ đêm?
\choice
{\True $12$ giờ đêm (24 giờ)}
{$11$ giờ đêm}
{$1$ giờ đêm}
{$10$ giờ đêm}
\loigiai{
$9 + 15 = 24$ giờ (tức 12 giờ đêm).\\
\textbf{Đáp án của câu hỏi là A}
}
\end{ex}

\textbf{Phần 2. Trắc nghiệm đúng sai}
\setcounterpageref{ex}{0}

\begin{ex} %[2H6H29]
Xét tính đúng sai về thời gian:
\choiceTF
{\True Một ngày có 24 giờ.}
{\True Một giờ có 60 phút.}
{14 giờ là 4 giờ chiều.}
{\True 14 giờ là 2 giờ chiều ($14 - 12 = 2$).}
\loigiai{
\begin{itemize}
    \item 1 ngày = 24 giờ: Đúng.
    \item 1 giờ = 60 phút: Đúng.
    \item 14 giờ = 2 giờ chiều (không phải 4 giờ). Ý 3 Sai, Ý 4 Đúng.
\end{itemize}
\begin{itemize}
    \item a) \textbf{Đúng}.
    \item b) \textbf{Đúng}.
    \item c) \textbf{Sai}.
    \item d) \textbf{Đúng}.
\end{itemize}
}
\end{ex}

\begin{ex} %[2H6V29]
Mai xem phim từ 19 giờ đến 20 giờ 30 phút. Xét tính đúng sai:
\choiceTF
{\True 19 giờ chính là 7 giờ tối.}
{\True Bộ phim kéo dài 1 giờ 30 phút.}
{20 giờ 30 phút chính là 9 giờ 30 phút tối.}
{\True 20 giờ 30 phút chính là 8 giờ 30 phút tối ($20 - 12 = 8$).}
\loigiai{
\begin{itemize}
    \item 19 giờ = 7 giờ tối: Đúng.
    \item Thời gian: 1 giờ 30 phút: Đúng.
    \item 20 giờ = 8 giờ tối (không phải 9 giờ). Ý 3 Sai, Ý 4 Đúng.
\end{itemize}
\begin{itemize}
    \item a) \textbf{Đúng}.
    \item b) \textbf{Đúng}.
    \item c) \textbf{Sai}.
    \item d) \textbf{Đúng}.
\end{itemize}
}
\end{ex}

\textbf{Phần 3. Trả lời ngắn}
\setcounterpageref{ex}{0}

\begin{ex} %[2H6H29]
Một ngày có bao nhiêu giờ?
\shortans{24}
\loigiai{
\textbf{Đáp số:} $24$\\
24 giờ.
}
\end{ex}

\begin{ex} %[2H6H29]
Một giờ có bao nhiêu phút?
\shortans{60}
\loigiai{
\textbf{Đáp số:} $60$\\
60 phút.
}
\end{ex}

\begin{ex} %[2H6V29]
Lúc 17 giờ chính là mấy giờ chiều?
\shortans{5}
\loigiai{
\textbf{Đáp số:} $5$\\
$17 - 12 = 5$ giờ chiều.
}
\end{ex}

\begin{ex} %[2H6V29]
Lúc 9 giờ tối tương ứng với bao nhiêu giờ trong ngày?
\shortans{21}
\loigiai{
\textbf{Đáp số:} $21$\\
$9 + 12 = 21$ giờ.
}
\end{ex}

\begin{ex} %[2H6V29]
Nam tập thể dục từ 6 giờ sáng đến 6 giờ 30 phút sáng. Nam đã tập thể dục bao nhiêu phút?
\shortans{30}
\loigiai{
\textbf{Đáp số:} $30$\\
30 phút.
}
\end{ex}

\begin{ex} %[2H6C29]
2 ngày có tất cả bao nhiêu giờ?
\shortans{48}
\loigiai{
\textbf{Đáp số:} $48$\\
$24 + 24 = 48$ giờ.
}
\end{ex}

\newpage

\TaoPhieuKiemTra{Chủ đề 6 -- Ngày -- giờ, giờ -- phút, ngày -- tháng}{30}{Giờ - Phút}{https://loptoanthaylinh.vn/thi-online/toan-2/bai-30}
\textbf{Phần 1. Trắc nghiệm lựa chọn}
\setcounterpageref{ex}{0}

\begin{ex} %[2H6N30]
Một tuần lễ có bao nhiêu ngày?
\choice
{\True $7$ ngày}
{$5$ ngày}
{$6$ ngày}
{$8$ ngày}
\loigiai{
Một tuần lễ có 7 ngày.\\
\textbf{Đáp án của câu hỏi là A}
}
\end{ex}

\begin{ex} %[2H6N30]
Một năm có bao nhiêu tháng?
\choice
{$10$ tháng}
{$11$ tháng}
{\True $12$ tháng}
{$24$ tháng}
\loigiai{
Một năm có 12 tháng.\\
\textbf{Đáp án của câu hỏi là C}
}
\end{ex}

\begin{ex} %[2H6H30]
Tháng 1 có bao nhiêu ngày?
\choice
{\True $31$ ngày}
{$28$ ngày}
{$30$ ngày}
{$29$ ngày}
\loigiai{
Tháng 1 có 31 ngày.\\
\textbf{Đáp án của câu hỏi là A}
}
\end{ex}

\begin{ex} %[2H6H30]
Tháng 4 có bao nhiêu ngày?
\choice
{$31$ ngày}
{$29$ ngày}
{$28$ ngày}
{\True $30$ ngày}
\loigiai{
Tháng 4 có 30 ngày.\\
\textbf{Đáp án của câu hỏi là D}
}
\end{ex}

\begin{ex} %[2H6H30]
Các tháng nào sau đây đều có 31 ngày?
\choice
{Tháng 4, tháng 6, tháng 9}
{Tháng 2, tháng 4, tháng 6}
{\True Tháng 1, tháng 3, tháng 5, tháng 7, tháng 8, tháng 10, tháng 12}
{Tháng 6, tháng 9, tháng 11}
\loigiai{
Các tháng có 31 ngày: 1, 3, 5, 7, 8, 10, 12.\\
\textbf{Đáp án của câu hỏi là C}
}
\end{ex}

\begin{ex} %[2H6H30]
Ngày Quốc tế Thiếu nhi là ngày nào trong năm?
\choice
{Ngày 8 tháng 3}
{\True Ngày 1 tháng 6}
{Ngày 20 tháng 11}
{Ngày 2 tháng 9}
\loigiai{
Quốc tế Thiếu nhi là ngày 1 tháng 6.\\
\textbf{Đáp án của câu hỏi là B}
}
\end{ex}

\begin{ex} %[2H6V30]
Nếu hôm nay là thứ Ba ngày 5 thì thứ Ba tuần sau là ngày mấy?
\choice
{$11$}
{$13$}
{$10$}
{\True $12$}
\loigiai{
$5 + 7 = 12$.\\
\textbf{Đáp án của câu hỏi là D}
}
\end{ex}

\begin{ex} %[2H6V30]
Nếu ngày 10 là thứ Năm thì ngày 17 của tháng đó là thứ mấy?
\choice
{Thứ Tư}
{\True Thứ Năm}
{Thứ Sáu}
{Thứ Bảy}
\loigiai{
$10 + 7 = 17$ nên cũng là thứ Năm.\\
\textbf{Đáp án của câu hỏi là B}
}
\end{ex}

\begin{ex} %[2H6V30]
Nếu hôm nay là thứ Sáu ngày 15 thì thứ Sáu tuần trước là ngày mấy?
\choice
{$7$}
{$9$}
{\True $8$}
{$10$}
\loigiai{
$15 - 7 = 8$.\\
\textbf{Đáp án của câu hỏi là C}
}
\end{ex}

\begin{ex} %[2H6V30]
Tháng 12 có bao nhiêu ngày?
\choice
{$30$ ngày}
{\True $31$ ngày}
{$28$ ngày}
{$29$ ngày}
\loigiai{
Tháng 12 có 31 ngày.\\
\textbf{Đáp án của câu hỏi là B}
}
\end{ex}

\begin{ex} %[2H6V30]
Ngày Nhà giáo Việt Nam là ngày 20 tháng 11. Tháng 11 có bao nhiêu ngày?
\choice
{\True $30$ ngày}
{$31$ ngày}
{$28$ ngày}
{$29$ ngày}
\loigiai{
Tháng 11 có 30 ngày.\\
\textbf{Đáp án của câu hỏi là A}
}
\end{ex}

\begin{ex} %[2H6C30]
Trong một tháng có 3 ngày Chủ nhật đều là ngày chẵn. Hỏi ngày Chủ nhật đầu tiên của tháng đó là ngày mùng mấy?
\choice
{Ngày mùng 4}
{Ngày mùng 6}
{Ngày mùng 8}
{\True Ngày mùng 2}
\loigiai{
Các ngày Chủ nhật: 2, 9, 16, 23, 30. Có 3 ngày chẵn là 2, 16, 30. Vậy ngày đầu tiên là mùng 2.\\
\textbf{Đáp án của câu hỏi là D}
}
\end{ex}

\textbf{Phần 2. Trắc nghiệm đúng sai}
\setcounterpageref{ex}{0}

\begin{ex} %[2H6H30]
Xét tính đúng sai về lịch và các tháng:
\choiceTF
{\True Một tuần lễ có 7 ngày.}
{\True Tháng 3 có 31 ngày.}
{Tháng 4 có 31 ngày.}
{\True Tháng 4 có 30 ngày.}
\loigiai{
\begin{itemize}
    \item 1 tuần = 7 ngày: Đúng.
    \item Tháng 3 có 31 ngày: Đúng.
    \item Tháng 4 có 30 ngày (không phải 31). Ý 3 Sai, Ý 4 Đúng.
\end{itemize}
\begin{itemize}
    \item a) \textbf{Đúng}.
    \item b) \textbf{Đúng}.
    \item c) \textbf{Sai}.
    \item d) \textbf{Đúng}.
\end{itemize}
}
\end{ex}

\begin{ex} %[2H6V30]
Nếu ngày 2 tháng 9 là thứ Năm. Xét tính đúng sai:
\choiceTF
{\True Ngày 9 tháng 9 là thứ Năm ($2 + 7 = 9$).}
{\True Ngày 3 tháng 9 là thứ Sáu.}
{Ngày 1 tháng 9 là thứ Sáu.}
{\True Ngày 1 tháng 9 là thứ Tư.}
\loigiai{
\begin{itemize}
    \item 9 tháng 9 là thứ Năm: Đúng.
    \item 3 tháng 9 là thứ Sáu: Đúng.
    \item 1 tháng 9 là thứ Tư (không phải thứ Sáu). Ý 3 Sai, Ý 4 Đúng.
\end{itemize}
\begin{itemize}
    \item a) \textbf{Đúng}.
    \item b) \textbf{Đúng}.
    \item c) \textbf{Sai}.
    \item d) \textbf{Đúng}.
\end{itemize}
}
\end{ex}

\textbf{Phần 3. Trả lời ngắn}
\setcounterpageref{ex}{0}

\begin{ex} %[2H6H30]
Một tuần lễ có bao nhiêu ngày?
\shortans{7}
\loigiai{
\textbf{Đáp số:} $7$\\
7 ngày.
}
\end{ex}

\begin{ex} %[2H6H30]
Tháng 5 có bao nhiêu ngày?
\shortans{31}
\loigiai{
\textbf{Đáp số:} $31$\\
31 ngày.
}
\end{ex}

\begin{ex} %[2H6V30]
Tháng 6 có bao nhiêu ngày?
\shortans{30}
\loigiai{
\textbf{Đáp số:} $30$\\
30 ngày.
}
\end{ex}

\begin{ex} %[2H6V30]
Nếu Chủ nhật là ngày 7 thì Chủ nhật tuần sau là ngày bao nhiêu?
\shortans{14}
\loigiai{
\textbf{Đáp số:} $14$\\
$7 + 7 = 14$.
}
\end{ex}

\begin{ex} %[2H6V30]
Nếu thứ Bảy là ngày 20 thì thứ Bảy tuần trước là ngày bao nhiêu?
\shortans{13}
\loigiai{
\textbf{Đáp số:} $13$\\
$20 - 7 = 13$.
}
\end{ex}

\begin{ex} %[2H6C30]
Một năm có bao nhiêu tháng có 31 ngày? (Gồm các tháng: 1, 3, 5, 7, 8, 10, 12)
\shortans{7}
\loigiai{
\textbf{Đáp số:} $7$\\
Có 7 tháng.
}
\end{ex}

\newpage

\TaoPhieuKiemTra{Chủ đề 6 -- Ngày -- giờ, giờ -- phút, ngày -- tháng}{31}{Ngày - Tháng}{https://loptoanthaylinh.vn/thi-online/toan-2/bai-31}
\textbf{Phần 1. Trắc nghiệm lựa chọn}
\setcounterpageref{ex}{0}

\begin{ex} %[2H6N31]
Kim dài của đồng hồ quay 1 vòng (từ số 12 đến số 12) là được bao nhiêu phút?
\choice
{$30$ phút}
{$12$ phút}
{$24$ phút}
{\True $60$ phút (1 giờ)}
\loigiai{
Kim dài quay 1 vòng là 60 phút.\\
\textbf{Đáp án của câu hỏi là D}
}
\end{ex}

\begin{ex} %[2H6N31]
Khi kim ngắn chỉ số 6, kim dài chỉ số 12 thì đồng hồ chỉ mấy giờ?
\choice
{$12$ giờ}
{\True $6$ giờ đúng}
{$6$ giờ $30$ phút}
{$12$ giờ $30$ phút}
\loigiai{
Lúc đó là 6 giờ đúng.\\
\textbf{Đáp án của câu hỏi là B}
}
\end{ex}

\begin{ex} %[2H6H31]
Đồng hồ điện tử hiển thị $16:30$. Lúc đó là
\choice
{$6$ giờ $30$ phút chiều}
{$3$ giờ $30$ phút chiều}
{\True $4$ giờ $30$ phút chiều (hoặc 4 giờ rưỡi chiều)}
{$5$ giờ $30$ phút chiều}
\loigiai{
$16 - 12 = 4$ giờ 30 phút chiều.\\
\textbf{Đáp án của câu hỏi là C}
}
\end{ex}

\begin{ex} %[2H6H31]
Hôm nay là thứ Năm ngày 10. Hôm qua là
\choice
{\True Thứ Tư ngày 9}
{Thứ Ba ngày 9}
{Thứ Sáu ngày 11}
{Thứ Tư ngày 8}
\loigiai{
Hôm qua là thứ Tư ngày 9.\\
\textbf{Đáp án của câu hỏi là A}
}
\end{ex}

\begin{ex} %[2H6H31]
Hôm nay là thứ Bảy ngày 20. Ngày mai là
\choice
{Chủ nhật ngày 19}
{\True Chủ nhật ngày 21}
{Thứ Sáu ngày 19}
{Thứ Hai ngày 22}
\loigiai{
Ngày mai là Chủ nhật ngày 21.\\
\textbf{Đáp án của câu hỏi là B}
}
\end{ex}

\begin{ex} %[2H6H31]
Em đi ngủ lúc 21 giờ 30 phút. Lúc đó còn gọi là
\choice
{$8$ giờ $30$ phút tối}
{$10$ giờ $30$ phút tối}
{\True $9$ giờ $30$ phút tối}
{$9$ giờ tối}
\loigiai{
$21 - 12 = 9$ giờ 30 phút tối.\\
\textbf{Đáp án của câu hỏi là C}
}
\end{ex}

\begin{ex} %[2H6V31]
Bạn An vẽ tranh từ 8 giờ đến 8 giờ 30 phút. An đã vẽ tranh trong
\choice
{\True $30$ phút (nửa giờ)}
{$15$ phút}
{$45$ phút}
{$1$ giờ}
\loigiai{
Thời gian vẽ: 30 phút.\\
\textbf{Đáp án của câu hỏi là A}
}
\end{ex}

\begin{ex} %[2H6V31]
Bố đi làm lúc 7 giờ sáng và về nhà lúc 17 giờ. Bố đã ở nơi làm việc trong số giờ là
\choice
{$8$ giờ}
{$12$ giờ}
{$9$ giờ}
{\True $10$ giờ}
\loigiai{
$17 - 7 = 10$ giờ.\\
\textbf{Đáp án của câu hỏi là D}
}
\end{ex}

\begin{ex} %[2H6V31]
Tháng 10 có 31 ngày. Ngày cuối cùng của tháng 10 là ngày
\choice
{$30$}
{\True $31$}
{$28$}
{$29$}
\loigiai{
Ngày 31 tháng 10.\\
\textbf{Đáp án của câu hỏi là B}
}
\end{ex}

\begin{ex} %[2H6V31]
Một trận bóng đá bắt đầu lúc 15 giờ và kết thúc sau 90 phút (1 giờ 30 phút). Trận bóng kết thúc lúc
\choice
{\True $16$ giờ $30$ phút}
{$16$ giờ}
{$17$ giờ}
{$17$ giờ $30$ phút}
\loigiai{
$15$ giờ $+ 1$ giờ $30$ phút = $16$ giờ $30$ phút.\\
\textbf{Đáp án của câu hỏi là A}
}
\end{ex}

\begin{ex} %[2H6V31]
Gia đình Lan đi du lịch từ ngày 12 đến hết ngày 15 của tháng. Chuyến du lịch kéo dài bao nhiêu ngày?
\choice
{$3$ ngày}
{$5$ ngày}
{$6$ ngày}
{\True $4$ ngày (ngày 12, 13, 14, 15)}
\loigiai{
$15 - 12 + 1 = 4$ ngày.\\
\textbf{Đáp án của câu hỏi là D}
}
\end{ex}

\begin{ex} %[2H6C31]
Một đồng hồ quả lắc đánh chuông theo số giờ (1 giờ điểm 1 tiếng, 2 giờ điểm 2 tiếng...). Từ lúc 1 giờ đến 4 giờ chiều, đồng hồ điểm tất cả bao nhiêu tiếng chuông?
\choice
{$8$ tiếng}
{$9$ tiếng}
{\True $10$ tiếng ($1 + 2 + 3 + 4 = 10$)}
{$12$ tiếng}
\loigiai{
$1 + 2 + 3 + 4 = 10$ tiếng chuông.\\
\textbf{Đáp án của câu hỏi là C}
}
\end{ex}

\textbf{Phần 2. Trắc nghiệm đúng sai}
\setcounterpageref{ex}{0}

\begin{ex} %[2H6H31]
Xét tính đúng sai khi xem đồng hồ:
\choiceTF
{\True Kim ngắn chỉ giờ, kim dài chỉ phút.}
{\True Khi kim dài chỉ số 6 là 30 phút.}
{Khi kim dài chỉ số 12 là 12 phút.}
{\True Khi kim dài chỉ số 12 là 00 phút (giờ đúng).}
\loigiai{
\begin{itemize}
    \item Kim ngắn chỉ giờ, dài chỉ phút: Đúng.
    \item Chỉ số 6 là 30 phút: Đúng.
    \item Chỉ số 12 là 0 phút (không phải 12 phút). Ý 3 Sai, Ý 4 Đúng.
\end{itemize}
\begin{itemize}
    \item a) \textbf{Đúng}.
    \item b) \textbf{Đúng}.
    \item c) \textbf{Sai}.
    \item d) \textbf{Đúng}.
\end{itemize}
}
\end{ex}

\begin{ex} %[2H6V31]
Hôm nay là thứ Tư ngày 14. Xét tính đúng sai:
\choiceTF
{\True Hôm qua là thứ Ba ngày 13.}
{\True Ngày mai là thứ Năm ngày 15.}
{Thứ Tư tuần trước là ngày 8.}
{\True Thứ Tư tuần trước là ngày 7 ($14 - 7 = 7$).}
\loigiai{
\begin{itemize}
    \item Hôm qua: Thứ Ba ngày 13: Đúng.
    \item Ngày mai: Thứ Năm ngày 15: Đúng.
    \item $14 - 7 = 7 \neq 8$. Ý 3 Sai, Ý 4 Đúng.
\end{itemize}
\begin{itemize}
    \item a) \textbf{Đúng}.
    \item b) \textbf{Đúng}.
    \item c) \textbf{Sai}.
    \item d) \textbf{Đúng}.
\end{itemize}
}
\end{ex}

\textbf{Phần 3. Trả lời ngắn}
\setcounterpageref{ex}{0}

\begin{ex} %[2H6H31]
Khi kim ngắn chỉ số 7, kim dài chỉ số 12 thì đồng hồ chỉ mấy giờ đúng?
\shortans{7}
\loigiai{
\textbf{Đáp số:} $7$\\
7 giờ.
}
\end{ex}

\begin{ex} %[2H6H31]
Khi kim dài chỉ số 6 thì tương ứng với bao nhiêu phút?
\shortans{30}
\loigiai{
\textbf{Đáp số:} $30$\\
30 phút.
}
\end{ex}

\begin{ex} %[2H6V31]
Lúc 14 giờ 30 phút còn gọi là mấy giờ rưỡi chiều?
\shortans{2}
\loigiai{
\textbf{Đáp số:} $2$\\
$14 - 12 = 2$ giờ rưỡi chiều.
}
\end{ex}

\begin{ex} %[2H6V31]
Nếu thứ Sáu là ngày 8 thì thứ Sáu tuần sau là ngày bao nhiêu?
\shortans{15}
\loigiai{
\textbf{Đáp số:} $15$\\
$8 + 7 = 15$.
}
\end{ex}

\begin{ex} %[2H6V31]
Một tiết mục biểu diễn văn nghệ kéo dài từ 8 giờ đến 8 giờ 30 phút. Hỏi tiết mục kéo dài bao nhiêu phút?
\shortans{30}
\loigiai{
\textbf{Đáp số:} $30$\\
30 phút.
}
\end{ex}

\begin{ex} %[2H6C31]
Một tuần lễ em đi học 5 ngày, nghỉ 2 ngày cuối tuần. Hỏi trong 4 tuần lễ em được nghỉ bao nhiêu ngày cuối tuần?
\shortans{8}
\loigiai{
\textbf{Đáp số:} $8$\\
$2 \times 4 = 8$ ngày.
}
\end{ex}

\newpage

\TaoPhieuKiemTra{Chủ đề 6 -- Ngày -- giờ, giờ -- phút, ngày -- tháng}{32}{Luyện tập chung thời gian}{https://loptoanthaylinh.vn/thi-online/toan-2/bai-32}
\textbf{Phần 1. Trắc nghiệm lựa chọn}
\setcounterpageref{ex}{0}

\begin{ex} %[2H6N32]
Một tuần lễ gồm có mấy ngày?
\choice
{\True $7$ ngày}
{$6$ ngày}
{$8$ ngày}
{$5$ ngày}
\loigiai{
Một tuần có 7 ngày.\\
\textbf{Đáp án của câu hỏi là A}
}
\end{ex}

\begin{ex} %[2H6N32]
Tháng 8 có bao nhiêu ngày?
\choice
{$30$ ngày}
{\True $31$ ngày}
{$28$ ngày}
{$29$ ngày}
\loigiai{
Tháng 8 có 31 ngày.\\
\textbf{Đáp án của câu hỏi là B}
}
\end{ex}

\begin{ex} %[2H6H32]
Đồng hồ chỉ 20 giờ nghĩa là
\choice
{$7$ giờ tối}
{$9$ giờ tối}
{\True $8$ giờ tối}
{$10$ giờ tối}
\loigiai{
$20 - 12 = 8$ giờ tối.\\
\textbf{Đáp án của câu hỏi là C}
}
\end{ex}

\begin{ex} %[2H6H32]
Tháng 9 có bao nhiêu ngày?
\choice
{$31$ ngày}
{$29$ ngày}
{$28$ ngày}
{\True $30$ ngày}
\loigiai{
Tháng 9 có 30 ngày.\\
\textbf{Đáp án của câu hỏi là D}
}
\end{ex}

\begin{ex} %[2H6H32]
Nếu ngày 3 là thứ Hai thì ngày 10 của tháng đó là
\choice
{\True Thứ Hai}
{Chủ nhật}
{Thứ Ba}
{Thứ Tư}
\loigiai{
$3 + 7 = 10$ nên cũng là thứ Hai.\\
\textbf{Đáp án của câu hỏi là A}
}
\end{ex}

\begin{ex} %[2H6H32]
Kim dài chỉ số 12, kim ngắn chỉ số 4. Đồng hồ chỉ
\choice
{$12$ giờ}
{$4$ giờ $30$ phút}
{\True $4$ giờ đúng}
{$12$ giờ $20$ phút}
\loigiai{
Đồng hồ chỉ 4 giờ đúng.\\
\textbf{Đáp án của câu hỏi là C}
}
\end{ex}

\begin{ex} %[2H6V32]
Một ngày có 24 giờ. Hỏi 2 ngày có tất cả bao nhiêu giờ?
\choice
{$40$ giờ}
{$50$ giờ}
{$72$ giờ}
{\True $48$ giờ}
\loigiai{
$24 + 24 = 48$ giờ.\\
\textbf{Đáp án của câu hỏi là D}
}
\end{ex}

\begin{ex} %[2H6V32]
Nam thức dậy lúc 6 giờ sáng và bắt đầu học bài lúc 7 giờ 30 phút sáng. Khoảng thời gian từ lúc thức dậy đến lúc học là
\choice
{\True $1$ giờ $30$ phút}
{$1$ giờ}
{$2$ giờ}
{$30$ phút}
\loigiai{
$7$ giờ $30$ phút - $6$ giờ = $1$ giờ $30$ phút.\\
\textbf{Đáp án của câu hỏi là A}
}
\end{ex}

\begin{ex} %[2H6V32]
Nếu hôm nay là thứ Tư ngày 20 thì thứ Tư tuần trước là ngày
\choice
{$12$}
{\True $13$}
{$14$}
{$15$}
\loigiai{
$20 - 7 = 13$.\\
\textbf{Đáp án của câu hỏi là B}
}
\end{ex}

\begin{ex} %[2H6V32]
Tháng 7 và tháng 8 có tất cả bao nhiêu ngày?
\choice
{$60$ ngày}
{$61$ ngày}
{\True $62$ ngày}
{$59$ ngày}
\loigiai{
$31 + 31 = 62$ ngày.\\
\textbf{Đáp án của câu hỏi là C}
}
\end{ex}

\begin{ex} %[2H6V32]
Hà đi ngủ lúc 21 giờ và thức dậy lúc 6 giờ sáng hôm sau. Hà đã ngủ bao nhiêu giờ?
\choice
{$8$ giờ}
{\True $9$ giờ}
{$10$ giờ}
{$7$ giờ}
\loigiai{
Từ 21 giờ đến 24 giờ là 3 giờ. Thêm 6 giờ sáng nữa là $3 + 6 = 9$ giờ.\\
\textbf{Đáp án của câu hỏi là B}
}
\end{ex}

\begin{ex} %[2H6C32]
Trong một tháng có 5 ngày thứ Bảy. Nếu ngày thứ Bảy đầu tiên là ngày mùng 1 thì ngày thứ Bảy cuối cùng của tháng đó là ngày bao nhiêu?
\choice
{$28$}
{$30$}
{$31$}
{\True $29$}
\loigiai{
Các thứ Bảy: 1, 8, 15, 22, 29. Ngày cuối cùng là ngày 29.\\
\textbf{Đáp án của câu hỏi là D}
}
\end{ex}

\textbf{Phần 2. Trắc nghiệm đúng sai}
\setcounterpageref{ex}{0}

\begin{ex} %[2H6H32]
Xét tính đúng sai về thời gian:
\choiceTF
{\True 22 giờ chính là 10 giờ đêm ($22 - 12 = 10$).}
{\True Tháng 7 có 31 ngày.}
{Tháng 6 có 31 ngày.}
{\True Tháng 6 chỉ có 30 ngày.}
\loigiai{
\begin{itemize}
    \item 22 giờ = 10 giờ đêm: Đúng.
    \item Tháng 7 có 31 ngày: Đúng.
    \item Tháng 6 có 30 ngày (không phải 31). Ý 3 Sai, Ý 4 Đúng.
\end{itemize}
\begin{itemize}
    \item a) \textbf{Đúng}.
    \item b) \textbf{Đúng}.
    \item c) \textbf{Sai}.
    \item d) \textbf{Đúng}.
\end{itemize}
}
\end{ex}

\begin{ex} %[2H6V32]
Hôm nay là Chủ nhật ngày 16. Xét tính đúng sai:
\choiceTF
{\True Chủ nhật tuần sau là ngày 23 ($16 + 7 = 23$).}
{\True Chủ nhật tuần trước là ngày 9 ($16 - 7 = 9$).}
{Thứ Hai ngày mai là ngày 18.}
{\True Thứ Hai ngày mai là ngày 17 ($16 + 1 = 17$).}
\loigiai{
\begin{itemize}
    \item Tuần sau: $16 + 7 = 23$: Đúng.
    \item Tuần trước: $16 - 7 = 9$: Đúng.
    \item $16 + 1 = 17 \neq 18$. Ý 3 Sai, Ý 4 Đúng.
\end{itemize}
\begin{itemize}
    \item a) \textbf{Đúng}.
    \item b) \textbf{Đúng}.
    \item c) \textbf{Sai}.
    \item d) \textbf{Đúng}.
\end{itemize}
}
\end{ex}

\textbf{Phần 3. Trả lời ngắn}
\setcounterpageref{ex}{0}

\begin{ex} %[2H6H32]
Một tuần lễ có bao nhiêu ngày?
\shortans{7}
\loigiai{
\textbf{Đáp số:} $7$\\
7 ngày.
}
\end{ex}

\begin{ex} %[2H6H32]
Tháng 10 có bao nhiêu ngày?
\shortans{31}
\loigiai{
\textbf{Đáp số:} $31$\\
31 ngày.
}
\end{ex}

\begin{ex} %[2H6V32]
Nếu thứ Ba là ngày 11 thì thứ Ba tuần sau là ngày bao nhiêu?
\shortans{18}
\loigiai{
\textbf{Đáp số:} $18$\\
$11 + 7 = 18$.
}
\end{ex}

\begin{ex} %[2H6V32]
Lúc 16 giờ chiều tương ứng với mấy giờ chiều?
\shortans{4}
\loigiai{
\textbf{Đáp số:} $4$\\
$16 - 12 = 4$ giờ chiều.
}
\end{ex}

\begin{ex} %[2H6V32]
Một buổi học kéo dài từ 7 giờ đến 11 giờ trưa. Hỏi buổi học kéo dài bao nhiêu giờ?
\shortans{4}
\loigiai{
\textbf{Đáp số:} $4$\\
$11 - 7 = 4$ giờ.
}
\end{ex}

\begin{ex} %[2H6C32]
Tổng số ngày của tháng 4 và tháng 5 là bao nhiêu ngày?
\shortans{61}
\loigiai{
\textbf{Đáp số:} $61$\\
$30 + 31 = 61$ ngày.
}
\end{ex}

\newpage

% =========================================================================
% CHỦ ĐỀ 7 - ÔN TẬP HỌC KÌ 1
% =========================================================================
\section{Ôn tập học kì 1}

\TaoPhieuKiemTra{Chủ đề 7 -- Ôn tập học kì 1}{33}{Ôn tập các số trong phạm vi 100}{https://loptoanthaylinh.vn/thi-online/toan-2/bai-33}
\textbf{Phần 1. Trắc nghiệm lựa chọn}
\setcounterpageref{ex}{0}

\begin{ex} %[2D7N33]
Kết quả của phép tính $9 + 8$ là
\choice
{\True $17$}
{$16$}
{$18$}
{$19$}
\loigiai{
$9 + 8 = 17$.\\
\textbf{Đáp án của câu hỏi là A}
}
\end{ex}

\begin{ex} %[2D7N33]
Kết quả của phép tính $16 - 9$ là
\choice
{$6$}
{\True $7$}
{$8$}
{$9$}
\loigiai{
$16 - 9 = 7$.\\
\textbf{Đáp án của câu hỏi là B}
}
\end{ex}

\begin{ex} %[2D7H33]
Kết quả của phép tính $47 + 38$ là
\choice
{\True $85$}
{$84$}
{$86$}
{$75$}
\loigiai{
$47 + 38 = 85$.\\
\textbf{Đáp án của câu hỏi là A}
}
\end{ex}

\begin{ex} %[2D7H33]
Kết quả của phép tính $83 - 46$ là
\choice
{$36$}
{$38$}
{\True $37$}
{$47$}
\loigiai{
$83 - 46 = 37$.\\
\textbf{Đáp án của câu hỏi là C}
}
\end{ex}

\begin{ex} %[2D7H33]
Tìm $x$ biết: $x - 35 = 28$.
\choice
{$62$}
{$64$}
{$53$}
{\True $63$}
\loigiai{
$x = 28 + 35 = 63$.\\
\textbf{Đáp án của câu hỏi là D}
}
\end{ex}

\begin{ex} %[2D7H33]
Điền dấu thích hợp: $58 + 26 \dots 90 - 6$.
\choice
{\True $=$}
{$>$}
{$<$}
{Không so sánh được}
\loigiai{
$58 + 26 = 84$ và $90 - 6 = 84$. Vậy $58 + 26 = 90 - 6$.\\
\textbf{Đáp án của câu hỏi là A}
}
\end{ex}

\begin{ex} %[2D7V33]
Một cuộn dây dài $95\text{ m}$, người ta cắt đi $38\text{ m}$. Cuộn dây còn lại dài là
\choice
{$56\text{ m}$}
{\True $57\text{ m}$}
{$58\text{ m}$}
{$59\text{ m}$}
\loigiai{
$95 - 38 = 57\text{ m}$.\\
\textbf{Đáp án của câu hỏi là B}
}
\end{ex}

\begin{ex} %[2D7V33]
Tính giá trị biểu thức: $46 + 28 - 37$.
\choice
{$36$}
{$38$}
{\True $37$}
{$39$}
\loigiai{
$74 - 37 = 37$.\\
\textbf{Đáp án của câu hỏi là C}
}
\end{ex}

\begin{ex} %[2D7V33]
Lớp 2A có 18 học sinh giỏi, lớp 2B có 19 học sinh giỏi. Cả hai lớp có tất cả số học sinh giỏi là
\choice
{$36$ bạn}
{\True $37$ bạn}
{$38$ bạn}
{$39$ bạn}
\loigiai{
$18 + 19 = 37$ bạn.\\
\textbf{Đáp án của câu hỏi là B}
}
\end{ex}

\begin{ex} %[2D7V33]
Tìm $y$ biết: $72 - y = 25 + 19$.
\choice
{$27$}
{$29$}
{$30$}
{\True $28$}
\loigiai{
$72 - y = 44 \Rightarrow y = 72 - 44 = 28$.\\
\textbf{Đáp án của câu hỏi là D}
}
\end{ex}

\begin{ex} %[2D7V33]
Tính nhanh: $38 + 27 + 12 + 23$.
\choice
{$98$}
{$102$}
{\True $100$}
{$95$}
\loigiai{
$(38 + 12) + (27 + 23) = 50 + 50 = 100$.\\
\textbf{Đáp án của câu hỏi là C}
}
\end{ex}

\begin{ex} %[2D7C33]
Hiệu của hai số là 45. Nếu giữ nguyên số trừ và tăng số bị trừ thêm 8 đơn vị thì hiệu mới là
\choice
{$37$}
{$45$}
{$50$}
{\True $53$}
\loigiai{
$45 + 8 = 53$.\\
\textbf{Đáp án của câu hỏi là D}
}
\end{ex}

\textbf{Phần 2. Trắc nghiệm đúng sai}
\setcounterpageref{ex}{0}

\begin{ex} %[2D7H33]
Xét tính đúng sai của các phép tính:
\choiceTF
{\True $67 + 28 = 95$.}
{\True $92 - 48 = 44$.}
{$15 - 8 = 6$.}
{\True $15 - 8 = 7$.}
\loigiai{
\begin{itemize}
    \item $67 + 28 = 95$: Đúng.
    \item $92 - 48 = 44$: Đúng.
    \item $15 - 8 = 7 \neq 6$. Ý 3 Sai, Ý 4 Đúng.
\end{itemize}
\begin{itemize}
    \item a) \textbf{Đúng}.
    \item b) \textbf{Đúng}.
    \item c) \textbf{Sai}.
    \item d) \textbf{Đúng}.
\end{itemize}
}
\end{ex}

\begin{ex} %[2D7V33]
Bác nông dân có 64 con vịt, bác bán đi 28 con vịt. Xét tính đúng sai:
\choiceTF
{\True Bác còn lại 36 con vịt ($64 - 28 = 36$).}
{\True Số vịt còn lại nhiều hơn số vịt đã bán là 8 con ($36 - 28 = 8$).}
{Nếu bác mua thêm 10 con thì có 40 con vịt.}
{\True Khi đó bác có $36 + 10 = 46$ con vịt.}
\loigiai{
\begin{itemize}
    \item Còn lại: $64 - 28 = 36$: Đúng.
    \item Còn lại hơn đã bán: $36 - 28 = 8$: Đúng.
    \item $36 + 10 = 46 \neq 40$. Ý 3 Sai, Ý 4 Đúng.
\end{itemize}
\begin{itemize}
    \item a) \textbf{Đúng}.
    \item b) \textbf{Đúng}.
    \item c) \textbf{Sai}.
    \item d) \textbf{Đúng}.
\end{itemize}
}
\end{ex}

\textbf{Phần 3. Trả lời ngắn}
\setcounterpageref{ex}{0}

\begin{ex} %[2D7H33]
Tính: $56 + 37$.
\shortans{93}
\loigiai{
\textbf{Đáp số:} $93$\\
$56 + 37 = 93$.
}
\end{ex}

\begin{ex} %[2D7H33]
Tính: $81 - 35$.
\shortans{46}
\loigiai{
\textbf{Đáp số:} $46$\\
$81 - 35 = 46$.
}
\end{ex}

\begin{ex} %[2D7V33]
Tính giá trị biểu thức: $18 + 29 + 13$.
\shortans{60}
\loigiai{
\textbf{Đáp số:} $60$\\
$47 + 13 = 60$.
}
\end{ex}

\begin{ex} %[2D7V33]
Tìm $x$ biết: $x + 39 = 75$.
\shortans{36}
\loigiai{
\textbf{Đáp số:} $36$\\
$x = 75 - 39 = 36$.
}
\end{ex}

\begin{ex} %[2D7V33]
Thùng thứ nhất có $45\text{ l}$ dầu, thùng thứ hai có ít hơn thùng thứ nhất $18\text{ l}$ dầu. Hỏi thùng thứ hai có bao nhiêu lít dầu?
\shortans{27}
\loigiai{
\textbf{Đáp số:} $27$\\
$45 - 18 = 27\text{ l}$.
}
\end{ex}

\begin{ex} %[2D7C33]
Tính nhanh: $96 - 38 - 18$.
\shortans{40}
\loigiai{
\textbf{Đáp số:} $40$\\
$96 - (38 + 18) = 96 - 56 = 40$.
}
\end{ex}

\newpage

\TaoPhieuKiemTra{Chủ đề 7 -- Ôn tập học kì 1}{34}{Ôn tập phép cộng, phép trừ trong phạm vi 100}{https://loptoanthaylinh.vn/thi-online/toan-2/bai-34}
\textbf{Phần 1. Trắc nghiệm lựa chọn}
\setcounterpageref{ex}{0}

\begin{ex} %[2H7N34]
Hình tứ giác có bao nhiêu cạnh?
\choice
{\True $4$ cạnh}
{$3$ cạnh}
{$5$ cạnh}
{$6$ cạnh}
\loigiai{
Hình tứ giác có 4 cạnh.\\
\textbf{Đáp án của câu hỏi là A}
}
\end{ex}

\begin{ex} %[2H7N34]
Đoạn thẳng $AB$ dài $14\text{ cm}$, đoạn thẳng $CD$ dài $16\text{ cm}$. Tổng độ dài 2 đoạn thẳng là
\choice
{$28\text{ cm}$}
{\True $30\text{ cm}$}
{$32\text{ cm}$}
{$20\text{ cm}$}
\loigiai{
$14 + 16 = 30\text{ cm}$.\\
\textbf{Đáp án của câu hỏi là B}
}
\end{ex}

\begin{ex} %[2H7H34]
Đường gấp khúc $ABCD$ gồm 3 đoạn thẳng có độ dài lần lượt là $4\text{ cm}, 6\text{ cm}, 8\text{ cm}$. Độ dài đường gấp khúc là
\choice
{$16\text{ cm}$}
{$20\text{ cm}$}
{$22\text{ cm}$}
{\True $18\text{ cm}$}
\loigiai{
$4 + 6 + 8 = 18\text{ cm}$.\\
\textbf{Đáp án của câu hỏi là D}
}
\end{ex}

\begin{ex} %[2H7H34]
Cho 3 điểm thẳng hàng $M, N, P$ với $N$ nằm giữa $M$ và $P$. Biết $MN = 9\text{ cm}, MP = 17\text{ cm}$. Độ dài $NP$ là
\choice
{$7\text{ cm}$}
{$9\text{ cm}$}
{\True $8\text{ cm}$}
{$26\text{ cm}$}
\loigiai{
$NP = 17 - 9 = 8\text{ cm}$.\\
\textbf{Đáp án của câu hỏi là C}
}
\end{ex}

\begin{ex} %[2H7H34]
Một hình tứ giác có độ dài các cạnh lần lượt là $3\text{ cm}, 5\text{ cm}, 7\text{ cm}, 9\text{ cm}$. Tổng độ dài 4 cạnh là
\choice
{$22\text{ cm}$}
{$26\text{ cm}$}
{$28\text{ cm}$}
{\True $24\text{ cm}$}
\loigiai{
$3 + 5 + 7 + 9 = 24\text{ cm}$.\\
\textbf{Đáp án của câu hỏi là D}
}
\end{ex}

\begin{ex} %[2H7H34]
Cắt một hình vuông theo đường chéo tạo thành
\choice
{2 hình vuông nhỏ}
{2 hình chữ nhật}
{\True 2 hình tam giác}
{2 hình tròn}
\loigiai{
Cắt theo đường chéo tạo thành 2 hình tam giác.\\
\textbf{Đáp án của câu hỏi là C}
}
\end{ex}

\begin{ex} %[2H7V34]
Đường gấp khúc gồm 4 đoạn thẳng bằng nhau, mỗi đoạn dài $6\text{ cm}$. Độ dài đường gấp khúc là
\choice
{\True $24\text{ cm}$}
{$20\text{ cm}$}
{$28\text{ cm}$}
{$18\text{ cm}$}
\loigiai{
$6 \times 4 = 24\text{ cm}$.\\
\textbf{Đáp án của câu hỏi là A}
}
\end{ex}

\begin{ex} %[2H7V34]
Có 12 que tính bằng nhau, xếp được nhiều nhất bao nhiêu hình tam giác rời nhau?
\choice
{$3$ hình}
{\True $4$ hình}
{$5$ hình}
{$6$ hình}
\loigiai{
$12 : 3 = 4$ hình tam giác.\\
\textbf{Đáp án của câu hỏi là B}
}
\end{ex}

\begin{ex} %[2H7V34]
Một đoạn dây đồng dài $36\text{ cm}$ được uốn thành một hình vuông. Độ dài mỗi cạnh hình vuông là
\choice
{$8\text{ cm}$}
{$10\text{ cm}$}
{\True $9\text{ cm}$}
{$12\text{ cm}$}
\loigiai{
$36 : 4 = 9\text{ cm}$.\\
\textbf{Đáp án của câu hỏi là C}
}
\end{ex}

\begin{ex} %[2H7V34]
Đường gấp khúc $ABC$ dài $28\text{ cm}$. Biết đoạn $AB = 15\text{ cm}$. Đoạn $BC$ ngắn hơn đoạn $AB$ là
\choice
{$1\text{ cm}$}
{\True $2\text{ cm}$}
{$3\text{ cm}$}
{$4\text{ cm}$}
\loigiai{
$BC = 28 - 15 = 13\text{ cm}$.\\
$AB - BC = 15 - 13 = 2\text{ cm}$.\\
\textbf{Đáp án của câu hỏi là B}
}
\end{ex}

\begin{ex} %[2H7V34]
Cho 4 điểm phân biệt, trong đó không có 3 điểm nào thẳng hàng. Nối từng cặp điểm ta được bao nhiêu đoạn thẳng?
\choice
{\True $6$ đoạn thẳng}
{$4$ đoạn thẳng}
{$5$ đoạn thẳng}
{$8$ đoạn thẳng}
\loigiai{
Có 6 đoạn thẳng.\\
\textbf{Đáp án của câu hỏi là A}
}
\end{ex}

\begin{ex} %[2H7C34]
Một hình tứ giác có tổng độ dài các cạnh là $32\text{ cm}$. Biết 3 cạnh có độ dài bằng nhau và bằng $7\text{ cm}$. Độ dài cạnh thứ tư là
\choice
{$10\text{ cm}$}
{$12\text{ cm}$}
{$9\text{ cm}$}
{\True $11\text{ cm}$}
\loigiai{
Tổng 3 cạnh: $7 \times 3 = 21\text{ cm}$.\\
Cạnh 4: $32 - 21 = 11\text{ cm}$.\\
\textbf{Đáp án của câu hỏi là D}
}
\end{ex}

\textbf{Phần 2. Trắc nghiệm đúng sai}
\setcounterpageref{ex}{0}

\begin{ex} %[2H7H34]
Xét tính đúng sai về hình học phẳng:
\choiceTF
{\True Hình vuông là một hình tứ giác đặc biệt có 4 cạnh bằng nhau.}
{\True Hình tam giác luôn có 3 đỉnh và 3 cạnh.}
{Đoạn thẳng là đường thẳng kéo dài mãi mãi.}
{\True Đoạn thẳng có điểm đầu và điểm cuối xác định.}
\loigiai{
\begin{itemize}
    \item Hình vuông là tứ giác: Đúng.
    \item Tam giác: 3 đỉnh, 3 cạnh: Đúng.
    \item Đoạn thẳng có giới hạn. Ý 3 Sai, Ý 4 Đúng.
\end{itemize}
\begin{itemize}
    \item a) \textbf{Đúng}.
    \item b) \textbf{Đúng}.
    \item c) \textbf{Sai}.
    \item d) \textbf{Đúng}.
\end{itemize}
}
\end{ex}

\begin{ex} %[2H7V34]
Cho đường gấp khúc $ABCD$ có $AB = 6\text{ cm}, BC = 8\text{ cm}, CD = 10\text{ cm}$. Xét tính đúng sai:
\choiceTF
{\True Độ dài đường gấp khúc $ABCD$ là $24\text{ cm}$ ($6 + 8 + 10 = 24$).}
{\True Đoạn $CD$ dài hơn đoạn $AB$ là $4\text{ cm}$ ($10 - 6 = 4$).}
{Đường gấp khúc $ABCD$ gồm 4 đoạn thẳng.}
{\True Đường gấp khúc $ABCD$ chỉ gồm 3 đoạn thẳng là $AB, BC, CD$.}
\loigiai{
\begin{itemize}
    \item Độ dài: $24\text{ cm}$: Đúng.
    \item $CD - AB = 4\text{ cm}$: Đúng.
    \item Gồm 3 đoạn thẳng. Ý 3 Sai, Ý 4 Đúng.
\end{itemize}
\begin{itemize}
    \item a) \textbf{Đúng}.
    \item b) \textbf{Đúng}.
    \item c) \textbf{Sai}.
    \item d) \textbf{Đúng}.
\end{itemize}
}
\end{ex}

\textbf{Phần 3. Trả lời ngắn}
\setcounterpageref{ex}{0}

\begin{ex} %[2H7H34]
Tính độ dài đường gấp khúc gồm 2 đoạn thẳng dài $15\text{ cm}$ và $18\text{ cm}$ theo xăng-ti-mét.
\shortans{33}
\loigiai{
\textbf{Đáp số:} $33$\\
$15 + 18 = 33\text{ cm}$.
}
\end{ex}

\begin{ex} %[2H7H34]
Một tứ giác có độ dài các cạnh là $6\text{ cm}, 7\text{ cm}, 8\text{ cm}, 9\text{ cm}$. Tính tổng độ dài 4 cạnh theo xăng-ti-mét.
\shortans{30}
\loigiai{
\textbf{Đáp số:} $30$\\
$6 + 7 + 8 + 9 = 30\text{ cm}$.
}
\end{ex}

\begin{ex} %[2H7V34]
Đoạn thẳng $AB = 25\text{ cm}$, điểm $M$ nằm trên $AB$ sao cho $AM = 14\text{ cm}$. Tính độ dài $MB$ theo xăng-ti-mét.
\shortans{11}
\loigiai{
\textbf{Đáp số:} $11$\\
$MB = 25 - 14 = 11\text{ cm}$.
}
\end{ex}

\begin{ex} %[2H7V34]
Có 16 que tính bằng nhau, xếp được nhiều nhất bao nhiêu hình vuông rời nhau?
\shortans{4}
\loigiai{
\textbf{Đáp số:} $4$\\
$16 : 4 = 4$ hình vuông.
}
\end{ex}

\begin{ex} %[2H7V34]
Đường gấp khúc gồm 3 đoạn thẳng có độ dài bằng nhau, tổng độ dài là $27\text{ cm}$. Tính độ dài mỗi đoạn thẳng theo xăng-ti-mét.
\shortans{9}
\loigiai{
\textbf{Đáp số:} $9$\\
$27 : 3 = 9\text{ cm}$.
}
\end{ex}

\begin{ex} %[2H7C34]
Một hình tứ giác có tổng độ dài các cạnh là $45\text{ cm}$. Biết tổng độ dài 3 cạnh đầu là $32\text{ cm}$. Tính độ dài cạnh thứ tư theo xăng-ti-mét.
\shortans{13}
\loigiai{
\textbf{Đáp số:} $13$\\
$45 - 32 = 13\text{ cm}$.
}
\end{ex}

\newpage

\TaoPhieuKiemTra{Chủ đề 7 -- Ôn tập học kì 1}{35}{Ôn tập hình học và đo lường}{https://loptoanthaylinh.vn/thi-online/toan-2/bai-35}
\textbf{Phần 1. Trắc nghiệm lựa chọn}
\setcounterpageref{ex}{0}

\begin{ex} %[2H7N35]
Đơn vị dùng để đo khối lượng là
\choice
{Xăng-ti-mét}
{Lít}
{\True Ki-lô-gam}
{Giờ}
\loigiai{
Đo khối lượng dùng ki-lô-gam.\\
\textbf{Đáp án của câu hỏi là C}
}
\end{ex}

\begin{ex} %[2H7N35]
Đơn vị dùng để đo dung tích (sức chứa chất lỏng) là
\choice
{Ki-lô-gam}
{\True Lít}
{Mét}
{Ngày}
\loigiai{
Đo dung tích dùng lít.\\
\textbf{Đáp án của câu hỏi là B}
}
\end{ex}

\begin{ex} %[2H7H35]
Tính: $38\text{ kg} + 27\text{ kg} = \dots$
\choice
{$64\text{ kg}$}
{$66\text{ kg}$}
{$67\text{ kg}$}
{\True $65\text{ kg}$}
\loigiai{
$38 + 27 = 65\text{ kg}$.\\
\textbf{Đáp án của câu hỏi là D}
}
\end{ex}

\begin{ex} %[2H7H35]
Tính: $72\text{ l} - 38\text{ l} = \dots$
\choice
{\True $34\text{ l}$}
{$33\text{ l}$}
{$35\text{ l}$}
{$36\text{ l}$}
\loigiai{
$72 - 38 = 34\text{ l}$.\\
\textbf{Đáp án của câu hỏi là A}
}
\end{ex}

\begin{ex} %[2H7H35]
Lúc 19 giờ 30 phút tối còn gọi là
\choice
{$6$ giờ $30$ phút tối}
{\True $7$ giờ $30$ phút tối (hoặc 7 giờ rưỡi tối)}
{$8$ giờ $30$ phút tối}
{$9$ giờ $30$ phút tối}
\loigiai{
$19 - 12 = 7$ giờ 30 phút tối.\\
\textbf{Đáp án của câu hỏi là B}
}
\end{ex}

\begin{ex} %[2H7H35]
Tháng 7 có bao nhiêu ngày?
\choice
{$30$ ngày}
{$28$ ngày}
{$29$ ngày}
{\True $31$ ngày}
\loigiai{
Tháng 7 có 31 ngày.\\
\textbf{Đáp án của câu hỏi là D}
}
\end{ex}

\begin{ex} %[2H7V35]
Bao gạo to nặng $45\text{ kg}$, bao gạo nhỏ nhẹ hơn bao gạo to $18\text{ kg}$. Bao gạo nhỏ nặng số ki-lô-gam là
\choice
{\True $27\text{ kg}$}
{$26\text{ kg}$}
{$28\text{ kg}$}
{$29\text{ kg}$}
\loigiai{
$45 - 18 = 27\text{ kg}$.\\
\textbf{Đáp án của câu hỏi là A}
}
\end{ex}

\begin{ex} %[2H7V35]
Một can chứa $25\text{ l}$ mật ong, đã bán đi $9\text{ l}$. Can còn lại số lít mật ong là
\choice
{$15\text{ l}$}
{$17\text{ l}$}
{\True $16\text{ l}$}
{$18\text{ l}$}
\loigiai{
$25 - 9 = 16\text{ l}$.\\
\textbf{Đáp án của câu hỏi là C}
}
\end{ex}

\begin{ex} %[2H7V35]
Nếu hôm nay là thứ Sáu ngày 9 thì thứ Sáu tuần sau là ngày mấy?
\choice
{$15$}
{\True $16$}
{$17$}
{$18$}
\loigiai{
$9 + 7 = 16$.\\
\textbf{Đáp án của câu hỏi là B}
}
\end{ex}

\begin{ex} %[2H7V35]
Tính giá trị biểu thức: $45\text{ kg} + 35\text{ kg} - 28\text{ kg}$.
\choice
{\True $52\text{ kg}$}
{$50\text{ kg}$}
{$54\text{ kg}$}
{$56\text{ kg}$}
\loigiai{
$80 - 28 = 52\text{ kg}$.\\
\textbf{Đáp án của câu hỏi là A}
}
\end{ex}

\begin{ex} %[2H7V35]
Một bể cá chứa $60\text{ l}$ nước, sau khi xả bớt $25\text{ l}$ thì trong bể còn lại số lít nước là
\choice
{$34\text{ l}$}
{$36\text{ l}$}
{$37\text{ l}$}
{\True $35\text{ l}$}
\loigiai{
$60 - 25 = 35\text{ l}$.\\
\textbf{Đáp án của câu hỏi là D}
}
\end{ex}

\begin{ex} %[2H7C35]
Thùng thứ nhất có $28\text{ l}$ dầu. Nếu đổ thêm $6\text{ l}$ vào thùng thứ nhất thì thùng thứ nhất vẫn ít hơn thùng thứ hai là $4\text{ l}$ dầu. Hỏi thùng thứ hai có bao nhiêu lít dầu?
\choice
{$36\text{ l}$}
{$40\text{ l}$}
{\True $38\text{ l}$}
{$34\text{ l}$}
\loigiai{
Thùng 1 sau khi thêm: $28 + 6 = 34\text{ l}$.\\
Thùng 2 có: $34 + 4 = 38\text{ l}$.\\
\textbf{Đáp án của câu hỏi là C}
}
\end{ex}

\textbf{Phần 2. Trắc nghiệm đúng sai}
\setcounterpageref{ex}{0}

\begin{ex} %[2H7H35]
Xét tính đúng sai về các đơn vị đo đại lượng:
\choiceTF
{\True $25\text{ kg} + 35\text{ kg} = 60\text{ kg}$.}
{\True $40\text{ l} - 15\text{ l} = 25\text{ l}$.}
{1 ngày có 12 giờ.}
{\True 1 ngày có 24 giờ.}
\loigiai{
\begin{itemize}
    \item $25 + 35 = 60\text{ kg}$: Đúng.
    \item $40 - 15 = 25\text{ l}$: Đúng.
    \item 1 ngày = 24 giờ (không phải 12 giờ). Ý 3 Sai, Ý 4 Đúng.
\end{itemize}
\begin{itemize}
    \item a) \textbf{Đúng}.
    \item b) \textbf{Đúng}.
    \item c) \textbf{Sai}.
    \item d) \textbf{Đúng}.
\end{itemize}
}
\end{ex}

\begin{ex} %[2H7V35]
Một can dầu chứa $18\text{ l}$, mẹ mua thêm $7\text{ l}$ dầu nữa. Xét tính đúng sai:
\choiceTF
{\True Tổng lượng dầu mẹ có là $25\text{ l}$ ($18 + 7 = 25$).}
{\True Nếu dùng hết $10\text{ l}$ dầu thì còn lại $15\text{ l}$ dầu ($25 - 10 = 15$).}
{Lượng dầu còn lại ít hơn $10\text{ l}$.}
{\True Lượng dầu còn lại là $15\text{ l}$ (nhiều hơn $10\text{ l}$).}
\loigiai{
\begin{itemize}
    \item Tổng: $18 + 7 = 25\text{ l}$: Đúng.
    \item Còn lại: $25 - 10 = 15\text{ l}$: Đúng.
    \item $15 > 10$ (không ít hơn). Ý 3 Sai, Ý 4 Đúng.
\end{itemize}
\begin{itemize}
    \item a) \textbf{Đúng}.
    \item b) \textbf{Đúng}.
    \item c) \textbf{Sai}.
    \item d) \textbf{Đúng}.
\end{itemize}
}
\end{ex}

\textbf{Phần 3. Trả lời ngắn}
\setcounterpageref{ex}{0}

\begin{ex} %[2H7H35]
Tính: $37\text{ kg} + 45\text{ kg}$ theo ki-lô-gam.
\shortans{82}
\loigiai{
\textbf{Đáp số:} $82$\\
$37 + 45 = 82\text{ kg}$.
}
\end{ex}

\begin{ex} %[2H7H35]
Tính: $65\text{ l} - 28\text{ l}$ theo lít.
\shortans{37}
\loigiai{
\textbf{Đáp số:} $37$\\
$65 - 28 = 37\text{ l}$.
}
\end{ex}

\begin{ex} %[2H7V35]
Bao gạo to nặng $50\text{ kg}$, bao gạo nhỏ nặng $25\text{ kg}$. Bao gạo to nặng hơn bao nhỏ bao nhiêu ki-lô-gam?
\shortans{25}
\loigiai{
\textbf{Đáp số:} $25$\\
$50 - 25 = 25\text{ kg}$.
}
\end{ex}

\begin{ex} %[2H7V35]
Một thùng có $45\text{ l}$ nước mắm, sau khi bán còn lại $19\text{ l}$. Hỏi đã bán bao nhiêu lít nước mắm?
\shortans{26}
\loigiai{
\textbf{Đáp số:} $26$\\
$45 - 19 = 26\text{ l}$.
}
\end{ex}

\begin{ex} %[2H7V35]
Tháng 12 có bao nhiêu ngày?
\shortans{31}
\loigiai{
\textbf{Đáp số:} $31$\\
31 ngày.
}
\end{ex}

\begin{ex} %[2H7C35]
Một cửa hàng buổi sáng bán được $35\text{ kg}$ gạo, buổi chiều bán được nhiều hơn buổi sáng $18\text{ kg}$ gạo. Hỏi cả ngày cửa hàng bán được bao nhiêu ki-lô-gam gạo?
\shortans{88}
\loigiai{
\textbf{Đáp số:} $88$\\
Buổi chiều: $35 + 18 = 53\text{ kg}$.\\
Cả ngày: $35 + 53 = 88\text{ kg}$.
}
\end{ex}

\newpage

\TaoPhieuKiemTra{Chủ đề 7 -- Ôn tập học kì 1}{36}{Ôn tập chung học kì 1}{https://loptoanthaylinh.vn/thi-online/toan-2/bai-36}
\textbf{Phần 1. Trắc nghiệm lựa chọn}
\setcounterpageref{ex}{0}

\begin{ex} %[2D7N36]
Số gồm $9$ chục và $4$ đơn vị viết là
\choice
{$49$}
{$904$}
{$90$}
{\True $94$}
\loigiai{
Số đó là $94$.\\
\textbf{Đáp án của câu hỏi là D}
}
\end{ex}

\begin{ex} %[2D7N36]
Số liền sau của số $99$ là
\choice
{$98$}
{\True $100$}
{$101$}
{$90$}
\loigiai{
$99 + 1 = 100$.\\
\textbf{Đáp án của câu hỏi là B}
}
\end{ex}

\begin{ex} %[2D7H36]
Kết quả của phép tính $56 + 37$ là
\choice
{$92$}
{$94$}
{\True $93$}
{$83$}
\loigiai{
$56 + 37 = 93$.\\
\textbf{Đáp án của câu hỏi là C}
}
\end{ex}

\begin{ex} %[2D7H36]
Kết quả của phép tính $91 - 47$ là
\choice
{\True $44$}
{$43$}
{$45$}
{$54$}
\loigiai{
$91 - 47 = 44$.\\
\textbf{Đáp án của câu hỏi là A}
}
\end{ex}

\begin{ex} %[2D7H36]
Đường gấp khúc gồm 3 đoạn thẳng có độ dài lần lượt là $8\text{ cm}, 10\text{ cm}, 12\text{ cm}$. Độ dài đường gấp khúc là
\choice
{$28\text{ cm}$}
{\True $30\text{ cm}$}
{$32\text{ cm}$}
{$34\text{ cm}$}
\loigiai{
$8 + 10 + 12 = 30\text{ cm}$.\\
\textbf{Đáp án của câu hỏi là B}
}
\end{ex}

\begin{ex} %[2D7H36]
Một can đựng $15\text{ l}$ dầu, một can khác đựng $20\text{ l}$ dầu. Cả hai can đựng số lít dầu là
\choice
{$30\text{ l}$}
{$40\text{ l}$}
{$45\text{ l}$}
{\True $35\text{ l}$}
\loigiai{
$15 + 20 = 35\text{ l}$.\\
\textbf{Đáp án của câu hỏi là D}
}
\end{ex}

\begin{ex} %[2D7V36]
Lớp 2A trồng được 38 cây, lớp 2B trồng được nhiều hơn lớp 2A là 15 cây. Lớp 2B trồng được số cây là
\choice
{$52$ cây}
{$54$ cây}
{\True $53$ cây}
{$55$ cây}
\loigiai{
$38 + 15 = 53$ cây.\\
\textbf{Đáp án của câu hỏi là C}
}
\end{ex}

\begin{ex} %[2D7V36]
Tìm $x$ biết: $x + 28 = 74$.
\choice
{\True $46$}
{$45$}
{$47$}
{$48$}
\loigiai{
$x = 74 - 28 = 46$.\\
\textbf{Đáp án của câu hỏi là A}
}
\end{ex}

\begin{ex} %[2D7V36]
Tính giá trị biểu thức: $75 - 28 + 19$.
\choice
{$64$}
{$68$}
{$70$}
{\True $66$}
\loigiai{
$47 + 19 = 66$.\\
\textbf{Đáp án của câu hỏi là D}
}
\end{ex}

\begin{ex} %[2D7V36]
Nếu hôm nay là thứ Năm ngày 12 thì thứ Năm tuần trước là ngày
\choice
{\True $5$}
{$4$}
{$6$}
{$7$}
\loigiai{
$12 - 7 = 5$.\\
\textbf{Đáp án của câu hỏi là A}
}
\end{ex}

\begin{ex} %[2D7V36]
Tính nhanh: $45 + 18 + 15 + 22$.
\choice
{$98$}
{\True $100$}
{$102$}
{$90$}
\loigiai{
$(45 + 15) + (18 + 22) = 60 + 40 = 100$.\\
\textbf{Đáp án của câu hỏi là B}
}
\end{ex}

\begin{ex} %[2D7C36]
Tìm số có hai chữ số có tổng các chữ số bằng 14, biết chữ số hàng chục hơn chữ số hàng đơn vị là 2.
\choice
{$77$}
{$95$}
{\True $86$}
{$68$}
\loigiai{
Chữ số hàng chục: $(14 + 2) : 2 = 8$.\\
Chữ số hàng đơn vị: $8 - 2 = 6$. Số đó là $86$.\\
\textbf{Đáp án của câu hỏi là C}
}
\end{ex}

\textbf{Phần 2. Trắc nghiệm đúng sai}
\setcounterpageref{ex}{0}

\begin{ex} %[2D7H36]
Xét tính đúng sai về kiến thức học kì 1:
\choiceTF
{\True Số lớn nhất có hai chữ số là 99.}
{\True $48 + 36 = 84$.}
{Hình tứ giác có 3 cạnh.}
{\True Hình tứ giác có 4 cạnh.}
\loigiai{
\begin{itemize}
    \item Số lớn nhất có 2 chữ số là 99: Đúng.
    \item $48 + 36 = 84$: Đúng.
    \item Tứ giác có 4 cạnh (không phải 3). Ý 3 Sai, Ý 4 Đúng.
\end{itemize}
\begin{itemize}
    \item a) \textbf{Đúng}.
    \item b) \textbf{Đúng}.
    \item c) \textbf{Sai}.
    \item d) \textbf{Đúng}.
\end{itemize}
}
\end{ex}

\begin{ex} %[2D7V36]
Một trang trại có 52 con bò và 38 con trâu. Xét tính đúng sai:
\choiceTF
{\True Trang trại có tất cả 90 con trâu và bò ($52 + 38 = 90$).}
{\True Số bò nhiều hơn số trâu là 14 con ($52 - 38 = 14$).}
{Nếu bán đi 12 con bò thì số bò còn lại ít hơn số trâu.}
{\True Khi đó còn $52 - 12 = 40$ con bò (vẫn nhiều hơn 38 con trâu).}
\loigiai{
\begin{itemize}
    \item Tổng: $52 + 38 = 90$: Đúng.
    \item Bò hơn trâu: $52 - 38 = 14$: Đúng.
    \item Bò còn 40 con $> 38$ con trâu. Ý 3 Sai, Ý 4 Đúng.
\end{itemize}
\begin{itemize}
    \item a) \textbf{Đúng}.
    \item b) \textbf{Đúng}.
    \item c) \textbf{Sai}.
    \item d) \textbf{Đúng}.
\end{itemize}
}
\end{ex}

\textbf{Phần 3. Trả lời ngắn}
\setcounterpageref{ex}{0}

\begin{ex} %[2D7H36]
Tính tổng: $49 + 38$.
\shortans{87}
\loigiai{
\textbf{Đáp số:} $87$\\
$49 + 38 = 87$.
}
\end{ex}

\begin{ex} %[2D7H36]
Tính hiệu: $85 - 39$.
\shortans{46}
\loigiai{
\textbf{Đáp số:} $46$\\
$85 - 39 = 46$.
}
\end{ex}

\begin{ex} %[2D7V36]
Tính độ dài đường gấp khúc gồm 3 đoạn thẳng có độ dài là $7\text{ cm}, 13\text{ cm}, 15\text{ cm}$ theo xăng-ti-mét.
\shortans{35}
\loigiai{
\textbf{Đáp số:} $35$\\
$7 + 13 + 15 = 35\text{ cm}$.
}
\end{ex}

\begin{ex} %[2D7V36]
Tìm $x$ biết: $x - 36 = 47$.
\shortans{83}
\loigiai{
\textbf{Đáp số:} $83$\\
$x = 47 + 36 = 83$.
}
\end{ex}

\begin{ex} %[2D7V36]
Một cửa hàng có $75\text{ kg}$ đường, đã bán đi $48\text{ kg}$. Hỏi cửa hàng còn lại bao nhiêu ki-lô-gam đường?
\shortans{27}
\loigiai{
\textbf{Đáp số:} $27$\\
$75 - 48 = 27\text{ kg}$.
}
\end{ex}

\begin{ex} %[2D7C36]
Tính nhanh tổng: $16 + 27 + 34 + 23$.
\shortans{100}
\loigiai{
\textbf{Đáp số:} $100$\\
$(16 + 34) + (27 + 23) = 50 + 50 = 100$.
}
\end{ex}

\newpage

% =========================================================================
% CHỦ ĐỀ 8 - PHÉP NHÂN, PHÉP CHIA
% =========================================================================
\section{Phép nhân, phép chia}

\TaoPhieuKiemTra{Chủ đề 8 -- Phép nhân, phép chia}{37}{Phép nhân. Thừa số - Tích}{https://loptoanthaylinh.vn/thi-online/toan-2/bai-37}
\textbf{Phần 1. Trắc nghiệm lựa chọn}
\setcounterpageref{ex}{0}

\begin{ex} %[2D8N37]
Tổng các số hạng bằng nhau $2 + 2 + 2 + 2$ được viết thành phép nhân là
\choice
{$2 \times 3$}
{$4 \times 4$}
{$2 \times 5$}
{\True $2 \times 4$}
\loigiai{
Số 2 được lấy 4 lần nên viết là $2 \times 4$.\\
\textbf{Đáp án của câu hỏi là D}
}
\end{ex}

\begin{ex} %[2D8N37]
Tổng $5 + 5 + 5$ viết thành phép nhân là
\choice
{\True $5 \times 3$}
{$5 \times 2$}
{$3 \times 5$}
{$5 \times 4$}
\loigiai{
$5$ lấy 3 lần là $5 \times 3$.\\
\textbf{Đáp án của câu hỏi là A}
}
\end{ex}

\begin{ex} %[2D8H37]
Phép nhân $3 \times 4$ được chuyển thành tổng các số hạng bằng nhau là
\choice
{$4 + 4 + 4$}
{\True $3 + 3 + 3 + 3$}
{$3 + 4$}
{$3 \times 3 \times 3 \times 3$}
\loigiai{
$3 \times 4 = 3 + 3 + 3 + 3$.\\
\textbf{Đáp án của câu hỏi là B}
}
\end{ex}

\begin{ex} %[2D8H37]
Tính kết quả phép nhân: $2 \times 5 = \dots$
\choice
{\True $10$}
{$8$}
{$12$}
{$15$}
\loigiai{
$2 + 2 + 2 + 2 + 2 = 10$.\\
\textbf{Đáp án của câu hỏi là A}
}
\end{ex}

\begin{ex} %[2D8H37]
Tính kết quả phép nhân: $5 \times 4 = \dots$
\choice
{$15$}
{$25$}
{\True $20$}
{$30$}
\loigiai{
$5 + 5 + 5 + 5 = 20$.\\
\textbf{Đáp án của câu hỏi là C}
}
\end{ex}

\begin{ex} %[2D8H37]
Mỗi bạn có 2 cái bút chì. Hỏi 6 bạn có tất cả bao nhiêu cái bút chì?
\choice
{$10$ cái}
{\True $12$ cái}
{$14$ cái}
{$8$ cái}
\loigiai{
$2 \times 6 = 12$ cái bút chì.\\
\textbf{Đáp án của câu hỏi là B}
}
\end{ex}

\begin{ex} %[2D8V37]
Mỗi đĩa có 5 quả cam. Hỏi 4 đĩa như thế có tất cả bao nhiêu quả cam?
\choice
{$15$ quả}
{$25$ quả}
{$10$ quả}
{\True $20$ quả}
\loigiai{
$5 \times 4 = 20$ quả cam.\\
\textbf{Đáp án của câu hỏi là D}
}
\end{ex}

\begin{ex} %[2D8V37]
Mỗi con gà có 2 cái chân. Hỏi 8 con gà có tất cả bao nhiêu cái chân?
\choice
{\True $16$ cái}
{$14$ cái}
{$18$ cái}
{$10$ cái}
\loigiai{
$2 \times 8 = 16$ cái chân.\\
\textbf{Đáp án của câu hỏi là A}
}
\end{ex}

\begin{ex} %[2D8V37]
Mỗi chuồng có 5 con thỏ. Hỏi 6 chuồng như thế có tất cả bao nhiêu con thỏ?
\choice
{$25$ con}
{$35$ con}
{\True $30$ con}
{$40$ con}
\loigiai{
$5 \times 6 = 30$ con thỏ.\\
\textbf{Đáp án của câu hỏi là C}
}
\end{ex}

\begin{ex} %[2D8V37]
Điền dấu thích hợp: $2 \times 7 \dots 5 \times 3$.
\choice
{$>$}
{$=$}
{Không so sánh được}
{\True $<$}
\loigiai{
$2 \times 7 = 14$ và $5 \times 3 = 15$. Vì $14 < 15$ nên điền dấu $<$.\\
\textbf{Đáp án của câu hỏi là D}
}
\end{ex}

\begin{ex} %[2D8V37]
Tính giá trị biểu thức: $5 \times 5 + 15$.
\choice
{$35$}
{$45$}
{\True $40$}
{$30$}
\loigiai{
$25 + 15 = 40$.\\
\textbf{Đáp án của câu hỏi là C}
}
\end{ex}

\begin{ex} %[2D8C37]
Có 4 gói kẹo, mỗi gói có 5 cái kẹo. Bé An ăn hết 3 cái kẹo. Hỏi còn lại bao nhiêu cái kẹo?
\choice
{$16$ cái}
{\True $17$ cái}
{$18$ cái}
{$15$ cái}
\loigiai{
Tổng số kẹo: $5 \times 4 = 20$ cái.\\
Còn lại: $20 - 3 = 17$ cái.\\
\textbf{Đáp án của câu hỏi là B}
}
\end{ex}

\textbf{Phần 2. Trắc nghiệm đúng sai}
\setcounterpageref{ex}{0}

\begin{ex} %[2D8H37]
Xét tính đúng sai về phép nhân:
\choiceTF
{\True Phép nhân là cách viết gọn của tổng các số hạng bằng nhau.}
{\True $2 + 2 + 2 = 2 \times 3 = 6$.}
{$5 + 5 + 5 + 5 = 5 \times 5$.}
{\True $5 + 5 + 5 + 5 = 5 \times 4 = 20$.}
\loigiai{
\begin{itemize}
    \item Phép nhân là tổng các số hạng bằng nhau: Đúng.
    \item $2 \times 3 = 6$: Đúng.
    \item $5$ lấy 4 lần là $5 \times 4$ (không phải $5 \times 5$). Ý 3 Sai, Ý 4 Đúng.
\end{itemize}
\begin{itemize}
    \item a) \textbf{Đúng}.
    \item b) \textbf{Đúng}.
    \item c) \textbf{Sai}.
    \item d) \textbf{Đúng}.
\end{itemize}
}
\end{ex}

\begin{ex} %[2D8V37]
Một bàn học có 2 bạn ngồi. Xét tính đúng sai:
\choiceTF
{\True 5 bàn học như thế có tất cả 10 bạn ($2 \times 5 = 10$).}
{\True 8 bàn học như thế có tất cả 16 bạn ($2 \times 8 = 16$).}
{10 bàn học có tất cả 18 bạn.}
{\True 10 bàn học có tất cả $2 \times 10 = 20$ bạn.}
\loigiai{
\begin{itemize}
    \item 5 bàn: $2 \times 5 = 10$: Đúng.
    \item 8 bàn: $2 \times 8 = 16$: Đúng.
    \item 10 bàn: $2 \times 10 = 20 \neq 18$. Ý 3 Sai, Ý 4 Đúng.
\end{itemize}
\begin{itemize}
    \item a) \textbf{Đúng}.
    \item b) \textbf{Đúng}.
    \item c) \textbf{Sai}.
    \item d) \textbf{Đúng}.
\end{itemize}
}
\end{ex}

\textbf{Phần 3. Trả lời ngắn}
\setcounterpageref{ex}{0}

\begin{ex} %[2D8H37]
Tính kết quả phép nhân: $2 \times 7$.
\shortans{14}
\loigiai{
\textbf{Đáp số:} $14$\\
$2 \times 7 = 14$.
}
\end{ex}

\begin{ex} %[2D8H37]
Tính kết quả phép nhân: $5 \times 6$.
\shortans{30}
\loigiai{
\textbf{Đáp số:} $30$\\
$5 \times 6 = 30$.
}
\end{ex}

\begin{ex} %[2D8V37]
Mỗi ô tô có 4 bánh xe. Hỏi 5 ô tô như thế có tất cả bao nhiêu bánh xe?
\shortans{20}
\loigiai{
\textbf{Đáp số:} $20$\\
$4 \times 5 = 20$ bánh xe.
}
\end{ex}

\begin{ex} %[2D8V37]
Mỗi hàng có 5 cây. Hỏi 7 hàng có tất cả bao nhiêu cây?
\shortans{35}
\loigiai{
\textbf{Đáp số:} $35$\\
$5 \times 7 = 35$ cây.
}
\end{ex}

\begin{ex} %[2D8V37]
Tính giá trị biểu thức: $2 \times 9 + 12$.
\shortans{30}
\loigiai{
\textbf{Đáp số:} $30$\\
$18 + 12 = 30$.
}
\end{ex}

\begin{ex} %[2D8C37]
Mỗi gói bánh có 5 cái bánh. Mẹ mua 8 gói bánh và cho các bạn 15 cái bánh. Hỏi mẹ còn lại bao nhiêu cái bánh?
\shortans{25}
\loigiai{
\textbf{Đáp số:} $25$\\
Tổng: $5 \times 8 = 40$ cái. Còn lại: $40 - 15 = 25$ cái bánh.
}
\end{ex}

\newpage

\TaoPhieuKiemTra{Chủ đề 8 -- Phép nhân, phép chia}{38}{Bảng nhân 2}{https://loptoanthaylinh.vn/thi-online/toan-2/bai-38}
\textbf{Phần 1. Trắc nghiệm lựa chọn}
\setcounterpageref{ex}{0}

\begin{ex} %[2D8N38]
Trong phép tính $2 \times 6 = 12$, số 12 được gọi là
\choice
{Thừa số}
{Số bị chia}
{\True Tích}
{Thương}
\loigiai{
Kết quả của phép nhân gọi là Tích.\\
\textbf{Đáp án của câu hỏi là C}
}
\end{ex}

\begin{ex} %[2D8N38]
Trong phép tính $5 \times 4 = 20$, các số 5 và 4 được gọi là
\choice
{Tích}
{\True Thừa số}
{Số hạng}
{Số bị trừ}
\loigiai{
5 và 4 là các thừa số.\\
\textbf{Đáp án của câu hỏi là B}
}
\end{ex}

\begin{ex} %[2D8H38]
Tìm tích biết hai thừa số lần lượt là 2 và 8:
\choice
{\True $16$}
{$14$}
{$18$}
{$10$}
\loigiai{
$2 \times 8 = 16$.\\
\textbf{Đáp án của câu hỏi là A}
}
\end{ex}

\begin{ex} %[2D8H38]
Tìm tích biết hai thừa số lần lượt là 5 và 7:
\choice
{$30$}
{$40$}
{$45$}
{\True $35$}
\loigiai{
$5 \times 7 = 35$.\\
\textbf{Đáp án của câu hỏi là D}
}
\end{ex}

\begin{ex} %[2D8H38]
Khi đổi chỗ hai thừa số trong một tích thì tích sẽ
\choice
{Tăng gấp đôi}
{Giảm đi một nửa}
{\True Không thay đổi}
{Bằng 0}
\loigiai{
Khi đổi chỗ các thừa số thì tích không thay đổi.\\
\textbf{Đáp án của câu hỏi là C}
}
\end{ex}

\begin{ex} %[2D8H38]
Tích của $2 \times 9$ có giá trị bằng tích nào sau đây?
\choice
{$9 + 2$}
{$2 + 9$}
{$2 \times 8$}
{\True $9 \times 2$}
\loigiai{
$2 \times 9 = 9 \times 2 = 18$.\\
\textbf{Đáp án của câu hỏi là D}
}
\end{ex}

\begin{ex} %[2D8V38]
Tích của hai thừa số là 18. Biết một thừa số là 2, thừa số kia là
\choice
{\True $9$}
{$8$}
{$7$}
{$6$}
\loigiai{
$18 : 2 = 9$.\\
\textbf{Đáp án của câu hỏi là A}
}
\end{ex}

\begin{ex} %[2D8V38]
Tích của hai thừa số là 40. Biết một thừa số là 5, thừa số kia là
\choice
{$7$}
{$9$}
{$6$}
{\True $8$}
\loigiai{
$40 : 5 = 8$.\\
\textbf{Đáp án của câu hỏi là D}
}
\end{ex}

\begin{ex} %[2D8V38]
Một thừa số là 5, thừa số kia là số liền sau số 8. Tích của hai số đó là
\choice
{$40$}
{\True $45$}
{$50$}
{$35$}
\loigiai{
Số liền sau của 8 là 9. Tích: $5 \times 9 = 45$.\\
\textbf{Đáp án của câu hỏi là B}
}
\end{ex}

\begin{ex} %[2D8V38]
Một thừa số là 2, thừa số kia là số lớn nhất có một chữ số. Tích là
\choice
{\True $18$}
{$16$}
{$20$}
{$14$}
\loigiai{
Số lớn nhất có 1 chữ số là 9. Tích: $2 \times 9 = 18$.\\
\textbf{Đáp án của câu hỏi là A}
}
\end{ex}

\begin{ex} %[2D8V38]
Tính giá trị biểu thức: $2 \times 8 + 5 \times 4$.
\choice
{$34$}
{$38$}
{\True $36$}
{$40$}
\loigiai{
$16 + 20 = 36$.\\
\textbf{Đáp án của câu hỏi là C}
}
\end{ex}

\begin{ex} %[2D8C38]
Tích của hai số là 10. Nếu giữ nguyên một thừa số và gấp thừa số kia lên 2 lần thì tích mới là
\choice
{$15$}
{\True $20$}
{$25$}
{$30$}
\loigiai{
$10 \times 2 = 20$.\\
\textbf{Đáp án của câu hỏi là B}
}
\end{ex}

\textbf{Phần 2. Trắc nghiệm đúng sai}
\setcounterpageref{ex}{0}

\begin{ex} %[2D8H38]
Xét tính đúng sai về thành phần của phép nhân:
\choiceTF
{\True Trong phép nhân $5 \times 6 = 30$, 5 và 6 là các thừa số.}
{\True 30 được gọi là tích.}
{$5 \times 6 = 6 + 5$.}
{\True $5 \times 6 = 6 \times 5 = 30$.}
\loigiai{
\begin{itemize}
    \item 5 và 6 là thừa số: Đúng.
    \item 30 là tích: Đúng.
    \item $5 \times 6 = 30 \neq 11$. Ý 3 Sai, Ý 4 Đúng.
\end{itemize}
\begin{itemize}
    \item a) \textbf{Đúng}.
    \item b) \textbf{Đúng}.
    \item c) \textbf{Sai}.
    \item d) \textbf{Đúng}.
\end{itemize}
}
\end{ex}

\begin{ex} %[2D8V38]
Cho phép nhân $2 \times 7 = 14$. Xét tính đúng sai:
\choiceTF
{\True Tích là 14.}
{\True Thừa số thứ nhất là 2, thừa số thứ hai là 7.}
{Nếu tăng thừa số 2 thêm 1 đơn vị thì tích tăng thêm 1 đơn vị.}
{\True Khi đó tích là $3 \times 7 = 21$ (tăng thêm 7 đơn vị).}
\loigiai{
\begin{itemize}
    \item Tích là 14: Đúng.
    \item Thừa số là 2 và 7: Đúng.
    \item Tích tăng thêm 7 (không phải 1). Ý 3 Sai, Ý 4 Đúng.
\end{itemize}
\begin{itemize}
    \item a) \textbf{Đúng}.
    \item b) \textbf{Đúng}.
    \item c) \textbf{Sai}.
    \item d) \textbf{Đúng}.
\end{itemize}
}
\end{ex}

\textbf{Phần 3. Trả lời ngắn}
\setcounterpageref{ex}{0}

\begin{ex} %[2D8H38]
Tìm tích của hai thừa số $2$ và $5$.
\shortans{10}
\loigiai{
\textbf{Đáp số:} $10$\\
$2 \times 5 = 10$.
}
\end{ex}

\begin{ex} %[2D8H38]
Tìm tích của hai thừa số $5$ và $8$.
\shortans{40}
\loigiai{
\textbf{Đáp số:} $40$\\
$5 \times 8 = 40$.
}
\end{ex}

\begin{ex} %[2D8V38]
Một thừa số là 2, thừa số kia là 9. Tính tích.
\shortans{18}
\loigiai{
\textbf{Đáp số:} $18$\\
$2 \times 9 = 18$.
}
\end{ex}

\begin{ex} %[2D8V38]
Một thừa số là 5, thừa số kia là 5. Tính tích.
\shortans{25}
\loigiai{
\textbf{Đáp số:} $25$\\
$5 \times 5 = 25$.
}
\end{ex}

\begin{ex} %[2D8V38]
Tính giá trị biểu thức: $2 \times 6 + 18$.
\shortans{30}
\loigiai{
\textbf{Đáp số:} $30$\\
$12 + 18 = 30$.
}
\end{ex}

\begin{ex} %[2D8C38]
Tìm tích của số 2 với số chẵn lớn nhất có một chữ số.
\shortans{16}
\loigiai{
\textbf{Đáp số:} $16$\\
Số chẵn lớn nhất có 1 chữ số là 8. Tích: $2 \times 8 = 16$.
}
\end{ex}

\newpage

\TaoPhieuKiemTra{Chủ đề 8 -- Phép nhân, phép chia}{39}{Bảng nhân 5}{https://loptoanthaylinh.vn/thi-online/toan-2/bai-39}
\textbf{Phần 1. Trắc nghiệm lựa chọn}
\setcounterpageref{ex}{0}

\begin{ex} %[2D8N39]
Kết quả của phép tính $2 \times 4$ là
\choice
{$6$}
{$10$}
{$12$}
{\True $8$}
\loigiai{
$2 \times 4 = 8$.\\
\textbf{Đáp án của câu hỏi là D}
}
\end{ex}

\begin{ex} %[2D8N39]
Kết quả của phép tính $2 \times 7$ là
\choice
{$12$}
{$16$}
{\True $14$}
{$18$}
\loigiai{
$2 \times 7 = 14$.\\
\textbf{Đáp án của câu hỏi là C}
}
\end{ex}

\begin{ex} %[2D8H39]
Phép tính nào sau đây có kết quả bằng 16?
\choice
{\True $2 \times 8$}
{$2 \times 7$}
{$2 \times 9$}
{$2 \times 6$}
\loigiai{
$2 \times 8 = 16$.\\
\textbf{Đáp án của câu hỏi là A}
}
\end{ex}

\begin{ex} %[2D8H39]
Phép tính nào trong bảng nhân 2 có kết quả lớn nhất?
\choice
{$2 \times 8$}
{$2 \times 9$}
{$2 \times 7$}
{\True $2 \times 10$}
\loigiai{
$2 \times 10 = 20$ (lớn nhất).\\
\textbf{Đáp án của câu hỏi là D}
}
\end{ex}

\begin{ex} %[2D8H39]
Trong bảng nhân 2, hai tích liên tiếp hơn kém nhau bao nhiêu đơn vị?
\choice
{\True $2$ đơn vị}
{$1$ đơn vị}
{$3$ đơn vị}
{$4$ đơn vị}
\loigiai{
Hai tích liên tiếp hơn kém nhau 2 đơn vị.\\
\textbf{Đáp án của câu hỏi là A}
}
\end{ex}

\begin{ex} %[2D8H39]
Điền số thích hợp: $2 \times \dots = 18$.
\choice
{$8$}
{$7$}
{\True $9$}
{$6$}
\loigiai{
$2 \times 9 = 18$.\\
\textbf{Đáp án của câu hỏi là C}
}
\end{ex}

\begin{ex} %[2D8V39]
Mỗi bạn học sinh có 2 hộp bút. Hỏi 8 bạn học sinh có tất cả bao nhiêu hộp bút?
\choice
{$14$ hộp}
{\True $16$ hộp}
{$18$ hộp}
{$12$ hộp}
\loigiai{
$2 \times 8 = 16$ hộp.\\
\textbf{Đáp án của câu hỏi là B}
}
\end{ex}

\begin{ex} %[2D8V39]
Mỗi đôi đũa có 2 chiếc đũa. Hỏi 9 đôi đũa có tất cả bao nhiêu chiếc đũa?
\choice
{$16$ chiếc}
{$20$ chiếc}
{$14$ chiếc}
{\True $18$ chiếc}
\loigiai{
$2 \times 9 = 18$ chiếc đũa.\\
\textbf{Đáp án của câu hỏi là D}
}
\end{ex}

\begin{ex} %[2D8V39]
Tính giá trị biểu thức: $2 \times 6 + 18$.
\choice
{$28$}
{\True $30$}
{$32$}
{$34$}
\loigiai{
$12 + 18 = 30$.\\
\textbf{Đáp án của câu hỏi là B}
}
\end{ex}

\begin{ex} %[2D8V39]
Điền dấu thích hợp: $2 \times 5 + 6 \dots 2 \times 8$.
\choice
{$>$}
{$<$}
{\True $=$}
{Không so sánh được}
\loigiai{
$2 \times 5 + 6 = 10 + 6 = 16$ và $2 \times 8 = 16$. Vậy bằng nhau.\\
\textbf{Đáp án của câu hỏi là C}
}
\end{ex}

\begin{ex} %[2D8V39]
Có 7 can nước mắm, mỗi can chứa $2\text{ l}$. Cả 7 can chứa số lít nước mắm là
\choice
{$12\text{ l}$}
{\True $14\text{ l}$}
{$16\text{ l}$}
{$18\text{ l}$}
\loigiai{
$2 \times 7 = 14\text{ l}$.\\
\textbf{Đáp án của câu hỏi là B}
}
\end{ex}

\begin{ex} %[2D8C39]
Tìm số thích hợp điền vào ô trống: $2 \times \dots > 15$ và $< 18$.
\choice
{\True $8$}
{$7$}
{$9$}
{$6$}
\loigiai{
$2 \times 8 = 16$ thỏa mãn $15 < 16 < 18$.\\
\textbf{Đáp án của câu hỏi là A}
}
\end{ex}

\textbf{Phần 2. Trắc nghiệm đúng sai}
\setcounterpageref{ex}{0}

\begin{ex} %[2D8H39]
Xét tính đúng sai của các phép tính trong bảng nhân 2:
\choiceTF
{\True $2 \times 3 = 6$.}
{\True $2 \times 7 = 14$.}
{$2 \times 9 = 19$.}
{\True $2 \times 9 = 18$.}
\loigiai{
\begin{itemize}
    \item $2 \times 3 = 6$: Đúng.
    \item $2 \times 7 = 14$: Đúng.
    \item $2 \times 9 = 18 \neq 19$. Ý 3 Sai, Ý 4 Đúng.
\end{itemize}
\begin{itemize}
    \item a) \textbf{Đúng}.
    \item b) \textbf{Đúng}.
    \item c) \textbf{Sai}.
    \item d) \textbf{Đúng}.
\end{itemize}
}
\end{ex}

\begin{ex} %[2D8V39]
Một đôi chim bồ câu có 2 con chim. Xét tính đúng sai:
\choiceTF
{\True 6 đôi chim có 12 con chim ($2 \times 6 = 12$).}
{\True 10 đôi chim có 20 con chim ($2 \times 10 = 20$).}
{8 đôi chim có 14 con chim.}
{\True 8 đôi chim có $2 \times 8 = 16$ con chim.}
\loigiai{
\begin{itemize}
    \item 6 đôi: $2 \times 6 = 12$: Đúng.
    \item 10 đôi: $2 \times 10 = 20$: Đúng.
    \item 8 đôi: $2 \times 8 = 16 \neq 14$. Ý 3 Sai, Ý 4 Đúng.
\end{itemize}
\begin{itemize}
    \item a) \textbf{Đúng}.
    \item b) \textbf{Đúng}.
    \item c) \textbf{Sai}.
    \item d) \textbf{Đúng}.
\end{itemize}
}
\end{ex}

\textbf{Phần 3. Trả lời ngắn}
\setcounterpageref{ex}{0}

\begin{ex} %[2D8H39]
Tính: $2 \times 5$.
\shortans{10}
\loigiai{
\textbf{Đáp số:} $10$\\
$2 \times 5 = 10$.
}
\end{ex}

\begin{ex} %[2D8H39]
Tính: $2 \times 8$.
\shortans{16}
\loigiai{
\textbf{Đáp số:} $16$\\
$2 \times 8 = 16$.
}
\end{ex}

\begin{ex} %[2D8V39]
Tính giá trị biểu thức: $2 \times 9 - 8$.
\shortans{10}
\loigiai{
\textbf{Đáp số:} $10$\\
$18 - 8 = 10$.
}
\end{ex}

\begin{ex} %[2D8V39]
Mỗi con vịt có 2 cái chân. Hỏi 10 con vịt có bao nhiêu cái chân?
\shortans{20}
\loigiai{
\textbf{Đáp số:} $20$\\
$2 \times 10 = 20$ cái chân.
}
\end{ex}

\begin{ex} %[2D8V39]
Điền số thích hợp: $2 \times ? = 14$.
\shortans{7}
\loigiai{
\textbf{Đáp số:} $7$\\
$14 : 2 = 7$.
}
\end{ex}

\begin{ex} %[2D8C39]
Có 6 túi kẹo, mỗi túi có 2 cái kẹo. Mẹ mua thêm 8 cái kẹo nữa. Hỏi có tất cả bao nhiêu cái kẹo?
\shortans{20}
\loigiai{
\textbf{Đáp số:} $20$\\
$2 \times 6 + 8 = 12 + 8 = 20$ cái.
}
\end{ex}

\newpage

\TaoPhieuKiemTra{Chủ đề 8 -- Phép nhân, phép chia}{40}{Phép chia. Số bị chia - Số chia - Thương}{https://loptoanthaylinh.vn/thi-online/toan-2/bai-40}
\textbf{Phần 1. Trắc nghiệm lựa chọn}
\setcounterpageref{ex}{0}

\begin{ex} %[2D8N40]
Kết quả của phép tính $5 \times 3$ là
\choice
{\True $15$}
{$10$}
{$20$}
{$25$}
\loigiai{
$5 \times 3 = 15$.\\
\textbf{Đáp án của câu hỏi là A}
}
\end{ex}

\begin{ex} %[2D8N40]
Kết quả của phép tính $5 \times 6$ là
\choice
{$25$}
{$35$}
{\True $30$}
{$40$}
\loigiai{
$5 \times 6 = 30$.\\
\textbf{Đáp án của câu hỏi là C}
}
\end{ex}

\begin{ex} %[2D8H40]
Phép tính nào sau đây có kết quả bằng 40?
\choice
{$5 \times 7$}
{\True $5 \times 8$}
{$5 \times 9$}
{$5 \times 6$}
\loigiai{
$5 \times 8 = 40$.\\
\textbf{Đáp án của câu hỏi là B}
}
\end{ex}

\begin{ex} %[2D8H40]
Tất cả các tích trong bảng nhân 5 đều có chữ số tận cùng là
\choice
{$0$ hoặc $2$}
{$5$ hoặc $9$}
{Chỉ là chữ số $5$}
{\True $0$ hoặc $5$}
\loigiai{
Các tích bảng nhân 5 luôn tận cùng là 0 hoặc 5.\\
\textbf{Đáp án của câu hỏi là D}
}
\end{ex}

\begin{ex} %[2D8H40]
Trong bảng nhân 5, hai tích liên tiếp hơn kém nhau bao nhiêu đơn vị?
\choice
{$2$ đơn vị}
{$10$ đơn vị}
{\True $5$ đơn vị}
{$1$ đơn vị}
\loigiai{
Hơn kém nhau 5 đơn vị.\\
\textbf{Đáp án của câu hỏi là C}
}
\end{ex}

\begin{ex} %[2D8H40]
Điền số thích hợp: $5 \times \dots = 45$.
\choice
{$8$}
{\True $9$}
{$7$}
{$10$}
\loigiai{
$5 \times 9 = 45$.\\
\textbf{Đáp án của câu hỏi là B}
}
\end{ex}

\begin{ex} %[2D8V40]
Mỗi bàn tay có 5 ngón tay. Hỏi 6 bàn tay có tất cả bao nhiêu ngón tay?
\choice
{\True $30$ ngón}
{$25$ ngón}
{$35$ ngón}
{$40$ ngón}
\loigiai{
$5 \times 6 = 30$ ngón tay.\\
\textbf{Đáp án của câu hỏi là A}
}
\end{ex}

\begin{ex} %[2D8V40]
Mỗi túi có 5 quả cam. Hỏi 8 túi như thế có tất cả bao nhiêu quả cam?
\choice
{$35$ quả}
{\True $40$ quả}
{$45$ quả}
{$50$ quả}
\loigiai{
$5 \times 8 = 40$ quả cam.\\
\textbf{Đáp án của câu hỏi là B}
}
\end{ex}

\begin{ex} %[2D8V40]
Tính giá trị biểu thức: $5 \times 7 + 15$.
\choice
{$45$}
{$55$}
{$60$}
{\True $50$}
\loigiai{
$35 + 15 = 50$.\\
\textbf{Đáp án của câu hỏi là D}
}
\end{ex}

\begin{ex} %[2D8V40]
Điền dấu thích hợp: $5 \times 4 \dots 2 \times 10$.
\choice
{$>$}
{$<$}
{\True $=$}
{Không so sánh được}
\loigiai{
$5 \times 4 = 20$ và $2 \times 10 = 20$. Vậy bằng nhau.\\
\textbf{Đáp án của câu hỏi là C}
}
\end{ex}

\begin{ex} %[2D8V40]
Một ngôi sao 5 cánh. Hỏi 9 ngôi sao như thế có tất cả bao nhiêu cánh?
\choice
{\True $45$ cánh}
{$40$ cánh}
{$50$ cánh}
{$35$ cánh}
\loigiai{
$5 \times 9 = 45$ cánh.\\
\textbf{Đáp án của câu hỏi là A}
}
\end{ex}

\begin{ex} %[2D8C40]
Có 7 can dầu loại $5\text{ l}$. Người ta đã dùng hết $18\text{ l}$ dầu. Hỏi còn lại bao nhiêu lít dầu?
\choice
{$16\text{ l}$}
{$18\text{ l}$}
{$19\text{ l}$}
{\True $17\text{ l}$}
\loigiai{
Tổng dầu: $5 \times 7 = 35\text{ l}$.\\
Còn lại: $35 - 18 = 17\text{ l}$.\\
\textbf{Đáp án của câu hỏi là D}
}
\end{ex}

\textbf{Phần 2. Trắc nghiệm đúng sai}
\setcounterpageref{ex}{0}

\begin{ex} %[2D8H40]
Xét tính đúng sai của các phép tính trong bảng nhân 5:
\choiceTF
{\True $5 \times 4 = 20$.}
{\True $5 \times 8 = 40$.}
{$5 \times 7 = 30$.}
{\True $5 \times 7 = 35$.}
\loigiai{
\begin{itemize}
    \item $5 \times 4 = 20$: Đúng.
    \item $5 \times 8 = 40$: Đúng.
    \item $5 \times 7 = 35 \neq 30$. Ý 3 Sai, Ý 4 Đúng.
\end{itemize}
\begin{itemize}
    \item a) \textbf{Đúng}.
    \item b) \textbf{Đúng}.
    \item c) \textbf{Sai}.
    \item d) \textbf{Đúng}.
\end{itemize}
}
\end{ex}

\begin{ex} %[2D8V40]
Một can đựng $5\text{ l}$ nước mắm. Xét tính đúng sai:
\choiceTF
{\True 5 can như thế đựng $25\text{ l}$ nước mắm ($5 \times 5 = 25$).}
{\True 9 can như thế đựng $45\text{ l}$ nước mắm ($5 \times 9 = 45$).}
{10 can như thế đựng $40\text{ l}$ nước mắm.}
{\True 10 can như thế đựng $5 \times 10 = 50\text{ l}$ nước mắm.}
\loigiai{
\begin{itemize}
    \item 5 can: $5 \times 5 = 25\text{ l}$: Đúng.
    \item 9 can: $5 \times 9 = 45\text{ l}$: Đúng.
    \item 10 can: $50\text{ l} \neq 40\text{ l}$. Ý 3 Sai, Ý 4 Đúng.
\end{itemize}
\begin{itemize}
    \item a) \textbf{Đúng}.
    \item b) \textbf{Đúng}.
    \item c) \textbf{Sai}.
    \item d) \textbf{Đúng}.
\end{itemize}
}
\end{ex}

\textbf{Phần 3. Trả lời ngắn}
\setcounterpageref{ex}{0}

\begin{ex} %[2D8H40]
Tính: $5 \times 5$.
\shortans{25}
\loigiai{
\textbf{Đáp số:} $25$\\
$5 \times 5 = 25$.
}
\end{ex}

\begin{ex} %[2D8H40]
Tính: $5 \times 9$.
\shortans{45}
\loigiai{
\textbf{Đáp số:} $45$\\
$5 \times 9 = 45$.
}
\end{ex}

\begin{ex} %[2D8V40]
Tính giá trị biểu thức: $5 \times 8 - 15$.
\shortans{25}
\loigiai{
\textbf{Đáp số:} $25$\\
$40 - 15 = 25$.
}
\end{ex}

\begin{ex} %[2D8V40]
Mỗi gói bánh có 5 cái bánh. Hỏi 6 gói bánh có bao nhiêu cái bánh?
\shortans{30}
\loigiai{
\textbf{Đáp số:} $30$\\
$5 \times 6 = 30$ cái bánh.
}
\end{ex}

\begin{ex} %[2D8V40]
Điền số thích hợp: $5 \times ? = 35$.
\shortans{7}
\loigiai{
\textbf{Đáp số:} $7$\\
$35 : 5 = 7$.
}
\end{ex}

\begin{ex} %[2D8C40]
Một nhóm có 8 bạn, mỗi bạn góp 5 quyển vở ủng hộ vùng cao. Sau đó cả nhóm mua thêm 10 quyển vở nữa. Hỏi nhóm đã ủng hộ tất cả bao nhiêu quyển vở?
\shortans{50}
\loigiai{
\textbf{Đáp số:} $50$\\
$5 \times 8 + 10 = 40 + 10 = 50$ quyển.
}
\end{ex}

\newpage

\TaoPhieuKiemTra{Chủ đề 8 -- Phép nhân, phép chia}{41}{Bảng chia 2}{https://loptoanthaylinh.vn/thi-online/toan-2/bai-41}
\textbf{Phần 1. Trắc nghiệm lựa chọn}
\setcounterpageref{ex}{0}

\begin{ex} %[2D8N41]
Chia đều 8 cái kẹo cho 2 bạn. Mỗi bạn được số cái kẹo là
\choice
{$3$ cái}
{$5$ cái}
{$2$ cái}
{\True $4$ cái}
\loigiai{
$8 : 2 = 4$ cái kẹo.\\
\textbf{Đáp án của câu hỏi là D}
}
\end{ex}

\begin{ex} %[2D8N41]
Từ phép nhân $2 \times 4 = 8$, ta viết được phép chia là
\choice
{$8 : 4 = 3$}
{\True $8 : 2 = 4$ và $8 : 4 = 2$}
{$8 : 8 = 1$}
{$4 : 2 = 2$}
\loigiai{
Từ $2 \times 4 = 8$ lập được $8 : 2 = 4$ và $8 : 4 = 2$.\\
\textbf{Đáp án của câu hỏi là B}
}
\end{ex}

\begin{ex} %[2D8H41]
Từ phép nhân $5 \times 3 = 15$, phép chia tương ứng là
\choice
{\True $15 : 5 = 3$ và $15 : 3 = 5$}
{$15 : 2 = 7$}
{$15 : 4 = 3$}
{$10 : 2 = 5$}
\loigiai{
$15 : 5 = 3$ và $15 : 3 = 5$.\\
\textbf{Đáp án của câu hỏi là A}
}
\end{ex}

\begin{ex} %[2D8H41]
Kết quả của phép chia $10 : 2$ là
\choice
{$4$}
{$6$}
{$8$}
{\True $5$}
\loigiai{
$10 : 2 = 5$.\\
\textbf{Đáp án của câu hỏi là D}
}
\end{ex}

\begin{ex} %[2D8H41]
Kết quả của phép chia $20 : 5$ là
\choice
{$3$}
{$5$}
{\True $4$}
{$6$}
\loigiai{
$20 : 5 = 4$.\\
\textbf{Đáp án của câu hỏi là C}
}
\end{ex}

\begin{ex} %[2D8H41]
Có 12 quả cam chia đều vào 2 đĩa. Mỗi đĩa có số quả cam là
\choice
{\True $6$ quả}
{$5$ quả}
{$7$ quả}
{$8$ quả}
\loigiai{
$12 : 2 = 6$ quả cam.\\
\textbf{Đáp án của câu hỏi là A}
}
\end{ex}

\begin{ex} %[2D8V41]
Có 15 bông hoa cắm đều vào 5 lọ. Mỗi lọ có số bông hoa là
\choice
{$2$ bông}
{\True $3$ bông}
{$4$ bông}
{$5$ bông}
\loigiai{
$15 : 5 = 3$ bông hoa.\\
\textbf{Đáp án của câu hỏi là B}
}
\end{ex}

\begin{ex} %[2D8V41]
Có 18 bạn xếp thành 2 hàng đều nhau. Mỗi hàng có số bạn là
\choice
{$8$ bạn}
{$10$ bạn}
{\True $9$ bạn}
{$7$ bạn}
\loigiai{
$18 : 2 = 9$ bạn.\\
\textbf{Đáp án của câu hỏi là C}
}
\end{ex}

\begin{ex} %[2D8V41]
Có 25 quyển vở chia đều cho 5 bạn. Mỗi bạn nhận được số quyển vở là
\choice
{\True $5$ quyển}
{$4$ quyển}
{$6$ quyển}
{$7$ quyển}
\loigiai{
$25 : 5 = 5$ quyển vở.\\
\textbf{Đáp án của câu hỏi là A}
}
\end{ex}

\begin{ex} %[2D8V41]
Điền dấu thích hợp: $14 : 2 \dots 30 : 5$.
\choice
{$<$}
{$=$}
{\True $>$}
{Không so sánh được}
\loigiai{
$14 : 2 = 7$ và $30 : 5 = 6$. Vì $7 > 6$ nên điền dấu $>$.\\
\textbf{Đáp án của câu hỏi là C}
}
\end{ex}

\begin{ex} %[2D8V41]
Tính giá trị biểu thức: $16 : 2 + 25$.
\choice
{$31$}
{$35$}
{$37$}
{\True $33$}
\loigiai{
$8 + 25 = 33$.\\
\textbf{Đáp án của câu hỏi là D}
}
\end{ex}

\begin{ex} %[2D8C41]
Có 35 cái kẹo chia đều vào 5 túi, lấy ra 2 túi kẹo. Số kẹo lấy ra là
\choice
{$12$ cái}
{\True $14$ cái}
{$15$ cái}
{$16$ cái}
\loigiai{
Mỗi túi có: $35 : 5 = 7$ cái kẹo.\\
2 túi có: $7 \times 2 = 14$ cái kẹo.\\
\textbf{Đáp án của câu hỏi là B}
}
\end{ex}

\textbf{Phần 2. Trắc nghiệm đúng sai}
\setcounterpageref{ex}{0}

\begin{ex} %[2D8H41]
Xét mối quan hệ giữa phép nhân và phép chia:
\choiceTF
{\True Từ phép nhân $2 \times 6 = 12$ ta có $12 : 2 = 6$.}
{\True Từ phép nhân $2 \times 6 = 12$ ta có $12 : 6 = 2$.}
{$12 : 2 = 5$.}
{\True $12 : 2 = 6$.}
\loigiai{
\begin{itemize}
    \item $12 : 2 = 6$: Đúng.
    \item $12 : 6 = 2$: Đúng.
    \item $12 : 2 = 6 \neq 5$. Ý 3 Sai, Ý 4 Đúng.
\end{itemize}
\begin{itemize}
    \item a) \textbf{Đúng}.
    \item b) \textbf{Đúng}.
    \item c) \textbf{Sai}.
    \item d) \textbf{Đúng}.
\end{itemize}
}
\end{ex}

\begin{ex} %[2D8V41]
Có 20 cái bánh chia đều vào các hộp:
\choiceTF
{\True Nếu chia vào 2 hộp thì mỗi hộp có 10 cái bánh ($20 : 2 = 10$).}
{\True Nếu chia vào 5 hộp thì mỗi hộp có 4 cái bánh ($20 : 5 = 4$).}
{Nếu chia vào 4 hộp thì mỗi hộp có 6 cái bánh.}
{\True Nếu chia vào 4 hộp thì mỗi hộp có $20 : 4 = 5$ cái bánh.}
\loigiai{
\begin{itemize}
    \item 2 hộp: $20 : 2 = 10$: Đúng.
    \item 5 hộp: $20 : 5 = 4$: Đúng.
    \item 4 hộp: $20 : 4 = 5 \neq 6$. Ý 3 Sai, Ý 4 Đúng.
\end{itemize}
\begin{itemize}
    \item a) \textbf{Đúng}.
    \item b) \textbf{Đúng}.
    \item c) \textbf{Sai}.
    \item d) \textbf{Đúng}.
\end{itemize}
}
\end{ex}

\textbf{Phần 3. Trả lời ngắn}
\setcounterpageref{ex}{0}

\begin{ex} %[2D8H41]
Tính: $14 : 2$.
\shortans{7}
\loigiai{
\textbf{Đáp số:} $7$\\
$14 : 2 = 7$.
}
\end{ex}

\begin{ex} %[2D8H41]
Tính: $35 : 5$.
\shortans{7}
\loigiai{
\textbf{Đáp số:} $7$\\
$35 : 5 = 7$.
}
\end{ex}

\begin{ex} %[2D8V41]
Có 16 lít dầu rót đều vào 2 can. Hỏi mỗi can có bao nhiêu lít dầu?
\shortans{8}
\loigiai{
\textbf{Đáp số:} $8$\\
$16 : 2 = 8\text{ l}$.
}
\end{ex}

\begin{ex} %[2D8V41]
Có 40 cái kẹo chia đều cho 5 bạn. Hỏi mỗi bạn được bao nhiêu cái kẹo?
\shortans{8}
\loigiai{
\textbf{Đáp số:} $8$\\
$40 : 5 = 8$ cái.
}
\end{ex}

\begin{ex} %[2D8V41]
Tính giá trị biểu thức: $45 : 5 + 16$.
\shortans{25}
\loigiai{
\textbf{Đáp số:} $25$\\
$9 + 16 = 25$.
}
\end{ex}

\begin{ex} %[2D8C41]
Một khúc gỗ dài $18\text{ dm}$ được cưa thành các đoạn bằng nhau dài $2\text{ dm}$. Hỏi cần bao nhiêu lần cưa?
\shortans{8}
\loigiai{
\textbf{Đáp số:} $8$\\
Số đoạn gỗ: $18 : 2 = 9$ đoạn.\\
Số lần cưa: $9 - 1 = 8$ lần.
}
\end{ex}

\newpage

\TaoPhieuKiemTra{Chủ đề 8 -- Phép nhân, phép chia}{42}{Bảng chia 5}{https://loptoanthaylinh.vn/thi-online/toan-2/bai-42}
\textbf{Phần 1. Trắc nghiệm lựa chọn}
\setcounterpageref{ex}{0}

\begin{ex} %[2D8N42]
Trong phép tính $18 : 2 = 9$, số 18 được gọi là
\choice
{Số chia}
{\True Số bị chia}
{Thương}
{Số hạng}
\loigiai{
18 là Số bị chia.\\
\textbf{Đáp án của câu hỏi là B}
}
\end{ex}

\begin{ex} %[2D8N42]
Trong phép tính $18 : 2 = 9$, số 2 được gọi là
\choice
{\True Số chia}
{Số bị chia}
{Thương}
{Tích}
\loigiai{
2 là Số chia.\\
\textbf{Đáp án của câu hỏi là A}
}
\end{ex}

\begin{ex} %[2D8H42]
Trong phép tính $18 : 2 = 9$, số 9 được gọi là
\choice
{Tổng}
{Số chia}
{\True Thương}
{Số bị chia}
\loigiai{
Kết quả của phép chia gọi là Thương.\\
\textbf{Đáp án của câu hỏi là C}
}
\end{ex}

\begin{ex} %[2D8H42]
Tìm thương biết số bị chia là 30, số chia là 5:
\choice
{$5$}
{$7$}
{$8$}
{\True $6$}
\loigiai{
$30 : 5 = 6$.\\
\textbf{Đáp án của câu hỏi là D}
}
\end{ex}

\begin{ex} %[2D8H42]
Tìm thương biết số bị chia là 16, số chia là 2:
\choice
{$6$}
{\True $8$}
{$7$}
{$9$}
\loigiai{
$16 : 2 = 8$.\\
\textbf{Đáp án của câu hỏi là B}
}
\end{ex}

\begin{ex} %[2D8H42]
Số bị chia là 45, số chia là 5. Thương là
\choice
{\True $9$}
{$8$}
{$7$}
{$6$}
\loigiai{
$45 : 5 = 9$.\\
\textbf{Đáp án của câu hỏi là A}
}
\end{ex}

\begin{ex} %[2D8V42]
Thương của phép chia là 7, số chia là 2. Số bị chia là
\choice
{$12$}
{$16$}
{$9$}
{\True $14$}
\loigiai{
Số bị chia: $7 \times 2 = 14$.\\
\textbf{Đáp án của câu hỏi là D}
}
\end{ex}

\begin{ex} %[2D8V42]
Thương của phép chia là 8, số chia là 5. Số bị chia là
\choice
{\True $40$}
{$35$}
{$45$}
{$30$}
\loigiai{
Số bị chia: $8 \times 5 = 40$.\\
\textbf{Đáp án của câu hỏi là A}
}
\end{ex}

\begin{ex} %[2D8V42]
Số bị chia là 20, thương là 4. Số chia là
\choice
{$4$}
{$6$}
{\True $5$}
{$10$}
\loigiai{
Số chia: $20 : 4 = 5$.\\
\textbf{Đáp án của câu hỏi là C}
}
\end{ex}

\begin{ex} %[2D8V42]
Số bị chia là 18, thương là 9. Số chia là
\choice
{$1$}
{$3$}
{$4$}
{\True $2$}
\loigiai{
Số chia: $18 : 9 = 2$.\\
\textbf{Đáp án của câu hỏi là D}
}
\end{ex}

\begin{ex} %[2D8V42]
Tính giá trị biểu thức: Thương của 35 và 5 cộng với 18.
\choice
{$23$}
{$27$}
{\True $25$}
{$29$}
\loigiai{
$35 : 5 + 18 = 7 + 18 = 25$.\\
\textbf{Đáp án của câu hỏi là C}
}
\end{ex}

\begin{ex} %[2D8C42]
Tìm một số biết rằng lấy số đó chia cho 5 thì được kết quả là số chẵn liền sau của 6.
\choice
{$35$}
{\True $40$}
{$45$}
{$30$}
\loigiai{
Số chẵn liền sau của 6 là 8. Số cần tìm: $8 \times 5 = 40$.\\
\textbf{Đáp án của câu hỏi là B}
}
\end{ex}

\textbf{Phần 2. Trắc nghiệm đúng sai}
\setcounterpageref{ex}{0}

\begin{ex} %[2D8H42]
Xét tính đúng sai về thành phần phép chia:
\choiceTF
{\True Trong phép chia $25 : 5 = 5$, 25 là số bị chia.}
{\True Số 5 đứng sau dấu chia gọi là số chia.}
{Kết quả 5 gọi là tích.}
{\True Kết quả của phép chia gọi là thương.}
\loigiai{
\begin{itemize}
    \item 25 là số bị chia: Đúng.
    \item 5 là số chia: Đúng.
    \item Kết quả gọi là thương (không phải tích). Ý 3 Sai, Ý 4 Đúng.
\end{itemize}
\begin{itemize}
    \item a) \textbf{Đúng}.
    \item b) \textbf{Đúng}.
    \item c) \textbf{Sai}.
    \item d) \textbf{Đúng}.
\end{itemize}
}
\end{ex}

\begin{ex} %[2D8V42]
Cho phép chia có số bị chia là 16, số chia là 2. Xét tính đúng sai:
\choiceTF
{\True Thương của phép chia là 8.}
{\True Nếu gấp số bị chia lên 2 lần (thành 32) thì thương mới là 16.}
{Nếu giảm số bị chia đi 2 lần thì thương không đổi.}
{\True Khi đó thương mới là $8 : 2 = 4$.}
\loigiai{
\begin{itemize}
    \item Thương là 8: Đúng.
    \item $32 : 2 = 16$: Đúng.
    \item Thương giảm đi 2 lần. Ý 3 Sai, Ý 4 Đúng.
\end{itemize}
\begin{itemize}
    \item a) \textbf{Đúng}.
    \item b) \textbf{Đúng}.
    \item c) \textbf{Sai}.
    \item d) \textbf{Đúng}.
\end{itemize}
}
\end{ex}

\textbf{Phần 3. Trả lời ngắn}
\setcounterpageref{ex}{0}

\begin{ex} %[2D8H42]
Tìm thương biết số bị chia là 14, số chia là 2.
\shortans{7}
\loigiai{
\textbf{Đáp số:} $7$\\
$14 : 2 = 7$.
}
\end{ex}

\begin{ex} %[2D8H42]
Tìm thương biết số bị chia là 40, số chia là 5.
\shortans{8}
\loigiai{
\textbf{Đáp số:} $8$\\
$40 : 5 = 8$.
}
\end{ex}

\begin{ex} %[2D8V42]
Tìm số bị chia biết số chia là 2, thương là 9.
\shortans{18}
\loigiai{
\textbf{Đáp số:} $18$\\
$9 \times 2 = 18$.
}
\end{ex}

\begin{ex} %[2D8V42]
Tìm số bị chia biết số chia là 5, thương là 6.
\shortans{30}
\loigiai{
\textbf{Đáp số:} $30$\\
$6 \times 5 = 30$.
}
\end{ex}

\begin{ex} %[2D8V42]
Tìm số chia biết số bị chia là 50, thương là 10.
\shortans{5}
\loigiai{
\textbf{Đáp số:} $5$\\
$50 : 10 = 5$.
}
\end{ex}

\begin{ex} %[2D8C42]
Tìm thương của số tròn chục nhỏ nhất có hai chữ số chia cho 2.
\shortans{5}
\loigiai{
\textbf{Đáp số:} $5$\\
Số tròn chục nhỏ nhất là 10. Thương: $10 : 2 = 5$.
}
\end{ex}

\newpage

\TaoPhieuKiemTra{Chủ đề 8 -- Phép nhân, phép chia}{43}{Số 1 và số 0 trong phép nhân, phép chia}{https://loptoanthaylinh.vn/thi-online/toan-2/bai-43}
\textbf{Phần 1. Trắc nghiệm lựa chọn}
\setcounterpageref{ex}{0}

\begin{ex} %[2D8N43]
Kết quả của phép tính $8 : 2$ là
\choice
{$3$}
{$5$}
{$6$}
{\True $4$}
\loigiai{
$8 : 2 = 4$.\\
\textbf{Đáp án của câu hỏi là D}
}
\end{ex}

\begin{ex} %[2D8N43]
Kết quả của phép tính $14 : 2$ là
\choice
{$6$}
{\True $7$}
{$8$}
{$9$}
\loigiai{
$14 : 2 = 7$.\\
\textbf{Đáp án của câu hỏi là B}
}
\end{ex}

\begin{ex} %[2D8H43]
Phép tính nào sau đây có kết quả bằng 9?
\choice
{\True $18 : 2$}
{$16 : 2$}
{$20 : 2$}
{$14 : 2$}
\loigiai{
$18 : 2 = 9$.\\
\textbf{Đáp án của câu hỏi là A}
}
\end{ex}

\begin{ex} %[2D8H43]
Phép tính nào trong bảng chia 2 có kết quả bé nhất?
\choice
{$4 : 2$}
{$6 : 2$}
{\True $2 : 2$}
{$8 : 2$}
\loigiai{
$2 : 2 = 1$ (bé nhất).\\
\textbf{Đáp án của câu hỏi là C}
}
\end{ex}

\begin{ex} %[2D8H43]
Trong bảng chia 2, các số bị chia đều là
\choice
{Các số lẻ}
{Các số tròn chục}
{Các số tận cùng là 5}
{\True Các số chẵn từ 2 đến 20}
\loigiai{
Các số bị chia là các số chẵn: 2, 4, 6, 8, ..., 20.\\
\textbf{Đáp án của câu hỏi là D}
}
\end{ex}

\begin{ex} %[2D8H43]
Điền số thích hợp: $16 : 2 = \dots$
\choice
{\True $8$}
{$7$}
{$9$}
{$6$}
\loigiai{
$16 : 2 = 8$.\\
\textbf{Đáp án của câu hỏi là A}
}
\end{ex}

\begin{ex} %[2D8V43]
Có 12 cái kẹo chia đều cho 2 bạn. Mỗi bạn nhận được số kẹo là
\choice
{$5$ cái}
{$7$ cái}
{\True $6$ cái}
{$8$ cái}
\loigiai{
$12 : 2 = 6$ cái kẹo.\\
\textbf{Đáp án của câu hỏi là C}
}
\end{ex}

\begin{ex} %[2D8V43]
Một sợi dây dài $20\text{ dm}$ cắt làm 2 đoạn bằng nhau. Mỗi đoạn dài là
\choice
{$8\text{ dm}$}
{$12\text{ dm}$}
{$15\text{ dm}$}
{\True $10\text{ dm}$}
\loigiai{
$20 : 2 = 10\text{ dm}$.\\
\textbf{Đáp án của câu hỏi là D}
}
\end{ex}

\begin{ex} %[2D8V43]
Tính giá trị biểu thức: $18 : 2 + 15$.
\choice
{$22$}
{\True $24$}
{$26$}
{$28$}
\loigiai{
$9 + 15 = 24$.\\
\textbf{Đáp án của câu hỏi là B}
}
\end{ex}

\begin{ex} %[2D8V43]
Điền dấu thích hợp: $14 : 2 \dots 16 : 2$.
\choice
{$>$}
{$=$}
{\True $<$}
{Không so sánh được}
\loigiai{
$14 : 2 = 7 < 16 : 2 = 8$.\\
\textbf{Đáp án của câu hỏi là C}
}
\end{ex}

\begin{ex} %[2D8V43]
Có 10 lít nước mắm rót đều vào 2 can. Mỗi can có số lít nước mắm là
\choice
{$4\text{ l}$}
{\True $5\text{ l}$}
{$6\text{ l}$}
{$7\text{ l}$}
\loigiai{
$10 : 2 = 5\text{ l}$.\\
\textbf{Đáp án của câu hỏi là B}
}
\end{ex}

\begin{ex} %[2D8C43]
Có 18 quả cam chia đều vào 2 rổ. Mẹ lấy ở một rổ ra 3 quả. Hỏi rổ đó còn lại bao nhiêu quả cam?
\choice
{\True $6$ quả}
{$5$ quả}
{$7$ quả}
{$8$ quả}
\loigiai{
Mỗi rổ có: $18 : 2 = 9$ quả.\\
Còn lại: $9 - 3 = 6$ quả.\\
\textbf{Đáp án của câu hỏi là A}
}
\end{ex}

\textbf{Phần 2. Trắc nghiệm đúng sai}
\setcounterpageref{ex}{0}

\begin{ex} %[2D8H43]
Xét tính đúng sai của các phép tính trong bảng chia 2:
\choiceTF
{\True $6 : 2 = 3$.}
{\True $12 : 2 = 6$.}
{$18 : 2 = 8$.}
{\True $18 : 2 = 9$.}
\loigiai{
\begin{itemize}
    \item $6 : 2 = 3$: Đúng.
    \item $12 : 2 = 6$: Đúng.
    \item $18 : 2 = 9 \neq 8$. Ý 3 Sai, Ý 4 Đúng.
\end{itemize}
\begin{itemize}
    \item a) \textbf{Đúng}.
    \item b) \textbf{Đúng}.
    \item c) \textbf{Sai}.
    \item d) \textbf{Đúng}.
\end{itemize}
}
\end{ex}

\begin{ex} %[2D8V43]
Có 16 học sinh chia đều thành 2 tổ. Xét tính đúng sai:
\choiceTF
{\True Mỗi tổ có 8 học sinh ($16 : 2 = 8$).}
{\True Số học sinh của hai tổ bằng nhau.}
{Nếu chuyển 2 bạn từ tổ 1 sang tổ 2 thì tổ 2 có 9 bạn.}
{\True Khi đó tổ 2 có $8 + 2 = 10$ bạn.}
\loigiai{
\begin{itemize}
    \item Mỗi tổ 8 bạn: Đúng.
    \item Hai tổ bằng nhau: Đúng.
    \item $8 + 2 = 10 \neq 9$. Ý 3 Sai, Ý 4 Đúng.
\end{itemize}
\begin{itemize}
    \item a) \textbf{Đúng}.
    \item b) \textbf{Đúng}.
    \item c) \textbf{Sai}.
    \item d) \textbf{Đúng}.
\end{itemize}
}
\end{ex}

\textbf{Phần 3. Trả lời ngắn}
\setcounterpageref{ex}{0}

\begin{ex} %[2D8H43]
Tính: $10 : 2$.
\shortans{5}
\loigiai{
\textbf{Đáp số:} $5$\\
$10 : 2 = 5$.
}
\end{ex}

\begin{ex} %[2D8H43]
Tính: $16 : 2$.
\shortans{8}
\loigiai{
\textbf{Đáp số:} $8$\\
$16 : 2 = 8$.
}
\end{ex}

\begin{ex} %[2D8V43]
Tính: $20 : 2$.
\shortans{10}
\loigiai{
\textbf{Đáp số:} $10$\\
$20 : 2 = 10$.
}
\end{ex}

\begin{ex} %[2D8V43]
Tính giá trị biểu thức: $14 : 2 + 18$.
\shortans{25}
\loigiai{
\textbf{Đáp số:} $25$\\
$7 + 18 = 25$.
}
\end{ex}

\begin{ex} %[2D8V43]
Có 12 cái bánh chia đều cho 2 bạn. Hỏi mỗi bạn được bao nhiêu cái bánh?
\shortans{6}
\loigiai{
\textbf{Đáp số:} $6$\\
$12 : 2 = 6$ cái bánh.
}
\end{ex}

\begin{ex} %[2D8C43]
Một đoạn dây dài $18\text{ cm}$ gập đôi lại. Hỏi độ dài mỗi nửa đoạn dây là bao nhiêu xăng-ti-mét?
\shortans{9}
\loigiai{
\textbf{Đáp số:} $9$\\
$18 : 2 = 9\text{ cm}$.
}
\end{ex}

\newpage

\TaoPhieuKiemTra{Chủ đề 8 -- Phép nhân, phép chia}{44}{Luyện tập phép nhân và phép chia}{https://loptoanthaylinh.vn/thi-online/toan-2/bai-44}
\textbf{Phần 1. Trắc nghiệm lựa chọn}
\setcounterpageref{ex}{0}

\begin{ex} %[2D8N44]
Kết quả của phép tính $15 : 5$ là
\choice
{$2$}
{$4$}
{\True $3$}
{$5$}
\loigiai{
$15 : 5 = 3$.\\
\textbf{Đáp án của câu hỏi là C}
}
\end{ex}

\begin{ex} %[2D8N44]
Kết quả của phép tính $25 : 5$ là
\choice
{$4$}
{$6$}
{$7$}
{\True $5$}
\loigiai{
$25 : 5 = 5$.\\
\textbf{Đáp án của câu hỏi là D}
}
\end{ex}

\begin{ex} %[2D8H44]
Phép tính nào sau đây có kết quả bằng 8?
\choice
{\True $40 : 5$}
{$35 : 5$}
{$45 : 5$}
{$30 : 5$}
\loigiai{
$40 : 5 = 8$.\\
\textbf{Đáp án của câu hỏi là A}
}
\end{ex}

\begin{ex} %[2D8H44]
Phép tính nào trong bảng chia 5 có kết quả lớn nhất?
\choice
{$45 : 5$}
{\True $50 : 5$}
{$40 : 5$}
{$35 : 5$}
\loigiai{
$50 : 5 = 10$ (lớn nhất).\\
\textbf{Đáp án của câu hỏi là B}
}
\end{ex}

\begin{ex} %[2D8H44]
Trong bảng chia 5, các số bị chia đều tận cùng là chữ số nào?
\choice
{$1$ hoặc $5$}
{$2$ hoặc $5$}
{\True $0$ hoặc $5$}
{Chỉ chữ số $0$}
\loigiai{
Tận cùng là 0 hoặc 5.\\
\textbf{Đáp án của câu hỏi là C}
}
\end{ex}

\begin{ex} %[2D8H44]
Điền số thích hợp: $35 : 5 = \dots$
\choice
{$6$}
{$8$}
{$9$}
{\True $7$}
\loigiai{
$35 : 5 = 7$.\\
\textbf{Đáp án của câu hỏi là D}
}
\end{ex}

\begin{ex} %[2D8V44]
Có 20 quyển vở chia đều cho 5 bạn. Mỗi bạn nhận được số quyển vở là
\choice
{$3$ quyển}
{\True $4$ quyển}
{$5$ quyển}
{$6$ quyển}
\loigiai{
$20 : 5 = 4$ quyển vở.\\
\textbf{Đáp án của câu hỏi là B}
}
\end{ex}

\begin{ex} %[2D8V44]
Một bao đường nặng $30\text{ kg}$ chia đều vào 5 túi. Mỗi túi nặng số ki-lô-gam là
\choice
{\True $6\text{ kg}$}
{$5\text{ kg}$}
{$7\text{ kg}$}
{$8\text{ kg}$}
\loigiai{
$30 : 5 = 6\text{ kg}$.\\
\textbf{Đáp án của câu hỏi là A}
}
\end{ex}

\begin{ex} %[2D8V44]
Tính giá trị biểu thức: $45 : 5 + 28$.
\choice
{$35$}
{$39$}
{$41$}
{\True $37$}
\loigiai{
$9 + 28 = 37$.\\
\textbf{Đáp án của câu hỏi là D}
}
\end{ex}

\begin{ex} %[2D8V44]
Điền dấu thích hợp: $40 : 5 \dots 18 : 2$.
\choice
{\True $<$}
{$>$}
{$=$}
{Không so sánh được}
\loigiai{
$40 : 5 = 8 < 18 : 2 = 9$.\\
\textbf{Đáp án của câu hỏi là A}
}
\end{ex}

\begin{ex} %[2D8V44]
Có 50 lít dầu rót đều vào 5 thùng. Mỗi thùng có số lít dầu là
\choice
{$8\text{ l}$}
{$12\text{ l}$}
{\True $10\text{ l}$}
{$15\text{ l}$}
\loigiai{
$50 : 5 = 10\text{ l}$.\\
\textbf{Đáp án của câu hỏi là C}
}
\end{ex}

\begin{ex} %[2D8C44]
Có 30 bông hoa cắm đều vào 5 bình hoa. Sau đó cắm thêm vào mỗi bình 2 bông hoa nữa. Hỏi mỗi bình có tất cả bao nhiêu bông hoa?
\choice
{$7$ bông}
{\True $8$ bông}
{$9$ bông}
{$10$ bông}
\loigiai{
Ban đầu mỗi bình: $30 : 5 = 6$ bông.\\
Lúc sau mỗi bình: $6 + 2 = 8$ bông hoa.\\
\textbf{Đáp án của câu hỏi là B}
}
\end{ex}

\textbf{Phần 2. Trắc nghiệm đúng sai}
\setcounterpageref{ex}{0}

\begin{ex} %[2D8H44]
Xét tính đúng sai của các phép tính trong bảng chia 5:
\choiceTF
{\True $10 : 5 = 2$.}
{\True $25 : 5 = 5$.}
{$45 : 5 = 8$.}
{\True $45 : 5 = 9$.}
\loigiai{
\begin{itemize}
    \item $10 : 5 = 2$: Đúng.
    \item $25 : 5 = 5$: Đúng.
    \item $45 : 5 = 9 \neq 8$. Ý 3 Sai, Ý 4 Đúng.
\end{itemize}
\begin{itemize}
    \item a) \textbf{Đúng}.
    \item b) \textbf{Đúng}.
    \item c) \textbf{Sai}.
    \item d) \textbf{Đúng}.
\end{itemize}
}
\end{ex}

\begin{ex} %[2D8V44]
Có 35 quả cam chia đều vào 5 đĩa. Xét tính đúng sai:
\choiceTF
{\True Mỗi đĩa có 7 quả cam ($35 : 5 = 7$).}
{\True 3 đĩa như thế có 21 quả cam ($7 \times 3 = 21$).}
{Nếu lấy bớt 1 đĩa thì còn lại 30 quả cam.}
{\True Khi đó còn lại $7 \times 4 = 28$ quả cam.}
\loigiai{
\begin{itemize}
    \item Mỗi đĩa 7 quả: Đúng.
    \item 3 đĩa có 21 quả: Đúng.
    \item Còn lại: $7 \times 4 = 28 \neq 30$. Ý 3 Sai, Ý 4 Đúng.
\end{itemize}
\begin{itemize}
    \item a) \textbf{Đúng}.
    \item b) \textbf{Đúng}.
    \item c) \textbf{Sai}.
    \item d) \textbf{Đúng}.
\end{itemize}
}
\end{ex}

\textbf{Phần 3. Trả lời ngắn}
\setcounterpageref{ex}{0}

\begin{ex} %[2D8H44]
Tính: $20 : 5$.
\shortans{4}
\loigiai{
\textbf{Đáp số:} $4$\\
$20 : 5 = 4$.
}
\end{ex}

\begin{ex} %[2D8H44]
Tính: $35 : 5$.
\shortans{7}
\loigiai{
\textbf{Đáp số:} $7$\\
$35 : 5 = 7$.
}
\end{ex}

\begin{ex} %[2D8V44]
Tính: $40 : 5$.
\shortans{8}
\loigiai{
\textbf{Đáp số:} $8$\\
$40 : 5 = 8$.
}
\end{ex}

\begin{ex} %[2D8V44]
Tính giá trị biểu thức: $50 : 5 - 6$.
\shortans{4}
\loigiai{
\textbf{Đáp số:} $4$\\
$10 - 6 = 4$.
}
\end{ex}

\begin{ex} %[2D8V44]
Có 45 cái kẹo chia đều cho 5 bạn. Hỏi mỗi bạn được bao nhiêu cái kẹo?
\shortans{9}
\loigiai{
\textbf{Đáp số:} $9$\\
$45 : 5 = 9$ cái kẹo.
}
\end{ex}

\begin{ex} %[2D8C44]
Một lớp có 35 học sinh chia đều thành 5 tổ. Tổ Một có 3 bạn nữ. Hỏi tổ Một có bao nhiêu bạn nam?
\shortans{4}
\loigiai{
\textbf{Đáp số:} $4$\\
Mỗi tổ có: $35 : 5 = 7$ bạn.\\
Số bạn nam: $7 - 3 = 4$ bạn.
}
\end{ex}

\newpage

\TaoPhieuKiemTra{Chủ đề 8 -- Phép nhân, phép chia}{45}{Luyện tập chung}{https://loptoanthaylinh.vn/thi-online/toan-2/bai-45}
\textbf{Phần 1. Trắc nghiệm lựa chọn}
\setcounterpageref{ex}{0}

\begin{ex} %[2D8N45]
Kết quả của phép tính $2 \times 8$ là
\choice
{$14$}
{$18$}
{\True $16$}
{$20$}
\loigiai{
$2 \times 8 = 16$.\\
\textbf{Đáp án của câu hỏi là C}
}
\end{ex}

\begin{ex} %[2D8N45]
Kết quả của phép tính $45 : 5$ là
\choice
{$8$}
{$7$}
{$6$}
{\True $9$}
\loigiai{
$45 : 5 = 9$.\\
\textbf{Đáp án của câu hỏi là D}
}
\end{ex}

\begin{ex} %[2D8H45]
Kết quả của phép tính $5 \times 7$ là
\choice
{\True $35$}
{$30$}
{$40$}
{$45$}
\loigiai{
$5 \times 7 = 35$.\\
\textbf{Đáp án của câu hỏi là A}
}
\end{ex}

\begin{ex} %[2D8H45]
Kết quả của phép tính $18 : 2$ là
\choice
{$8$}
{\True $9$}
{$7$}
{$6$}
\loigiai{
$18 : 2 = 9$.\\
\textbf{Đáp án của câu hỏi là B}
}
\end{ex}

\begin{ex} %[2D8H45]
Điền dấu thích hợp: $2 \times 9 \dots 5 \times 4$.
\choice
{\True $<$}
{$>$}
{$=$}
{Không so sánh được}
\loigiai{
$2 \times 9 = 18 < 5 \times 4 = 20$.\\
\textbf{Đáp án của câu hỏi là A}
}
\end{ex}

\begin{ex} %[2D8H45]
Tính giá trị biểu thức: $5 \times 6 + 18 : 2$.
\choice
{$37$}
{$41$}
{$43$}
{\True $39$}
\loigiai{
$30 + 9 = 39$.\\
\textbf{Đáp án của câu hỏi là D}
}
\end{ex}

\begin{ex} %[2D8V45]
Có 16 lít mật ong rót đều vào 2 can. Mỗi can có số lít mật ong là
\choice
{$7\text{ l}$}
{\True $8\text{ l}$}
{$9\text{ l}$}
{$10\text{ l}$}
\loigiai{
$16 : 2 = 8\text{ l}$.\\
\textbf{Đáp án của câu hỏi là B}
}
\end{ex}

\begin{ex} %[2D8V45]
Mỗi gói bánh có 5 chiếc bánh. Hỏi 9 gói bánh có tất cả bao nhiêu chiếc bánh?
\choice
{$40$ chiếc}
{$50$ chiếc}
{\True $45$ chiếc}
{$35$ chiếc}
\loigiai{
$5 \times 9 = 45$ chiếc bánh.\\
\textbf{Đáp án của câu hỏi là C}
}
\end{ex}

\begin{ex} %[2D8V45]
Tìm $x$ biết: $x \times 2 = 14$.
\choice
{\True $7$}
{$6$}
{$8$}
{$9$}
\loigiai{
$x = 14 : 2 = 7$.\\
\textbf{Đáp án của câu hỏi là A}
}
\end{ex}

\begin{ex} %[2D8V45]
Tìm $y$ biết: $y : 5 = 8$.
\choice
{$35$}
{$45$}
{$30$}
{\True $40$}
\loigiai{
$y = 8 \times 5 = 40$.\\
\textbf{Đáp án của câu hỏi là D}
}
\end{ex}

\begin{ex} %[2D8V45]
Tính giá trị biểu thức: $2 \times 10 - 25 : 5$.
\choice
{$14$}
{$16$}
{\True $15$}
{$17$}
\loigiai{
$20 - 5 = 15$.\\
\textbf{Đáp án của câu hỏi là C}
}
\end{ex}

\begin{ex} %[2D8C45]
Có 4 tổ học sinh, mỗi tổ có 5 bạn. Nếu chuyển 3 bạn đi làm nhiệm vụ thì còn lại bao nhiêu bạn?
\choice
{$16$ bạn}
{\True $17$ bạn}
{$18$ bạn}
{$19$ bạn}
\loigiai{
Tổng số bạn: $5 \times 4 = 20$ bạn.\\
Còn lại: $20 - 3 = 17$ bạn.\\
\textbf{Đáp án của câu hỏi là B}
}
\end{ex}

\textbf{Phần 2. Trắc nghiệm đúng sai}
\setcounterpageref{ex}{0}

\begin{ex} %[2D8H45]
Xét tính đúng sai của các phép toán nhân, chia:
\choiceTF
{\True $2 \times 7 = 14$.}
{\True $14 : 2 = 7$.}
{$5 \times 8 = 45$.}
{\True $5 \times 8 = 40$.}
\loigiai{
\begin{itemize}
    \item $2 \times 7 = 14$: Đúng.
    \item $14 : 2 = 7$: Đúng.
    \item $5 \times 8 = 40 \neq 45$. Ý 3 Sai, Ý 4 Đúng.
\end{itemize}
\begin{itemize}
    \item a) \textbf{Đúng}.
    \item b) \textbf{Đúng}.
    \item c) \textbf{Sai}.
    \item d) \textbf{Đúng}.
\end{itemize}
}
\end{ex}

\begin{ex} %[2D8V45]
Có 30 quyển vở chia đều cho 5 bạn. Xét tính đúng sai:
\choiceTF
{\True Mỗi bạn được 6 quyển vở ($30 : 5 = 6$).}
{\True 2 bạn được tất cả 12 quyển vở ($6 \times 2 = 12$).}
{Nếu có 40 quyển vở thì mỗi bạn được 7 quyển.}
{\True Khi đó mỗi bạn nhận được $40 : 5 = 8$ quyển vở.}
\loigiai{
\begin{itemize}
    \item Mỗi bạn 6 quyển: Đúng.
    \item 2 bạn có 12 quyển: Đúng.
    \item $40 : 5 = 8 \neq 7$. Ý 3 Sai, Ý 4 Đúng.
\end{itemize}
\begin{itemize}
    \item a) \textbf{Đúng}.
    \item b) \textbf{Đúng}.
    \item c) \textbf{Sai}.
    \item d) \textbf{Đúng}.
\end{itemize}
}
\end{ex}

\textbf{Phần 3. Trả lời ngắn}
\setcounterpageref{ex}{0}

\begin{ex} %[2D8H45]
Tính: $2 \times 9$.
\shortans{18}
\loigiai{
\textbf{Đáp số:} $18$\\
$2 \times 9 = 18$.
}
\end{ex}

\begin{ex} %[2D8H45]
Tính: $5 \times 8$.
\shortans{40}
\loigiai{
\textbf{Đáp số:} $40$\\
$5 \times 8 = 40$.
}
\end{ex}

\begin{ex} %[2D8V45]
Tính: $16 : 2$.
\shortans{8}
\loigiai{
\textbf{Đáp số:} $8$\\
$16 : 2 = 8$.
}
\end{ex}

\begin{ex} %[2D8V45]
Tính: $35 : 5$.
\shortans{7}
\loigiai{
\textbf{Đáp số:} $7$\\
$35 : 5 = 7$.
}
\end{ex}

\begin{ex} %[2D8V45]
Tìm $x$ biết: $x \times 5 = 45$.
\shortans{9}
\loigiai{
\textbf{Đáp số:} $9$\\
$x = 45 : 5 = 9$.
}
\end{ex}

\begin{ex} %[2D8C45]
Có 5 túi bi, mỗi túi có 5 viên bi. Sau khi cho bạn 8 viên bi thì còn lại bao nhiêu viên bi?
\shortans{17}
\loigiai{
\textbf{Đáp số:} $17$\\
$5 \times 5 - 8 = 25 - 8 = 17$ viên bi.
}
\end{ex}

\newpage

% =========================================================================
% CHỦ ĐỀ 9 - LÀM QUEN VỚI HÌNH KHỐI
% =========================================================================
\section{Làm quen với hình khối}

\TaoPhieuKiemTra{Chủ đề 9 -- Làm quen với hình khối}{46}{Khối trụ, khối cầu}{https://loptoanthaylinh.vn/thi-online/toan-2/bai-46}
\textbf{Phần 1. Trắc nghiệm lựa chọn}
\setcounterpageref{ex}{0}

\begin{ex} %[2H9N46]
Vật nào sau đây có dạng khối cầu?
\choice
{Hộp sữa bột dạng lon}
{Viên gạch}
{Quyển sách}
{\True Quả bóng đá}
\loigiai{
Quả bóng đá có dạng khối cầu.\\
\textbf{Đáp án của câu hỏi là D}
}
\end{ex}

\begin{ex} %[2H9N46]
Vật nào sau đây có dạng khối trụ?
\choice
{Quả địa cầu}
{\True Lon nước ngọt}
{Hộp phấn viết bảng}
{Quả bóng bàn}
\loigiai{
Lon nước ngọt có dạng khối trụ.\\
\textbf{Đáp án của câu hỏi là B}
}
\end{ex}

\begin{ex} %[2H9H46]
Khối cầu có đặc điểm nào sau đây?
\choice
{Có các mặt phẳng vuông vức}
{Chỉ lăn được về hai phía}
{Có 6 mặt bằng nhau}
{\True Lăn được về mọi phía}
\loigiai{
Khối cầu lăn được về mọi phía.\\
\textbf{Đáp án của câu hỏi là D}
}
\end{ex}

\begin{ex} %[2H9H46]
Khối trụ có 2 đáy là hai hình gì?
\choice
{\True Hình tròn bằng nhau}
{Hình vuông}
{Hình tam giác}
{Hình chữ nhật}
\loigiai{
Khối trụ có 2 đáy là 2 hình tròn bằng nhau.\\
\textbf{Đáp án của câu hỏi là A}
}
\end{ex}

\begin{ex} %[2H9H46]
Khi đặt khối trụ đứng trên mặt bàn, khối trụ có thể
\choice
{Lăn tròn}
{Tự quay tròn}
{\True Trượt hoặc đặt chồng lên nhau}
{Lăn về mọi phía}
\loigiai{
Khi đặt đứng thì khối trụ trượt được và xếp chồng lên nhau được.\\
\textbf{Đáp án của câu hỏi là C}
}
\end{ex}

\begin{ex} %[2H9H46]
Khi đặt khối trụ nằm ngang trên mặt bàn, khối trụ có thể
\choice
{Đặt chồng vững chắc}
{\True Lăn được}
{Không di chuyển được}
{Lăn về mọi phía như khối cầu}
\loigiai{
Khi đặt nằm ngang khối trụ có thể lăn được.\\
\textbf{Đáp án của câu hỏi là B}
}
\end{ex}

\begin{ex} %[2H9V46]
Trong các vật sau: quả cam, hộp chè hình trụ, viên bi, hộp bánh hình lập phương. Có bao nhiêu vật có dạng khối cầu?
\choice
{$1$ vật}
{$3$ vật}
{\True $2$ vật (quả cam, viên bi)}
{$4$ vật}
\loigiai{
Có 2 vật: quả cam và viên bi.\\
\textbf{Đáp án của câu hỏi là C}
}
\end{ex}

\begin{ex} %[2H9V46]
Trong các vật sau: lon sữa đặc, ống bơ đong gạo, khối ru-bích, khúc gỗ tròn. Có bao nhiêu vật có dạng khối trụ?
\choice
{\True $3$ vật (lon sữa, ống bơ, khúc gỗ tròn)}
{$2$ vật}
{$4$ vật}
{$1$ vật}
\loigiai{
Có 3 vật có dạng khối trụ.\\
\textbf{Đáp án của câu hỏi là A}
}
\end{ex}

\begin{ex} %[2H9V46]
Để xếp chồng được lên nhau một cách vững chắc nhất, người ta nên dùng khối nào?
\choice
{Khối cầu}
{Quả bóng}
{Viên bi}
{\True Khối trụ (đặt đứng) hoặc khối lập phương}
\loigiai{
Khối trụ đặt đứng hoặc lập phương có mặt phẳng để xếp chồng.\\
\textbf{Đáp án của câu hỏi là D}
}
\end{ex}

\begin{ex} %[2H9V46]
Bạn An có 4 khối cầu và 5 khối trụ. An có tất cả số khối hình là
\choice
{$8$ khối}
{$10$ khối}
{\True $9$ khối}
{$11$ khối}
\loigiai{
$4 + 5 = 9$ khối.\\
\textbf{Đáp án của câu hỏi là C}
}
\end{ex}

\begin{ex} %[2H9V46]
Vật nào sau đây KHÔNG có dạng khối cầu?
\choice
{Quả bóng rổ}
{\True Chiếc nón lá}
{Hòn bi ve}
{Quả bóng bàn}
\loigiai{
Chiếc nón lá có dạng hình nón (không phải khối cầu).\\
\textbf{Đáp án của câu hỏi là B}
}
\end{ex}

\begin{ex} %[2H9C46]
Xếp 3 khối trụ giống nhau chồng lên nhau theo phương thẳng đứng. Hình thu được vẫn là một
\choice
{\True Khối trụ cao hơn}
{Khối cầu}
{Khối lập phương}
{Khối chữ nhật}
\loigiai{
Xếp chồng các khối trụ cùng đáy tạo thành 1 khối trụ cao hơn.\\
\textbf{Đáp án của câu hỏi là A}
}
\end{ex}

\textbf{Phần 2. Trắc nghiệm đúng sai}
\setcounterpageref{ex}{0}

\begin{ex} %[2H9H46]
Xét tính đúng sai về hình khối:
\choiceTF
{\True Khối cầu không có cạnh và không có đỉnh.}
{\True Khối trụ có hai mặt đáy phẳng hình tròn.}
{Khối cầu có thể xếp chồng lên nhau một cách dễ dàng.}
{\True Khối cầu dễ bị lăn nên rất khó xếp chồng trực tiếp lên nhau.}
\loigiai{
\begin{itemize}
    \item Khối cầu không có cạnh, đỉnh: Đúng.
    \item Khối trụ có 2 đáy tròn: Đúng.
    \item Khối cầu trơn lăn khó xếp chồng. Ý 3 Sai, Ý 4 Đúng.
\end{itemize}
\begin{itemize}
    \item a) \textbf{Đúng}.
    \item b) \textbf{Đúng}.
    \item c) \textbf{Sai}.
    \item d) \textbf{Đúng}.
\end{itemize}
}
\end{ex}

\begin{ex} %[2H9V46]
Cho tập hợp các đồ vật: quả bóng đá, lon cocacola, bình thủy tinh hình trụ, viên bi ve. Xét tính đúng sai:
\choiceTF
{\True Có 2 vật có dạng khối cầu (quả bóng, viên bi).}
{\True Có 2 vật có dạng khối trụ (lon cocacola, bình thủy tinh).}
{Tất cả các vật trên đều có dạng khối cầu.}
{\True Các vật trên gồm 2 khối cầu và 2 khối trụ.}
\loigiai{
\begin{itemize}
    \item 2 khối cầu: Đúng.
    \item 2 khối trụ: Đúng.
    \item Không phải tất cả là khối cầu. Ý 3 Sai, Ý 4 Đúng.
\end{itemize}
\begin{itemize}
    \item a) \textbf{Đúng}.
    \item b) \textbf{Đúng}.
    \item c) \textbf{Sai}.
    \item d) \textbf{Đúng}.
\end{itemize}
}
\end{ex}

\textbf{Phần 3. Trả lời ngắn}
\setcounterpageref{ex}{0}

\begin{ex} %[2H9H46]
Khối trụ có mấy mặt đáy hình tròn?
\shortans{2}
\loigiai{
\textbf{Đáp số:} $2$\\
2 mặt đáy.
}
\end{ex}

\begin{ex} %[2H9H46]
Quả bóng đá có dạng khối gì? (Viết: khối cầu)
\shortans{khối cầu}
\loigiai{
\textbf{Đáp số:} khối cầu\\
Khối cầu.
}
\end{ex}

\begin{ex} %[2H9V46]
Lon nước ngọt dạng hình gì? (Viết: khối trụ)
\shortans{khối trụ}
\loigiai{
\textbf{Đáp số:} khối trụ\\
Khối trụ.
}
\end{ex}

\begin{ex} %[2H9V46]
Có 6 khối cầu và 8 khối trụ. Hỏi có tất cả bao nhiêu khối hình?
\shortans{14}
\loigiai{
\textbf{Đáp số:} $14$\\
$6 + 8 = 14$ khối.
}
\end{ex}

\begin{ex} %[2H9V46]
Có 15 khối trụ, chia đều thành 5 cột bằng nhau. Hỏi mỗi cột có bao nhiêu khối trụ?
\shortans{3}
\loigiai{
\textbf{Đáp số:} $3$\\
$15 : 5 = 3$ khối.
}
\end{ex}

\begin{ex} %[2H9C46]
Bé Nam dùng 2 khối lập phương và 3 khối trụ để xếp một tòa tháp. Hỏi tòa tháp được tạo từ bao nhiêu khối hình?
\shortans{5}
\loigiai{
\textbf{Đáp số:} $5$\\
$2 + 3 = 5$ khối.
}
\end{ex}

\newpage

\TaoPhieuKiemTra{Chủ đề 9 -- Làm quen với hình khối}{47}{Nhận biết hình khối thực tế}{https://loptoanthaylinh.vn/thi-online/toan-2/bai-47}
\textbf{Phần 1. Trắc nghiệm lựa chọn}
\setcounterpageref{ex}{0}

\begin{ex} %[2H9N47]
Khối lập phương có bao nhiêu mặt?
\choice
{$4$ mặt}
{\True $6$ mặt vuông bằng nhau}
{$8$ mặt}
{$12$ mặt}
\loigiai{
Khối lập phương có 6 mặt vuông bằng nhau.\\
\textbf{Đáp án của câu hỏi là B}
}
\end{ex}

\begin{ex} %[2H9N47]
Vật nào sau đây lăn được về mọi phía?
\choice
{Hộp bánh hình chữ nhật}
{Khối rubik}
{\True Quả bóng bàn (khối cầu)}
{Lon sữa bò}
\loigiai{
Quả bóng bàn có dạng khối cầu lăn về mọi phía.\\
\textbf{Đáp án của câu hỏi là C}
}
\end{ex}

\begin{ex} %[2H9H47]
Khối trụ khi đặt nằm ngang thì có thể
\choice
{Không di chuyển}
{\True Lăn được}
{Xếp chồng vững chắc}
{Lăn về mọi hướng}
\loigiai{
Khi nằm ngang thì khối trụ lăn được.\\
\textbf{Đáp án của câu hỏi là B}
}
\end{ex}

\begin{ex} %[2H9H47]
Trong các khối: khối cầu, khối trụ, khối hộp chữ nhật, khối lập phương. Khối nào KHÔNG có mặt phẳng nào?
\choice
{\True Khối cầu}
{Khối trụ}
{Khối lập phương}
{Khối hộp chữ nhật}
\loigiai{
Khối cầu không có mặt phẳng nào.\\
\textbf{Đáp án của câu hỏi là A}
}
\end{ex}

\begin{ex} %[2H9H47]
Để xếp một ngôi nhà đồ chơi gồm phần thân là khối hộp chữ nhật và mái nhà là một khối tam giác, cần mấy khối hình?
\choice
{$1$ khối}
{$3$ khối}
{\True $2$ khối}
{$4$ khối}
\loigiai{
Cần 2 khối hình ghép lại.\\
\textbf{Đáp án của câu hỏi là C}
}
\end{ex}

\begin{ex} %[2H9H47]
Một hộp phấn có dạng hình khối nào?
\choice
{\True Khối hộp chữ nhật hoặc lập phương}
{Khối cầu}
{Khối trụ}
{Hình tròn}
\loigiai{
Hộp phấn có dạng khối hộp chữ nhật hoặc lập phương.\\
\textbf{Đáp án của câu hỏi là A}
}
\end{ex}

\begin{ex} %[2H9V47]
Bác thợ mộc có 8 khúc gỗ hình khối trụ. Bác cưa đôi mỗi khúc gỗ. Hỏi bác thu được bao nhiêu khúc gỗ hình khối trụ nhỏ?
\choice
{$14$ khúc}
{$18$ khúc}
{\True $16$ khúc}
{$12$ khúc}
\loigiai{
$8 \times 2 = 16$ khúc.\\
\textbf{Đáp án của câu hỏi là C}
}
\end{ex}

\begin{ex} %[2H9V47]
Trong một giỏ đồ chơi có 5 quả bóng, 4 hộp sữa bột hình trụ và 3 khối rubik. Tổng số khối hình trong giỏ là
\choice
{$11$ khối}
{$13$ khối}
{$14$ khối}
{\True $12$ khối}
\loigiai{
$5 + 4 + 3 = 12$ khối.\\
\textbf{Đáp án của câu hỏi là D}
}
\end{ex}

\begin{ex} %[2H9V47]
Một thùng có 20 lon nước ngọt xếp đều thành 4 hàng. Mỗi hàng có số lon nước ngọt là
\choice
{$4$ lon}
{\True $5$ lon}
{$6$ lon}
{$7$ lon}
\loigiai{
$20 : 4 = 5$ lon.\\
\textbf{Đáp án của câu hỏi là B}
}
\end{ex}

\begin{ex} %[2H9V47]
Muốn ghép 8 khối lập phương nhỏ cạnh $1\text{ cm}$ thành một khối lập phương lớn hơn, cạnh của khối lập phương lớn là
\choice
{$1\text{ cm}$}
{$3\text{ cm}$}
{$4\text{ cm}$}
{\True $2\text{ cm}$}
\loigiai{
Cạnh khối lập phương lớn: $1 + 1 = 2\text{ cm}$.\\
\textbf{Đáp án của câu hỏi là D}
}
\end{ex}

\begin{ex} %[2H9V47]
Hình nào sau đây có thể vừa lăn được vừa trượt được?
\choice
{\True Khối trụ}
{Khối cầu}
{Khối lập phương}
{Khối hộp chữ nhật}
\loigiai{
Khối trụ đặt đứng thì trượt được, đặt nằm thì lăn được.\\
\textbf{Đáp án của câu hỏi là A}
}
\end{ex}

\begin{ex} %[2H9C47]
Có 12 khối lập phương nhỏ. Có thể xếp thành một khối hộp chữ nhật gồm mấy tầng, mỗi tầng 4 khối?
\choice
{$2$ tầng}
{$4$ tầng}
{$5$ tầng}
{\True $3$ tầng}
\loigiai{
$12 : 4 = 3$ tầng.\\
\textbf{Đáp án của câu hỏi là D}
}
\end{ex}

\textbf{Phần 2. Trắc nghiệm đúng sai}
\setcounterpageref{ex}{0}

\begin{ex} %[2H9H47]
Xét tính đúng sai về các hình khối đã học:
\choiceTF
{\True Khối cầu lăn được về mọi phía.}
{\True Khối trụ có 2 mặt phẳng đối diện là hai hình tròn bằng nhau.}
{Khối lập phương có các mặt là hình chữ nhật.}
{\True Khối lập phương có 6 mặt đều là hình vuông bằng nhau.}
\loigiai{
\begin{itemize}
    \item Khối cầu lăn mọi phía: Đúng.
    \item Khối trụ có 2 mặt đáy tròn: Đúng.
    \item Các mặt của khối lập phương là hình vuông. Ý 3 Sai, Ý 4 Đúng.
\end{itemize}
\begin{itemize}
    \item a) \textbf{Đúng}.
    \item b) \textbf{Đúng}.
    \item c) \textbf{Sai}.
    \item d) \textbf{Đúng}.
\end{itemize}
}
\end{ex}

\begin{ex} %[2H9V47]
Bạn Mai có 6 khối cầu và 9 khối trụ. Xét tính đúng sai:
\choiceTF
{\True Tổng số khối hình Mai có là 15 khối ($6 + 9 = 15$).}
{\True Số khối trụ nhiều hơn số khối cầu là 3 khối ($9 - 6 = 3$).}
{Mai có số khối cầu nhiều hơn số khối trụ.}
{\True Mai có số khối cầu ít hơn số khối trụ là 3 khối.}
\loigiai{
\begin{itemize}
    \item Tổng: $6 + 9 = 15$: Đúng.
    \item Trụ hơn cầu: $9 - 6 = 3$: Đúng.
    \item Cầu ít hơn trụ. Ý 3 Sai, Ý 4 Đúng.
\end{itemize}
\begin{itemize}
    \item a) \textbf{Đúng}.
    \item b) \textbf{Đúng}.
    \item c) \textbf{Sai}.
    \item d) \textbf{Đúng}.
\end{itemize}
}
\end{ex}

\textbf{Phần 3. Trả lời ngắn}
\setcounterpageref{ex}{0}

\begin{ex} %[2H9H47]
Khối lập phương có bao nhiêu mặt?
\shortans{6}
\loigiai{
\textbf{Đáp số:} $6$\\
6 mặt.
}
\end{ex}

\begin{ex} %[2H9H47]
Khối trụ có bao nhiêu mặt đáy phẳng hình tròn?
\shortans{2}
\loigiai{
\textbf{Đáp số:} $2$\\
2 mặt đáy.
}
\end{ex}

\begin{ex} %[2H9V47]
Có 7 khối cầu và 8 khối trụ. Hỏi có tất cả bao nhiêu khối hình?
\shortans{15}
\loigiai{
\textbf{Đáp số:} $15$\\
$7 + 8 = 15$ khối.
}
\end{ex}

\begin{ex} %[2H9V47]
Một hộp có 18 khối gỗ gồm khối trụ và khối cầu, trong đó có 9 khối cầu. Hỏi có bao nhiêu khối trụ?
\shortans{9}
\loigiai{
\textbf{Đáp số:} $9$\\
$18 - 9 = 9$ khối.
}
\end{ex}

\begin{ex} %[2H9V47]
Có 20 khối trụ chia đều vào 5 hộp. Hỏi mỗi hộp có bao nhiêu khối trụ?
\shortans{4}
\loigiai{
\textbf{Đáp số:} $4$\\
$20 : 5 = 4$ khối.
}
\end{ex}

\begin{ex} %[2H9C47]
Dùng các khối lập phương nhỏ xếp thành một khối hộp chữ nhật có 3 tầng, mỗi tầng có 6 khối nhỏ. Hỏi cần tất cả bao nhiêu khối lập phương nhỏ?
\shortans{18}
\loigiai{
\textbf{Đáp số:} $18$\\
$6 \times 3 = 18$ khối.
}
\end{ex}

\newpage

% =========================================================================
% CHỦ ĐỀ 10 - CÁC SỐ TRONG PHẠM VI 1 000
% =========================================================================
\section{Các số trong phạm vi 1 000}

\TaoPhieuKiemTra{Chủ đề 10 -- Các số trong phạm vi 1 000}{48}{Đơn vị đo: Đề-xi-mét}{https://loptoanthaylinh.vn/thi-online/toan-2/bai-48}
\textbf{Phần 1. Trắc nghiệm lựa chọn}
\setcounterpageref{ex}{0}

\begin{ex} %[2D10N48]
10 đơn vị bằng mấy chục?
\choice
{$2$ chục}
{$10$ chục}
{$1$ trăm}
{\True $1$ chục}
\loigiai{
10 đơn vị = 1 chục.\\
\textbf{Đáp án của câu hỏi là D}
}
\end{ex}

\begin{ex} %[2D10N48]
10 chục bằng mấy trăm?
\choice
{$10$ trăm}
{\True $1$ trăm}
{$1$ nghìn}
{$2$ trăm}
\loigiai{
10 chục = 1 trăm (100).\\
\textbf{Đáp án của câu hỏi là B}
}
\end{ex}

\begin{ex} %[2D10H48]
10 trăm bằng mấy nghìn?
\choice
{$10$ nghìn}
{$100$ nghìn}
{\True $1$ nghìn ($1\,000$)}
{$2$ nghìn}
\loigiai{
10 trăm = 1 nghìn ($1000$).\\
\textbf{Đáp án của câu hỏi là C}
}
\end{ex}

\begin{ex} %[2D10H48]
Số một nghìn được viết là
\choice
{$100$}
{$10\,000$}
{$10$}
{\True $1\,000$}
\loigiai{
Một nghìn viết là $1000$.\\
\textbf{Đáp án của câu hỏi là D}
}
\end{ex}

\begin{ex} %[2D10H48]
Số 100 gồm mấy chữ số?
\choice
{\True $3$ chữ số}
{$2$ chữ số}
{$4$ chữ số}
{$1$ chữ số}
\loigiai{
Số 100 có 3 chữ số.\\
\textbf{Đáp án của câu hỏi là A}
}
\end{ex}

\begin{ex} %[2D10H48]
Số $1\,000$ gồm mấy chữ số?
\choice
{$3$ chữ số}
{\True $4$ chữ số}
{$5$ chữ số}
{$2$ chữ số}
\loigiai{
Số 1000 có 4 chữ số (gồm chữ số 1 và ba chữ số 0).\\
\textbf{Đáp án của câu hỏi là B}
}
\end{ex}

\begin{ex} %[2D10V48]
Số gồm 4 trăm, 0 chục và 0 đơn vị viết là
\choice
{\True $400$}
{$40$}
{$4000$}
{$4$}
\loigiai{
Số đó là $400$.\\
\textbf{Đáp án của câu hỏi là A}
}
\end{ex}

\begin{ex} %[2D10V48]
Số gồm 1 nghìn, 0 trăm, 0 chục, 0 đơn vị là
\choice
{$100$}
{$10$}
{$10\,000$}
{\True $1\,000$}
\loigiai{
Số đó là $1000$.\\
\textbf{Đáp án của câu hỏi là D}
}
\end{ex}

\begin{ex} %[2D10V48]
Điền số thích hợp: 30 chục = $\dots$ trăm.
\choice
{$2$}
{$4$}
{\True $3$}
{$30$}
\loigiai{
30 chục = 3 trăm (300).\\
\textbf{Đáp án của câu hỏi là C}
}
\end{ex}

\begin{ex} %[2D10V48]
Điền số thích hợp: 50 trăm = $\dots$ nghìn.
\choice
{$4$}
{\True $5$}
{$6$}
{$50$}
\loigiai{
50 trăm = 5 nghìn.\\
\textbf{Đáp án của câu hỏi là B}
}
\end{ex}

\begin{ex} %[2D10V48]
Một hộp có 100 chiếc khẩu trang. 10 hộp như thế có tất cả số chiếc khẩu trang là
\choice
{\True $1\,000$ chiếc}
{$100$ chiếc}
{$500$ chiếc}
{$10\,000$ chiếc}
\loigiai{
$100 \times 10 = 1000$ chiếc.\\
\textbf{Đáp án của câu hỏi là A}
}
\end{ex}

\begin{ex} %[2D10C48]
Số bé nhất có bốn chữ số là
\choice
{$999$}
{$1001$}
{\True $1\,000$}
{$1111$}
\loigiai{
Số bé nhất có 4 chữ số là $1000$.\\
\textbf{Đáp án của câu hỏi là C}
}
\end{ex}

\textbf{Phần 2. Trắc nghiệm đúng sai}
\setcounterpageref{ex}{0}

\begin{ex} %[2D10H48]
Xét mối quan hệ giữa các hàng đơn vị, chục, trăm, nghìn:
\choiceTF
{\True 10 đơn vị = 1 chục.}
{\True 10 chục = 1 trăm.}
{10 trăm = 10 nghìn.}
{\True 10 trăm = 1 nghìn ($1\,000$).}
\loigiai{
\begin{itemize}
    \item 10 đơn vị = 1 chục: Đúng.
    \item 10 chục = 1 trăm: Đúng.
    \item 10 trăm = 1 nghìn (không phải 10 nghìn). Ý 3 Sai, Ý 4 Đúng.
\end{itemize}
\begin{itemize}
    \item a) \textbf{Đúng}.
    \item b) \textbf{Đúng}.
    \item c) \textbf{Sai}.
    \item d) \textbf{Đúng}.
\end{itemize}
}
\end{ex}

\begin{ex} %[2D10V48]
Cho số 1 000. Xét tính đúng sai:
\choiceTF
{\True 1 000 là số có bốn chữ số.}
{\True 1 000 là số liền sau của số 999.}
{1 000 là số lớn nhất có 3 chữ số.}
{\True 1 000 là số bé nhất có 4 chữ số.}
\loigiai{
\begin{itemize}
    \item 1000 có 4 chữ số: Đúng.
    \item Liền sau của 999: Đúng.
    \item 1000 là số có 4 chữ số (không phải 3 chữ số). Ý 3 Sai, Ý 4 Đúng.
\end{itemize}
\begin{itemize}
    \item a) \textbf{Đúng}.
    \item b) \textbf{Đúng}.
    \item c) \textbf{Sai}.
    \item d) \textbf{Đúng}.
\end{itemize}
}
\end{ex}

\textbf{Phần 3. Trả lời ngắn}
\setcounterpageref{ex}{0}

\begin{ex} %[2D10H48]
10 chục bằng bao nhiêu trăm?
\shortans{1}
\loigiai{
\textbf{Đáp số:} $1$\\
1 trăm.
}
\end{ex}

\begin{ex} %[2D10H48]
10 trăm bằng bao nhiêu nghìn?
\shortans{1}
\loigiai{
\textbf{Đáp số:} $1$\\
1 nghìn.
}
\end{ex}

\begin{ex} %[2D10V48]
Viết số một nghìn bằng chữ số.
\shortans{1000}
\loigiai{
\textbf{Đáp số:} $1000$\\
1000.
}
\end{ex}

\begin{ex} %[2D10V48]
Tìm số liền trước của số $1\,000$.
\shortans{999}
\loigiai{
\textbf{Đáp số:} $999$\\
$1000 - 1 = 999$.
}
\end{ex}

\begin{ex} %[2D10V48]
Tìm số liền sau của số $999$.
\shortans{1000}
\loigiai{
\textbf{Đáp số:} $1000$\\
$999 + 1 = 1000$.
}
\end{ex}

\begin{ex} %[2D10C48]
Có bao nhiêu chữ số 0 trong số một nghìn ($1\,000$)?
\shortans{3}
\loigiai{
\textbf{Đáp số:} $3$\\
Có 3 chữ số 0.
}
\end{ex}

\newpage

\TaoPhieuKiemTra{Chủ đề 10 -- Các số trong phạm vi 1 000}{49}{Đơn vị đo: Mét}{https://loptoanthaylinh.vn/thi-online/toan-2/bai-49}
\textbf{Phần 1. Trắc nghiệm lựa chọn}
\setcounterpageref{ex}{0}

\begin{ex} %[2D10N49]
Số nào sau đây là số tròn trăm?
\choice
{$150$}
{$205$}
{\True $200$}
{$99$}
\loigiai{
Số tròn trăm tận cùng là hai chữ số 0: 200.\\
\textbf{Đáp án của câu hỏi là C}
}
\end{ex}

\begin{ex} %[2D10N49]
Số nào sau đây là số tròn chục?
\choice
{\True $320$}
{$325$}
{$305$}
{$321$}
\loigiai{
Số tròn chục tận cùng là chữ số 0: 320.\\
\textbf{Đáp án của câu hỏi là A}
}
\end{ex}

\begin{ex} %[2D10H49]
Số tròn trăm liền sau số 400 là
\choice
{$401$}
{\True $500$}
{$300$}
{$410$}
\loigiai{
Số tròn trăm liền sau của 400 là 500.\\
\textbf{Đáp án của câu hỏi là B}
}
\end{ex}

\begin{ex} %[2D10H49]
Số tròn chục liền trước số 150 là
\choice
{$149$}
{$160$}
{\True $140$}
{$100$}
\loigiai{
Số tròn chục liền trước là $150 - 10 = 140$.\\
\textbf{Đáp án của câu hỏi là C}
}
\end{ex}

\begin{ex} %[2D10H49]
Các số tròn trăm từ 100 đến 500 là
\choice
{$100, 200, 400, 500$}
{$100, 150, 200, 250$}
{$200, 300, 400, 500$}
{\True $100, 200, 300, 400, 500$}
\loigiai{
Gồm 5 số: 100, 200, 300, 400, 500.\\
\textbf{Đáp án của câu hỏi là D}
}
\end{ex}

\begin{ex} %[2D10H49]
Tính: $300 + 200 = \dots$
\choice
{\True $500$}
{$400$}
{$600$}
{$50$}
\loigiai{
$300 + 200 = 500$.\\
\textbf{Đáp án của câu hỏi là A}
}
\end{ex}

\begin{ex} %[2D10V49]
Tính: $800 - 300 = \dots$
\choice
{$400$}
{$600$}
{\True $500$}
{$50$}
\loigiai{
$800 - 300 = 500$.\\
\textbf{Đáp án của câu hỏi là C}
}
\end{ex}

\begin{ex} %[2D10V49]
Tính: $200 + 400 + 100 = \dots$
\choice
{\True $700$}
{$600$}
{$800$}
{$900$}
\loigiai{
$200 + 400 + 100 = 700$.\\
\textbf{Đáp án của câu hỏi là A}
}
\end{ex}

\begin{ex} %[2D10V49]
Số tròn chục lớn nhất có ba chữ số là
\choice
{$900$}
{\True $990$}
{$999$}
{$1000$}
\loigiai{
Số tròn chục lớn nhất có 3 chữ số là 990.\\
\textbf{Đáp án của câu hỏi là B}
}
\end{ex}

\begin{ex} %[2D10V49]
Số tròn trăm lớn nhất có ba chữ số là
\choice
{$990$}
{$999$}
{$1000$}
{\True $900$}
\loigiai{
Số tròn trăm lớn nhất có 3 chữ số là 900.\\
\textbf{Đáp án của câu hỏi là D}
}
\end{ex}

\begin{ex} %[2D10V49]
Có bao nhiêu số tròn trăm có ba chữ số?
\choice
{$8$ số}
{\True $9$ số (từ $100$ đến $900$)}
{$10$ số}
{$11$ số}
\loigiai{
Từ 100 đến 900 có 9 số tròn trăm.\\
\textbf{Đáp án của câu hỏi là B}
}
\end{ex}

\begin{ex} %[2D10C49]
Tìm một số tròn trăm biết rằng lấy số đó cộng với 300 rồi trừ đi 200 thì được 600.
\choice
{$400$}
{$600$}
{$700$}
{\True $500$}
\loigiai{
$600 + 200 - 300 = 500$.\\
\textbf{Đáp án của câu hỏi là D}
}
\end{ex}

\textbf{Phần 2. Trắc nghiệm đúng sai}
\setcounterpageref{ex}{0}

\begin{ex} %[2D10H49]
Xét tính đúng sai về số tròn trăm:
\choiceTF
{\True Số tròn trăm có chữ số hàng chục và hàng đơn vị đều là 0.}
{\True $100$ là số tròn trăm bé nhất có ba chữ số.}
{$990$ là số tròn trăm lớn nhất có ba chữ số.}
{\True $900$ là số tròn trăm lớn nhất có ba chữ số.}
\loigiai{
\begin{itemize}
    \item Tận cùng là 00: Đúng.
    \item 100 là số tròn trăm bé nhất: Đúng.
    \item 990 là số tròn chục (không phải tròn trăm). Ý 3 Sai, Ý 4 Đúng.
\end{itemize}
\begin{itemize}
    \item a) \textbf{Đúng}.
    \item b) \textbf{Đúng}.
    \item c) \textbf{Sai}.
    \item d) \textbf{Đúng}.
\end{itemize}
}
\end{ex}

\begin{ex} %[2D10V49]
Cho các phép tính với số tròn trăm:
\choiceTF
{\True $400 + 300 = 700$.}
{\True $900 - 500 = 400$.}
{$200 + 600 = 900$.}
{\True $200 + 600 = 800$.}
\loigiai{
\begin{itemize}
    \item $400 + 300 = 700$: Đúng.
    \item $900 - 500 = 400$: Đúng.
    \item $200 + 600 = 800 \neq 900$. Ý 3 Sai, Ý 4 Đúng.
\end{itemize}
\begin{itemize}
    \item a) \textbf{Đúng}.
    \item b) \textbf{Đúng}.
    \item c) \textbf{Sai}.
    \item d) \textbf{Đúng}.
\end{itemize}
}
\end{ex}

\textbf{Phần 3. Trả lời ngắn}
\setcounterpageref{ex}{0}

\begin{ex} %[2D10H49]
Tính: $400 + 500$.
\shortans{900}
\loigiai{
\textbf{Đáp số:} $900$\\
$400 + 500 = 900$.
}
\end{ex}

\begin{ex} %[2D10H49]
Tính: $700 - 300$.
\shortans{400}
\loigiai{
\textbf{Đáp số:} $400$\\
$700 - 300 = 400$.
}
\end{ex}

\begin{ex} %[2D10V49]
Tìm số tròn trăm liền trước số $600$.
\shortans{500}
\loigiai{
\textbf{Đáp số:} $500$\\
$600 - 100 = 500$.
}
\end{ex}

\begin{ex} %[2D10V49]
Tìm số tròn trăm liền sau số $700$.
\shortans{800}
\loigiai{
\textbf{Đáp số:} $800$\\
$700 + 100 = 800$.
}
\end{ex}

\begin{ex} %[2D10V49]
Tính giá trị biểu thức: $1000 - 400 - 200$.
\shortans{400}
\loigiai{
\textbf{Đáp số:} $400$\\
$600 - 200 = 400$.
}
\end{ex}

\begin{ex} %[2D10C49]
Có bao nhiêu số tròn trăm nằm giữa số 200 và số 700? (Gồm các số: 300, 400, 500, 600)
\shortans{4}
\loigiai{
\textbf{Đáp số:} $4$\\
Có 4 số.
}
\end{ex}

\newpage

\TaoPhieuKiemTra{Chủ đề 10 -- Các số trong phạm vi 1 000}{50}{Đơn vị đo: Ki-lô-mét}{https://loptoanthaylinh.vn/thi-online/toan-2/bai-50}
\textbf{Phần 1. Trắc nghiệm lựa chọn}
\setcounterpageref{ex}{0}

\begin{ex} %[2D10N50]
Điền dấu thích hợp: $300 \dots 500$.
\choice
{$>$}
{$=$}
{Không so sánh được}
{\True $<$}
\loigiai{
$300 < 500$.\\
\textbf{Đáp án của câu hỏi là D}
}
\end{ex}

\begin{ex} %[2D10N50]
Điền dấu thích hợp: $700 \dots 600$.
\choice
{$<$}
{$=$}
{\True $>$}
{Không so sánh được}
\loigiai{
$700 > 600$.\\
\textbf{Đáp án của câu hỏi là C}
}
\end{ex}

\begin{ex} %[2D10H50]
Điền dấu thích hợp: $400 + 200 \dots 600$.
\choice
{\True $=$}
{$>$}
{$<$}
{Không so sánh được}
\loigiai{
$400 + 200 = 600$.\\
\textbf{Đáp án của câu hỏi là A}
}
\end{ex}

\begin{ex} %[2D10H50]
Trong các số $300, 700, 500, 900$, số lớn nhất là
\choice
{$700$}
{\True $900$}
{$500$}
{$300$}
\loigiai{
Số lớn nhất là 900.\\
\textbf{Đáp án của câu hỏi là B}
}
\end{ex}

\begin{ex} %[2D10H50]
Trong các số $400, 200, 800, 600$, số bé nhất là
\choice
{$400$}
{$600$}
{\True $200$}
{$800$}
\loigiai{
Số bé nhất là 200.\\
\textbf{Đáp án của câu hỏi là C}
}
\end{ex}

\begin{ex} %[2D10H50]
Sắp xếp các số $500, 200, 800, 300$ theo thứ tự từ bé đến lớn:
\choice
{$800, 500, 300, 200$}
{$200, 500, 300, 800$}
{$300, 200, 500, 800$}
{\True $200, 300, 500, 800$}
\loigiai{
Thứ tự từ bé đến lớn: $200 < 300 < 500 < 800$.\\
\textbf{Đáp án của câu hỏi là D}
}
\end{ex}

\begin{ex} %[2D10V50]
Điền dấu thích hợp: $800 - 300 \dots 400 + 200$.
\choice
{$>$}
{\True $<$}
{$=$}
{Không so sánh được}
\loigiai{
$800 - 300 = 500 < 400 + 200 = 600$.\\
\textbf{Đáp án của câu hỏi là B}
}
\end{ex}

\begin{ex} %[2D10V50]
Tìm số tròn trăm $x$ biết: $400 < x < 600$.
\choice
{$450$}
{$550$}
{$600$}
{\True $500$}
\loigiai{
Số tròn trăm thỏa mãn là 500.\\
\textbf{Đáp án của câu hỏi là D}
}
\end{ex}

\begin{ex} %[2D10V50]
Sắp xếp các số tròn trăm $900, 400, 700, 100$ theo thứ tự giảm dần:
\choice
{$100, 400, 700, 900$}
{$900, 400, 700, 100$}
{\True $900, 700, 400, 100$}
{$700, 900, 400, 100$}
\loigiai{
Thứ tự giảm dần: $900 > 700 > 400 > 100$.\\
\textbf{Đáp án của câu hỏi là C}
}
\end{ex}

\begin{ex} %[2D10V50]
Có bao nhiêu số tròn trăm lớn hơn 300 và bé hơn 800?
\choice
{\True $4$ số ($400, 500, 600, 700$)}
{$3$ số}
{$5$ số}
{$6$ số}
\loigiai{
Các số: 400, 500, 600, 700 (có 4 số).\\
\textbf{Đáp án của câu hỏi là A}
}
\end{ex}

\begin{ex} %[2D10V50]
Điền dấu thích hợp: $1\,000 \dots 900 + 100$.
\choice
{$>$}
{\True $=$}
{$<$}
{Không so sánh được}
\loigiai{
$900 + 100 = 1000$.\\
\textbf{Đáp án của câu hỏi là B}
}
\end{ex}

\begin{ex} %[2D10C50]
Tìm số tròn trăm lớn nhất có ba chữ số bé hơn tổng $500 + 300$.
\choice
{\True $700$}
{$600$}
{$800$}
{$900$}
\loigiai{
$500 + 300 = 800$. Số tròn trăm lớn nhất bé hơn 800 là 700.\\
\textbf{Đáp án của câu hỏi là A}
}
\end{ex}

\textbf{Phần 2. Trắc nghiệm đúng sai}
\setcounterpageref{ex}{0}

\begin{ex} %[2D10H50]
Xét tính đúng sai của các phép so sánh:
\choiceTF
{\True $600 > 400$.}
{\True $300 < 800$.}
{$500 = 50$.}
{\True $500 > 50$.}
\loigiai{
\begin{itemize}
    \item $600 > 400$: Đúng.
    \item $300 < 800$: Đúng.
    \item $500 > 50 \neq 50$. Ý 3 Sai, Ý 4 Đúng.
\end{itemize}
\begin{itemize}
    \item a) \textbf{Đúng}.
    \item b) \textbf{Đúng}.
    \item c) \textbf{Sai}.
    \item d) \textbf{Đúng}.
\end{itemize}
}
\end{ex}

\begin{ex} %[2D10V50]
Cho 4 số: $200, 500, 700, 900$. Xét tính đúng sai:
\choiceTF
{\True Số lớn nhất là 900.}
{\True Số bé nhất là 200.}
{Tổng số bé nhất và số lớn nhất là 1 000.}
{\True Tổng số bé nhất và số lớn nhất là $200 + 900 = 1\,100$.}
\loigiai{
\begin{itemize}
    \item Lớn nhất là 900: Đúng.
    \item Bé nhất là 200: Đúng.
    \item $200 + 900 = 1100 \neq 1000$. Ý 3 Sai, Ý 4 Đúng.
\end{itemize}
\begin{itemize}
    \item a) \textbf{Đúng}.
    \item b) \textbf{Đúng}.
    \item c) \textbf{Sai}.
    \item d) \textbf{Đúng}.
\end{itemize}
}
\end{ex}

\textbf{Phần 3. Trả lời ngắn}
\setcounterpageref{ex}{0}

\begin{ex} %[2D10H50]
Trong các số $400, 700, 200, 600$, tìm số lớn nhất.
\shortans{700}
\loigiai{
\textbf{Đáp số:} $700$\\
700.
}
\end{ex}

\begin{ex} %[2D10H50]
Trong các số $500, 800, 300, 900$, tìm số bé nhất.
\shortans{300}
\loigiai{
\textbf{Đáp số:} $300$\\
300.
}
\end{ex}

\begin{ex} %[2D10V50]
Tìm số tròn trăm $x$ biết $600 < x < 800$.
\shortans{700}
\loigiai{
\textbf{Đáp số:} $700$\\
700.
}
\end{ex}

\begin{ex} %[2D10V50]
Có bao nhiêu số tròn trăm nằm giữa 100 và 500?
\shortans{3}
\loigiai{
\textbf{Đáp số:} $3$\\
Gồm 200, 300, 400 (có 3 số).
}
\end{ex}

\begin{ex} %[2D10V50]
Tìm hiệu giữa số tròn trăm lớn nhất có 3 chữ số và số tròn trăm bé nhất có 3 chữ số.
\shortans{800}
\loigiai{
\textbf{Đáp số:} $800$\\
$900 - 100 = 800$.
}
\end{ex}

\begin{ex} %[2D10C50]
Tìm tổng của tất cả các số tròn trăm có ba chữ số lớn hơn 600 và nhỏ hơn 1 000. (Gồm 700, 800, 900)
\shortans{2400}
\loigiai{
\textbf{Đáp số:} $2400$\\
$700 + 800 + 900 = 2400$.
}
\end{ex}

\newpage

\TaoPhieuKiemTra{Chủ đề 10 -- Các số trong phạm vi 1 000}{51}{Đơn vị đo: Mi-li-mét}{https://loptoanthaylinh.vn/thi-online/toan-2/bai-51}
\textbf{Phần 1. Trắc nghiệm lựa chọn}
\setcounterpageref{ex}{0}

\begin{ex} %[2D10N51]
Số gồm $3$ trăm, $4$ chục và $5$ đơn vị viết là
\choice
{$354$}
{$543$}
{\True $345$}
{$3405$}
\loigiai{
Số đó là $345$.\\
\textbf{Đáp án của câu hỏi là C}
}
\end{ex}

\begin{ex} %[2D10N51]
Số $528$ đọc là
\choice
{\True Năm trăm hai mươi tám}
{Năm trăm hai mươi}
{Năm trăm hai tám}
{Năm hai tám}
\loigiai{
Đọc là Năm trăm hai mươi tám.\\
\textbf{Đáp án của câu hỏi là A}
}
\end{ex}

\begin{ex} %[2D10H51]
Trong số $764$, chữ số 7 ở hàng nào?
\choice
{Hàng đơn vị}
{\True Hàng trăm}
{Hàng chục}
{Hàng nghìn}
\loigiai{
Chữ số 7 ở hàng trăm.\\
\textbf{Đáp án của câu hỏi là B}
}
\end{ex}

\begin{ex} %[2D10H51]
Trong số $839$, giá trị của chữ số 3 là
\choice
{$3$}
{$300$}
{$39$}
{\True $30$}
\loigiai{
Chữ số 3 ở hàng chục có giá trị là 30.\\
\textbf{Đáp án của câu hỏi là D}
}
\end{ex}

\begin{ex} %[2D10H51]
Số "Bảy trăm linh năm" viết là
\choice
{$750$}
{\True $705$}
{$75$}
{$7005$}
\loigiai{
Viết là $705$.\\
\textbf{Đáp án của câu hỏi là B}
}
\end{ex}

\begin{ex} %[2D10H51]
Số liền sau của số $459$ là
\choice
{$458$}
{$461$}
{$470$}
{\True $460$}
\loigiai{
$459 + 1 = 460$.\\
\textbf{Đáp án của câu hỏi là D}
}
\end{ex}

\begin{ex} %[2D10V51]
Số liền trước của số $600$ là
\choice
{$601$}
{$590$}
{\True $599$}
{$500$}
\loigiai{
$600 - 1 = 599$.\\
\textbf{Đáp án của câu hỏi là C}
}
\end{ex}

\begin{ex} %[2D10V51]
Số bé nhất có ba chữ số khác nhau là
\choice
{$100$}
{\True $102$}
{$120$}
{$123$}
\loigiai{
Số bé nhất có 3 chữ số khác nhau là 102.\\
\textbf{Đáp án của câu hỏi là B}
}
\end{ex}

\begin{ex} %[2D10V51]
Số lớn nhất có ba chữ số khác nhau là
\choice
{\True $987$}
{$999$}
{$989$}
{$897$}
\loigiai{
Số lớn nhất có 3 chữ số khác nhau là 987.\\
\textbf{Đáp án của câu hỏi là A}
}
\end{ex}

\begin{ex} %[2D10V51]
Từ 3 chữ số $1, 2, 3$ lập được bao nhiêu số có ba chữ số khác nhau?
\choice
{$4$ số}
{$8$ số}
{\True $6$ số}
{$3$ số}
\loigiai{
Các số: 123, 132, 213, 231, 312, 321 (có 6 số).\\
\textbf{Đáp án của câu hỏi là C}
}
\end{ex}

\begin{ex} %[2D10V51]
Số gồm 6 trăm và 8 đơn vị viết là
\choice
{$68$}
{$680$}
{$6008$}
{\True $608$}
\loigiai{
Số đó là $608$.\\
\textbf{Đáp án của câu hỏi là D}
}
\end{ex}

\begin{ex} %[2D10C51]
Tìm số có ba chữ số biết chữ số hàng trăm là 4, chữ số hàng chục gấp đôi chữ số hàng trăm và chữ số hàng đơn vị bằng 0.
\choice
{\True $480$}
{$420$}
{$440$}
{$408$}
\loigiai{
Hàng trăm là 4. Hàng chục là $4 \times 2 = 8$. Hàng đơn vị là 0. Số đó là $480$.\\
\textbf{Đáp án của câu hỏi là A}
}
\end{ex}

\textbf{Phần 2. Trắc nghiệm đúng sai}
\setcounterpageref{ex}{0}

\begin{ex} %[2D10H51]
Xét tính đúng sai về cấu tạo số có ba chữ số:
\choiceTF
{\True Số 100 là số bé nhất có ba chữ số.}
{\True Số 999 là số lớn nhất có ba chữ số.}
{Số 305 đọc là Ba trăm năm mươi.}
{\True Số 305 đọc là Ba trăm linh năm.}
\loigiai{
\begin{itemize}
    \item Bé nhất có 3 chữ số là 100: Đúng.
    \item Lớn nhất có 3 chữ số là 999: Đúng.
    \item 305 đọc là Ba trăm linh năm (không phải năm mươi). Ý 3 Sai, Ý 4 Đúng.
\end{itemize}
\begin{itemize}
    \item a) \textbf{Đúng}.
    \item b) \textbf{Đúng}.
    \item c) \textbf{Sai}.
    \item d) \textbf{Đúng}.
\end{itemize}
}
\end{ex}

\begin{ex} %[2D10V51]
Cho số 746. Xét tính đúng sai:
\choiceTF
{\True Chữ số 7 có giá trị là 700.}
{\True Chữ số 4 có giá trị là 40.}
{Chữ số 6 có giá trị là 60.}
{\True Chữ số 6 ở hàng đơn vị nên có giá trị là 6.}
\loigiai{
\begin{itemize}
    \item Chữ số 7 là 700: Đúng.
    \item Chữ số 4 là 40: Đúng.
    \item Chữ số 6 ở hàng đơn vị giá trị 6 (không phải 60). Ý 3 Sai, Ý 4 Đúng.
\end{itemize}
\begin{itemize}
    \item a) \textbf{Đúng}.
    \item b) \textbf{Đúng}.
    \item c) \textbf{Sai}.
    \item d) \textbf{Đúng}.
\end{itemize}
}
\end{ex}

\textbf{Phần 3. Trả lời ngắn}
\setcounterpageref{ex}{0}

\begin{ex} %[2D10H51]
Viết số gồm 5 trăm, 7 chục và 2 đơn vị.
\shortans{572}
\loigiai{
\textbf{Đáp số:} $572$\\
572.
}
\end{ex}

\begin{ex} %[2D10H51]
Viết số gồm 8 trăm và 4 chục.
\shortans{840}
\loigiai{
\textbf{Đáp số:} $840$\\
840.
}
\end{ex}

\begin{ex} %[2D10V51]
Tìm số liền trước của số $350$.
\shortans{349}
\loigiai{
\textbf{Đáp số:} $349$\\
$350 - 1 = 349$.
}
\end{ex}

\begin{ex} %[2D10V51]
Tìm số liền sau của số $789$.
\shortans{790}
\loigiai{
\textbf{Đáp số:} $790$\\
$789 + 1 = 790$.
}
\end{ex}

\begin{ex} %[2D10V51]
Tìm số lớn nhất có ba chữ số khác nhau.
\shortans{987}
\loigiai{
\textbf{Đáp số:} $987$\\
987.
}
\end{ex}

\begin{ex} %[2D10C51]
Tìm tổng các chữ số của số bé nhất có ba chữ số khác nhau ($102$).
\shortans{3}
\loigiai{
\textbf{Đáp số:} $3$\\
$1 + 0 + 2 = 3$.
}
\end{ex}

\newpage

\TaoPhieuKiemTra{Chủ đề 10 -- Các số trong phạm vi 1 000}{52}{Thực hành đo và ước lượng độ dài}{https://loptoanthaylinh.vn/thi-online/toan-2/bai-52}
\textbf{Phần 1. Trắc nghiệm lựa chọn}
\setcounterpageref{ex}{0}

\begin{ex} %[2D10N52]
Số $456$ viết thành tổng các trăm, chục, đơn vị là
\choice
{$40 + 50 + 6$}
{$400 + 5 + 6$}
{\True $400 + 50 + 6$}
{$4 + 5 + 6$}
\loigiai{
$456 = 400 + 50 + 6$.\\
\textbf{Đáp án của câu hỏi là C}
}
\end{ex}

\begin{ex} %[2D10N52]
Số $708$ viết thành tổng là
\choice
{\True $700 + 8$}
{$700 + 80$}
{$70 + 8$}
{$700 + 80 + 0$}
\loigiai{
$708 = 700 + 8$.\\
\textbf{Đáp án của câu hỏi là A}
}
\end{ex}

\begin{ex} %[2D10H52]
Tổng $500 + 60 + 3$ là dạng phân tích của số nào?
\choice
{$536$}
{$653$}
{$5063$}
{\True $563$}
\loigiai{
$500 + 60 + 3 = 563$.\\
\textbf{Đáp án của câu hỏi là D}
}
\end{ex}

\begin{ex} %[2D10H52]
Tổng $800 + 50$ viết thành số là
\choice
{\True $850$}
{$805$}
{$85$}
{$8005$}
\loigiai{
$800 + 50 = 850$.\\
\textbf{Đáp án của câu hỏi là A}
}
\end{ex}

\begin{ex} %[2D10H52]
Số $990$ viết thành tổng các trăm và chục là
\choice
{$900 + 9$}
{$90 + 90$}
{\True $900 + 90$}
{$900 + 900$}
\loigiai{
$990 = 900 + 90$.\\
\textbf{Đáp án của câu hỏi là C}
}
\end{ex}

\begin{ex} %[2D10H52]
Điền số thích hợp: $635 = 600 + \dots + 5$.
\choice
{$3$}
{\True $30$}
{$300$}
{$35$}
\loigiai{
$635 = 600 + 30 + 5$.\\
\textbf{Đáp án của câu hỏi là B}
}
\end{ex}

\begin{ex} %[2D10V52]
Điền số thích hợp: $742 = \dots + 40 + 2$.
\choice
{$7$}
{$70$}
{$740$}
{\True $700$}
\loigiai{
$742 = 700 + 40 + 2$.\\
\textbf{Đáp án của câu hỏi là D}
}
\end{ex}

\begin{ex} %[2D10V52]
Cho số $\overline{a45} = 300 + 40 + 5$. Chữ số $a$ là
\choice
{$4$}
{\True $3$}
{$5$}
{$2$}
\loigiai{
Chữ số hàng trăm là 3.\\
\textbf{Đáp án của câu hỏi là B}
}
\end{ex}

\begin{ex} %[2D10V52]
Tính giá trị của tổng: $600 + 80 + 9$.
\choice
{\True $689$}
{$698$}
{$869$}
{$896$}
\loigiai{
$600 + 80 + 9 = 689$.\\
\textbf{Đáp án của câu hỏi là A}
}
\end{ex}

\begin{ex} %[2D10V52]
Số gồm $9$ trăm và $9$ đơn vị viết thành tổng là
\choice
{$900 + 90$}
{$90 + 9$}
{\True $900 + 9$}
{$9 + 9$}
\loigiai{
Số đó là $909 = 900 + 9$.\\
\textbf{Đáp án của câu hỏi là C}
}
\end{ex}

\begin{ex} %[2D10V52]
Điền dấu thích hợp: $300 + 40 + 5 \dots 354$.
\choice
{$>$}
{$=$}
{Không so sánh được}
{\True $<$}
\loigiai{
$300 + 40 + 5 = 345 < 354$.\\
\textbf{Đáp án của câu hỏi là D}
}
\end{ex}

\begin{ex} %[2D10C52]
Tìm số có 3 chữ số biết khi viết thành tổng các trăm, chục, đơn vị ta được $a00 + b0 + c$ với $a = 2 \times b$ và $b = 2 \times c = 4$.
\choice
{$442$}
{\True $842$}
{$824$}
{$482$}
\loigiai{
$c = 4 : 2 = 2; b = 4; a = 2 \times 4 = 8$. Số đó là $842$.\\
\textbf{Đáp án của câu hỏi là B}
}
\end{ex}

\textbf{Phần 2. Trắc nghiệm đúng sai}
\setcounterpageref{ex}{0}

\begin{ex} %[2D10H52]
Xét tính đúng sai khi phân tích số:
\choiceTF
{\True $548 = 500 + 40 + 8$.}
{\True $306 = 300 + 6$.}
{$720 = 700 + 2$.}
{\True $720 = 700 + 20$.}
\loigiai{
\begin{itemize}
    \item $548 = 500 + 40 + 8$: Đúng.
    \item $306 = 300 + 6$: Đúng.
    \item $720 = 700 + 20 \neq 700 + 2$. Ý 3 Sai, Ý 4 Đúng.
\end{itemize}
\begin{itemize}
    \item a) \textbf{Đúng}.
    \item b) \textbf{Đúng}.
    \item c) \textbf{Sai}.
    \item d) \textbf{Đúng}.
\end{itemize}
}
\end{ex}

\begin{ex} %[2D10V52]
Cho số $M = 800 + 70 + 5$ và $N = 800 + 50 + 7$. Xét tính đúng sai:
\choiceTF
{\True $M = 875$.}
{\True $N = 857$.}
{$M < N$.}
{\True $M > N$ ($875 > 857$).}
\loigiai{
\begin{itemize}
    \item $M = 875$: Đúng.
    \item $N = 857$: Đúng.
    \item $M > N$ (không bé hơn). Ý 3 Sai, Ý 4 Đúng.
\end{itemize}
\begin{itemize}
    \item a) \textbf{Đúng}.
    \item b) \textbf{Đúng}.
    \item c) \textbf{Sai}.
    \item d) \textbf{Đúng}.
\end{itemize}
}
\end{ex}

\textbf{Phần 3. Trả lời ngắn}
\setcounterpageref{ex}{0}

\begin{ex} %[2D10H52]
Viết kết quả của tổng: $400 + 70 + 3$.
\shortans{473}
\loigiai{
\textbf{Đáp số:} $473$\\
473.
}
\end{ex}

\begin{ex} %[2D10H52]
Viết kết quả của tổng: $600 + 9$.
\shortans{609}
\loigiai{
\textbf{Đáp số:} $609$\\
609.
}
\end{ex}

\begin{ex} %[2D10V52]
Viết kết quả của tổng: $700 + 50$.
\shortans{750}
\loigiai{
\textbf{Đáp số:} $750$\\
750.
}
\end{ex}

\begin{ex} %[2D10V52]
Điền số thích hợp: $852 = 800 + ? + 2$.
\shortans{50}
\loigiai{
\textbf{Đáp số:} $50$\\
50.
}
\end{ex}

\begin{ex} %[2D10V52]
Điền số thích hợp: $936 = ? + 30 + 6$.
\shortans{900}
\loigiai{
\textbf{Đáp số:} $900$\\
900.
}
\end{ex}

\begin{ex} %[2D10C52]
Tìm giá trị của biểu thức: $(500 + 60 + 4) - 64$.
\shortans{500}
\loigiai{
\textbf{Đáp số:} $500$\\
$564 - 64 = 500$.
}
\end{ex}

\newpage

\TaoPhieuKiemTra{Chủ đề 10 -- Các số trong phạm vi 1 000}{53}{Tiền Việt Nam}{https://loptoanthaylinh.vn/thi-online/toan-2/bai-53}
\textbf{Phần 1. Trắc nghiệm lựa chọn}
\setcounterpageref{ex}{0}

\begin{ex} %[2D10N53]
Điền dấu thích hợp: $456 \dots 465$.
\choice
{$>$}
{\True $<$}
{$=$}
{Không so sánh được}
\loigiai{
Hàng trăm bằng 4. Hàng chục: $5 < 6$ nên $456 < 465$.\\
\textbf{Đáp án của câu hỏi là B}
}
\end{ex}

\begin{ex} %[2D10N53]
Điền dấu thích hợp: $789 \dots 780$.
\choice
{\True $>$}
{$<$}
{$=$}
{Không so sánh được}
\loigiai{
$789 > 780$.\\
\textbf{Đáp án của câu hỏi là A}
}
\end{ex}

\begin{ex} %[2D10H53]
Trong các số $567, 765, 657, 576$, số lớn nhất là
\choice
{$657$}
{$576$}
{\True $765$}
{$567$}
\loigiai{
Số lớn nhất là 765.\\
\textbf{Đáp án của câu hỏi là C}
}
\end{ex}

\begin{ex} %[2D10H53]
Trong các số $432, 234, 342, 243$, số bé nhất là
\choice
{$243$}
{\True $234$}
{$342$}
{$432$}
\loigiai{
Số bé nhất là 234.\\
\textbf{Đáp án của câu hỏi là B}
}
\end{ex}

\begin{ex} %[2D10H53]
Sắp xếp các số $324, 423, 234, 342$ theo thứ tự từ bé đến lớn:
\choice
{\True $234, 324, 342, 423$}
{$423, 342, 324, 234$}
{$234, 342, 324, 423$}
{$324, 234, 342, 423$}
\loigiai{
$234 < 324 < 342 < 423$.\\
\textbf{Đáp án của câu hỏi là A}
}
\end{ex}

\begin{ex} %[2D10H53]
Điền dấu thích hợp: $500 + 40 \dots 504$.
\choice
{$<$}
{$=$}
{Không so sánh được}
{\True $>$}
\loigiai{
$500 + 40 = 540 > 504$.\\
\textbf{Đáp án của câu hỏi là D}
}
\end{ex}

\begin{ex} %[2D10V53]
Tìm chữ số $x$ thích hợp để: $\overline{5x8} < 518$.
\choice
{$1$}
{\True $0$}
{$2$}
{$3$}
\loigiai{
Chữ số hàng chục phải bé hơn 1 nên $x = 0$.\\
\textbf{Đáp án của câu hỏi là B}
}
\end{ex}

\begin{ex} %[2D10V53]
Tìm chữ số $y$ thích hợp để: $\overline{78y} > 788$.
\choice
{\True $9$}
{$7$}
{$8$}
{$6$}
\loigiai{
Chữ số hàng đơn vị phải lớn hơn 8 nên $y = 9$.\\
\textbf{Đáp án của câu hỏi là A}
}
\end{ex}

\begin{ex} %[2D10V53]
Sắp xếp các số $890, 908, 809, 980$ theo thứ tự từ lớn đến bé:
\choice
{$809, 890, 908, 980$}
{$980, 890, 908, 809$}
{\True $980, 908, 890, 809$}
{$908, 980, 890, 809$}
\loigiai{
Thứ tự giảm dần: $980 > 908 > 890 > 809$.\\
\textbf{Đáp án của câu hỏi là C}
}
\end{ex}

\begin{ex} %[2D10V53]
Có bao nhiêu số tự nhiên có ba chữ số lớn hơn 995? (Gồm 996, 997, 998, 999)
\choice
{$3$ số}
{$5$ số}
{$6$ số}
{\True $4$ số}
\loigiai{
Các số: 996, 997, 998, 999 (có 4 số có 3 chữ số).\\
\textbf{Đáp án của câu hỏi là D}
}
\end{ex}

\begin{ex} %[2D10V53]
Điền dấu thích hợp: $600 + 70 + 8 \dots 600 + 80 + 7$.
\choice
{$>$}
{$=$}
{\True $<$}
{Không so sánh được}
\loigiai{
$678 < 687$.\\
\textbf{Đáp án của câu hỏi là C}
}
\end{ex}

\begin{ex} %[2D10C53]
Cho các chữ số $2, 5, 8$. Lập số lớn nhất và số bé nhất có ba chữ số khác nhau từ 3 chữ số trên. Hiệu của hai số đó là
\choice
{$594$}
{$600$}
{$624$}
{\True $594$ ($852 - 258 = 594$)}
\loigiai{
Lớn nhất: 852. Bé nhất: 258. Hiệu: $852 - 258 = 594$.\\
\textbf{Đáp án của câu hỏi là D}
}
\end{ex}

\textbf{Phần 2. Trắc nghiệm đúng sai}
\setcounterpageref{ex}{0}

\begin{ex} %[2D10H53]
Xét tính đúng sai của các so sánh số có 3 chữ số:
\choiceTF
{\True $435 < 453$.}
{\True $702 < 720$.}
{$890 < 809$.}
{\True $890 > 809$.}
\loigiai{
\begin{itemize}
    \item $435 < 453$: Đúng.
    \item $702 < 720$: Đúng.
    \item $890 > 809 \neq 809$. Ý 3 Sai, Ý 4 Đúng.
\end{itemize}
\begin{itemize}
    \item a) \textbf{Đúng}.
    \item b) \textbf{Đúng}.
    \item c) \textbf{Sai}.
    \item d) \textbf{Đúng}.
\end{itemize}
}
\end{ex}

\begin{ex} %[2D10V53]
Cho dãy số: $350, 355, 360, 365, 370$. Xét tính đúng sai:
\choiceTF
{\True Đây là dãy số tăng dần đều 5 đơn vị.}
{\True Số lớn nhất trong dãy là 370.}
{Số bé nhất trong dãy là 360.}
{\True Số bé nhất trong dãy là 350.}
\loigiai{
\begin{itemize}
    \item Dãy tăng 5 đơn vị: Đúng.
    \item Lớn nhất là 370: Đúng.
    \item Bé nhất là 350 (không phải 360). Ý 3 Sai, Ý 4 Đúng.
\end{itemize}
\begin{itemize}
    \item a) \textbf{Đúng}.
    \item b) \textbf{Đúng}.
    \item c) \textbf{Sai}.
    \item d) \textbf{Đúng}.
\end{itemize}
}
\end{ex}

\textbf{Phần 3. Trả lời ngắn}
\setcounterpageref{ex}{0}

\begin{ex} %[2D10H53]
Trong các số $672, 726, 267, 627$, tìm số lớn nhất.
\shortans{726}
\loigiai{
\textbf{Đáp số:} $726$\\
726.
}
\end{ex}

\begin{ex} %[2D10H53]
Trong các số $581, 185, 815, 158$, tìm số bé nhất.
\shortans{158}
\loigiai{
\textbf{Đáp số:} $158$\\
158.
}
\end{ex}

\begin{ex} %[2D10V53]
Tìm số lớn nhất có 3 chữ số mà tổng các chữ số của nó bằng 3 ($300$).
\shortans{300}
\loigiai{
\textbf{Đáp số:} $300$\\
300.
}
\end{ex}

\begin{ex} %[2D10V53]
Tìm số bé nhất có 3 chữ số mà tổng các chữ số của nó bằng 3 ($102$).
\shortans{102}
\loigiai{
\textbf{Đáp số:} $102$\\
102.
}
\end{ex}

\begin{ex} %[2D10V53]
Có bao nhiêu số tự nhiên có ba chữ số lớn hơn 990 và nhỏ hơn 996?
\shortans{5}
\loigiai{
\textbf{Đáp số:} $5$\\
Gồm 991, 992, 993, 994, 995 (có 5 số).
}
\end{ex}

\begin{ex} %[2D10C53]
Tìm số liền trước của số lớn nhất có ba chữ số khác nhau ($987$).
\shortans{986}
\loigiai{
\textbf{Đáp số:} $986$\\
$987 - 1 = 986$.
}
\end{ex}

\newpage

\TaoPhieuKiemTra{Chủ đề 10 -- Các số trong phạm vi 1 000}{54}{Luyện tập chung đại lượng}{https://loptoanthaylinh.vn/thi-online/toan-2/bai-54}
\textbf{Phần 1. Trắc nghiệm lựa chọn}
\setcounterpageref{ex}{0}

\begin{ex} %[2D10N54]
Số gồm $7$ trăm, $2$ chục và $0$ đơn vị là
\choice
{\True $720$}
{$702$}
{$270$}
{$72$}
\loigiai{
Số đó là $720$.\\
\textbf{Đáp án của câu hỏi là A}
}
\end{ex}

\begin{ex} %[2D10N54]
Số $805$ đọc là
\choice
{Tám trăm năm mươi}
{Tám không năm}
{\True Tám trăm linh năm}
{Tám trăm năm}
\loigiai{
Đọc là Tám trăm linh năm.\\
\textbf{Đáp án của câu hỏi là C}
}
\end{ex}

\begin{ex} %[2D10H54]
Số liền sau của số $899$ là
\choice
{\True $900$}
{$898$}
{$901$}
{$890$}
\loigiai{
$899 + 1 = 900$.\\
\textbf{Đáp án của câu hỏi là A}
}
\end{ex}

\begin{ex} %[2D10H54]
Số liền trước của số $1\,000$ là
\choice
{$1001$}
{$990$}
{$900$}
{\True $999$}
\loigiai{
$1000 - 1 = 999$.\\
\textbf{Đáp án của câu hỏi là D}
}
\end{ex}

\begin{ex} %[2D10H54]
Tính: $500 + 300 = \dots$
\choice
{$700$}
{\True $800$}
{$900$}
{$600$}
\loigiai{
$500 + 300 = 800$.\\
\textbf{Đáp án của câu hỏi là B}
}
\end{ex}

\begin{ex} %[2D10H54]
Số $642$ được viết thành tổng là
\choice
{\True $600 + 40 + 2$}
{$600 + 4 + 2$}
{$60 + 40 + 2$}
{$640 + 20$}
\loigiai{
$642 = 600 + 40 + 2$.\\
\textbf{Đáp án của câu hỏi là A}
}
\end{ex}

\begin{ex} %[2D10V54]
Điền dấu thích hợp: $789 \dots 798$.
\choice
{$>$}
{$=$}
{\True $<$}
{Không so sánh được}
\loigiai{
$789 < 798$.\\
\textbf{Đáp án của câu hỏi là C}
}
\end{ex}

\begin{ex} %[2D10V54]
Tìm số tròn chục lớn nhất có 3 chữ số:
\choice
{$900$}
{$999$}
{$1000$}
{\True $990$}
\loigiai{
Số đó là 990.\\
\textbf{Đáp án của câu hỏi là D}
}
\end{ex}

\begin{ex} %[2D10V54]
Sắp xếp các số $456, 654, 546, 465$ theo thứ tự từ bé đến lớn:
\choice
{$654, 546, 465, 456$}
{\True $456, 465, 546, 654$}
{$456, 546, 465, 654$}
{$465, 456, 546, 654$}
\loigiai{
$456 < 465 < 546 < 654$.\\
\textbf{Đáp án của câu hỏi là B}
}
\end{ex}

\begin{ex} %[2D10V54]
Tính giá trị biểu thức: $900 - 400 + 200$.
\choice
{$600$}
{$800$}
{$500$}
{\True $700$}
\loigiai{
$500 + 200 = 700$.\\
\textbf{Đáp án của câu hỏi là D}
}
\end{ex}

\begin{ex} %[2D10V54]
Số bé nhất có ba chữ số khác nhau mà chữ số hàng trăm là 5 là
\choice
{$500$}
{\True $501$}
{$510$}
{$512$}
\loigiai{
Số đó là 501.\\
\textbf{Đáp án của câu hỏi là B}
}
\end{ex}

\begin{ex} %[2D10C54]
Tìm số có 3 chữ số biết chữ số hàng trăm là 8, chữ số hàng đơn vị là 5 và chữ số hàng chục bằng hiệu của chữ số hàng trăm và hàng đơn vị.
\choice
{$853$}
{$825$}
{\True $835$}
{$845$}
\loigiai{
Hàng chục: $8 - 5 = 3$. Số đó là $835$.\\
\textbf{Đáp án của câu hỏi là C}
}
\end{ex}

\textbf{Phần 2. Trắc nghiệm đúng sai}
\setcounterpageref{ex}{0}

\begin{ex} %[2D10H54]
Xét tính đúng sai về số trong phạm vi 1 000:
\choiceTF
{\True Số 1 000 gồm 10 trăm.}
{\True $999$ là số lớn nhất có ba chữ số.}
{Số 605 đọc là Sáu trăm năm mươi.}
{\True Số 605 đọc là Sáu trăm linh năm.}
\loigiai{
\begin{itemize}
    \item 1000 = 10 trăm: Đúng.
    \item 999 là số lớn nhất có 3 chữ số: Đúng.
    \item 605 đọc là Sáu trăm linh năm. Ý 3 Sai, Ý 4 Đúng.
\end{itemize}
\begin{itemize}
    \item a) \textbf{Đúng}.
    \item b) \textbf{Đúng}.
    \item c) \textbf{Sai}.
    \item d) \textbf{Đúng}.
\end{itemize}
}
\end{ex}

\begin{ex} %[2D10V54]
Cho 4 số: $512, 521, 215, 251$. Xét tính đúng sai:
\choiceTF
{\True Số lớn nhất là 521.}
{\True Số bé nhất là 215.}
{Xếp tăng dần: $251, 215, 512, 521$.}
{\True Xếp tăng dần đúng là $215, 251, 512, 521$.}
\loigiai{
\begin{itemize}
    \item Lớn nhất 521: Đúng.
    \item Bé nhất 215: Đúng.
    \item Thứ tự đúng là $215 < 251$. Ý 3 Sai, Ý 4 Đúng.
\end{itemize}
\begin{itemize}
    \item a) \textbf{Đúng}.
    \item b) \textbf{Đúng}.
    \item c) \textbf{Sai}.
    \item d) \textbf{Đúng}.
\end{itemize}
}
\end{ex}

\textbf{Phần 3. Trả lời ngắn}
\setcounterpageref{ex}{0}

\begin{ex} %[2D10H54]
Viết số gồm 6 trăm, 3 chục và 8 đơn vị.
\shortans{638}
\loigiai{
\textbf{Đáp số:} $638$\\
638.
}
\end{ex}

\begin{ex} %[2D10H54]
Tìm số liền trước của số $500$.
\shortans{499}
\loigiai{
\textbf{Đáp số:} $499$\\
$500 - 1 = 499$.
}
\end{ex}

\begin{ex} %[2D10V54]
Tìm số liền sau của số $699$.
\shortans{700}
\loigiai{
\textbf{Đáp số:} $700$\\
$699 + 1 = 700$.
}
\end{ex}

\begin{ex} %[2D10V54]
Tính giá trị của tổng: $700 + 40 + 9$.
\shortans{749}
\loigiai{
\textbf{Đáp số:} $749$\\
749.
}
\end{ex}

\begin{ex} %[2D10V54]
Tính: $800 - 500 + 300$.
\shortans{600}
\loigiai{
\textbf{Đáp số:} $600$\\
$300 + 300 = 600$.
}
\end{ex}

\begin{ex} %[2D10C54]
Tìm số tròn chục bé nhất có ba chữ số khác nhau ($120$).
\shortans{120}
\loigiai{
\textbf{Đáp số:} $120$\\
120.
}
\end{ex}

\newpage

% =========================================================================
% CHỦ ĐỀ 11 - ĐỘ DÀI VÀ ĐƠN VỊ ĐO ĐỘ DÀI. TIỀN VIỆT NAM
% =========================================================================
\section{Độ dài và đơn vị đo độ dài. tiền việt nam}

\TaoPhieuKiemTra{Chủ đề 11 -- Độ dài và đơn vị đo độ dài. tiền việt nam}{55}{Đơn vị, chục, trăm, nghìn}{https://loptoanthaylinh.vn/thi-online/toan-2/bai-55}
\textbf{Phần 1. Trắc nghiệm lựa chọn}
\setcounterpageref{ex}{0}

\begin{ex} %[2H11N55]
Đề-xi-mét được viết tắt là
\choice
{$\text{cm}$}
{$\text{m}$}
{\True $\text{dm}$}
{$\text{km}$}
\loigiai{
Đề-xi-mét viết tắt là $\text{dm}$.\\
\textbf{Đáp án của câu hỏi là C}
}
\end{ex}

\begin{ex} %[2H11N55]
Mét được viết tắt là
\choice
{$\text{dm}$}
{$\text{km}$}
{$\text{mm}$}
{\True $\text{m}$}
\loigiai{
Mét viết tắt là $\text{m}$.\\
\textbf{Đáp án của câu hỏi là D}
}
\end{ex}

\begin{ex} %[2H11N55]
Ki-lô-mét được viết tắt là
\choice
{$\text{kg}$}
{\True $\text{km}$}
{$\text{m}$}
{$\text{dm}$}
\loigiai{
Ki-lô-mét viết tắt là $\text{km}$.\\
\textbf{Đáp án của câu hỏi là B}
}
\end{ex}

\begin{ex} %[2H11H55]
$1\text{ m}$ bằng bao nhiêu xăng-ti-mét?
\choice
{$10\text{ cm}$}
{$1000\text{ cm}$}
{\True $100\text{ cm}$}
{$1\text{ cm}$}
\loigiai{
$1\text{ m} = 100\text{ cm}$.\\
\textbf{Đáp án của câu hỏi là C}
}
\end{ex}

\begin{ex} %[2H11H55]
$1\text{ m}$ bằng bao nhiêu đề-xi-mét?
\choice
{\True $10\text{ dm}$}
{$100\text{ dm}$}
{$1\text{ dm}$}
{$1000\text{ dm}$}
\loigiai{
$1\text{ m} = 10\text{ dm}$.\\
\textbf{Đáp án của câu hỏi là A}
}
\end{ex}

\begin{ex} %[2H11H55]
$1\text{ km}$ bằng bao nhiêu mét?
\choice
{$100\text{ m}$}
{$10\text{ m}$}
{$10\,000\text{ m}$}
{\True $1\,000\text{ m}$}
\loigiai{
$1\text{ km} = 1000\text{ m}$.\\
\textbf{Đáp án của câu hỏi là D}
}
\end{ex}

\begin{ex} %[2H11V55]
Tính: $35\text{ m} + 28\text{ m} = \dots$
\choice
{$62\text{ m}$}
{$64\text{ m}$}
{\True $63\text{ m}$}
{$65\text{ m}$}
\loigiai{
$35 + 28 = 63\text{ m}$.\\
\textbf{Đáp án của câu hỏi là C}
}
\end{ex}

\begin{ex} %[2H11V55]
Tính: $80\text{ km} - 35\text{ km} = \dots$
\choice
{$44\text{ km}$}
{\True $45\text{ km}$}
{$46\text{ km}$}
{$55\text{ km}$}
\loigiai{
$80 - 35 = 45\text{ km}$.\\
\textbf{Đáp án của câu hỏi là B}
}
\end{ex}

\begin{ex} %[2H11V55]
Một cuộn vải dài $50\text{ m}$, người ta cắt bán đi $18\text{ m}$. Cuộn vải còn lại dài là
\choice
{\True $32\text{ m}$}
{$31\text{ m}$}
{$33\text{ m}$}
{$34\text{ m}$}
\loigiai{
$50 - 18 = 32\text{ m}$.\\
\textbf{Đáp án của câu hỏi là A}
}
\end{ex}

\begin{ex} %[2H11V55]
Quãng đường từ Hà Nội đến Hải Phòng dài khoảng
\choice
{$100\text{ m}$}
{$100\text{ cm}$}
{$100\text{ dm}$}
{\True $100\text{ km}$}
\loigiai{
Khoảng cách giữa các tỉnh/thành phố đo bằng ki-lô-mét ($100\text{ km}$).\\
\textbf{Đáp án của câu hỏi là D}
}
\end{ex}

\begin{ex} %[2H11V55]
Điền dấu thích hợp: $2\text{ m} \dots 180\text{ cm}$.
\choice
{\True $>$}
{$<$}
{$=$}
{Không so sánh được}
\loigiai{
$2\text{ m} = 200\text{ cm} > 180\text{ cm}$.\\
\textbf{Đáp án của câu hỏi là A}
}
\end{ex}

\begin{ex} %[2H11C55]
Một cây tre dài $6\text{ m}$, người ta chặt thành các đoạn dài $2\text{ m}$. Hỏi được bao nhiêu đoạn tre?
\choice
{$2$ đoạn}
{\True $3$ đoạn}
{$4$ đoạn}
{$5$ đoạn}
\loigiai{
$6 : 2 = 3$ đoạn.\\
\textbf{Đáp án của câu hỏi là B}
}
\end{ex}

\textbf{Phần 2. Trắc nghiệm đúng sai}
\setcounterpageref{ex}{0}

\begin{ex} %[2H11H55]
Xét tính đúng sai về đổi đơn vị đo độ dài:
\choiceTF
{\True $1\text{ m} = 10\text{ dm} = 100\text{ cm}$.}
{\True $1\text{ km} = 1\,000\text{ m}$.}
{$1\text{ dm} = 100\text{ cm}$.}
{\True $1\text{ dm} = 10\text{ cm}$.}
\loigiai{
\begin{itemize}
    \item $1\text{ m} = 10\text{ dm} = 100\text{ cm}$: Đúng.
    \item $1\text{ km} = 1000\text{ m}$: Đúng.
    \item $1\text{ dm} = 10\text{ cm} \neq 100\text{ cm}$. Ý 3 Sai, Ý 4 Đúng.
\end{itemize}
\begin{itemize}
    \item a) \textbf{Đúng}.
    \item b) \textbf{Đúng}.
    \item c) \textbf{Sai}.
    \item d) \textbf{Đúng}.
\end{itemize}
}
\end{ex}

\begin{ex} %[2H11V55]
Quãng đường $AB$ dài $45\text{ km}$, quãng đường $BC$ dài $38\text{ km}$ ($B$ nằm giữa $A$ và $C$). Xét tính đúng sai:
\choiceTF
{\True Quãng đường $AC$ dài $83\text{ km}$ ($45 + 38 = 83$).}
{\True Quãng đường $AB$ dài hơn quãng đường $BC$ là $7\text{ km}$ ($45 - 38 = 7$).}
{Quãng đường $AC$ ngắn hơn $80\text{ km}$.}
{\True Quãng đường $AC$ dài hơn $80\text{ km}$ (cụ thể là $83\text{ km}$).}
\loigiai{
\begin{itemize}
    \item $AC = 45 + 38 = 83\text{ km}$: Đúng.
    \item $AB - BC = 7\text{ km}$: Đúng.
    \item $83 > 80$ (không ngắn hơn). Ý 3 Sai, Ý 4 Đúng.
\end{itemize}
\begin{itemize}
    \item a) \textbf{Đúng}.
    \item b) \textbf{Đúng}.
    \item c) \textbf{Sai}.
    \item d) \textbf{Đúng}.
\end{itemize}
}
\end{ex}

\textbf{Phần 3. Trả lời ngắn}
\setcounterpageref{ex}{0}

\begin{ex} %[2H11H55]
$1\text{ m}$ bằng bao nhiêu xăng-ti-mét?
\shortans{100}
\loigiai{
\textbf{Đáp số:} $100$\\
100 cm.
}
\end{ex}

\begin{ex} %[2H11H55]
$1\text{ km}$ bằng bao nhiêu mét?
\shortans{1000}
\loigiai{
\textbf{Đáp số:} $1000$\\
1000 m.
}
\end{ex}

\begin{ex} %[2H11V55]
Tính: $48\text{ m} + 35\text{ m}$ theo mét.
\shortans{83}
\loigiai{
\textbf{Đáp số:} $83$\\
$48 + 35 = 83\text{ m}$.
}
\end{ex}

\begin{ex} %[2H11V55]
Tính: $92\text{ km} - 47\text{ km}$ theo ki-lô-mét.
\shortans{45}
\loigiai{
\textbf{Đáp số:} $45$\\
$92 - 47 = 45\text{ km}$.
}
\end{ex}

\begin{ex} %[2H11V55]
Cuộn dây thứ nhất dài $25\text{ m}$, cuộn dây thứ hai dài $35\text{ m}$. Cả hai cuộn dây dài bao nhiêu mét?
\shortans{60}
\loigiai{
\textbf{Đáp số:} $60$\\
$25 + 35 = 60\text{ m}$.
}
\end{ex}

\begin{ex} %[2H11C55]
Bác thợ may có mảnh vải dài $10\text{ m}$, bác may 5 bộ quần áo như nhau thì vừa hết. Hỏi mỗi bộ quần áo may hết bao nhiêu mét vải?
\shortans{2}
\loigiai{
\textbf{Đáp số:} $2$\\
$10 : 5 = 2\text{ m}$.
}
\end{ex}

\newpage

\TaoPhieuKiemTra{Chủ đề 11 -- Độ dài và đơn vị đo độ dài. tiền việt nam}{56}{Các số tròn chục, tròn trăm}{https://loptoanthaylinh.vn/thi-online/toan-2/bai-56}
\textbf{Phần 1. Trắc nghiệm lựa chọn}
\setcounterpageref{ex}{0}

\begin{ex} %[2H11N56]
Đơn vị tiền tệ của nước Việt Nam là
\choice
{Đô-la}
{Yên}
{\True Đồng}
{Ơ-rô}
\loigiai{
Đơn vị tiền tệ là Đồng.\\
\textbf{Đáp án của câu hỏi là C}
}
\end{ex}

\begin{ex} %[2H11N56]
Tờ tiền nào sau đây có mệnh giá lớn nhất?
\choice
{$100$ đồng}
{\True $1\,000$ đồng}
{$200$ đồng}
{$500$ đồng}
\loigiai{
Tờ 1 000 đồng lớn nhất.\\
\textbf{Đáp án của câu hỏi là B}
}
\end{ex}

\begin{ex} %[2H11H56]
2 tờ 500 đồng có giá trị bằng
\choice
{\True $1\,000$ đồng}
{$500$ đồng}
{$700$ đồng}
{$2\,000$ đồng}
\loigiai{
$500 + 500 = 1000$ đồng.\\
\textbf{Đáp án của câu hỏi là A}
}
\end{ex}

\begin{ex} %[2H11H56]
5 tờ 200 đồng có tổng giá trị là
\choice
{$800$ đồng}
{\True $1\,000$ đồng}
{$700$ đồng}
{$500$ đồng}
\loigiai{
$200 \times 5 = 1000$ đồng.\\
\textbf{Đáp án của câu hỏi là B}
}
\end{ex}

\begin{ex} %[2H11H56]
Mai có 1 tờ 500 đồng và 2 tờ 200 đồng. Mai có tất cả số tiền là
\choice
{\True $900$ đồng}
{$800$ đồng}
{$1000$ đồng}
{$700$ đồng}
\loigiai{
$500 + 200 + 200 = 900$ đồng.\\
\textbf{Đáp án của câu hỏi là A}
}
\end{ex}

\begin{ex} %[2H11H56]
Một chiếc kẹo mút giá 500 đồng, một cái bánh que giá 300 đồng. Mua cả hai thứ hết số tiền là
\choice
{$700$ đồng}
{$900$ đồng}
{$1000$ đồng}
{\True $800$ đồng}
\loigiai{
$500 + 300 = 800$ đồng.\\
\textbf{Đáp án của câu hỏi là D}
}
\end{ex}

\begin{ex} %[2H11V56]
Nam có 1 000 đồng, Nam mua một cục tẩy hết 600 đồng. Nam còn lại số tiền là
\choice
{$300$ đồng}
{$500$ đồng}
{\True $400$ đồng}
{$200$ đồng}
\loigiai{
$1000 - 600 = 400$ đồng.\\
\textbf{Đáp án của câu hỏi là C}
}
\end{ex}

\begin{ex} %[2H11V56]
Để có đủ 1 000 đồng từ các tờ 200 đồng, cần mấy tờ 200 đồng?
\choice
{$4$ tờ}
{$6$ tờ}
{$10$ tờ}
{\True $5$ tờ}
\loigiai{
$1000 : 200 = 5$ tờ.\\
\textbf{Đáp án của câu hỏi là D}
}
\end{ex}

\begin{ex} %[2H11V56]
Bình có 3 tờ 200 đồng và 1 tờ 100 đồng. Bình có tất cả số tiền là
\choice
{\True $700$ đồng}
{$600$ đồng}
{$800$ đồng}
{$900$ đồng}
\loigiai{
$200 \times 3 + 100 = 700$ đồng.\\
\textbf{Đáp án của câu hỏi là A}
}
\end{ex}

\begin{ex} %[2H11V56]
Một chiếc bút chì giá 800 đồng, chiếc thước kẻ rẻ hơn chiếc bút chì 200 đồng. Chiếc thước kẻ giá
\choice
{$500$ đồng}
{\True $600$ đồng}
{$700$ đồng}
{$1000$ đồng}
\loigiai{
$800 - 200 = 600$ đồng.\\
\textbf{Đáp án của câu hỏi là B}
}
\end{ex}

\begin{ex} %[2H11V56]
Lan có 1 000 đồng, Lan mua 2 chiếc kẹo giá 400 đồng một chiếc. Lan còn lại số tiền là
\choice
{$100$ đồng}
{$300$ đồng}
{\True $200$ đồng}
{$400$ đồng}
\loigiai{
Mua 2 kẹo: $400 \times 2 = 800$ đồng.\\
Còn lại: $1000 - 800 = 200$ đồng.\\
\textbf{Đáp án của câu hỏi là C}
}
\end{ex}

\begin{ex} %[2H11C56]
Mẹ đưa cho cô bán hàng tờ 1 000 đồng để mua một mớ rau giá 700 đồng. Cô bán hàng trả lại tiền thừa bằng các tờ tiền nào sau đây?
\choice
{1 tờ 200 đồng}
{2 tờ 200 đồng}
{3 tờ 200 đồng}
{\True 1 tờ 200 đồng và 1 tờ 100 đồng}
\loigiai{
Tiền thừa: $1000 - 700 = 300$ đồng = $200 + 100$ đồng.\\
\textbf{Đáp án của câu hỏi là D}
}
\end{ex}

\textbf{Phần 2. Trắc nghiệm đúng sai}
\setcounterpageref{ex}{0}

\begin{ex} %[2H11H56]
Xét tính đúng sai về mệnh giá tiền Việt Nam:
\choiceTF
{\True 2 tờ 500 đồng có giá trị bằng 1 000 đồng.}
{\True 5 tờ 100 đồng có giá trị bằng 500 đồng.}
{1 tờ 1 000 đồng có giá trị bằng 3 tờ 200 đồng.}
{\True 1 tờ 1 000 đồng có giá trị bằng 5 tờ 200 đồng.}
\loigiai{
\begin{itemize}
    \item $500 \times 2 = 1000$ đồng: Đúng.
    \item $100 \times 5 = 500$ đồng: Đúng.
    \item 1000 đồng = 5 tờ 200 đồng (không phải 3 tờ). Ý 3 Sai, Ý 4 Đúng.
\end{itemize}
\begin{itemize}
    \item a) \textbf{Đúng}.
    \item b) \textbf{Đúng}.
    \item c) \textbf{Sai}.
    \item d) \textbf{Đúng}.
\end{itemize}
}
\end{ex}

\begin{ex} %[2H11V56]
Hà có 1 000 đồng, Hà mua 1 gói kẹo giá 500 đồng và 1 phong bì giá 200 đồng. Xét tính đúng sai:
\choiceTF
{\True Tổng số tiền Hà đã mua là 700 đồng ($500 + 200 = 700$).}
{\True Số tiền Hà còn lại là 300 đồng ($1000 - 700 = 300$).}
{Hà đã tiêu hết toàn bộ 1 000 đồng.}
{\True Hà vẫn còn thừa 300 đồng.}
\loigiai{
\begin{itemize}
    \item Mua hết 700 đồng: Đúng.
    \item Còn lại 300 đồng: Đúng.
    \item Chưa tiêu hết. Ý 3 Sai, Ý 4 Đúng.
\end{itemize}
\begin{itemize}
    \item a) \textbf{Đúng}.
    \item b) \textbf{Đúng}.
    \item c) \textbf{Sai}.
    \item d) \textbf{Đúng}.
\end{itemize}
}
\end{ex}

\textbf{Phần 3. Trả lời ngắn}
\setcounterpageref{ex}{0}

\begin{ex} %[2H11H56]
2 tờ 500 đồng đổi được bao nhiêu đồng?
\shortans{1000}
\loigiai{
\textbf{Đáp số:} $1000$\\
1000 đồng.
}
\end{ex}

\begin{ex} %[2H11H56]
Có 4 tờ 200 đồng, hỏi có tất cả bao nhiêu đồng?
\shortans{800}
\loigiai{
\textbf{Đáp số:} $800$\\
$200 \times 4 = 800$ đồng.
}
\end{ex}

\begin{ex} %[2H11V56]
Lan có 1 000 đồng, mua kem hết 700 đồng. Hỏi Lan còn lại bao nhiêu đồng?
\shortans{300}
\loigiai{
\textbf{Đáp số:} $300$\\
$1000 - 700 = 300$ đồng.
}
\end{ex}

\begin{ex} %[2H11V56]
Một cái nhãn vở giá 200 đồng. Hỏi mua 4 cái nhãn vở hết bao nhiêu đồng?
\shortans{800}
\loigiai{
\textbf{Đáp số:} $800$\\
$200 \times 4 = 800$ đồng.
}
\end{ex}

\begin{ex} %[2H11V56]
Bình có 1 tờ 500 đồng và 2 tờ 200 đồng. Hỏi Bình có tất cả bao nhiêu đồng?
\shortans{900}
\loigiai{
\textbf{Đáp số:} $900$\\
$500 + 400 = 900$ đồng.
}
\end{ex}

\begin{ex} %[2H11C56]
Cần bao nhiêu tờ 200 đồng để đổi được 1 tờ 1 000 đồng?
\shortans{5}
\loigiai{
\textbf{Đáp số:} $5$\\
$1000 : 200 = 5$ tờ.
}
\end{ex}

\newpage

\TaoPhieuKiemTra{Chủ đề 11 -- Độ dài và đơn vị đo độ dài. tiền việt nam}{57}{Các số có ba chữ số}{https://loptoanthaylinh.vn/thi-online/toan-2/bai-57}
\textbf{Phần 1. Trắc nghiệm lựa chọn}
\setcounterpageref{ex}{0}

\begin{ex} %[2H11N57]
Để đo chiều dài của phòng học, người ta nên dùng đơn vị nào?
\choice
{Mi-li-mét}
{Ki-lô-mét}
{Đề-xi-mét}
{\True Mét}
\loigiai{
Đo phòng học dùng đơn vị mét.\\
\textbf{Đáp án của câu hỏi là D}
}
\end{ex}

\begin{ex} %[2H11N57]
Để đo khoảng cách giữa hai thành phố, người ta dùng đơn vị
\choice
{\True Ki-lô-mét}
{Mét}
{Xăng-ti-mét}
{Đề-xi-mét}
\loigiai{
Đo khoảng cách giữa 2 thành phố dùng ki-lô-mét.\\
\textbf{Đáp án của câu hỏi là A}
}
\end{ex}

\begin{ex} %[2H11H57]
Bảng lớp dài khoảng bao nhiêu?
\choice
{$3\text{ cm}$}
{$3\text{ km}$}
{$30\text{ m}$}
{\True $3\text{ m}$}
\loigiai{
Bảng lớp dài khoảng $3\text{ m}$.\\
\textbf{Đáp án của câu hỏi là D}
}
\end{ex}

\begin{ex} %[2H11H57]
Bút chì của em dài khoảng bao nhiêu?
\choice
{$15\text{ m}$}
{$15\text{ dm}$}
{\True $15\text{ cm}$}
{$15\text{ km}$}
\loigiai{
Bút chì dài khoảng $15\text{ cm}$.\\
\textbf{Đáp án của câu hỏi là C}
}
\end{ex}

\begin{ex} %[2H11H57]
Cột cờ trường học cao khoảng
\choice
{$8\text{ cm}$}
{$8\text{ km}$}
{$80\text{ m}$}
{\True $8\text{ m}$}
\loigiai{
Cột cờ cao khoảng $8\text{ m}$.\\
\textbf{Đáp án của câu hỏi là D}
}
\end{ex}

\begin{ex} %[2H11H57]
Một gang tay của học sinh lớp 2 dài khoảng
\choice
{$15\text{ m}$}
{\True $15\text{ cm}$ (hoặc hơn 1 dm)}
{$15\text{ km}$}
{$1\text{ cm}$}
\loigiai{
Gang tay dài khoảng $15\text{ cm}$.\\
\textbf{Đáp án của câu hỏi là B}
}
\end{ex}

\begin{ex} %[2H11V57]
Đoạn thẳng $AB$ dài $1\text{ m}$, đoạn thẳng $CD$ dài $80\text{ cm}$. Đoạn $AB$ dài hơn đoạn $CD$ là
\choice
{\True $20\text{ cm}$}
{$10\text{ cm}$}
{$30\text{ cm}$}
{$80\text{ cm}$}
\loigiai{
$1\text{ m} = 100\text{ cm}$.\\
$100 - 80 = 20\text{ cm}$.\\
\textbf{Đáp án của câu hỏi là A}
}
\end{ex}

\begin{ex} %[2H11V57]
Sợi dây thứ nhất dài $2\text{ m}$, sợi dây thứ hai dài $15\text{ dm}$. Cả hai sợi dây dài tất cả là
\choice
{$30\text{ dm}$}
{$40\text{ dm}$}
{\True $35\text{ dm}$}
{$25\text{ dm}$}
\loigiai{
$2\text{ m} = 20\text{ dm}$.\\
$20 + 15 = 35\text{ dm}$.\\
\textbf{Đáp án của câu hỏi là C}
}
\end{ex}

\begin{ex} %[2H11V57]
Bạn Nam bước đi 10 bước, mỗi bước dài khoảng $4\text{ dm}$. Nam đã đi được quãng đường dài
\choice
{$30\text{ dm}$}
{\True $40\text{ dm}$ ($4\text{ m}$)}
{$50\text{ dm}$}
{$4\text{ dm}$}
\loigiai{
$4 \times 10 = 40\text{ dm} = 4\text{ m}$.\\
\textbf{Đáp án của câu hỏi là B}
}
\end{ex}

\begin{ex} %[2H11V57]
Một mảnh vườn hình chữ nhật có chiều dài $20\text{ m}$, chiều rộng $10\text{ m}$. Nửa chu vi mảnh vườn là
\choice
{\True $30\text{ m}$}
{$25\text{ m}$}
{$35\text{ m}$}
{$40\text{ m}$}
\loigiai{
$20 + 10 = 30\text{ m}$.\\
\textbf{Đáp án của câu hỏi là A}
}
\end{ex}

\begin{ex} %[2H11V57]
Bác Ba trồng 5 cây cau thẳng hàng cách đều nhau, khoảng cách giữa 2 cây liên tiếp là $2\text{ m}$. Khoảng cách từ cây đầu tiên đến cây cuối cùng là
\choice
{$10\text{ m}$}
{\True $8\text{ m}$ (có 4 khoảng cách)}
{$6\text{ m}$}
{$12\text{ m}$}
\loigiai{
5 cây có $5 - 1 = 4$ khoảng cách.\\
$2 \times 4 = 8\text{ m}$.\\
\textbf{Đáp án của câu hỏi là B}
}
\end{ex}

\begin{ex} %[2H11C57]
Một con ốc sên bò lên một cây cột cao $5\text{ m}$. Ban ngày bò lên được $3\text{ m}$, ban đêm bị tụt xuống $1\text{ m}$. Hỏi sau bao nhiêu ngày đêm ốc sên bò lên đến đỉnh cột?
\choice
{$3$ ngày đêm}
{$4$ ngày đêm}
{\True $2$ ngày đêm}
{$5$ ngày đêm}
\loigiai{
Ngày 1 lên 3m, đêm 1 tụt 1m còn 2m.\\
Ngày 2 bò tiếp 3m: $2 + 3 = 5\text{ m}$ (đạt đỉnh luôn trong ngày thứ 2).\\
\textbf{Đáp án của câu hỏi là C}
}
\end{ex}

\textbf{Phần 2. Trắc nghiệm đúng sai}
\setcounterpageref{ex}{0}

\begin{ex} %[2H11H57]
Xét tính thực tế của các đơn vị đo độ dài:
\choiceTF
{\True Quyển sách Toán 2 dài khoảng $24\text{ cm}$.}
{\True Sân trường dài khoảng vài chục mét.}
{Khoảng cách từ Hà Nội đến Đà Nẵng dài khoảng $800\text{ m}$.}
{\True Khoảng cách từ Hà Nội đến Đà Nẵng dài khoảng gần $800\text{ km}$.}
\loigiai{
\begin{itemize}
    \item Quyển sách dài $24\text{ cm}$: Đúng.
    \item Sân trường dài vài chục mét: Đúng.
    \item Khoảng cách xa phải dùng $\text{km}$ (không phải $\text{m}$). Ý 3 Sai, Ý 4 Đúng.
\end{itemize}
\begin{itemize}
    \item a) \textbf{Đúng}.
    \item b) \textbf{Đúng}.
    \item c) \textbf{Sai}.
    \item d) \textbf{Đúng}.
\end{itemize}
}
\end{ex}

\begin{ex} %[2H11V57]
Sợi dây $A$ dài $1\text{ m}$, sợi dây $B$ dài $70\text{ cm}$. Xét tính đúng sai:
\choiceTF
{\True Sợi dây $A$ dài $100\text{ cm}$.}
{\True Sợi dây $A$ dài hơn sợi dây $B$ là $30\text{ cm}$ ($100 - 70 = 30$).}
{Tổng độ dài hai sợi dây là $80\text{ cm}$.}
{\True Tổng độ dài hai sợi dây là $170\text{ cm}$ ($100 + 70 = 170$).}
\loigiai{
\begin{itemize}
    \item $1\text{ m} = 100\text{ cm}$: Đúng.
    \item Hiệu: $100 - 70 = 30\text{ cm}$: Đúng.
    \item Tổng: $170\text{ cm} \neq 80\text{ cm}$. Ý 3 Sai, Ý 4 Đúng.
\end{itemize}
\begin{itemize}
    \item a) \textbf{Đúng}.
    \item b) \textbf{Đúng}.
    \item c) \textbf{Sai}.
    \item d) \textbf{Đúng}.
\end{itemize}
}
\end{ex}

\textbf{Phần 3. Trả lời ngắn}
\setcounterpageref{ex}{0}

\begin{ex} %[2H11H57]
Đổi: $2\text{ m} = \dots\text{ cm}$.
\shortans{200}
\loigiai{
\textbf{Đáp số:} $200$\\
200 cm.
}
\end{ex}

\begin{ex} %[2H11H57]
Đổi: $3\text{ m} = \dots\text{ dm}$.
\shortans{30}
\loigiai{
\textbf{Đáp số:} $30$\\
30 dm.
}
\end{ex}

\begin{ex} %[2H11V57]
Một đoạn đường dài $45\text{ m}$, đã rải nhựa được $25\text{ m}$. Hỏi còn bao nhiêu mét đường chưa rải nhựa?
\shortans{20}
\loigiai{
\textbf{Đáp số:} $20$\\
$45 - 25 = 20\text{ m}$.
}
\end{ex}

\begin{ex} %[2H11V57]
Quãng đường từ nhà An đến trường dài $800\text{ m}$, từ trường đến nhà bà ngoại dài thêm $200\text{ m}$. Tính quãng đường từ nhà An qua trường đến nhà bà ngoại theo mét.
\shortans{1000}
\loigiai{
\textbf{Đáp số:} $1000$\\
$800 + 200 = 1000\text{ m}$.
}
\end{ex}

\begin{ex} %[2H11V57]
Có 4 đoạn dây bằng nhau, mỗi đoạn dài $5\text{ m}$. Tính tổng độ dài 4 đoạn dây theo mét.
\shortans{20}
\loigiai{
\textbf{Đáp số:} $20$\\
$5 \times 4 = 20\text{ m}$.
}
\end{ex}

\begin{ex} %[2H11C57]
Bác thợ mộc có cây gỗ dài $8\text{ m}$, bác cưa thành các đoạn dài $2\text{ m}$. Hỏi bác phải cưa bao nhiêu lần?
\shortans{3}
\loigiai{
\textbf{Đáp số:} $3$\\
Số đoạn: $8 : 2 = 4$ đoạn.\\
Số lần cưa: $4 - 1 = 3$ lần.
}
\end{ex}

\newpage

\TaoPhieuKiemTra{Chủ đề 11 -- Độ dài và đơn vị đo độ dài. tiền việt nam}{58}{Viết số thành tổng các trăm, chục, đơn vị}{https://loptoanthaylinh.vn/thi-online/toan-2/bai-58}
\textbf{Phần 1. Trắc nghiệm lựa chọn}
\setcounterpageref{ex}{0}

\begin{ex} %[2H11N58]
Đơn vị nào sau đây dùng để đo độ dài lớn nhất?
\choice
{Mét}
{Đề-xi-mét}
{\True Ki-lô-mét}
{Xăng-ti-mét}
\loigiai{
Ki-lô-mét là đơn vị lớn nhất trong các đơn vị trên.\\
\textbf{Đáp án của câu hỏi là C}
}
\end{ex}

\begin{ex} %[2H11N58]
Tờ tiền mệnh giá 1 000 đồng bằng mấy tờ 500 đồng?
\choice
{\True $2$ tờ}
{$1$ tờ}
{$3$ tờ}
{$4$ tờ}
\loigiai{
$500 \times 2 = 1000$ đồng.\\
\textbf{Đáp án của câu hỏi là A}
}
\end{ex}

\begin{ex} %[2H11H58]
Tính: $45\text{ m} + 35\text{ m} = \dots$
\choice
{$70\text{ m}$}
{\True $80\text{ m}$}
{$90\text{ m}$}
{$85\text{ m}$}
\loigiai{
$45 + 35 = 80\text{ m}$.\\
\textbf{Đáp án của câu hỏi là B}
}
\end{ex}

\begin{ex} %[2H11H58]
Tính: $90\text{ km} - 45\text{ km} = \dots$
\choice
{$35\text{ km}$}
{$55\text{ km}$}
{\True $45\text{ km}$}
{$40\text{ km}$}
\loigiai{
$90 - 45 = 45\text{ km}$.\\
\textbf{Đáp án của câu hỏi là C}
}
\end{ex}

\begin{ex} %[2H11H58]
Mai có 2 tờ 200 đồng và 1 tờ 500 đồng. Mai có tất cả số tiền là
\choice
{$800$ đồng}
{$1000$ đồng}
{$700$ đồng}
{\True $900$ đồng}
\loigiai{
$200 \times 2 + 500 = 900$ đồng.\\
\textbf{Đáp án của câu hỏi là D}
}
\end{ex}

\begin{ex} %[2H11H58]
Điền dấu thích hợp: $1\text{ km} \dots 900\text{ m}$.
\choice
{\True $>$}
{$<$}
{$=$}
{Không so sánh được}
\loigiai{
$1\text{ km} = 1000\text{ m} > 900\text{ m}$.\\
\textbf{Đáp án của câu hỏi là A}
}
\end{ex}

\begin{ex} %[2H11V58]
Một cây cầu dài $350\text{ m}$, một cây cầu khác dài $450\text{ m}$. Cả hai cây cầu dài tổng cộng là
\choice
{$700\text{ m}$}
{$900\text{ m}$}
{$1000\text{ m}$}
{\True $800\text{ m}$}
\loigiai{
$350 + 450 = 800\text{ m}$.\\
\textbf{Đáp án của câu hỏi là D}
}
\end{ex}

\begin{ex} %[2H11V58]
Nam có 1 000 đồng, Nam mua 1 cái bánh giá 600 đồng. Số tiền Nam còn lại là
\choice
{$300$ đồng}
{\True $400$ đồng}
{$500$ đồng}
{$200$ đồng}
\loigiai{
$1000 - 600 = 400$ đồng.\\
\textbf{Đáp án của câu hỏi là B}
}
\end{ex}

\begin{ex} %[2H11V58]
Cuộn dây thứ nhất dài $48\text{ m}$, cuộn dây thứ hai ngắn hơn cuộn thứ nhất $15\text{ m}$. Cuộn dây thứ hai dài là
\choice
{$32\text{ m}$}
{$34\text{ m}$}
{\True $33\text{ m}$}
{$35\text{ m}$}
\loigiai{
$48 - 15 = 33\text{ m}$.\\
\textbf{Đáp án của câu hỏi là C}
}
\end{ex}

\begin{ex} %[2H11V58]
Tính giá trị biểu thức: $15\text{ m} \times 2 + 20\text{ m}$.
\choice
{$45\text{ m}$}
{$55\text{ m}$}
{$60\text{ m}$}
{\True $50\text{ m}$}
\loigiai{
$30 + 20 = 50\text{ m}$.\\
\textbf{Đáp án của câu hỏi là D}
}
\end{ex}

\begin{ex} %[2H11V58]
Mỗi chiếc kem giá 200 đồng. Mua 5 chiếc kem như thế hết số tiền là
\choice
{$800$ đồng}
{\True $1\,000$ đồng}
{$700$ đồng}
{$900$ đồng}
\loigiai{
$200 \times 5 = 1000$ đồng.\\
\textbf{Đáp án của câu hỏi là B}
}
\end{ex}

\begin{ex} %[2H11C58]
Hùng có 1 000 đồng, Dũng có 600 đồng. Hùng phải cho Dũng bao nhiêu đồng để số tiền của hai bạn bằng nhau?
\choice
{\True $200$ đồng}
{$400$ đồng}
{$300$ đồng}
{$100$ đồng}
\loigiai{
Hùng hơn Dũng: $1000 - 600 = 400$ đồng.\\
Hùng cho Dũng: $400 : 2 = 200$ đồng.\\
\textbf{Đáp án của câu hỏi là A}
}
\end{ex}

\textbf{Phần 2. Trắc nghiệm đúng sai}
\setcounterpageref{ex}{0}

\begin{ex} %[2H11H58]
Xét tính đúng sai về đại lượng và tiền tệ:
\choiceTF
{\True $1\text{ km} = 1\,000\text{ m}$.}
{\True 2 tờ 500 đồng = 1 000 đồng.}
{$1\text{ m} = 10\text{ cm}$.}
{\True $1\text{ m} = 100\text{ cm}$.}
\loigiai{
\begin{itemize}
    \item $1\text{ km} = 1000\text{ m}$: Đúng.
    \item 2 tờ 500đ = 1000đ: Đúng.
    \item $1\text{ m} = 100\text{ cm} \neq 10\text{ cm}$. Ý 3 Sai, Ý 4 Đúng.
\end{itemize}
\begin{itemize}
    \item a) \textbf{Đúng}.
    \item b) \textbf{Đúng}.
    \item c) \textbf{Sai}.
    \item d) \textbf{Đúng}.
\end{itemize}
}
\end{ex}

\begin{ex} %[2H11V58]
Quãng đường từ nhà đến trường dài $2\text{ km}$. Mỗi ngày An đi học cả đi lẫn về. Xét tính đúng sai:
\choiceTF
{\True Mỗi ngày An đi quãng đường là $4\text{ km}$ ($2 \times 2 = 4$).}
{\True Trong 5 ngày đi học, An đi tất cả $20\text{ km}$ ($4 \times 5 = 20$).}
{Trong 5 ngày An đi tất cả 10 km.}
{\True Nếu chỉ tính 1 chiều đi thì 5 ngày là 10 km, nhưng cả đi và về là 20 km.}
\loigiai{
\begin{itemize}
    \item 1 ngày cả đi về: 4 km: Đúng.
    \item 5 ngày: 20 km: Đúng.
    \item Cả đi và về là 20 km (không phải 10 km). Ý 3 Sai, Ý 4 Đúng.
\end{itemize}
\begin{itemize}
    \item a) \textbf{Đúng}.
    \item b) \textbf{Đúng}.
    \item c) \textbf{Sai}.
    \item d) \textbf{Đúng}.
\end{itemize}
}
\end{ex}

\textbf{Phần 3. Trả lời ngắn}
\setcounterpageref{ex}{0}

\begin{ex} %[2H11H58]
Tính: $36\text{ m} + 24\text{ m}$ theo mét.
\shortans{60}
\loigiai{
\textbf{Đáp số:} $60$\\
$36 + 24 = 60\text{ m}$.
}
\end{ex}

\begin{ex} %[2H11H58]
Tính: $100\text{ km} - 65\text{ km}$ theo ki-lô-mét.
\shortans{35}
\loigiai{
\textbf{Đáp số:} $35$\\
$100 - 65 = 35\text{ km}$.
}
\end{ex}

\begin{ex} %[2H11V58]
Có 3 tờ 200 đồng và 2 tờ 100 đồng. Hỏi có tất cả bao nhiêu đồng?
\shortans{800}
\loigiai{
\textbf{Đáp số:} $800$\\
$600 + 200 = 800$ đồng.
}
\end{ex}

\begin{ex} %[2H11V58]
Lan có 1 000 đồng, mua vở hết 800 đồng. Hỏi Lan còn lại bao nhiêu đồng?
\shortans{200}
\loigiai{
\textbf{Đáp số:} $200$\\
$1000 - 800 = 200$ đồng.
}
\end{ex}

\begin{ex} %[2H11V58]
Sợi dây dài $30\text{ dm}$ cắt thành 5 đoạn bằng nhau. Hỏi mỗi đoạn dài bao nhiêu đề-xi-mét?
\shortans{6}
\loigiai{
\textbf{Đáp số:} $6$\\
$30 : 5 = 6\text{ dm}$.
}
\end{ex}

\begin{ex} %[2H11C58]
Tính nhanh: $450\text{ m} + 280\text{ m} + 250\text{ m} + 20\text{ m}$ theo mét.
\shortans{1000}
\loigiai{
\textbf{Đáp số:} $1000$\\
$(450 + 250) + (280 + 20) = 700 + 300 = 1000\text{ m}$.
}
\end{ex}

\newpage

% =========================================================================
% CHỦ ĐỀ 12 - PHÉP CỘNG, PHÉP TRỪ TRONG PHẠM VI 1 000
% =========================================================================
\section{Phép cộng, phép trừ trong phạm vi 1 000}

\TaoPhieuKiemTra{Chủ đề 12 -- Phép cộng, phép trừ trong phạm vi 1 000}{59}{So sánh các số có ba chữ số}{https://loptoanthaylinh.vn/thi-online/toan-2/bai-59}
\textbf{Phần 1. Trắc nghiệm lựa chọn}
\setcounterpageref{ex}{0}

\begin{ex} %[2D12N59]
Kết quả của phép tính $324 + 143$ là
\choice
{\True $467$}
{$457$}
{$477$}
{$468$}
\loigiai{
$324 + 143 = 467$.\\
\textbf{Đáp án của câu hỏi là A}
}
\end{ex}

\begin{ex} %[2D12N59]
Kết quả của phép tính $500 + 300$ là
\choice
{$700$}
{$900$}
{$600$}
{\True $800$}
\loigiai{
$500 + 300 = 800$.\\
\textbf{Đáp án của câu hỏi là D}
}
\end{ex}

\begin{ex} %[2D12H59]
Kết quả của phép tính $250 + 130$ là
\choice
{$370$}
{$390$}
{\True $380$}
{$400$}
\loigiai{
$250 + 130 = 380$.\\
\textbf{Đáp án của câu hỏi là C}
}
\end{ex}

\begin{ex} %[2D12H59]
Kết quả của phép tính $416 + 232$ là
\choice
{$638$}
{\True $648$}
{$658$}
{$748$}
\loigiai{
$416 + 232 = 648$.\\
\textbf{Đáp án của câu hỏi là B}
}
\end{ex}

\begin{ex} %[2D12H59]
Tìm $x$ biết: $x - 213 = 342$.
\choice
{\True $555$}
{$545$}
{$565$}
{$535$}
\loigiai{
$x = 342 + 213 = 555$.\\
\textbf{Đáp án của câu hỏi là A}
}
\end{ex}

\begin{ex} %[2D12H59]
Điền dấu thích hợp: $352 + 124 \dots 241 + 235$.
\choice
{$<$}
{$>$}
{Không so sánh được}
{\True $=$}
\loigiai{
$352 + 124 = 476$ và $241 + 235 = 476$. Vậy bằng nhau.\\
\textbf{Đáp án của câu hỏi là D}
}
\end{ex}

\begin{ex} %[2D12V59]
Một trường tiểu học có 425 học sinh nam và 432 học sinh nữ. Trường có tất cả số học sinh là
\choice
{$847$ bạn}
{$867$ bạn}
{\True $857$ bạn}
{$877$ bạn}
\loigiai{
$425 + 432 = 857$ học sinh.\\
\textbf{Đáp án của câu hỏi là C}
}
\end{ex}

\begin{ex} %[2D12V59]
Thư viện có 320 quyển sách toán và 450 quyển sách tiếng việt. Thư viện có tất cả số quyển sách là
\choice
{$760$ quyển}
{$780$ quyển}
{$790$ quyển}
{\True $770$ quyển}
\loigiai{
$320 + 450 = 770$ quyển.\\
\textbf{Đáp án của câu hỏi là D}
}
\end{ex}

\begin{ex} %[2D12V59]
Tính giá trị biểu thức: $312 + 245 + 131$.
\choice
{$678$}
{\True $688$}
{$698$}
{$708$}
\loigiai{
$557 + 131 = 688$.\\
\textbf{Đáp án của câu hỏi là B}
}
\end{ex}

\begin{ex} %[2D12V59]
Đoạn đường thứ nhất dài $340\text{ m}$, đoạn đường thứ hai dài $420\text{ m}$. Cả hai đoạn đường dài tất cả là
\choice
{$750\text{ m}$}
{$770\text{ m}$}
{\True $760\text{ m}$}
{$780\text{ m}$}
\loigiai{
$340 + 420 = 760\text{ m}$.\\
\textbf{Đáp án của câu hỏi là C}
}
\end{ex}

\begin{ex} %[2D12V59]
Tìm tổng của số tròn trăm bé nhất có 3 chữ số và số 452.
\choice
{$452$}
{\True $552$}
{$652$}
{$752$}
\loigiai{
Số tròn trăm bé nhất có 3 chữ số là 100. Tổng: $100 + 452 = 552$.\\
\textbf{Đáp án của câu hỏi là B}
}
\end{ex}

\begin{ex} %[2D12C59]
Điền chữ số thích hợp vào dấu *: $3*5 + 241 = 576$.
\choice
{\True $3$}
{$2$}
{$4$}
{$5$}
\loigiai{
$* + 4 = 7 \Rightarrow * = 3$.\\
\textbf{Đáp án của câu hỏi là A}
}
\end{ex}

\textbf{Phần 2. Trắc nghiệm đúng sai}
\setcounterpageref{ex}{0}

\begin{ex} %[2D12H59]
Xét tính đúng sai của các phép tính cộng không nhớ:
\choiceTF
{\True $423 + 231 = 654$.}
{\True $510 + 340 = 850$.}
{$602 + 135 = 747$.}
{\True $602 + 135 = 737$.}
\loigiai{
\begin{itemize}
    \item $423 + 231 = 654$: Đúng.
    \item $510 + 340 = 850$: Đúng.
    \item $602 + 135 = 737 \neq 747$. Ý 3 Sai, Ý 4 Đúng.
\end{itemize}
\begin{itemize}
    \item a) \textbf{Đúng}.
    \item b) \textbf{Đúng}.
    \item c) \textbf{Sai}.
    \item d) \textbf{Đúng}.
\end{itemize}
}
\end{ex}

\begin{ex} %[2D12V59]
Một nông trại nuôi 320 con gà và 250 con vịt. Xét tính đúng sai:
\choiceTF
{\True Tổng số gia cầm là 570 con ($320 + 250 = 570$).}
{\True Số gà nhiều hơn số vịt là 70 con ($320 - 250 = 70$).}
{Nếu mua thêm 100 con gà thì đàn gà có 400 con.}
{\True Khi đó đàn gà có $320 + 100 = 420$ con.}
\loigiai{
\begin{itemize}
    \item Tổng: $320 + 250 = 570$: Đúng.
    \item Gà hơn vịt 70 con: Đúng.
    \item $320 + 100 = 420 \neq 400$. Ý 3 Sai, Ý 4 Đúng.
\end{itemize}
\begin{itemize}
    \item a) \textbf{Đúng}.
    \item b) \textbf{Đúng}.
    \item c) \textbf{Sai}.
    \item d) \textbf{Đúng}.
\end{itemize}
}
\end{ex}

\textbf{Phần 3. Trả lời ngắn}
\setcounterpageref{ex}{0}

\begin{ex} %[2D12H59]
Tính: $354 + 213$.
\shortans{567}
\loigiai{
\textbf{Đáp số:} $567$\\
$354 + 213 = 567$.
}
\end{ex}

\begin{ex} %[2D12H59]
Tính: $620 + 170$.
\shortans{790}
\loigiai{
\textbf{Đáp số:} $790$\\
$620 + 170 = 790$.
}
\end{ex}

\begin{ex} %[2D12V59]
Tính: $405 + 372$.
\shortans{777}
\loigiai{
\textbf{Đáp số:} $777$\\
$405 + 372 = 777$.
}
\end{ex}

\begin{ex} %[2D12V59]
Tìm $x$ biết: $x - 321 = 456$.
\shortans{777}
\loigiai{
\textbf{Đáp số:} $777$\\
$x = 456 + 321 = 777$.
}
\end{ex}

\begin{ex} %[2D12V59]
Một vườn cây có 230 cây cam và 340 cây bưởi. Hỏi vườn có tất cả bao nhiêu cây cam và bưởi?
\shortans{570}
\loigiai{
\textbf{Đáp số:} $570$\\
$230 + 340 = 570$ cây.
}
\end{ex}

\begin{ex} %[2D12C59]
Tính nhanh: $215 + 340 + 85$.
\shortans{640}
\loigiai{
\textbf{Đáp số:} $640$\\
$(215 + 85) + 340 = 300 + 340 = 640$.
}
\end{ex}

\newpage

\TaoPhieuKiemTra{Chủ đề 12 -- Phép cộng, phép trừ trong phạm vi 1 000}{60}{Luyện tập các số trong phạm vi 1000}{https://loptoanthaylinh.vn/thi-online/toan-2/bai-60}
\textbf{Phần 1. Trắc nghiệm lựa chọn}
\setcounterpageref{ex}{0}

\begin{ex} %[2D12N60]
Kết quả của phép tính $348 + 225$ là
\choice
{$563$}
{$583$}
{$574$}
{\True $573$}
\loigiai{
$348 + 225 = 573$.\\
\textbf{Đáp án của câu hỏi là D}
}
\end{ex}

\begin{ex} %[2D12N60]
Kết quả của phép tính $456 + 172$ là
\choice
{\True $628$}
{$618$}
{$638$}
{$528$}
\loigiai{
$456 + 172 = 628$.\\
\textbf{Đáp án của câu hỏi là A}
}
\end{ex}

\begin{ex} %[2D12H60]
Kết quả của phép tính $537 + 286$ là
\choice
{$813$}
{\True $823$}
{$833$}
{$824$}
\loigiai{
$537 + 286 = 823$.\\
\textbf{Đáp án của câu hỏi là B}
}
\end{ex}

\begin{ex} %[2D12H60]
Kết quả của phép tính $684 + 158$ là
\choice
{$832$}
{$852$}
{\True $842$}
{$742$}
\loigiai{
$684 + 158 = 842$.\\
\textbf{Đáp án của câu hỏi là C}
}
\end{ex}

\begin{ex} %[2D12H60]
Tìm $x$ biết: $x - 275 = 348$.
\choice
{\True $623$}
{$613$}
{$633$}
{$624$}
\loigiai{
$x = 348 + 275 = 623$.\\
\textbf{Đáp án của câu hỏi là A}
}
\end{ex}

\begin{ex} %[2D12H60]
Điền dấu thích hợp: $368 + 245 \dots 412 + 195$.
\choice
{$<$}
{$=$}
{Không so sánh được}
{\True $>$}
\loigiai{
$368 + 245 = 613 > 412 + 195 = 607$.\\
\textbf{Đáp án của câu hỏi là D}
}
\end{ex}

\begin{ex} %[2D12V60]
Một trang trại có 385 con vịt và 248 con ngan. Trang trại có tất cả số con vịt và ngan là
\choice
{$623$ con}
{\True $633$ con}
{$643$ con}
{$634$ con}
\loigiai{
$385 + 248 = 633$ con.\\
\textbf{Đáp án của câu hỏi là B}
}
\end{ex}

\begin{ex} %[2D12V60]
Một cửa hàng buổi sáng bán được $468\text{ kg}$ gạo, buổi chiều bán được $375\text{ kg}$ gạo. Cả ngày cửa hàng bán được số ki-lô-gam gạo là
\choice
{$833\text{ kg}$}
{$853\text{ kg}$}
{\True $843\text{ kg}$}
{$844\text{ kg}$}
\loigiai{
$468 + 375 = 843\text{ kg}$.\\
\textbf{Đáp án của câu hỏi là C}
}
\end{ex}

\begin{ex} %[2D12V60]
Tính giá trị biểu thức: $275 + 186 + 139$.
\choice
{$590$}
{\True $600$}
{$610$}
{$620$}
\loigiai{
$461 + 139 = 600$.\\
\textbf{Đáp án của câu hỏi là B}
}
\end{ex}

\begin{ex} %[2D12V60]
Quãng đường thứ nhất dài $378\text{ m}$, quãng đường thứ hai dài $456\text{ m}$. Cả hai quãng đường dài tất cả là
\choice
{\True $834\text{ m}$}
{$824\text{ m}$}
{$844\text{ m}$}
{$835\text{ m}$}
\loigiai{
$378 + 456 = 834\text{ m}$.\\
\textbf{Đáp án của câu hỏi là A}
}
\end{ex}

\begin{ex} %[2D12V60]
Tính nhanh: $368 + 275 + 132 + 125$.
\choice
{$890$}
{$910$}
{$850$}
{\True $900$}
\loigiai{
$(368 + 132) + (275 + 125) = 500 + 400 = 900$.\\
\textbf{Đáp án của câu hỏi là D}
}
\end{ex}

\begin{ex} %[2D12C60]
Tổng của số bé nhất có ba chữ số khác nhau ($102$) và số lớn nhất có ba chữ số khác nhau ($987$) là
\choice
{$1079$}
{$1099$}
{\True $1089$}
{$1088$}
\loigiai{
$102 + 987 = 1089$.\\
\textbf{Đáp án của câu hỏi là C}
}
\end{ex}

\textbf{Phần 2. Trắc nghiệm đúng sai}
\setcounterpageref{ex}{0}

\begin{ex} %[2D12H60]
Xét tính đúng sai của các phép cộng có nhớ:
\choiceTF
{\True $376 + 258 = 634$.}
{\True $489 + 175 = 664$.}
{$528 + 394 = 912$.}
{\True $528 + 394 = 922$.}
\loigiai{
\begin{itemize}
    \item $376 + 258 = 634$: Đúng.
    \item $489 + 175 = 664$: Đúng.
    \item $528 + 394 = 922 \neq 912$. Ý 3 Sai, Ý 4 Đúng.
\end{itemize}
\begin{itemize}
    \item a) \textbf{Đúng}.
    \item b) \textbf{Đúng}.
    \item c) \textbf{Sai}.
    \item d) \textbf{Đúng}.
\end{itemize}
}
\end{ex}

\begin{ex} %[2D12V60]
Một thư viện có 458 cuốn truyện thiếu nhi và 365 cuốn sách giáo khoa. Xét tính đúng sai:
\choiceTF
{\True Tổng số sách truyện là 823 cuốn ($458 + 365 = 823$).}
{\True Truyện thiếu nhi nhiều hơn sách giáo khoa là 93 cuốn ($458 - 365 = 93$).}
{Tổng số sách ít hơn 800 cuốn.}
{\True Tổng số sách là 823 cuốn (nhiều hơn 800 cuốn).}
\loigiai{
\begin{itemize}
    \item Tổng: 823 cuốn: Đúng.
    \item Hiệu: 93 cuốn: Đúng.
    \item $823 > 800$ (không ít hơn). Ý 3 Sai, Ý 4 Đúng.
\end{itemize}
\begin{itemize}
    \item a) \textbf{Đúng}.
    \item b) \textbf{Đúng}.
    \item c) \textbf{Sai}.
    \item d) \textbf{Đúng}.
\end{itemize}
}
\end{ex}

\textbf{Phần 3. Trả lời ngắn}
\setcounterpageref{ex}{0}

\begin{ex} %[2D12H60]
Tính tổng: $468 + 257$.
\shortans{725}
\loigiai{
\textbf{Đáp số:} $725$\\
$468 + 257 = 725$.
}
\end{ex}

\begin{ex} %[2D12H60]
Tính tổng: $579 + 346$.
\shortans{925}
\loigiai{
\textbf{Đáp số:} $925$\\
$579 + 346 = 925$.
}
\end{ex}

\begin{ex} %[2D12V60]
Tìm $x$ biết: $x - 389 = 456$.
\shortans{845}
\loigiai{
\textbf{Đáp số:} $845$\\
$x = 456 + 389 = 845$.
}
\end{ex}

\begin{ex} %[2D12V60]
Một đội công nhân tuần đầu trải được $375\text{ m}$ đường, tuần sau trải được $485\text{ m}$ đường. Cả hai tuần trải được bao nhiêu mét đường?
\shortans{860}
\loigiai{
\textbf{Đáp số:} $860$\\
$375 + 485 = 860\text{ m}$.
}
\end{ex}

\begin{ex} %[2D12V60]
Tính giá trị biểu thức: $348 + 265 + 187$.
\shortans{800}
\loigiai{
\textbf{Đáp số:} $800$\\
$613 + 187 = 800$.
}
\end{ex}

\begin{ex} %[2D12C60]
Tính nhanh tổng: $185 + 246 + 215 + 154$.
\shortans{800}
\loigiai{
\textbf{Đáp số:} $800$\\
$(185 + 215) + (246 + 154) = 400 + 400 = 800$.
}
\end{ex}

\newpage

\TaoPhieuKiemTra{Chủ đề 12 -- Phép cộng, phép trừ trong phạm vi 1 000}{61}{Phép cộng (không nhớ) trong phạm vi 1000}{https://loptoanthaylinh.vn/thi-online/toan-2/bai-61}
\textbf{Phần 1. Trắc nghiệm lựa chọn}
\setcounterpageref{ex}{0}

\begin{ex} %[2D12N61]
Kết quả của phép tính $567 - 234$ là
\choice
{$323$}
{$343$}
{\True $333$}
{$334$}
\loigiai{
$567 - 234 = 333$.\\
\textbf{Đáp án của câu hỏi là C}
}
\end{ex}

\begin{ex} %[2D12N61]
Kết quả của phép tính $800 - 300$ là
\choice
{$400$}
{\True $500$}
{$600$}
{$50$}
\loigiai{
$800 - 300 = 500$.\\
\textbf{Đáp án của câu hỏi là B}
}
\end{ex}

\begin{ex} %[2D12H61]
Kết quả của phép tính $789 - 456$ là
\choice
{\True $333$}
{$323$}
{$343$}
{$233$}
\loigiai{
$789 - 456 = 333$.\\
\textbf{Đáp án của câu hỏi là A}
}
\end{ex}

\begin{ex} %[2D12H61]
Kết quả của phép tính $650 - 230$ là
\choice
{$410$}
{$430$}
{$440$}
{\True $420$}
\loigiai{
$650 - 230 = 420$.\\
\textbf{Đáp án của câu hỏi là D}
}
\end{ex}

\begin{ex} %[2D12H61]
Tìm $x$ biết: $x + 231 = 654$.
\choice
{$413$}
{\True $423$}
{$433$}
{$424$}
\loigiai{
$x = 654 - 231 = 423$.\\
\textbf{Đáp án của câu hỏi là B}
}
\end{ex}

\begin{ex} %[2D12H61]
Điền dấu thích hợp: $758 - 324 \dots 689 - 255$.
\choice
{\True $=$}
{$>$}
{$<$}
{Không so sánh được}
\loigiai{
$758 - 324 = 434$ và $689 - 255 = 434$. Vậy bằng nhau.\\
\textbf{Đáp án của câu hỏi là A}
}
\end{ex}

\begin{ex} %[2D12V61]
Một cửa hàng có $680\text{ kg}$ gạo, đã bán đi $250\text{ kg}$. Cửa hàng còn lại số ki-lô-gam gạo là
\choice
{$420\text{ kg}$}
{$440\text{ kg}$}
{$450\text{ kg}$}
{\True $430\text{ kg}$}
\loigiai{
$680 - 250 = 430\text{ kg}$.\\
\textbf{Đáp án của câu hỏi là D}
}
\end{ex}

\begin{ex} %[2D12V61]
Một cuộn vải dài $485\text{ m}$, người ta đã cắt bán đi $162\text{ m}$. Cuộn vải còn lại dài là
\choice
{$313\text{ m}$}
{$333\text{ m}$}
{\True $323\text{ m}$}
{$324\text{ m}$}
\loigiai{
$485 - 162 = 323\text{ m}$.\\
\textbf{Đáp án của câu hỏi là C}
}
\end{ex}

\begin{ex} %[2D12V61]
Tính giá trị biểu thức: $896 - 342 - 231$.
\choice
{$313$}
{\True $323$}
{$333$}
{$343$}
\loigiai{
$554 - 231 = 323$.\\
\textbf{Đáp án của câu hỏi là B}
}
\end{ex}

\begin{ex} %[2D12V61]
Tìm $y$ biết: $975 - y = 432$.
\choice
{$533$}
{$553$}
{$544$}
{\True $543$}
\loigiai{
$y = 975 - 432 = 543$.\\
\textbf{Đáp án của câu hỏi là D}
}
\end{ex}

\begin{ex} %[2D12V61]
Khối Hai có 368 học sinh, trong đó có 145 học sinh nam. Khối Hai có số học sinh nữ là
\choice
{\True $223$ bạn}
{$213$ bạn}
{$233$ bạn}
{$224$ bạn}
\loigiai{
$368 - 145 = 223$ bạn.\\
\textbf{Đáp án của câu hỏi là A}
}
\end{ex}

\begin{ex} %[2D12C61]
Hiệu của số lớn nhất có 3 chữ số ($999$) và số tròn trăm lớn nhất có 3 chữ số ($900$) là
\choice
{$90$}
{$9$}
{\True $99$}
{$100$}
\loigiai{
$999 - 900 = 99$.\\
\textbf{Đáp án của câu hỏi là C}
}
\end{ex}

\textbf{Phần 2. Trắc nghiệm đúng sai}
\setcounterpageref{ex}{0}

\begin{ex} %[2D12H61]
Xét tính đúng sai của các phép trừ không nhớ:
\choiceTF
{\True $678 - 345 = 333$.}
{\True $890 - 450 = 440$.}
{$756 - 213 = 542$.}
{\True $756 - 213 = 543$.}
\loigiai{
\begin{itemize}
    \item $678 - 345 = 333$: Đúng.
    \item $890 - 450 = 440$: Đúng.
    \item $756 - 213 = 543 \neq 542$. Ý 3 Sai, Ý 4 Đúng.
\end{itemize}
\begin{itemize}
    \item a) \textbf{Đúng}.
    \item b) \textbf{Đúng}.
    \item c) \textbf{Sai}.
    \item d) \textbf{Đúng}.
\end{itemize}
}
\end{ex}

\begin{ex} %[2D12V61]
Một bể chứa $850\text{ l}$ nước, đã dùng hết $320\text{ l}$. Xét tính đúng sai:
\choiceTF
{\True Trong bể còn lại $530\text{ l}$ nước ($850 - 320 = 530$).}
{\True Lượng nước còn lại nhiều hơn lượng nước đã dùng là $210\text{ l}$ ($530 - 320 = 210$).}
{Nếu xả thêm $100\text{ l}$ nữa thì bể còn $400\text{ l}$.}
{\True Khi đó trong bể còn $530 - 100 = 430\text{ l}$ nước.}
\loigiai{
\begin{itemize}
    \item Còn lại: $530\text{ l}$: Đúng.
    \item Hiệu: $210\text{ l}$: Đúng.
    \item $530 - 100 = 430 \neq 400$. Ý 3 Sai, Ý 4 Đúng.
\end{itemize}
\begin{itemize}
    \item a) \textbf{Đúng}.
    \item b) \textbf{Đúng}.
    \item c) \textbf{Sai}.
    \item d) \textbf{Đúng}.
\end{itemize}
}
\end{ex}

\textbf{Phần 3. Trả lời ngắn}
\setcounterpageref{ex}{0}

\begin{ex} %[2D12H61]
Tính: $687 - 245$.
\shortans{442}
\loigiai{
\textbf{Đáp số:} $442$\\
$687 - 245 = 442$.
}
\end{ex}

\begin{ex} %[2D12H61]
Tính: $890 - 360$.
\shortans{530}
\loigiai{
\textbf{Đáp số:} $530$\\
$890 - 360 = 530$.
}
\end{ex}

\begin{ex} %[2D12V61]
Tính giá trị biểu thức: $768 - 234 - 112$.
\shortans{422}
\loigiai{
\textbf{Đáp số:} $422$\\
$534 - 112 = 422$.
}
\end{ex}

\begin{ex} %[2D12V61]
Tìm $x$ biết: $x + 314 = 859$.
\shortans{545}
\loigiai{
\textbf{Đáp số:} $545$\\
$x = 859 - 314 = 545$.
}
\end{ex}

\begin{ex} %[2D12V61]
Một nông trường có 670 con bò, trong đó có 340 con bò sữa. Hỏi có bao nhiêu con bò thịt?
\shortans{330}
\loigiai{
\textbf{Đáp số:} $330$\\
$670 - 340 = 330$ con.
}
\end{ex}

\begin{ex} %[2D12C61]
Tính nhanh: $875 - 230 - 175$.
\shortans{470}
\loigiai{
\textbf{Đáp số:} $470$\\
$(875 - 175) - 230 = 700 - 230 = 470$.
}
\end{ex}

\newpage

\TaoPhieuKiemTra{Chủ đề 12 -- Phép cộng, phép trừ trong phạm vi 1 000}{62}{Phép trừ (không nhớ) trong phạm vi 1000}{https://loptoanthaylinh.vn/thi-online/toan-2/bai-62}
\textbf{Phần 1. Trắc nghiệm lựa chọn}
\setcounterpageref{ex}{0}

\begin{ex} %[2D12N62]
Kết quả của phép tính $542 - 218$ là
\choice
{$314$}
{$334$}
{\True $324$}
{$325$}
\loigiai{
$542 - 218 = 324$.\\
\textbf{Đáp án của câu hỏi là C}
}
\end{ex}

\begin{ex} %[2D12N62]
Kết quả của phép tính $635 - 281$ là
\choice
{$344$}
{\True $354$}
{$364$}
{$355$}
\loigiai{
$635 - 281 = 354$.\\
\textbf{Đáp án của câu hỏi là B}
}
\end{ex}

\begin{ex} %[2D12H62]
Kết quả của phép tính $724 - 368$ là
\choice
{$346$}
{$366$}
{\True $356$}
{$358$}
\loigiai{
$724 - 368 = 356$.\\
\textbf{Đáp án của câu hỏi là C}
}
\end{ex}

\begin{ex} %[2D12H62]
Kết quả của phép tính $800 - 345$ là
\choice
{$445$}
{$465$}
{$555$}
{\True $455$}
\loigiai{
$800 - 345 = 455$.\\
\textbf{Đáp án của câu hỏi là D}
}
\end{ex}

\begin{ex} %[2D12H62]
Tìm $x$ biết: $x + 286 = 631$.
\choice
{\True $345$}
{$335$}
{$355$}
{$346$}
\loigiai{
$x = 631 - 286 = 345$.\\
\textbf{Đáp án của câu hỏi là A}
}
\end{ex}

\begin{ex} %[2D12H62]
Điền dấu thích hợp: $612 - 248 \dots 705 - 349$.
\choice
{$>$}
{\True $=$}
{$<$}
{Không so sánh được}
\loigiai{
$612 - 248 = 364$ và $705 - 341 = 364$. Do đó $364 = 364$.\\
\textbf{Đáp án của câu hỏi là B}
}
\end{ex}

\begin{ex} %[2D12V62]
Một kho thóc có $920\text{ kg}$ thóc, người ta xuất kho $475\text{ kg}$. Trong kho còn lại số ki-lô-gam thóc là
\choice
{$435\text{ kg}$}
{$455\text{ kg}$}
{$446\text{ kg}$}
{\True $445\text{ kg}$}
\loigiai{
$920 - 475 = 445\text{ kg}$.\\
\textbf{Đáp án của câu hỏi là D}
}
\end{ex}

\begin{ex} %[2D12V62]
Một cuộn cáp dài $1\,000\text{ m}$, người ta đã dùng hết $625\text{ m}$. Cuộn cáp còn lại dài là
\choice
{$365\text{ m}$}
{$385\text{ m}$}
{\True $375\text{ m}$}
{$376\text{ m}$}
\loigiai{
$1000 - 625 = 375\text{ m}$.\\
\textbf{Đáp án của câu hỏi là C}
}
\end{ex}

\begin{ex} %[2D12V62]
Tính giá trị biểu thức: $750 - 285 - 165$.
\choice
{\True $300$}
{$290$}
{$310$}
{$320$}
\loigiai{
$465 - 165 = 300$.\\
\textbf{Đáp án của câu hỏi là A}
}
\end{ex}

\begin{ex} %[2D12V62]
Tìm $y$ biết: $814 - y = 359$.
\choice
{$445$}
{\True $455$}
{$465$}
{$456$}
\loigiai{
$y = 814 - 359 = 455$.\\
\textbf{Đáp án của câu hỏi là B}
}
\end{ex}

\begin{ex} %[2D12V62]
Thùng thứ nhất có $625\text{ l}$ dầu, thùng thứ hai có ít hơn thùng thứ nhất $180\text{ l}$ dầu. Thùng thứ hai có số lít dầu là
\choice
{$435\text{ l}$}
{$455\text{ l}$}
{$446\text{ l}$}
{\True $445\text{ l}$}
\loigiai{
$625 - 180 = 445\text{ l}$.\\
\textbf{Đáp án của câu hỏi là D}
}
\end{ex}

\begin{ex} %[2D12C62]
Hiệu của $1\,000$ và số lớn nhất có ba chữ số khác nhau ($987$) là
\choice
{\True $13$}
{$12$}
{$14$}
{$15$}
\loigiai{
$1000 - 987 = 13$.\\
\textbf{Đáp án của câu hỏi là A}
}
\end{ex}

\textbf{Phần 2. Trắc nghiệm đúng sai}
\setcounterpageref{ex}{0}

\begin{ex} %[2D12H62]
Xét tính đúng sai của các phép trừ có nhớ:
\choiceTF
{\True $625 - 348 = 277$.}
{\True $800 - 462 = 338$.}
{$714 - 258 = 466$.}
{\True $714 - 258 = 456$.}
\loigiai{
\begin{itemize}
    \item $625 - 348 = 277$: Đúng.
    \item $800 - 462 = 338$: Đúng.
    \item $714 - 258 = 456 \neq 466$. Ý 3 Sai, Ý 4 Đúng.
\end{itemize}
\begin{itemize}
    \item a) \textbf{Đúng}.
    \item b) \textbf{Đúng}.
    \item c) \textbf{Sai}.
    \item d) \textbf{Đúng}.
\end{itemize}
}
\end{ex}

\begin{ex} %[2D12V62]
Cửa hàng có $840\text{ kg}$ xi măng, đã bán $375\text{ kg}$. Xét tính đúng sai:
\choiceTF
{\True Cửa hàng còn lại $465\text{ kg}$ xi măng ($840 - 375 = 465$).}
{\True Số xi măng còn lại nhiều hơn số đã bán là $90\text{ kg}$ ($465 - 375 = 90$).}
{Nếu bán thêm $65\text{ kg}$ nữa thì còn lại $410\text{ kg}$.}
{\True Khi đó còn lại $465 - 65 = 400\text{ kg}$.}
\loigiai{
\begin{itemize}
    \item Còn lại: $465\text{ kg}$: Đúng.
    \item Hiệu: $90\text{ kg}$: Đúng.
    \item $465 - 65 = 400 \neq 410$. Ý 3 Sai, Ý 4 Đúng.
\end{itemize}
\begin{itemize}
    \item a) \textbf{Đúng}.
    \item b) \textbf{Đúng}.
    \item c) \textbf{Sai}.
    \item d) \textbf{Đúng}.
\end{itemize}
}
\end{ex}

\textbf{Phần 3. Trả lời ngắn}
\setcounterpageref{ex}{0}

\begin{ex} %[2D12H62]
Tính: $634 - 278$.
\shortans{356}
\loigiai{
\textbf{Đáp số:} $356$\\
$634 - 278 = 356$.
}
\end{ex}

\begin{ex} %[2D12H62]
Tính: $900 - 456$.
\shortans{444}
\loigiai{
\textbf{Đáp số:} $444$\\
$900 - 456 = 444$.
}
\end{ex}

\begin{ex} %[2D12V62]
Tính: $1000 - 648$.
\shortans{352}
\loigiai{
\textbf{Đáp số:} $352$\\
$1000 - 648 = 352$.
}
\end{ex}

\begin{ex} %[2D12V62]
Tìm $x$ biết: $x + 367 = 812$.
\shortans{445}
\loigiai{
\textbf{Đáp số:} $445$\\
$x = 812 - 367 = 445$.
}
\end{ex}

\begin{ex} %[2D12V62]
Một đoạn đường dài $850\text{ m}$, đội công nhân đã làm xong $465\text{ m}$. Hỏi còn bao nhiêu mét đường chưa làm?
\shortans{385}
\loigiai{
\textbf{Đáp số:} $385$\\
$850 - 465 = 385\text{ m}$.
}
\end{ex}

\begin{ex} %[2D12C62]
Tính nhanh: $945 - 280 - 220$.
\shortans{445}
\loigiai{
\textbf{Đáp số:} $445$\\
$945 - (280 + 220) = 945 - 500 = 445$.
}
\end{ex}

\newpage

\TaoPhieuKiemTra{Chủ đề 12 -- Phép cộng, phép trừ trong phạm vi 1 000}{63}{Phép cộng (có nhớ) trong phạm vi 1000}{https://loptoanthaylinh.vn/thi-online/toan-2/bai-63}
\textbf{Phần 1. Trắc nghiệm lựa chọn}
\setcounterpageref{ex}{0}

\begin{ex} %[2D12N63]
Kết quả của phép tính $456 + 238$ là
\choice
{\True $694$}
{$684$}
{$704$}
{$695$}
\loigiai{
$456 + 238 = 694$.\\
\textbf{Đáp án của câu hỏi là A}
}
\end{ex}

\begin{ex} %[2D12N63]
Kết quả của phép tính $782 - 345$ là
\choice
{$427$}
{$447$}
{\True $437$}
{$438$}
\loigiai{
$782 - 345 = 437$.\\
\textbf{Đáp án của câu hỏi là C}
}
\end{ex}

\begin{ex} %[2D12H63]
Kết quả của phép tính $539 + 284$ là
\choice
{$813$}
{$833$}
{$824$}
{\True $823$}
\loigiai{
$539 + 284 = 823$.\\
\textbf{Đáp án của câu hỏi là D}
}
\end{ex}

\begin{ex} %[2D12H63]
Kết quả của phép tính $900 - 468$ là
\choice
{$422$}
{$442$}
{\True $432$}
{$433$}
\loigiai{
$900 - 468 = 432$.\\
\textbf{Đáp án của câu hỏi là C}
}
\end{ex}

\begin{ex} %[2D12H63]
Tìm $x$ biết: $x - 348 = 275$.
\choice
{$613$}
{\True $623$}
{$633$}
{$624$}
\loigiai{
$x = 275 + 348 = 623$.\\
\textbf{Đáp án của câu hỏi là B}
}
\end{ex}

\begin{ex} %[2D12H63]
Tính giá trị biểu thức: $450 + 280 - 360$.
\choice
{\True $370$}
{$360$}
{$380$}
{$390$}
\loigiai{
$730 - 360 = 370$.\\
\textbf{Đáp án của câu hỏi là A}
}
\end{ex}

\begin{ex} %[2D12V63]
Một trang trại nuôi 485 con gà đẻ trứng và 365 con gà thịt. Tổng số gà trang trại có là
\choice
{$840$ con}
{\True $850$ con}
{$860$ con}
{$870$ con}
\loigiai{
$485 + 365 = 850$ con.\\
\textbf{Đáp án của câu hỏi là B}
}
\end{ex}

\begin{ex} %[2D12V63]
Kho hàng có $950\text{ kg}$ đường, đã xuất đi $485\text{ kg}$. Trong kho còn lại số ki-lô-gam đường là
\choice
{$455\text{ kg}$}
{$475\text{ kg}$}
{\True $465\text{ kg}$}
{$466\text{ kg}$}
\loigiai{
$950 - 485 = 465\text{ kg}$.\\
\textbf{Đáp án của câu hỏi là C}
}
\end{ex}

\begin{ex} %[2D12V63]
Đoạn đường từ tỉnh $A$ đến tỉnh $B$ dài $145\text{ km}$, từ tỉnh $B$ đến tỉnh $C$ dài $180\text{ km}$. Quãng đường từ $A$ qua $B$ đến $C$ dài là
\choice
{$315\text{ km}$}
{$335\text{ km}$}
{$345\text{ km}$}
{\True $325\text{ km}$}
\loigiai{
$145 + 180 = 325\text{ km}$.\\
\textbf{Đáp án của câu hỏi là D}
}
\end{ex}

\begin{ex} %[2D12V63]
Tìm $y$ biết: $820 - y = 345 + 185$.
\choice
{\True $290$}
{$280$}
{$300$}
{$310$}
\loigiai{
$820 - y = 530 \Rightarrow y = 820 - 530 = 290$.\\
\textbf{Đáp án của câu hỏi là A}
}
\end{ex}

\begin{ex} %[2D12V63]
Tính nhanh: $275 + 184 + 125 + 216$.
\choice
{$790$}
{\True $800$}
{$810$}
{$820$}
\loigiai{
$(275 + 125) + (184 + 216) = 400 + 400 = 800$.\\
\textbf{Đáp án của câu hỏi là B}
}
\end{ex}

\begin{ex} %[2D12C63]
Tổng của hai số là 850. Nếu giữ nguyên số thứ nhất và bớt ở số thứ hai đi 60 đơn vị thì tổng mới là
\choice
{$910$}
{$800$}
{$850$}
{\True $790$}
\loigiai{
$850 - 60 = 790$.\\
\textbf{Đáp án của câu hỏi là D}
}
\end{ex}

\textbf{Phần 2. Trắc nghiệm đúng sai}
\setcounterpageref{ex}{0}

\begin{ex} %[2D12H63]
Xét tính đúng sai của các phép toán trong phạm vi 1 000:
\choiceTF
{\True $480 + 350 = 830$.}
{\True $830 - 480 = 350$.}
{$900 - 350 = 650$.}
{\True $900 - 350 = 550$.}
\loigiai{
\begin{itemize}
    \item $480 + 350 = 830$: Đúng.
    \item $830 - 480 = 350$: Đúng.
    \item $900 - 350 = 550 \neq 650$. Ý 3 Sai, Ý 4 Đúng.
\end{itemize}
\begin{itemize}
    \item a) \textbf{Đúng}.
    \item b) \textbf{Đúng}.
    \item c) \textbf{Sai}.
    \item d) \textbf{Đúng}.
\end{itemize}
}
\end{ex}

\begin{ex} %[2D12V63]
Một xí nghiệp may được 520 cái áo và 380 cái quần. Xét tính đúng sai:
\choiceTF
{\True Tổng số sản phẩm may được là 900 cái ($520 + 380 = 900$).}
{\True Số áo nhiều hơn số quần là 140 cái ($520 - 380 = 140$).}
{Xí nghiệp đã may được 1 000 sản phẩm.}
{\True Xí nghiệp còn thiếu 100 sản phẩm nữa mới đủ 1 000 cái.}
\loigiai{
\begin{itemize}
    \item Tổng: 900 cái: Đúng.
    \item Áo hơn quần: 140 cái: Đúng.
    \item Chưa đủ 1000 cái. Ý 3 Sai, Ý 4 Đúng.
\end{itemize}
\begin{itemize}
    \item a) \textbf{Đúng}.
    \item b) \textbf{Đúng}.
    \item c) \textbf{Sai}.
    \item d) \textbf{Đúng}.
\end{itemize}
}
\end{ex}

\textbf{Phần 3. Trả lời ngắn}
\setcounterpageref{ex}{0}

\begin{ex} %[2D12H63]
Tính: $567 + 285$.
\shortans{852}
\loigiai{
\textbf{Đáp số:} $852$\\
$567 + 285 = 852$.
}
\end{ex}

\begin{ex} %[2D12H63]
Tính: $824 - 458$.
\shortans{366}
\loigiai{
\textbf{Đáp số:} $366$\\
$824 - 458 = 366$.
}
\end{ex}

\begin{ex} %[2D12V63]
Tính giá trị biểu thức: $350 + 480 - 230$.
\shortans{600}
\loigiai{
\textbf{Đáp số:} $600$\\
$830 - 230 = 600$.
}
\end{ex}

\begin{ex} %[2D12V63]
Tìm $x$ biết: $x + 285 = 730$.
\shortans{445}
\loigiai{
\textbf{Đáp số:} $445$\\
$x = 730 - 285 = 445$.
}
\end{ex}

\begin{ex} %[2D12V63]
Một bể nước có $1\,000\text{ l}$, đã dùng hết $480\text{ l}$. Hỏi còn lại bao nhiêu lít nước?
\shortans{520}
\loigiai{
\textbf{Đáp số:} $520$\\
$1000 - 480 = 520\text{ l}$.
}
\end{ex}

\begin{ex} %[2D12C63]
Tính nhanh: $980 - 275 - 225$.
\shortans{480}
\loigiai{
\textbf{Đáp số:} $480$\\
$980 - (275 + 225) = 980 - 500 = 480$.
}
\end{ex}

\newpage

% =========================================================================
% CHỦ ĐỀ 13 - LÀM QUEN VỚI YẾU TỐ THỐNG KÊ, XÁC SUẤT
% =========================================================================
\section{Làm quen với yếu tố thống kê, xác suất}

\TaoPhieuKiemTra{Chủ đề 13 -- Làm quen với yếu tố thống kê, xác suất}{64}{Phép trừ (có nhớ) trong phạm vi 1000}{https://loptoanthaylinh.vn/thi-online/toan-2/bai-64}
\textbf{Phần 1. Trắc nghiệm lựa chọn}
\setcounterpageref{ex}{0}

\begin{ex} %[2X13N64]
Việc đếm số học sinh nam và số học sinh nữ trong lớp học được gọi là
\choice
{Vẽ tranh}
{\True Thu thập, phân loại và kiểm đếm số liệu}
{Giải toán có lời văn}
{Tính nhẩm}
\loigiai{
Đó là hoạt động thu thập, phân loại và kiểm đếm số liệu.\\
\textbf{Đáp án của câu hỏi là B}
}
\end{ex}

\begin{ex} %[2X13N64]
Trong một hộp có 5 viên bi đỏ, 4 viên bi xanh và 3 viên bi vàng. Các viên bi được phân loại theo
\choice
{Kích thước}
{Khối lượng}
{\True Màu sắc}
{Chất liệu}
\loigiai{
Phân loại theo màu sắc (đỏ, xanh, vàng).\\
\textbf{Đáp án của câu hỏi là C}
}
\end{ex}

\begin{ex} %[2X13H64]
Kiểm đếm các đồ dùng học tập của Lan: 4 cái bút chì, 2 cục tẩy, 1 cái thước kẻ. Lan có tất cả số đồ dùng là
\choice
{\True $7$ cái}
{$6$ cái}
{$8$ cái}
{$9$ cái}
\loigiai{
$4 + 2 + 1 = 7$ đồ dùng.\\
\textbf{Đáp án của câu hỏi là A}
}
\end{ex}

\begin{ex} %[2X13H64]
Kiểm đếm số cây trong vườn: 15 cây cam, 12 cây xoài, 8 cây bưởi. Loại cây nào có số lượng nhiều nhất?
\choice
{Cây xoài}
{Cây bưởi}
{Các loại bằng nhau}
{\True Cây cam ($15$ cây)}
\loigiai{
Cây cam có 15 cây là nhiều nhất.\\
\textbf{Đáp án của câu hỏi là D}
}
\end{ex}

\begin{ex} %[2X13H64]
Kiểm đếm số cây trong vườn: 15 cây cam, 12 cây xoài, 8 cây bưởi. Loại cây nào có số lượng ít nhất?
\choice
{\True Cây bưởi ($8$ cây)}
{Cây cam}
{Cây xoài}
{Không có}
\loigiai{
Cây bưởi có 8 cây là ít nhất.\\
\textbf{Đáp án của câu hỏi là A}
}
\end{ex}

\begin{ex} %[2X13H64]
Số cây cam nhiều hơn số cây bưởi là
\choice
{$5$ cây}
{$6$ cây}
{$8$ cây}
{\True $7$ cây}
\loigiai{
$15 - 8 = 7$ cây.\\
\textbf{Đáp án của câu hỏi là D}
}
\end{ex}

\begin{ex} %[2X13V64]
Tổ 1 có 4 bạn thích bơi, 3 bạn thích đá bóng, 5 bạn thích cầu lông. Tổng số bạn trong tổ 1 là
\choice
{\True $12$ bạn}
{$11$ bạn}
{$13$ bạn}
{$14$ bạn}
\loigiai{
$4 + 3 + 5 = 12$ bạn.\\
\textbf{Đáp án của câu hỏi là A}
}
\end{ex}

\begin{ex} %[2X13V64]
Trong một đĩa trái cây có: 6 quả chuối, 4 quả cam, 5 quả táo. Số quả chuối nhiều hơn số quả cam là
\choice
{$1$ quả}
{$3$ quả}
{\True $2$ quả}
{$4$ quả}
\loigiai{
$6 - 4 = 2$ quả.\\
\textbf{Đáp án của câu hỏi là C}
}
\end{ex}

\begin{ex} %[2X13V64]
Thống kê điểm 10 của lớp 2A trong tuần: Thứ Hai: 8 điểm 10; Thứ Ba: 10 điểm 10; Thứ Tư: 7 điểm 10; Thứ Năm: 9 điểm 10; Thứ Sáu: 11 điểm 10. Ngày có nhiều điểm 10 nhất là
\choice
{Thứ Ba}
{\True Thứ Sáu}
{Thứ Hai}
{Thứ Năm}
\loigiai{
Thứ Sáu có 11 điểm 10 (nhiều nhất).\\
\textbf{Đáp án của câu hỏi là B}
}
\end{ex}

\begin{ex} %[2X13V64]
Tổng số điểm 10 của lớp 2A trong cả tuần đó là
\choice
{$43$ điểm}
{$47$ điểm}
{$50$ điểm}
{\True $45$ điểm}
\loigiai{
$8 + 10 + 7 + 9 + 11 = 45$ điểm 10.\\
\textbf{Đáp án của câu hỏi là D}
}
\end{ex}

\begin{ex} %[2X13V64]
Một cửa hàng bán xe đạp ghi chép số xe bán trong tuần: Xe màu đỏ: 8 chiếc, xe màu xanh: 6 chiếc, xe màu đen: 10 chiếc. Số xe màu đen nhiều hơn xe màu xanh là
\choice
{$3$ chiếc}
{$5$ chiếc}
{\True $4$ chiếc}
{$2$ chiếc}
\loigiai{
$10 - 6 = 4$ chiếc.\\
\textbf{Đáp án của câu hỏi là C}
}
\end{ex}

\begin{ex} %[2X13C64]
Trong hộp có các viên bi màu xanh, đỏ, vàng. Biết số bi đỏ là 8 viên, số bi xanh gấp đôi bi đỏ, số bi vàng ít hơn bi xanh 4 viên. Số bi vàng có trong hộp là
\choice
{$10$ viên}
{\True $12$ viên}
{$14$ viên}
{$16$ viên}
\loigiai{
Bi xanh: $8 \times 2 = 16$ viên.\\
Bi vàng: $16 - 4 = 12$ viên.\\
\textbf{Đáp án của câu hỏi là B}
}
\end{ex}

\textbf{Phần 2. Trắc nghiệm đúng sai}
\setcounterpageref{ex}{0}

\begin{ex} %[2X13H64]
Một giá sách có: 18 quyển sách Toán, 15 quyển sách Tiếng Việt, 12 quyển truyện tranh. Xét tính đúng sai:
\choiceTF
{\True Sách Toán có số lượng nhiều nhất (18 quyển).}
{\True Truyện tranh có số lượng ít nhất (12 quyển).}
{Giá sách có tất cả 40 quyển.}
{\True Giá sách có tất cả 45 quyển ($18 + 15 + 12 = 45$).}
\loigiai{
\begin{itemize}
    \item Toán nhiều nhất: Đúng.
    \item Truyện ít nhất: Đúng.
    \item Tổng: $18 + 15 + 12 = 45 \neq 40$. Ý 3 Sai, Ý 4 Đúng.
\end{itemize}
\begin{itemize}
    \item a) \textbf{Đúng}.
    \item b) \textbf{Đúng}.
    \item c) \textbf{Sai}.
    \item d) \textbf{Đúng}.
\end{itemize}
}
\end{ex}

\begin{ex} %[2X13V64]
Kiểm đếm các phương tiện giao thông đi qua ngã tư trong 5 phút: 25 xe máy, 10 ô tô, 8 xe đạp. Xét tính đúng sai:
\choiceTF
{\True Xe máy đi qua nhiều nhất.}
{\True Xe đạp đi qua ít nhất.}
{Số ô tô nhiều hơn số xe máy.}
{\True Số xe máy nhiều hơn số ô tô là 15 chiếc ($25 - 10 = 15$).}
\loigiai{
\begin{itemize}
    \item Xe máy nhiều nhất: Đúng.
    \item Xe đạp ít nhất: Đúng.
    \item Xe máy (25) nhiều hơn ô tô (10). Ý 3 Sai, Ý 4 Đúng.
\end{itemize}
\begin{itemize}
    \item a) \textbf{Đúng}.
    \item b) \textbf{Đúng}.
    \item c) \textbf{Sai}.
    \item d) \textbf{Đúng}.
\end{itemize}
}
\end{ex}

\textbf{Phần 3. Trả lời ngắn}
\setcounterpageref{ex}{0}

\begin{ex} %[2X13H64]
Kiểm đếm số vật nuôi của nhà An: 12 con gà, 8 con vịt, 4 con thỏ. Tính tổng số vật nuôi nhà An.
\shortans{24}
\loigiai{
\textbf{Đáp số:} $24$\\
$12 + 8 + 4 = 24$ con.
}
\end{ex}

\begin{ex} %[2X13H64]
Số gà nhiều hơn số thỏ bao nhiêu con?
\shortans{8}
\loigiai{
\textbf{Đáp số:} $8$\\
$12 - 4 = 8$ con.
}
\end{ex}

\begin{ex} %[2X13V64]
Kiểm đếm số học sinh tham gia câu lạc bộ: Cờ vua có 8 bạn, Võ thuật có 12 bạn, Mĩ thuật có 10 bạn. Tính tổng số bạn tham gia 3 câu lạc bộ.
\shortans{30}
\loigiai{
\textbf{Đáp số:} $30$\\
$8 + 12 + 10 = 30$ bạn.
}
\end{ex}

\begin{ex} %[2X13V64]
Số bạn tham gia câu lạc bộ Võ thuật nhiều hơn Cờ vua bao nhiêu bạn?
\shortans{4}
\loigiai{
\textbf{Đáp số:} $4$\\
$12 - 8 = 4$ bạn.
}
\end{ex}

\begin{ex} %[2X13V64]
Trong tủ có 15 áo sơ mi trắng, 10 áo thun xanh và 5 áo khoác. Hỏi có tất cả bao nhiêu cái áo?
\shortans{30}
\loigiai{
\textbf{Đáp số:} $30$\\
$15 + 10 + 5 = 30$ cái áo.
}
\end{ex}

\begin{ex} %[2X13C64]
Trong một hộp bút có 3 loại bút: bút mực, bút chì và bút dạ. Biết có 6 bút mực, số bút chì gấp đôi bút mực, số bút dạ bằng một nửa số bút mực. Tính tổng số bút trong hộp.
\shortans{21}
\loigiai{
\textbf{Đáp số:} $21$\\
Bút chì: $6 \times 2 = 12$. Bút dạ: $6 : 2 = 3$.\\
Tổng: $6 + 12 + 3 = 21$ cái bút.
}
\end{ex}

\newpage

\TaoPhieuKiemTra{Chủ đề 13 -- Làm quen với yếu tố thống kê, xác suất}{65}{Luyện tập chung phép tính 1000}{https://loptoanthaylinh.vn/thi-online/toan-2/bai-65}
\textbf{Phần 1. Trắc nghiệm lựa chọn}
\setcounterpageref{ex}{0}

\begin{ex} %[2X13N65]
Biểu đồ tranh dùng cái gì để biểu diễn số liệu?
\choice
{Các đoạn thẳng nối tiếp}
{Các hình khối chữ nhật}
{\True Các hình ảnh hoặc biểu tượng tranh vẽ}
{Các con số đơn thuần}
\loigiai{
Biểu đồ tranh dùng các hình ảnh hoặc biểu tượng.\\
\textbf{Đáp án của câu hỏi là C}
}
\end{ex}

\begin{ex} %[2X13N65]
Nếu mỗi biểu tượng ngôi sao $\bigstar$ biểu diễn cho 1 bạn học sinh, có 5 ngôi sao $\bigstar$ tương ứng với
\choice
{$1$ bạn}
{\True $5$ bạn học sinh}
{$10$ bạn}
{$15$ bạn}
\loigiai{
5 ngôi sao tương ứng với 5 bạn học sinh.\\
\textbf{Đáp án của câu hỏi là B}
}
\end{ex}

\begin{ex} %[2X13H65]
Nếu mỗi biểu tượng quả bóng biểu diễn cho 2 quả bóng thực tế, có 4 biểu tượng tương ứng với
\choice
{$6$ quả bóng}
{$10$ quả bóng}
{\True $8$ quả bóng}
{$4$ quả bóng}
\loigiai{
$2 \times 4 = 8$ quả bóng.\\
\textbf{Đáp án của câu hỏi là C}
}
\end{ex}

\begin{ex} %[2X13H65]
Nếu mỗi bông hoa biểu diễn cho 5 bông hoa thật, có 3 bông hoa biểu tượng tương ứng với
\choice
{\True $15$ bông hoa}
{$10$ bông}
{$20$ bông}
{$25$ bông}
\loigiai{
$5 \times 3 = 15$ bông hoa.\\
\textbf{Đáp án của câu hỏi là A}
}
\end{ex}

\begin{ex} %[2X13H65]
Quan sát biểu đồ tranh về số quả táo của 3 bạn: Mai có 4 quả táo, Lan có 6 quả táo, Hà có 3 quả táo (mỗi hình vẽ là 1 quả). Bạn có nhiều táo nhất là
\choice
{Mai}
{Hà}
{Các bạn bằng nhau}
{\True Lan ($6$ quả)}
\loigiai{
Lan có 6 quả táo (nhiều nhất).\\
\textbf{Đáp án của câu hỏi là D}
}
\end{ex}

\begin{ex} %[2X13H65]
Lan có nhiều hơn Hà số quả táo là
\choice
{\True $3$ quả}
{$2$ quả}
{$4$ quả}
{$5$ quả}
\loigiai{
$6 - 3 = 3$ quả.\\
\textbf{Đáp án của câu hỏi là A}
}
\end{ex}

\begin{ex} %[2X13V65]
Tổng số quả táo của cả ba bạn Mai, Lan, Hà là
\choice
{$12$ quả}
{\True $13$ quả}
{$14$ quả}
{$15$ quả}
\loigiai{
$4 + 6 + 3 = 13$ quả táo.\\
\textbf{Đáp án của câu hỏi là B}
}
\end{ex}

\begin{ex} %[2X13V65]
Biểu đồ tranh số cây trồng của các tổ (mỗi hình cây $= 2$ cây thật): Tổ 1 có 3 hình cây, Tổ 2 có 5 hình cây. Tổ 2 trồng được số cây thật là
\choice
{$8$ cây}
{$12$ cây}
{$15$ cây}
{\True $10$ cây}
\loigiai{
$2 \times 5 = 10$ cây.\\
\textbf{Đáp án của câu hỏi là D}
}
\end{ex}

\begin{ex} %[2X13V65]
Tổ 2 trồng nhiều hơn Tổ 1 số cây thật là
\choice
{\True $4$ cây}
{$2$ cây}
{$6$ cây}
{$8$ cây}
\loigiai{
Tổ 1: $2 \times 3 = 6$ cây. Tổ 2: 10 cây.\\
$10 - 6 = 4$ cây.\\
\textbf{Đáp án của câu hỏi là A}
}
\end{ex}

\begin{ex} %[2X13V65]
Biểu đồ tranh số thú nuôi trong trang trại: Bò có 4 con biểu tượng, Cừu có 6 con biểu tượng, Lợn có 5 con biểu tượng (mỗi biểu tượng $= 5$ con thật). Số cừu trong trang trại là
\choice
{$25$ con}
{$35$ con}
{$40$ con}
{\True $30$ con}
\loigiai{
$5 \times 6 = 30$ con cừu.\\
\textbf{Đáp án của câu hỏi là D}
}
\end{ex}

\begin{ex} %[2X13V65]
Tổng số con vật có trong trang trại đó là
\choice
{$70$ con}
{$80$ con}
{\True $75$ con}
{$85$ con}
\loigiai{
Tổng số biểu tượng: $4 + 6 + 5 = 15$.\\
Tổng số con: $5 \times 15 = 75$ con.\\
\textbf{Đáp án của câu hỏi là C}
}
\end{ex}

\begin{ex} %[2X13C65]
Trên biểu đồ tranh, bạn An có 4 hình trái tim, bạn Bình có 7 hình trái tim. Mỗi hình trái tim biểu diễn cho 2 cuốn sách. Hỏi Bình có nhiều hơn An bao nhiêu cuốn sách?
\choice
{$4$ cuốn}
{\True $6$ cuốn}
{$8$ cuốn}
{$10$ cuốn}
\loigiai{
Bình hơn An: $7 - 4 = 3$ hình trái tim.\\
Số sách nhiều hơn: $2 \times 3 = 6$ cuốn.\\
\textbf{Đáp án của câu hỏi là B}
}
\end{ex}

\textbf{Phần 2. Trắc nghiệm đúng sai}
\setcounterpageref{ex}{0}

\begin{ex} %[2X13H65]
Biểu đồ tranh ghi số ô tô bán được trong tuần (mỗi ô tô vẽ $= 2$ xe thật): Thứ Hai có 3 hình vẽ, Thứ Ba có 4 hình vẽ. Xét tính đúng sai:
\choiceTF
{\True Thứ Hai bán được 6 xe ô tô thật ($2 \times 3 = 6$).}
{\True Thứ Ba bán được 8 xe ô tô thật ($2 \times 4 = 8$).}
{Thứ Ba bán nhiều hơn Thứ Hai 1 xe thật.}
{\True Thứ Ba bán nhiều hơn Thứ Hai 2 xe thật ($8 - 6 = 2$).}
\loigiai{
\begin{itemize}
    \item Thứ Hai: 6 xe: Đúng.
    \item Thứ Ba: 8 xe: Đúng.
    \item Hơn 2 xe (không phải 1 xe). Ý 3 Sai, Ý 4 Đúng.
\end{itemize}
\begin{itemize}
    \item a) \textbf{Đúng}.
    \item b) \textbf{Đúng}.
    \item c) \textbf{Sai}.
    \item d) \textbf{Đúng}.
\end{itemize}
}
\end{ex}

\begin{ex} %[2X13V65]
Biểu đồ tranh về số quả bóng (mỗi biểu tượng $= 5$ quả): Lớp 2A có 4 biểu tượng, Lớp 2B có 3 biểu tượng. Xét tính đúng sai:
\choiceTF
{\True Lớp 2A có 20 quả bóng ($5 \times 4 = 20$).}
{\True Lớp 2B có 15 quả bóng ($5 \times 3 = 15$).}
{Cả hai lớp có tất cả 40 quả bóng.}
{\True Cả hai lớp có tất cả 35 quả bóng ($20 + 15 = 35$).}
\loigiai{
\begin{itemize}
    \item Lớp 2A có 20 quả: Đúng.
    \item Lớp 2B có 15 quả: Đúng.
    \item Tổng: $20 + 15 = 35 \neq 40$. Ý 3 Sai, Ý 4 Đúng.
\end{itemize}
\begin{itemize}
    \item a) \textbf{Đúng}.
    \item b) \textbf{Đúng}.
    \item c) \textbf{Sai}.
    \item d) \textbf{Đúng}.
\end{itemize}
}
\end{ex}

\textbf{Phần 3. Trả lời ngắn}
\setcounterpageref{ex}{0}

\begin{ex} %[2X13H65]
Trên biểu đồ tranh, mỗi ngôi sao biểu diễn cho 2 điểm mười. Bạn Nam có 5 ngôi sao. Hỏi Nam có bao nhiêu điểm mười?
\shortans{10}
\loigiai{
\textbf{Đáp số:} $10$\\
$2 \times 5 = 10$ điểm mười.
}
\end{ex}

\begin{ex} %[2X13H65]
Bạn Mai có 7 ngôi sao (mỗi ngôi sao $= 2$ điểm mười). Hỏi Mai có bao nhiêu điểm mười?
\shortans{14}
\loigiai{
\textbf{Đáp số:} $14$\\
$2 \times 7 = 14$ điểm mười.
}
\end{ex}

\begin{ex} %[2X13V65]
Trên biểu đồ, mỗi hình cái cây biểu diễn cho 5 cây thật. Lớp 2A có 6 hình cái cây. Hỏi lớp 2A trồng được bao nhiêu cây?
\shortans{30}
\loigiai{
\textbf{Đáp số:} $30$\\
$5 \times 6 = 30$ cây.
}
\end{ex}

\begin{ex} %[2X13V65]
Lớp 2B có 8 hình cái cây (mỗi hình $= 5$ cây). Hỏi lớp 2B trồng được bao nhiêu cây?
\shortans{40}
\loigiai{
\textbf{Đáp số:} $40$\\
$5 \times 8 = 40$ cây.
}
\end{ex}

\begin{ex} %[2X13V65]
Lớp 2B trồng nhiều hơn lớp 2A bao nhiêu cây?
\shortans{10}
\loigiai{
\textbf{Đáp số:} $10$\\
$40 - 30 = 10$ cây.
}
\end{ex}

\begin{ex} %[2X13C65]
Có 45 quyển sách cần biểu diễn lên biểu đồ tranh. Biết mỗi biểu tượng quyển sách biểu diễn cho 5 cuốn sách thật. Hỏi cần vẽ bao nhiêu biểu tượng?
\shortans{9}
\loigiai{
\textbf{Đáp số:} $9$\\
$45 : 5 = 9$ biểu tượng.
}
\end{ex}

\newpage

\TaoPhieuKiemTra{Chủ đề 13 -- Làm quen với yếu tố thống kê, xác suất}{66}{Thu thập, phân loại, kiểm đếm số liệu}{https://loptoanthaylinh.vn/thi-online/toan-2/bai-66}
\textbf{Phần 1. Trắc nghiệm lựa chọn}
\setcounterpageref{ex}{0}

\begin{ex} %[2X13N66]
Trong một hộp chỉ có các viên bi màu đỏ. Nhắm mắt lấy ra 1 viên bi, sự kiện "Lấy được viên bi màu đỏ" là
\choice
{Có thể}
{\True Chắc chắn}
{Không thể}
{Chắc chắn không xảy ra}
\loigiai{
Trong hộp chỉ có bi đỏ nên chắc chắn lấy được bi đỏ.\\
\textbf{Đáp án của câu hỏi là B}
}
\end{ex}

\begin{ex} %[2X13N66]
Trong một hộp chỉ có các viên bi màu đỏ. Nhắm mắt lấy ra 1 viên bi, sự kiện "Lấy được viên bi màu xanh" là
\choice
{\True Không thể}
{Chắc chắn}
{Có thể}
{Có thể hoặc không}
\loigiai{
Không có bi xanh nên không thể lấy được bi xanh.\\
\textbf{Đáp án của câu hỏi là A}
}
\end{ex}

\begin{ex} %[2X13H66]
Trong một hộp có cả bi đỏ và bi xanh. Nhắm mắt lấy ra 1 viên bi, sự kiện "Lấy được viên bi màu đỏ" là
\choice
{Chắc chắn}
{Không thể}
{Không bao giờ}
{\True Có thể}
\loigiai{
Có cả bi đỏ và bi xanh nên có thể lấy được bi đỏ.\\
\textbf{Đáp án của câu hỏi là D}
}
\end{ex}

\begin{ex} %[2X13H66]
Gieo một con xúc xắc 6 mặt (có các số chấm từ 1 đến 6). Sự kiện "Xuất hiện mặt có 7 chấm" là
\choice
{Chắc chắn}
{\True Không thể}
{Có thể}
{Có khả năng cao}
\loigiai{
Xúc xắc chỉ có từ 1 đến 6 chấm nên không thể có 7 chấm.\\
\textbf{Đáp án của câu hỏi là B}
}
\end{ex}

\begin{ex} %[2X13H66]
Gieo một con xúc xắc 6 mặt. Sự kiện "Xuất hiện mặt có số chấm bé hơn 7" là
\choice
{\True Chắc chắn}
{Không thể}
{Có thể}
{Chưa biết}
\loigiai{
Mọi mặt xúc xắc đều có số chấm từ 1 đến 6 (bé hơn 7) nên là chắc chắn.\\
\textbf{Đáp án của câu hỏi là A}
}
\end{ex}

\begin{ex} %[2X13H66]
Gieo một con xúc xắc 6 mặt. Sự kiện "Xuất hiện mặt có 3 chấm" là
\choice
{Chắc chắn}
{Không thể}
{\True Có thể}
{Hoàn toàn không}
\loigiai{
Có thể xuất hiện mặt 3 chấm.\\
\textbf{Đáp án của câu hỏi là C}
}
\end{ex}

\begin{ex} %[2X13V66]
Một túi đựng 3 quả bóng xanh và 2 quả bóng đỏ. Nam nhắm mắt lấy ra cùng lúc 3 quả bóng. Khẳng định nào sau đây là ĐÚNG?
\choice
{\True Có thể lấy được 3 quả bóng xanh}
{Chắc chắn lấy được 3 quả bóng đỏ}
{Không thể lấy được quả bóng đỏ nào}
{Chắc chắn lấy được 3 quả bóng xanh}
\loigiai{
Vì có 3 bóng xanh nên có thể lấy được 3 bóng xanh.\\
\textbf{Đáp án của câu hỏi là A}
}
\end{ex}

\begin{ex} %[2X13V66]
Một túi đựng 5 cái kẹo dâu. An lấy ngẫu nhiên 2 cái kẹo. Sự kiện "An lấy được 2 cái kẹo dâu" là
\choice
{Có thể}
{Không thể}
{\True Chắc chắn}
{Chắc chắn không xảy ra}
\loigiai{
Tất cả kẹo đều vị dâu nên chắc chắn lấy được 2 kẹo dâu.\\
\textbf{Đáp án của câu hỏi là C}
}
\end{ex}

\begin{ex} %[2X13V66]
Hôm nay là thứ Năm. Ngày mai là thứ Sáu là sự kiện
\choice
{Có thể}
{Không thể}
{Tùy từng tuần}
{\True Chắc chắn}
\loigiai{
Thứ Năm luôn đến thứ Sáu nên là chắc chắn.\\
\textbf{Đáp án của câu hỏi là D}
}
\end{ex}

\begin{ex} %[2X13V66]
Ngày mai trời sẽ có mưa là sự kiện
\choice
{Chắc chắn}
{Không thể}
{\True Có thể}
{Chắc chắn không}
\loigiai{
Thời tiết ngày mai là sự kiện có thể xảy ra.\\
\textbf{Đáp án của câu hỏi là C}
}
\end{ex}

\begin{ex} %[2X13V66]
Trong một hộp có 4 chiếc bút chì đen. Hà lấy ra 1 chiếc bút. Khẳng định nào sau đây là KHÔNG ĐÚNG?
\choice
{Chắc chắn lấy được bút chì đen}
{\True Có thể lấy được bút chì màu đỏ}
{Không thể lấy được bút chì màu đỏ}
{Chắc chắn chiếc bút lấy ra là màu đen}
\loigiai{
Không có bút đỏ nên không thể nói "có thể lấy được bút đỏ".\\
\textbf{Đáp án của câu hỏi là B}
}
\end{ex}

\begin{ex} %[2X13C66]
Trong hộp có 3 viên bi đỏ và 3 viên bi xanh. Lấy ra ngẫu nhiên 4 viên bi. Khẳng định nào sau đây là CHẮC CHẮN?
\choice
{Chắc chắn lấy được 4 viên bi đỏ}
{Chắc chắn lấy được 4 viên bi xanh}
{Chắc chắn lấy được 3 viên bi vàng}
{\True Chắc chắn trong 4 viên lấy ra có ít nhất 1 viên bi đỏ và 1 viên bi xanh}
\loigiai{
Vì chỉ có 3 bi mỗi màu, lấy 4 viên thì chắc chắn phải có cả 2 màu đỏ và xanh.\\
\textbf{Đáp án của câu hỏi là D}
}
\end{ex}

\textbf{Phần 2. Trắc nghiệm đúng sai}
\setcounterpageref{ex}{0}

\begin{ex} %[2X13H66]
Trong hộp có 5 viên bi vàng và 3 viên bi xanh. Xét tính đúng sai của các sự kiện khi lấy ra 1 viên bi:
\choiceTF
{\True Có thể lấy được viên bi màu vàng.}
{\True Có thể lấy được viên bi màu xanh.}
{Chắc chắn lấy được viên bi màu vàng.}
{\True Không thể lấy được viên bi màu đỏ.}
\loigiai{
\begin{itemize}
    \item Có thể lấy bi vàng: Đúng.
    \item Có thể lấy bi xanh: Đúng.
    \item Không chắc chắn vì có thể ra bi xanh. Ý 3 Sai, Ý 4 Đúng.
\end{itemize}
\begin{itemize}
    \item a) \textbf{Đúng}.
    \item b) \textbf{Đúng}.
    \item c) \textbf{Sai}.
    \item d) \textbf{Đúng}.
\end{itemize}
}
\end{ex}

\begin{ex} %[2X13V66]
Quay một vòng quay có 2 màu: Đỏ và Xanh. Xét tính đúng sai:
\choiceTF
{\True Kim quay có thể dừng ở phần màu đỏ.}
{\True Kim quay có thể dừng ở phần màu xanh.}
{Kim quay chắc chắn dừng ở phần màu đỏ.}
{\True Kim quay không thể dừng ở phần màu vàng (vì không có màu vàng).}
\loigiai{
\begin{itemize}
    \item Có thể dừng ở đỏ: Đúng.
    \item Có thể dừng ở xanh: Đúng.
    \item Không chắc chắn đỏ (vì có thể vào xanh). Ý 3 Sai, Ý 4 Đúng.
\end{itemize}
\begin{itemize}
    \item a) \textbf{Đúng}.
    \item b) \textbf{Đúng}.
    \item c) \textbf{Sai}.
    \item d) \textbf{Đúng}.
\end{itemize}
}
\end{ex}

\textbf{Phần 3. Trả lời ngắn}
\setcounterpageref{ex}{0}

\begin{ex} %[2X13H66]
Trong túi có 5 quả bóng đỏ. Lấy ra 1 quả bóng, sự kiện "Lấy được bóng đỏ" là gì? (Viết: chắc chắn, có thể, hoặc không thể)
\shortans{chắc chắn}
\loigiai{
\textbf{Đáp số:} chắc chắn\\
Chắc chắn.
}
\end{ex}

\begin{ex} %[2X13H66]
Trong túi có 5 quả bóng đỏ. Lấy ra 1 quả bóng, sự kiện "Lấy được bóng vàng" là gì? (Viết: chắc chắn, có thể, hoặc không thể)
\shortans{không thể}
\loigiai{
\textbf{Đáp số:} không thể\\
Không thể.
}
\end{ex}

\begin{ex} %[2X13V66]
Trong túi có 3 bóng đỏ và 3 bóng xanh. Lấy ra 1 quả bóng, sự kiện "Lấy được bóng đỏ" là gì? (Viết: chắc chắn, có thể, hoặc không thể)
\shortans{có thể}
\loigiai{
\textbf{Đáp số:} có thể\\
Có thể.
}
\end{ex}

\begin{ex} %[2X13V66]
Một con cá biết bay trên trời là sự kiện gì? (Viết: chắc chắn, có thể, hoặc không thể)
\shortans{không thể}
\loigiai{
\textbf{Đáp số:} không thể\\
Không thể.
}
\end{ex}

\begin{ex} %[2X13V66]
Mặt trời mọc ở hướng Đông là sự kiện gì? (Viết: chắc chắn, có thể, hoặc không thể)
\shortans{chắc chắn}
\loigiai{
\textbf{Đáp số:} chắc chắn\\
Chắc chắn.
}
\end{ex}

\begin{ex} %[2X13C66]
Tung một đồng xu có 2 mặt Sấp và Ngửa. Sự kiện "Đồng xu rơi xuống mặt Sấp" là gì? (Viết: chắc chắn, có thể, hoặc không thể)
\shortans{có thể}
\loigiai{
\textbf{Đáp số:} có thể\\
Có thể.
}
\end{ex}

\newpage

\TaoPhieuKiemTra{Chủ đề 13 -- Làm quen với yếu tố thống kê, xác suất}{67}{Biểu đồ tranh}{https://loptoanthaylinh.vn/thi-online/toan-2/bai-67}
\textbf{Phần 1. Trắc nghiệm lựa chọn}
\setcounterpageref{ex}{0}

\begin{ex} %[2X13N67]
Để kiểm đếm số lượng các bạn trong lớp sinh vào các tháng, ta tiến hành
\choice
{Đo chiều cao}
{\True Thu thập ngày tháng sinh, phân loại theo tháng và đếm}
{Đo cân nặng}
{Làm bài kiểm tra}
\loigiai{
Thu thập ngày sinh, phân loại theo tháng và đếm số lượng.\\
\textbf{Đáp án của câu hỏi là B}
}
\end{ex}

\begin{ex} %[2X13N67]
Trong một giờ học thực hành, nhóm của Lan tung đồng xu 10 lần và ghi lại số lần xuất hiện mặt Sấp và mặt Ngửa. Hoạt động này là
\choice
{Đo độ dài}
{Xếp hình khối}
{\True Thực hành kiểm đếm và ghi chép số liệu}
{Vẽ đoạn thẳng}
\loigiai{
Đó là thực hành kiểm đếm và ghi chép số liệu.\\
\textbf{Đáp án của câu hỏi là C}
}
\end{ex}

\begin{ex} %[2X13H67]
Tung đồng xu 10 lần thu được 6 lần mặt Sấp. Số lần xuất hiện mặt Ngửa là
\choice
{\True $4$ lần}
{$3$ lần}
{$5$ lần}
{$6$ lần}
\loigiai{
$10 - 6 = 4$ lần.\\
\textbf{Đáp án của câu hỏi là A}
}
\end{ex}

\begin{ex} %[2X13H67]
Gieo xúc xắc 20 lần ghi nhận được: Mặt 1 chấm xuất hiện 3 lần, mặt 2 chấm: 4 lần, mặt 3 chấm: 5 lần, mặt 4 chấm: 2 lần, mặt 5 chấm: 3 lần. Số lần xuất hiện mặt 6 chấm là
\choice
{$2$ lần}
{$4$ lần}
{\True $3$ lần}
{$5$ lần}
\loigiai{
Tổng số lần mặt 1 đến 5: $3 + 4 + 5 + 2 + 3 = 17$ lần.\\
Mặt 6 chấm: $20 - 17 = 3$ lần.\\
\textbf{Đáp án của câu hỏi là C}
}
\end{ex}

\begin{ex} %[2X13H67]
Trong bảng số liệu trên, mặt nào xuất hiện nhiều lần nhất?
\choice
{\True Mặt 3 chấm ($5$ lần)}
{Mặt 2 chấm}
{Mặt 1 chấm}
{Mặt 5 chấm}
\loigiai{
Mặt 3 chấm xuất hiện 5 lần (nhiều nhất).\\
\textbf{Đáp án của câu hỏi là A}
}
\end{ex}

\begin{ex} %[2X13H67]
Mặt nào xuất hiện ít lần nhất?
\choice
{Mặt 1 chấm}
{Mặt 6 chấm}
{Mặt 5 chấm}
{\True Mặt 4 chấm ($2$ lần)}
\loigiai{
Mặt 4 chấm xuất hiện 2 lần (ít nhất).\\
\textbf{Đáp án của câu hỏi là D}
}
\end{ex}

\begin{ex} %[2X13V67]
Khảo sát món ăn yêu thích của 30 bạn học sinh: 12 bạn thích gà rán, 10 bạn thích xúc xích, còn lại thích pizza. Số bạn thích pizza là
\choice
{\True $8$ bạn}
{$7$ bạn}
{$9$ bạn}
{$10$ bạn}
\loigiai{
$30 - (12 + 10) = 30 - 22 = 8$ bạn.\\
\textbf{Đáp án của câu hỏi là A}
}
\end{ex}

\begin{ex} %[2X13V67]
Số bạn thích gà rán nhiều hơn số bạn thích pizza là
\choice
{$3$ bạn}
{$5$ bạn}
{$2$ bạn}
{\True $4$ bạn}
\loigiai{
$12 - 8 = 4$ bạn.\\
\textbf{Đáp án của câu hỏi là D}
}
\end{ex}

\begin{ex} %[2X13V67]
Khảo sát phương tiện đi học của lớp 2B (có 32 bạn): 18 bạn đi xe máy (bố mẹ chở), 8 bạn đi bộ, còn lại đi xe đạp. Số bạn đi xe đạp là
\choice
{$5$ bạn}
{\True $6$ bạn}
{$7$ bạn}
{$8$ bạn}
\loigiai{
$32 - (18 + 8) = 32 - 26 = 6$ bạn.\\
\textbf{Đáp án của câu hỏi là B}
}
\end{ex}

\begin{ex} %[2X13V67]
Số bạn đi xe máy nhiều hơn số bạn đi bộ là
\choice
{$8$ bạn}
{$12$ bạn}
{\True $10$ bạn}
{$14$ bạn}
\loigiai{
$18 - 8 = 10$ bạn.\\
\textbf{Đáp án của câu hỏi là C}
}
\end{ex}

\begin{ex} %[2X13V67]
Thực hành quay chong chóng 15 lần: màu đỏ quay được 8 lần, màu xanh quay được 7 lần. Sự kiện xuất hiện nhiều hơn là
\choice
{Màu xanh}
{Hai màu bằng nhau}
{Không xác định được}
{\True Màu đỏ}
\loigiai{
Màu đỏ xuất hiện 8 lần (nhiều hơn).\\
\textbf{Đáp án của câu hỏi là D}
}
\end{ex}

\begin{ex} %[2X13C67]
Trong một trò chơi rút thăm có thưởng, có 20 lá thăm gồm: 5 thăm trúng gấu bông, 7 thăm trúng hộp bút, còn lại là thăm chúc may mắn lần sau. Số thăm chúc may mắn lần sau là
\choice
{$7$ lá}
{\True $8$ lá}
{$9$ lá}
{$10$ lá}
\loigiai{
$20 - (5 + 7) = 20 - 12 = 8$ lá thăm.\\
\textbf{Đáp án của câu hỏi là B}
}
\end{ex}

\textbf{Phần 2. Trắc nghiệm đúng sai}
\setcounterpageref{ex}{0}

\begin{ex} %[2X13H67]
Khảo sát loài hoa yêu thích của 25 bạn: Hoa hồng: 10 bạn, Hoa hướng dương: 8 bạn, Hoa sen: 7 bạn. Xét tính đúng sai:
\choiceTF
{\True Hoa hồng được nhiều bạn yêu thích nhất (10 bạn).}
{\True Hoa sen có ít bạn yêu thích nhất (7 bạn).}
{Hoa hồng nhiều hơn hoa hướng dương 3 bạn.}
{\True Hoa hồng nhiều hơn hoa hướng dương 2 bạn ($10 - 8 = 2$).}
\loigiai{
\begin{itemize}
    \item Hồng nhiều nhất: Đúng.
    \item Sen ít nhất: Đúng.
    \item Hơn 2 bạn (không phải 3). Ý 3 Sai, Ý 4 Đúng.
\end{itemize}
\begin{itemize}
    \item a) \textbf{Đúng}.
    \item b) \textbf{Đúng}.
    \item c) \textbf{Sai}.
    \item d) \textbf{Đúng}.
\end{itemize}
}
\end{ex}

\begin{ex} %[2X13V67]
Gieo đồng xu 20 lần thu được 12 lần mặt Sấp và 8 lần mặt Ngửa. Xét tính đúng sai:
\choiceTF
{\True Mặt Sấp xuất hiện nhiều hơn mặt Ngửa.}
{\True Mặt Sấp nhiều hơn mặt Ngửa 4 lần ($12 - 8 = 4$).}
{Tổng số lần gieo là 24 lần.}
{\True Tổng số lần gieo là $12 + 8 = 20$ lần.}
\loigiai{
\begin{itemize}
    \item Sấp nhiều hơn: Đúng.
    \item Sấp hơn Ngửa 4 lần: Đúng.
    \item Tổng là 20 lần (không phải 24). Ý 3 Sai, Ý 4 Đúng.
\end{itemize}
\begin{itemize}
    \item a) \textbf{Đúng}.
    \item b) \textbf{Đúng}.
    \item c) \textbf{Sai}.
    \item d) \textbf{Đúng}.
\end{itemize}
}
\end{ex}

\textbf{Phần 3. Trả lời ngắn}
\setcounterpageref{ex}{0}

\begin{ex} %[2X13H67]
Khảo sát số con vật yêu thích của 20 bạn: Chó có 9 bạn, Mèo có 7 bạn, Thỏ có 4 bạn. Có bao nhiêu bạn thích Chó?
\shortans{9}
\loigiai{
\textbf{Đáp số:} $9$\\
9 bạn.
}
\end{ex}

\begin{ex} %[2X13H67]
Số bạn thích Mèo nhiều hơn số bạn thích Thỏ bao nhiêu bạn?
\shortans{3}
\loigiai{
\textbf{Đáp số:} $3$\\
$7 - 4 = 3$ bạn.
}
\end{ex}

\begin{ex} %[2X13V67]
Lớp 2A có 35 bạn, trong đó có 18 bạn nam. Hỏi lớp 2A có bao nhiêu bạn nữ?
\shortans{17}
\loigiai{
\textbf{Đáp số:} $17$\\
$35 - 18 = 17$ bạn.
}
\end{ex}

\begin{ex} %[2X13V67]
Tung xúc xắc 15 lần, mặt chẵn xuất hiện 8 lần. Hỏi mặt lẻ xuất hiện bao nhiêu lần?
\shortans{7}
\loigiai{
\textbf{Đáp số:} $7$\\
$15 - 8 = 7$ lần.
}
\end{ex}

\begin{ex} %[2X13V67]
Trong giỏ có 12 quả táo, 8 quả lê và 5 quả xoài. Tính tổng số quả trong giỏ.
\shortans{25}
\loigiai{
\textbf{Đáp số:} $25$\\
$12 + 8 + 5 = 25$ quả.
}
\end{ex}

\begin{ex} %[2X13C67]
Khảo sát 40 học sinh về môn thể thao yêu thích: Bóng đá có 18 bạn, Bơi lội có 12 bạn, Cầu lông có 6 bạn, còn lại thích Bóng rổ. Hỏi có bao nhiêu bạn thích Bóng rổ?
\shortans{4}
\loigiai{
\textbf{Đáp số:} $4$\\
$40 - (18 + 12 + 6) = 40 - 36 = 4$ bạn.
}
\end{ex}

\newpage

% =========================================================================
% CHỦ ĐỀ 14 - ÔN TẬP CUỐI NĂM
% =========================================================================
\section{Ôn tập cuối năm}

\TaoPhieuKiemTra{Chủ đề 14 -- Ôn tập cuối năm}{68}{Chắc chắn, có thể, không thể}{https://loptoanthaylinh.vn/thi-online/toan-2/bai-68}
\textbf{Phần 1. Trắc nghiệm lựa chọn}
\setcounterpageref{ex}{0}

\begin{ex} %[2D14N68]
Số gồm $5$ trăm, $8$ chục và $3$ đơn vị viết là
\choice
{\True $583$}
{$538$}
{$853$}
{$385$}
\loigiai{
Số đó là $583$.\\
\textbf{Đáp án của câu hỏi là A}
}
\end{ex}

\begin{ex} %[2D14N68]
Số $905$ đọc là
\choice
{Chín trăm năm mươi}
{\True Chín trăm linh năm}
{Chín trăm năm}
{Chín không năm}
\loigiai{
Đọc là Chín trăm linh năm.\\
\textbf{Đáp án của câu hỏi là B}
}
\end{ex}

\begin{ex} %[2D14H68]
Số liền trước của số $800$ là
\choice
{$801$}
{$790$}
{\True $799$}
{$700$}
\loigiai{
$800 - 1 = 799$.\\
\textbf{Đáp án của câu hỏi là C}
}
\end{ex}

\begin{ex} %[2D14H68]
Số liền sau của số $999$ là
\choice
{$998$}
{\True $1\,000$}
{$1001$}
{$990$}
\loigiai{
$999 + 1 = 1000$.\\
\textbf{Đáp án của câu hỏi là B}
}
\end{ex}

\begin{ex} %[2D14H68]
Số $672$ được viết thành tổng các trăm, chục, đơn vị là
\choice
{$60 + 70 + 2$}
{$600 + 7 + 2$}
{\True $600 + 70 + 2$}
{$600 + 72$}
\loigiai{
$672 = 600 + 70 + 2$.\\
\textbf{Đáp án của câu hỏi là C}
}
\end{ex}

\begin{ex} %[2D14H68]
Điền dấu thích hợp: $856 \dots 865$.
\choice
{\True $<$}
{$>$}
{$=$}
{Không so sánh được}
\loigiai{
$856 < 865$.\\
\textbf{Đáp án của câu hỏi là A}
}
\end{ex}

\begin{ex} %[2D14V68]
Trong các số $645, 564, 456, 654$, số lớn nhất là
\choice
{$645$}
{\True $654$}
{$564$}
{$456$}
\loigiai{
Số lớn nhất là 654.\\
\textbf{Đáp án của câu hỏi là B}
}
\end{ex}

\begin{ex} %[2D14V68]
Trong các số $789, 879, 798, 897$, số bé nhất là
\choice
{$798$}
{$879$}
{$897$}
{\True $789$}
\loigiai{
Số bé nhất là 789.\\
\textbf{Đáp án của câu hỏi là D}
}
\end{ex}

\begin{ex} %[2D14V68]
Sắp xếp các số $312, 123, 231, 132$ theo thứ tự từ bé đến lớn:
\choice
{\True $123, 132, 231, 312$}
{$312, 231, 132, 123$}
{$123, 231, 132, 312$}
{$132, 123, 231, 312$}
\loigiai{
$123 < 132 < 231 < 312$.\\
\textbf{Đáp án của câu hỏi là A}
}
\end{ex}

\begin{ex} %[2D14V68]
Số chẵn lớn nhất có 3 chữ số khác nhau là
\choice
{$998$}
{$988$}
{$978$}
{\True $986$}
\loigiai{
Số đó là 986.\\
\textbf{Đáp án của câu hỏi là D}
}
\end{ex}

\begin{ex} %[2D14V68]
Số tròn chục bé nhất có ba chữ số khác nhau là
\choice
{$100$}
{$102$}
{\True $120$}
{$110$}
\loigiai{
Số đó là 120.\\
\textbf{Đáp án của câu hỏi là C}
}
\end{ex}

\begin{ex} %[2D14C68]
Tìm số có 3 chữ số biết chữ số hàng trăm là số lớn nhất có một chữ số, chữ số hàng chục bằng một phần ba chữ số hàng trăm và chữ số hàng đơn vị bằng 0.
\choice
{$960$}
{$910$}
{$900$}
{\True $930$}
\loigiai{
Hàng trăm: 9. Hàng chục: $9 : 3 = 3$. Hàng đơn vị: 0. Số đó là $930$.\\
\textbf{Đáp án của câu hỏi là D}
}
\end{ex}

\textbf{Phần 2. Trắc nghiệm đúng sai}
\setcounterpageref{ex}{0}

\begin{ex} %[2D14H68]
Xét tính đúng sai về số trong phạm vi 1 000:
\choiceTF
{\True $1\,000$ là số bé nhất có bốn chữ số.}
{\True Số liền trước của số 600 là 599.}
{Số 408 viết thành tổng là $400 + 80$.}
{\True Số 408 viết thành tổng là $400 + 8$.}
\loigiai{
\begin{itemize}
    \item 1000 bé nhất có 4 chữ số: Đúng.
    \item Liền trước 600 là 599: Đúng.
    \item $408 = 400 + 8 \neq 400 + 80$. Ý 3 Sai, Ý 4 Đúng.
\end{itemize}
\begin{itemize}
    \item a) \textbf{Đúng}.
    \item b) \textbf{Đúng}.
    \item c) \textbf{Sai}.
    \item d) \textbf{Đúng}.
\end{itemize}
}
\end{ex}

\begin{ex} %[2D14V68]
Cho các số: $450, 540, 405, 504$. Xét tính đúng sai:
\choiceTF
{\True Số lớn nhất là 540.}
{\True Số bé nhất là 405.}
{Sắp xếp theo thứ tự giảm dần là $540 > 450 > 504 > 405$.}
{\True Sắp xếp giảm dần đúng là $540 > 504 > 450 > 405$.}
\loigiai{
\begin{itemize}
    \item Lớn nhất 540: Đúng.
    \item Bé nhất 405: Đúng.
    \item $504 > 450$. Ý 3 Sai, Ý 4 Đúng.
\end{itemize}
\begin{itemize}
    \item a) \textbf{Đúng}.
    \item b) \textbf{Đúng}.
    \item c) \textbf{Sai}.
    \item d) \textbf{Đúng}.
\end{itemize}
}
\end{ex}

\textbf{Phần 3. Trả lời ngắn}
\setcounterpageref{ex}{0}

\begin{ex} %[2D14H68]
Viết số gồm 7 trăm, 0 chục và 9 đơn vị.
\shortans{709}
\loigiai{
\textbf{Đáp số:} $709$\\
709.
}
\end{ex}

\begin{ex} %[2D14H68]
Tìm số liền trước của số $900$.
\shortans{899}
\loigiai{
\textbf{Đáp số:} $899$\\
$900 - 1 = 899$.
}
\end{ex}

\begin{ex} %[2D14V68]
Tìm số liền sau của số $889$.
\shortans{890}
\loigiai{
\textbf{Đáp số:} $890$\\
$889 + 1 = 890$.
}
\end{ex}

\begin{ex} %[2D14V68]
Viết kết quả của tổng: $800 + 60 + 4$.
\shortans{864}
\loigiai{
\textbf{Đáp số:} $864$\\
864.
}
\end{ex}

\begin{ex} %[2D14V68]
Tìm số lớn nhất có 3 chữ số khác nhau.
\shortans{987}
\loigiai{
\textbf{Đáp số:} $987$\\
987.
}
\end{ex}

\begin{ex} %[2D14C68]
Từ 3 chữ số $0, 4, 7$, viết số bé nhất có 3 chữ số khác nhau.
\shortans{407}
\loigiai{
\textbf{Đáp số:} $407$\\
407.
}
\end{ex}

\newpage

\TaoPhieuKiemTra{Chủ đề 14 -- Ôn tập cuối năm}{69}{Luyện tập chung thống kê và xác suất}{https://loptoanthaylinh.vn/thi-online/toan-2/bai-69}
\textbf{Phần 1. Trắc nghiệm lựa chọn}
\setcounterpageref{ex}{0}

\begin{ex} %[2D14N69]
Tính: $36 + 28 = \dots$
\choice
{\True $64$}
{$54$}
{$74$}
{$65$}
\loigiai{
$36 + 28 = 64$.\\
\textbf{Đáp án của câu hỏi là A}
}
\end{ex}

\begin{ex} %[2D14N69]
Tính: $72 - 35 = \dots$
\choice
{$36$}
{$38$}
{\True $37$}
{$47$}
\loigiai{
$72 - 35 = 37$.\\
\textbf{Đáp án của câu hỏi là C}
}
\end{ex}

\begin{ex} %[2D14H69]
Kết quả của phép tính $54 + 46$ là
\choice
{\True $100$}
{$90$}
{$99$}
{$101$}
\loigiai{
$54 + 46 = 100$.\\
\textbf{Đáp án của câu hỏi là A}
}
\end{ex}

\begin{ex} %[2D14H69]
Kết quả của phép tính $100 - 43$ là
\choice
{$56$}
{\True $57$}
{$58$}
{$67$}
\loigiai{
$100 - 43 = 57$.\\
\textbf{Đáp án của câu hỏi là B}
}
\end{ex}

\begin{ex} %[2D14H69]
Tìm $x$ biết: $x - 27 = 48$.
\choice
{$65$}
{$85$}
{$76$}
{\True $75$}
\loigiai{
$x = 48 + 27 = 75$.\\
\textbf{Đáp án của câu hỏi là D}
}
\end{ex}

\begin{ex} %[2D14H69]
Tìm $y$ biết: $85 - y = 39$.
\choice
{\True $46$}
{$45$}
{$47$}
{$56$}
\loigiai{
$y = 85 - 39 = 46$.\\
\textbf{Đáp án của câu hỏi là A}
}
\end{ex}

\begin{ex} %[2D14V69]
Lớp 2A có 38 học sinh, lớp 2B có 35 học sinh. Cả hai lớp có tất cả số học sinh là
\choice
{$72$ bạn}
{$74$ bạn}
{$75$ bạn}
{\True $73$ bạn}
\loigiai{
$38 + 35 = 73$ học sinh.\\
\textbf{Đáp án của câu hỏi là D}
}
\end{ex}

\begin{ex} %[2D14V69]
Một bến xe có 62 ô tô, lúc sau có 28 ô tô rời bến. Bến xe còn lại số ô tô là
\choice
{$33$ chiếc}
{\True $34$ chiếc}
{$35$ chiếc}
{$44$ chiếc}
\loigiai{
$62 - 28 = 34$ chiếc.\\
\textbf{Đáp án của câu hỏi là B}
}
\end{ex}

\begin{ex} %[2D14V69]
Tính giá trị biểu thức: $45 + 38 - 29$.
\choice
{$53$}
{$55$}
{\True $54$}
{$56$}
\loigiai{
$83 - 29 = 54$.\\
\textbf{Đáp án của câu hỏi là C}
}
\end{ex}

\begin{ex} %[2D14V69]
Điền dấu thích hợp: $56 + 27 \dots 48 + 35$.
\choice
{$>$}
{$<$}
{Không so sánh được}
{\True $=$}
\loigiai{
$56 + 27 = 83$ và $48 + 35 = 83$. Vậy bằng nhau.\\
\textbf{Đáp án của câu hỏi là D}
}
\end{ex}

\begin{ex} %[2D14V69]
Mẹ mua 45 quả trứng gà và 35 quả trứng vịt. Mẹ đã dùng 28 quả trứng. Hỏi mẹ còn lại bao nhiêu quả trứng?
\choice
{$51$ quả}
{$53$ quả}
{\True $52$ quả}
{$54$ quả}
\loigiai{
Tổng số trứng: $45 + 35 = 80$ quả.\\
Còn lại: $80 - 28 = 52$ quả.\\
\textbf{Đáp án của câu hỏi là C}
}
\end{ex}

\begin{ex} %[2D14C69]
Tổng của hai số là 84. Nếu thêm vào số thứ nhất 8 đơn vị và bớt ở số thứ hai 12 đơn vị thì tổng mới là
\choice
{$88$}
{\True $80$}
{$82$}
{$86$}
\loigiai{
$84 + 8 - 12 = 80$.\\
\textbf{Đáp án của câu hỏi là B}
}
\end{ex}

\textbf{Phần 2. Trắc nghiệm đúng sai}
\setcounterpageref{ex}{0}

\begin{ex} %[2D14H69]
Xét tính đúng sai của các phép toán trong phạm vi 100:
\choiceTF
{\True $47 + 38 = 85$.}
{\True $91 - 46 = 45$.}
{$63 - 28 = 45$.}
{\True $63 - 28 = 35$.}
\loigiai{
\begin{itemize}
    \item $47 + 38 = 85$: Đúng.
    \item $91 - 46 = 45$: Đúng.
    \item $63 - 28 = 35 \neq 45$. Ý 3 Sai, Ý 4 Đúng.
\end{itemize}
\begin{itemize}
    \item a) \textbf{Đúng}.
    \item b) \textbf{Đúng}.
    \item c) \textbf{Sai}.
    \item d) \textbf{Đúng}.
\end{itemize}
}
\end{ex}

\begin{ex} %[2D14V69]
Bạn Mai có 45 cái kẹo, bạn Lan có ít hơn bạn Mai 12 cái kẹo. Xét tính đúng sai:
\choiceTF
{\True Lan có 33 cái kẹo ($45 - 12 = 33$).}
{\True Cả hai bạn có tất cả 78 cái kẹo ($45 + 33 = 78$).}
{Cả hai bạn có tất cả 80 cái kẹo.}
{\True Nếu Lan được cho thêm 2 cái kẹo thì cả hai bạn có 80 cái kẹo.}
\loigiai{
\begin{itemize}
    \item Lan có 33 cái: Đúng.
    \item Cả 2 bạn có 78 cái: Đúng.
    \item $78 \neq 80$. Ý 3 Sai, Ý 4 Đúng.
\end{itemize}
\begin{itemize}
    \item a) \textbf{Đúng}.
    \item b) \textbf{Đúng}.
    \item c) \textbf{Sai}.
    \item d) \textbf{Đúng}.
\end{itemize}
}
\end{ex}

\textbf{Phần 3. Trả lời ngắn}
\setcounterpageref{ex}{0}

\begin{ex} %[2D14H69]
Tính: $48 + 37$.
\shortans{85}
\loigiai{
\textbf{Đáp số:} $85$\\
$48 + 37 = 85$.
}
\end{ex}

\begin{ex} %[2D14H69]
Tính: $82 - 45$.
\shortans{37}
\loigiai{
\textbf{Đáp số:} $37$\\
$82 - 45 = 37$.
}
\end{ex}

\begin{ex} %[2D14V69]
Tính giá trị biểu thức: $100 - 36 - 24$.
\shortans{40}
\loigiai{
\textbf{Đáp số:} $40$\\
$64 - 24 = 40$.
}
\end{ex}

\begin{ex} %[2D14V69]
Tìm $x$ biết: $x + 38 = 91$.
\shortans{53}
\loigiai{
\textbf{Đáp số:} $53$\\
$x = 91 - 38 = 53$.
}
\end{ex}

\begin{ex} %[2D14V69]
Một cửa hàng có $80\text{ kg}$ gạo, buổi sáng bán được $35\text{ kg}$, buổi chiều bán được $28\text{ kg}$. Hỏi cửa hàng còn lại bao nhiêu ki-lô-gam gạo?
\shortans{17}
\loigiai{
\textbf{Đáp số:} $17$\\
$80 - (35 + 28) = 80 - 63 = 17\text{ kg}$.
}
\end{ex}

\begin{ex} %[2D14C69]
Tính nhanh: $27 + 38 + 13 + 12$.
\shortans{90}
\loigiai{
\textbf{Đáp số:} $90$\\
$(27 + 13) + (38 + 12) = 40 + 50 = 90$.
}
\end{ex}

\newpage

\TaoPhieuKiemTra{Chủ đề 14 -- Ôn tập cuối năm}{70}{Ôn tập các số trong phạm vi 1000}{https://loptoanthaylinh.vn/thi-online/toan-2/bai-70}
\textbf{Phần 1. Trắc nghiệm lựa chọn}
\setcounterpageref{ex}{0}

\begin{ex} %[2D14N70]
Kết quả của phép tính $450 + 320$ là
\choice
{$760$}
{$780$}
{\True $770$}
{$790$}
\loigiai{
$450 + 320 = 770$.\\
\textbf{Đáp án của câu hỏi là C}
}
\end{ex}

\begin{ex} %[2D14N70]
Kết quả của phép tính $890 - 450$ là
\choice
{$430$}
{\True $440$}
{$450$}
{$460$}
\loigiai{
$890 - 450 = 440$.\\
\textbf{Đáp án của câu hỏi là B}
}
\end{ex}

\begin{ex} %[2D14H70]
Kết quả của phép tính $538 + 276$ là
\choice
{\True $814$}
{$804$}
{$824$}
{$815$}
\loigiai{
$538 + 276 = 814$.\\
\textbf{Đáp án của câu hỏi là A}
}
\end{ex}

\begin{ex} %[2D14H70]
Kết quả của phép tính $725 - 368$ là
\choice
{$347$}
{$367$}
{\True $357$}
{$358$}
\loigiai{
$725 - 368 = 357$.\\
\textbf{Đáp án của câu hỏi là C}
}
\end{ex}

\begin{ex} %[2D14H70]
Tìm $x$ biết: $x - 345 = 485$.
\choice
{$820$}
{$840$}
{$831$}
{\True $830$}
\loigiai{
$x = 485 + 345 = 830$.\\
\textbf{Đáp án của câu hỏi là D}
}
\end{ex}

\begin{ex} %[2D14H70]
Tìm $y$ biết: $900 - y = 428$.
\choice
{\True $472$}
{$462$}
{$482$}
{$473$}
\loigiai{
$y = 900 - 428 = 472$.\\
\textbf{Đáp án của câu hỏi là A}
}
\end{ex}

\begin{ex} %[2D14V70]
Một trường tiểu học có 468 học sinh nam và 452 học sinh nữ. Toàn trường có tất cả số học sinh là
\choice
{$910$ bạn}
{$930$ bạn}
{$900$ bạn}
{\True $920$ bạn}
\loigiai{
$468 + 452 = 920$ học sinh.\\
\textbf{Đáp án của câu hỏi là D}
}
\end{ex}

\begin{ex} %[2D14V70]
Kho thóc có $1\,000\text{ kg}$ thóc, đã xuất đi $645\text{ kg}$. Trong kho còn lại số ki-lô-gam thóc là
\choice
{$345\text{ kg}$}
{\True $355\text{ kg}$}
{$365\text{ kg}$}
{$356\text{ kg}$}
\loigiai{
$1000 - 645 = 355\text{ kg}$.\\
\textbf{Đáp án của câu hỏi là B}
}
\end{ex}

\begin{ex} %[2D14V70]
Tính giá trị biểu thức: $425 + 285 - 360$.
\choice
{\True $350$}
{$340$}
{$360$}
{$370$}
\loigiai{
$710 - 360 = 350$.\\
\textbf{Đáp án của câu hỏi là A}
}
\end{ex}

\begin{ex} %[2D14V70]
Điền dấu thích hợp: $536 + 284 \dots 900 - 80$.
\choice
{$>$}
{\True $=$}
{$<$}
{Không so sánh được}
\loigiai{
$536 + 284 = 820$ và $900 - 80 = 820$. Vậy bằng nhau.\\
\textbf{Đáp án của câu hỏi là B}
}
\end{ex}

\begin{ex} %[2D14V70]
Tính nhanh: $345 + 180 + 155 + 120$.
\choice
{$790$}
{$810$}
{\True $800$}
{$820$}
\loigiai{
$(345 + 155) + (180 + 120) = 500 + 300 = 800$.\\
\textbf{Đáp án của câu hỏi là C}
}
\end{ex}

\begin{ex} %[2D14C70]
Hiệu của số tròn trăm lớn nhất có ba chữ số ($900$) và số bé nhất có ba chữ số khác nhau ($102$) là
\choice
{$788$}
{$808$}
{$799$}
{\True $798$}
\loigiai{
$900 - 102 = 798$.\\
\textbf{Đáp án của câu hỏi là D}
}
\end{ex}

\textbf{Phần 2. Trắc nghiệm đúng sai}
\setcounterpageref{ex}{0}

\begin{ex} %[2D14H70]
Xét tính đúng sai của các phép toán trong phạm vi 1 000:
\choiceTF
{\True $468 + 352 = 820$.}
{\True $820 - 468 = 352$.}
{$1\,000 - 450 = 650$.}
{\True $1\,000 - 450 = 550$.}
\loigiai{
\begin{itemize}
    \item $468 + 352 = 820$: Đúng.
    \item $820 - 468 = 352$: Đúng.
    \item $1000 - 450 = 550 \neq 650$. Ý 3 Sai, Ý 4 Đúng.
\end{itemize}
\begin{itemize}
    \item a) \textbf{Đúng}.
    \item b) \textbf{Đúng}.
    \item c) \textbf{Sai}.
    \item d) \textbf{Đúng}.
\end{itemize}
}
\end{ex}

\begin{ex} %[2D14V70]
Đội Một trồng được 350 cây, Đội Hai trồng được nhiều hơn Đội Một 80 cây. Xét tính đúng sai:
\choiceTF
{\True Đội Hai trồng được 430 cây ($350 + 80 = 430$).}
{\True Cả hai đội trồng được 780 cây ($350 + 430 = 780$).}
{Cả hai đội trồng được 800 cây.}
{\True Hai đội cần trồng thêm 20 cây nữa để tròn 800 cây.}
\loigiai{
\begin{itemize}
    \item Đội 2: 430 cây: Đúng.
    \item Cả 2 đội: 780 cây: Đúng.
    \item $780 \neq 800$. Ý 3 Sai, Ý 4 Đúng.
\end{itemize}
\begin{itemize}
    \item a) \textbf{Đúng}.
    \item b) \textbf{Đúng}.
    \item c) \textbf{Sai}.
    \item d) \textbf{Đúng}.
\end{itemize}
}
\end{ex}

\textbf{Phần 3. Trả lời ngắn}
\setcounterpageref{ex}{0}

\begin{ex} %[2D14H70]
Tính: $475 + 286$.
\shortans{761}
\loigiai{
\textbf{Đáp số:} $761$\\
$475 + 286 = 761$.
}
\end{ex}

\begin{ex} %[2D14H70]
Tính: $834 - 369$.
\shortans{465}
\loigiai{
\textbf{Đáp số:} $465$\\
$834 - 369 = 465$.
}
\end{ex}

\begin{ex} %[2D14V70]
Tính giá trị biểu thức: $1000 - 350 - 250$.
\shortans{400}
\loigiai{
\textbf{Đáp số:} $400$\\
$650 - 250 = 400$.
}
\end{ex}

\begin{ex} %[2D14V70]
Tìm $x$ biết: $x + 278 = 615$.
\shortans{337}
\loigiai{
\textbf{Đáp số:} $337$\\
$x = 615 - 278 = 337$.
}
\end{ex}

\begin{ex} %[2D14V70]
Một bể chứa $900\text{ l}$ nước, đã tháo ra $465\text{ l}$. Hỏi trong bể còn lại bao nhiêu lít nước?
\shortans{435}
\loigiai{
\textbf{Đáp số:} $435$\\
$900 - 465 = 435\text{ l}$.
}
\end{ex}

\begin{ex} %[2D14C70]
Tính nhanh: $875 - 190 - 210$.
\shortans{475}
\loigiai{
\textbf{Đáp số:} $475$\\
$875 - (190 + 210) = 875 - 400 = 475$.
}
\end{ex}

\newpage

\TaoPhieuKiemTra{Chủ đề 14 -- Ôn tập cuối năm}{71}{Ôn tập phép cộng, phép trừ trong phạm vi 1000}{https://loptoanthaylinh.vn/thi-online/toan-2/bai-71}
\textbf{Phần 1. Trắc nghiệm lựa chọn}
\setcounterpageref{ex}{0}

\begin{ex} %[2D14N71]
Kết quả của phép tính $2 \times 7$ là
\choice
{$12$}
{$16$}
{\True $14$}
{$18$}
\loigiai{
$2 \times 7 = 14$.\\
\textbf{Đáp án của câu hỏi là C}
}
\end{ex}

\begin{ex} %[2D14N71]
Kết quả của phép tính $5 \times 8$ là
\choice
{$35$}
{$45$}
{$50$}
{\True $40$}
\loigiai{
$5 \times 8 = 40$.\\
\textbf{Đáp án của câu hỏi là D}
}
\end{ex}

\begin{ex} %[2D14H71]
Kết quả của phép tính $18 : 2$ là
\choice
{\True $9$}
{$8$}
{$7$}
{$6$}
\loigiai{
$18 : 2 = 9$.\\
\textbf{Đáp án của câu hỏi là A}
}
\end{ex}

\begin{ex} %[2D14H71]
Kết quả của phép tính $45 : 5$ là
\choice
{$8$}
{\True $9$}
{$7$}
{$6$}
\loigiai{
$45 : 5 = 9$.\\
\textbf{Đáp án của câu hỏi là B}
}
\end{ex}

\begin{ex} %[2D14H71]
Trong phép nhân $5 \times 9 = 45$, số 45 được gọi là
\choice
{Thừa số}
{Số bị chia}
{Thương}
{\True Tích}
\loigiai{
45 là Tích.\\
\textbf{Đáp án của câu hỏi là D}
}
\end{ex}

\begin{ex} %[2D14H71]
Trong phép chia $20 : 2 = 10$, số 2 được gọi là
\choice
{Số bị chia}
{\True Số chia}
{Thương}
{Tích}
\loigiai{
2 là Số chia.\\
\textbf{Đáp án của câu hỏi là B}
}
\end{ex}

\begin{ex} %[2D14V71]
Mỗi bạn có 5 cái nhãn vở. Hỏi 7 bạn có tất cả bao nhiêu cái nhãn vở?
\choice
{\True $35$ cái}
{$30$ cái}
{$40$ cái}
{$45$ cái}
\loigiai{
$5 \times 7 = 35$ cái.\\
\textbf{Đáp án của câu hỏi là A}
}
\end{ex}

\begin{ex} %[2D14V71]
Có 16 cái bánh chia đều cho 2 bạn. Mỗi bạn được số cái bánh là
\choice
{$7$ cái}
{$9$ cái}
{\True $8$ cái}
{$10$ cái}
\loigiai{
$16 : 2 = 8$ cái.\\
\textbf{Đáp án của câu hỏi là C}
}
\end{ex}

\begin{ex} %[2D14V71]
Có 35 lít nước mắm rót đều vào 5 can. Mỗi can có số lít nước mắm là
\choice
{$6\text{ l}$}
{$8\text{ l}$}
{$9\text{ l}$}
{\True $7\text{ l}$}
\loigiai{
$35 : 5 = 7\text{ l}$.\\
\textbf{Đáp án của câu hỏi là D}
}
\end{ex}

\begin{ex} %[2D14V71]
Tính giá trị biểu thức: $2 \times 9 + 32$.
\choice
{\True $50$}
{$48$}
{$52$}
{$54$}
\loigiai{
$18 + 32 = 50$.\\
\textbf{Đáp án của câu hỏi là A}
}
\end{ex}

\begin{ex} %[2D14V71]
Tính giá trị biểu thức: $50 : 5 - 4$.
\choice
{$5$}
{\True $6$}
{$7$}
{$8$}
\loigiai{
$10 - 4 = 6$.\\
\textbf{Đáp án của câu hỏi là B}
}
\end{ex}

\begin{ex} %[2D14C71]
Một lớp có 30 học sinh chia đều thành 5 tổ. Mỗi tổ có 2 bạn nữ. Hỏi cả lớp có bao nhiêu bạn nam?
\choice
{$18$ bạn}
{$22$ bạn}
{\True $20$ bạn}
{$24$ bạn}
\loigiai{
Số bạn nữ cả lớp: $2 \times 5 = 10$ bạn.\\
Số bạn nam cả lớp: $30 - 10 = 20$ bạn.\\
\textbf{Đáp án của câu hỏi là C}
}
\end{ex}

\textbf{Phần 2. Trắc nghiệm đúng sai}
\setcounterpageref{ex}{0}

\begin{ex} %[2D14H71]
Xét tính đúng sai của các phép toán nhân chia:
\choiceTF
{\True $2 \times 8 = 16$.}
{\True $5 \times 6 = 30$.}
{$14 : 2 = 8$.}
{\True $14 : 2 = 7$.}
\loigiai{
\begin{itemize}
    \item $2 \times 8 = 16$: Đúng.
    \item $5 \times 6 = 30$: Đúng.
    \item $14 : 2 = 7 \neq 8$. Ý 3 Sai, Ý 4 Đúng.
\end{itemize}
\begin{itemize}
    \item a) \textbf{Đúng}.
    \item b) \textbf{Đúng}.
    \item c) \textbf{Sai}.
    \item d) \textbf{Đúng}.
\end{itemize}
}
\end{ex}

\begin{ex} %[2D14V71]
Có 40 cái kẹo chia đều vào 5 hộp. Xét tính đúng sai:
\choiceTF
{\True Mỗi hộp có 8 cái kẹo ($40 : 5 = 8$).}
{\True 3 hộp như thế có 24 cái kẹo ($8 \times 3 = 24$).}
{Nếu bớt đi 1 hộp thì còn lại 30 cái kẹo.}
{\True Khi bớt 1 hộp thì còn lại $8 \times 4 = 32$ cái kẹo.}
\loigiai{
\begin{itemize}
    \item Mỗi hộp 8 cái: Đúng.
    \item 3 hộp có 24 cái: Đúng.
    \item $8 \times 4 = 32 \neq 30$. Ý 3 Sai, Ý 4 Đúng.
\end{itemize}
\begin{itemize}
    \item a) \textbf{Đúng}.
    \item b) \textbf{Đúng}.
    \item c) \textbf{Sai}.
    \item d) \textbf{Đúng}.
\end{itemize}
}
\end{ex}

\textbf{Phần 3. Trả lời ngắn}
\setcounterpageref{ex}{0}

\begin{ex} %[2D14H71]
Tính: $2 \times 9$.
\shortans{18}
\loigiai{
\textbf{Đáp số:} $18$\\
$2 \times 9 = 18$.
}
\end{ex}

\begin{ex} %[2D14H71]
Tính: $5 \times 9$.
\shortans{45}
\loigiai{
\textbf{Đáp số:} $45$\\
$5 \times 9 = 45$.
}
\end{ex}

\begin{ex} %[2D14V71]
Tính: $18 : 2$.
\shortans{9}
\loigiai{
\textbf{Đáp số:} $9$\\
$18 : 2 = 9$.
}
\end{ex}

\begin{ex} %[2D14V71]
Tính: $40 : 5$.
\shortans{8}
\loigiai{
\textbf{Đáp số:} $8$\\
$40 : 5 = 8$.
}
\end{ex}

\begin{ex} %[2D14V71]
Tính giá trị biểu thức: $5 \times 7 + 25$.
\shortans{60}
\loigiai{
\textbf{Đáp số:} $60$\\
$35 + 25 = 60$.
}
\end{ex}

\begin{ex} %[2D14C71]
Có 5 túi bi, mỗi túi 4 viên bi. Sau khi cho bạn 6 viên bi thì còn lại bao nhiêu viên bi?
\shortans{14}
\loigiai{
\textbf{Đáp số:} $14$\\
$5 \times 4 - 6 = 20 - 6 = 14$ viên bi.
}
\end{ex}

\newpage

\TaoPhieuKiemTra{Chủ đề 14 -- Ôn tập cuối năm}{72}{Ôn tập phép nhân, phép chia}{https://loptoanthaylinh.vn/thi-online/toan-2/bai-72}
\textbf{Phần 1. Trắc nghiệm lựa chọn}
\setcounterpageref{ex}{0}

\begin{ex} %[2H14N72]
Hình nào sau đây có 3 cạnh và 3 đỉnh?
\choice
{Hình vuông}
{Hình tứ giác}
{\True Hình tam giác}
{Hình chữ nhật}
\loigiai{
Hình tam giác có 3 cạnh và 3 đỉnh.\\
\textbf{Đáp án của câu hỏi là C}
}
\end{ex}

\begin{ex} %[2H14N72]
Hình tứ giác có bao nhiêu cạnh?
\choice
{$3$ cạnh}
{$5$ cạnh}
{$6$ cạnh}
{\True $4$ cạnh}
\loigiai{
Hình tứ giác có 4 cạnh.\\
\textbf{Đáp án của câu hỏi là D}
}
\end{ex}

\begin{ex} %[2H14H72]
Đường gấp khúc $ABCD$ gồm mấy đoạn thẳng?
\choice
{$2$ đoạn thẳng}
{$4$ đoạn thẳng}
{\True $3$ đoạn thẳng ($AB, BC, CD$)}
{$5$ đoạn thẳng}
\loigiai{
Gồm 3 đoạn thẳng: $AB, BC, CD$.\\
\textbf{Đáp án của câu hỏi là C}
}
\end{ex}

\begin{ex} %[2H14H72]
Độ dài đường gấp khúc gồm hai đoạn thẳng dài $15\text{ cm}$ và $25\text{ cm}$ là
\choice
{\True $40\text{ cm}$}
{$35\text{ cm}$}
{$45\text{ cm}$}
{$50\text{ cm}$}
\loigiai{
$15 + 25 = 40\text{ cm}$.\\
\textbf{Đáp án của câu hỏi là A}
}
\end{ex}

\begin{ex} %[2H14H72]
Khối nào sau đây lăn được về mọi phía?
\choice
{Khối trụ}
{Khối lập phương}
{Khối hộp chữ nhật}
{\True Khối cầu}
\loigiai{
Khối cầu lăn được về mọi phía.\\
\textbf{Đáp án của câu hỏi là D}
}
\end{ex}

\begin{ex} %[2H14H72]
Khối lập phương có bao nhiêu mặt?
\choice
{$4$ mặt}
{$8$ mặt}
{\True $6$ mặt vuông bằng nhau}
{$12$ mặt}
\loigiai{
Khối lập phương có 6 mặt vuông bằng nhau.\\
\textbf{Đáp án của câu hỏi là C}
}
\end{ex}

\begin{ex} %[2H14V72]
Đường gấp khúc gồm 3 đoạn thẳng có độ dài lần lượt là $12\text{ cm}, 18\text{ cm}, 20\text{ cm}$. Độ dài đường gấp khúc đó là
\choice
{$40\text{ cm}$}
{\True $50\text{ cm}$}
{$60\text{ cm}$}
{$45\text{ cm}$}
\loigiai{
$12 + 18 + 20 = 50\text{ cm}$.\\
\textbf{Đáp án của câu hỏi là B}
}
\end{ex}

\begin{ex} %[2H14V72]
Một hình tam giác có độ dài các cạnh là $15\text{ cm}, 20\text{ cm}, 25\text{ cm}$. Chu vi hình tam giác đó là
\choice
{$50\text{ cm}$}
{$70\text{ cm}$}
{$55\text{ cm}$}
{\True $60\text{ cm}$}
\loigiai{
$15 + 20 + 25 = 60\text{ cm}$.\\
\textbf{Đáp án của câu hỏi là D}
}
\end{ex}

\begin{ex} %[2H14V72]
Một hình tứ giác có độ dài 4 cạnh đều bằng $5\text{ cm}$. Chu vi của hình tứ giác đó là
\choice
{$15\text{ cm}$}
{\True $20\text{ cm}$}
{$25\text{ cm}$}
{$10\text{ cm}$}
\loigiai{
$5 \times 4 = 20\text{ cm}$.\\
\textbf{Đáp án của câu hỏi là B}
}
\end{ex}

\begin{ex} %[2H14V72]
Trong các vật sau: hộp sữa chua hình trụ, quả bóng bàn, viên gạch, quả cam. Có bao nhiêu vật có dạng khối cầu?
\choice
{\True $2$ vật (quả bóng bàn, quả cam)}
{$1$ vật}
{$3$ vật}
{$4$ vật}
\loigiai{
Có 2 vật: quả bóng bàn và quả cam.\\
\textbf{Đáp án của câu hỏi là A}
}
\end{ex}

\begin{ex} %[2H14V72]
Một đường gấp khúc gồm 4 đoạn thẳng bằng nhau, mỗi đoạn dài $5\text{ dm}$. Độ dài đường gấp khúc theo mét là
\choice
{$1\text{ m}$}
{\True $2\text{ m}$ ($20\text{ dm}$)}
{$3\text{ m}$}
{$4\text{ m}$}
\loigiai{
$5 \times 4 = 20\text{ dm} = 2\text{ m}$.\\
\textbf{Đáp án của câu hỏi là B}
}
\end{ex}

\begin{ex} %[2H14C72]
Một sợi dây đồng uốn thành một hình tam giác có 3 cạnh bằng nhau, mỗi cạnh dài $8\text{ cm}$. Nếu duỗi thẳng sợi dây rồi uốn thành hình tứ giác có 4 cạnh bằng nhau thì mỗi cạnh của hình tứ giác dài là
\choice
{\True $6\text{ cm}$}
{$5\text{ cm}$}
{$7\text{ cm}$}
{$8\text{ cm}$}
\loigiai{
Độ dài dây: $8 \times 3 = 24\text{ cm}$.\\
Cạnh hình tứ giác: $24 : 4 = 6\text{ cm}$.\\
\textbf{Đáp án của câu hỏi là A}
}
\end{ex}

\textbf{Phần 2. Trắc nghiệm đúng sai}
\setcounterpageref{ex}{0}

\begin{ex} %[2H14H72]
Xét tính đúng sai về các hình hình học:
\choiceTF
{\True Ba điểm thẳng hàng là ba điểm cùng nằm trên một đường thẳng.}
{\True Đường gấp khúc $MNPQ$ gồm 3 đoạn thẳng $MN, NP, PQ$.}
{Khối trụ có hai mặt đáy là hai hình vuông.}
{\True Khối trụ có hai mặt đáy là hai hình tròn bằng nhau.}
\loigiai{
\begin{itemize}
    \item 3 điểm thẳng hàng: Đúng.
    \item Đường gấp khúc MNPQ có 3 đoạn: Đúng.
    \item Đáy khối trụ là hình tròn (không phải hình vuông). Ý 3 Sai, Ý 4 Đúng.
\end{itemize}
\begin{itemize}
    \item a) \textbf{Đúng}.
    \item b) \textbf{Đúng}.
    \item c) \textbf{Sai}.
    \item d) \textbf{Đúng}.
\end{itemize}
}
\end{ex}

\begin{ex} %[2H14V72]
Một hình tứ giác có độ dài các cạnh lần lượt là $10\text{ cm}, 12\text{ cm}, 15\text{ cm}, 13\text{ cm}$. Xét tính đúng sai:
\choiceTF
{\True Chu vi hình tứ giác là $50\text{ cm}$ ($10 + 12 + 15 + 13 = 50$).}
{\True Chu vi hình tứ giác đổi ra đề-xi-mét là $5\text{ dm}$.}
{Chu vi hình tứ giác là $60\text{ cm}$.}
{\True $50\text{ cm} = 5\text{ dm} \neq 60\text{ cm}$.}
\loigiai{
\begin{itemize}
    \item Chu vi = 50 cm: Đúng.
    \item $50\text{ cm} = 5\text{ dm}$: Đúng.
    \item $50 \neq 60$. Ý 3 Sai, Ý 4 Đúng.
\end{itemize}
\begin{itemize}
    \item a) \textbf{Đúng}.
    \item b) \textbf{Đúng}.
    \item c) \textbf{Sai}.
    \item d) \textbf{Đúng}.
\end{itemize}
}
\end{ex}

\textbf{Phần 3. Trả lời ngắn}
\setcounterpageref{ex}{0}

\begin{ex} %[2H14H72]
Hình tứ giác có mấy cạnh?
\shortans{4}
\loigiai{
\textbf{Đáp số:} $4$\\
4 cạnh.
}
\end{ex}

\begin{ex} %[2H14H72]
Khối lập phương có mấy mặt?
\shortans{6}
\loigiai{
\textbf{Đáp số:} $6$\\
6 mặt.
}
\end{ex}

\begin{ex} %[2H14V72]
Đường gấp khúc gồm 2 đoạn thẳng dài $18\text{ cm}$ và $22\text{ cm}$. Tính độ dài đường gấp khúc theo xăng-ti-mét.
\shortans{40}
\loigiai{
\textbf{Đáp số:} $40$\\
$18 + 22 = 40\text{ cm}$.
}
\end{ex}

\begin{ex} %[2H14V72]
Một hình tam giác có 3 cạnh đều dài $7\text{ cm}$. Tính chu vi hình tam giác đó theo xăng-ti-mét.
\shortans{21}
\loigiai{
\textbf{Đáp số:} $21$\\
$7 \times 3 = 21\text{ cm}$.
}
\end{ex}

\begin{ex} %[2H14V72]
Có 8 khối cầu và 7 khối trụ. Hỏi có tất cả bao nhiêu khối hình?
\shortans{15}
\loigiai{
\textbf{Đáp số:} $15$\\
$8 + 7 = 15$ khối.
}
\end{ex}

\begin{ex} %[2H14C72]
Đường gấp khúc gồm 5 đoạn thẳng bằng nhau có tổng độ dài là $35\text{ cm}$. Hỏi mỗi đoạn thẳng dài bao nhiêu xăng-ti-mét?
\shortans{7}
\loigiai{
\textbf{Đáp số:} $7$\\
$35 : 5 = 7\text{ cm}$.
}
\end{ex}

\newpage

\TaoPhieuKiemTra{Chủ đề 14 -- Ôn tập cuối năm}{73}{Ôn tập hình học và đo lường}{https://loptoanthaylinh.vn/thi-online/toan-2/bai-73}
\textbf{Phần 1. Trắc nghiệm lựa chọn}
\setcounterpageref{ex}{0}

\begin{ex} %[2H14N73]
$1\text{ m}$ bằng bao nhiêu xăng-ti-mét?
\choice
{$10\text{ cm}$}
{$1000\text{ cm}$}
{$1\text{ cm}$}
{\True $100\text{ cm}$}
\loigiai{
$1\text{ m} = 100\text{ cm}$.\\
\textbf{Đáp án của câu hỏi là D}
}
\end{ex}

\begin{ex} %[2H14N73]
$1\text{ km}$ bằng bao nhiêu mét?
\choice
{\True $1\,000\text{ m}$}
{$100\text{ m}$}
{$10\text{ m}$}
{$10\,000\text{ m}$}
\loigiai{
$1\text{ km} = 1000\text{ m}$.\\
\textbf{Đáp án của câu hỏi là A}
}
\end{ex}

\begin{ex} %[2H14H73]
1 ngày có bao nhiêu giờ?
\choice
{$12$ giờ}
{\True $24$ giờ}
{$30$ giờ}
{$60$ giờ}
\loigiai{
1 ngày = 24 giờ.\\
\textbf{Đáp án của câu hỏi là B}
}
\end{ex}

\begin{ex} %[2H14H73]
1 giờ có bao nhiêu phút?
\choice
{$30$ phút}
{$100$ phút}
{$24$ phút}
{\True $60$ phút}
\loigiai{
1 giờ = 60 phút.\\
\textbf{Đáp án của câu hỏi là D}
}
\end{ex}

\begin{ex} %[2H14H73]
Tháng 1 có bao nhiêu ngày?
\choice
{\True $31$ ngày}
{$30$ ngày}
{$28$ ngày}
{$29$ ngày}
\loigiai{
Tháng 1 có 31 ngày.\\
\textbf{Đáp án của câu hỏi là A}
}
\end{ex}

\begin{ex} %[2H14H73]
Đồng hồ chỉ 8 giờ 15 phút, kim phút đang chỉ vào số nào?
\choice
{Số 12}
{Số 6}
{\True Số 3}
{Số 9}
\loigiai{
15 phút tương ứng với kim phút chỉ số 3.\\
\textbf{Đáp án của câu hỏi là C}
}
\end{ex}

\begin{ex} %[2H14V73]
Tính: $35\text{ kg} + 28\text{ kg} = \dots$
\choice
{\True $63\text{ kg}$}
{$62\text{ kg}$}
{$64\text{ kg}$}
{$65\text{ kg}$}
\loigiai{
$35 + 28 = 63\text{ kg}$.\\
\textbf{Đáp án của câu hỏi là A}
}
\end{ex}

\begin{ex} %[2H14V73]
Tính: $50\text{ l} - 18\text{ l} = \dots$
\choice
{$31\text{ l}$}
{\True $32\text{ l}$}
{$33\text{ l}$}
{$34\text{ l}$}
\loigiai{
$50 - 18 = 32\text{ l}$.\\
\textbf{Đáp án của câu hỏi là B}
}
\end{ex}

\begin{ex} %[2H14V73]
Mai có 2 tờ 500 đồng. Mai có tất cả số tiền là
\choice
{$500$ đồng}
{$800$ đồng}
{\True $1\,000$ đồng}
{$700$ đồng}
\loigiai{
$500 \times 2 = 1000$ đồng.\\
\textbf{Đáp án của câu hỏi là C}
}
\end{ex}

\begin{ex} %[2H14V73]
Một cuộn vải dài $45\text{ m}$, người ta đã bán $18\text{ m}$. Cuộn vải còn lại dài là
\choice
{$26\text{ m}$}
{\True $27\text{ m}$}
{$28\text{ m}$}
{$29\text{ m}$}
\loigiai{
$45 - 18 = 27\text{ m}$.\\
\textbf{Đáp án của câu hỏi là B}
}
\end{ex}

\begin{ex} %[2H14V73]
Bình bắt đầu học bài lúc 19 giờ và học xong lúc 20 giờ. Bình đã học bài trong thời gian là
\choice
{$30$ phút}
{$2$ giờ}
{\True $1$ giờ ($60$ phút)}
{$15$ phút}
\loigiai{
$20 - 19 = 1$ giờ.\\
\textbf{Đáp án của câu hỏi là C}
}
\end{ex}

\begin{ex} %[2H14C73]
Nếu ngày 20 tháng 11 là thứ Tư thì ngày 27 tháng 11 cùng năm đó là thứ mấy?
\choice
{Thứ Ba}
{Thứ Năm}
{Thứ Sáu}
{\True Thứ Tư}
\loigiai{
Cách nhau đúng 7 ngày ($27 - 20 = 7$) nên vẫn là ngày Thứ Tư.\\
\textbf{Đáp án của câu hỏi là D}
}
\end{ex}

\textbf{Phần 2. Trắc nghiệm đúng sai}
\setcounterpageref{ex}{0}

\begin{ex} %[2H14H73]
Xét tính đúng sai về các đơn vị đo đại lượng:
\choiceTF
{\True $1\text{ m} = 10\text{ dm} = 100\text{ cm}$.}
{\True 1 ngày = 24 giờ, 1 giờ = 60 phút.}
{Tháng 4 có 31 ngày.}
{\True Tháng 4 có 30 ngày.}
\loigiai{
\begin{itemize}
    \item $1\text{ m} = 10\text{ dm} = 100\text{ cm}$: Đúng.
    \item 1 ngày = 24h, 1h = 60p: Đúng.
    \item Tháng 4 có 30 ngày (không phải 31). Ý 3 Sai, Ý 4 Đúng.
\end{itemize}
\begin{itemize}
    \item a) \textbf{Đúng}.
    \item b) \textbf{Đúng}.
    \item c) \textbf{Sai}.
    \item d) \textbf{Đúng}.
\end{itemize}
}
\end{ex}

\begin{ex} %[2H14V73]
Một can dầu có $20\text{ l}$, rót đều vào 4 chai nhỏ. Xét tính đúng sai:
\choiceTF
{\True Mỗi chai nhỏ có $5\text{ l}$ dầu ($20 : 4 = 5$).}
{\True 3 chai nhỏ có tất cả $15\text{ l}$ dầu ($5 \times 3 = 15$).}
{Nếu rót vào 5 chai thì mỗi chai có $6\text{ l}$.}
{\True Nếu rót vào 5 chai thì mỗi chai có $20 : 5 = 4\text{ l}$.}
\loigiai{
\begin{itemize}
    \item Mỗi chai 5l: Đúng.
    \item 3 chai 15l: Đúng.
    \item $20 : 5 = 4\text{ l} \neq 6\text{ l}$. Ý 3 Sai, Ý 4 Đúng.
\end{itemize}
\begin{itemize}
    \item a) \textbf{Đúng}.
    \item b) \textbf{Đúng}.
    \item c) \textbf{Sai}.
    \item d) \textbf{Đúng}.
\end{itemize}
}
\end{ex}

\textbf{Phần 3. Trả lời ngắn}
\setcounterpageref{ex}{0}

\begin{ex} %[2H14H73]
$1\text{ m}$ bằng bao nhiêu xăng-ti-mét?
\shortans{100}
\loigiai{
\textbf{Đáp số:} $100$\\
100 cm.
}
\end{ex}

\begin{ex} %[2H14H73]
1 giờ có bao nhiêu phút?
\shortans{60}
\loigiai{
\textbf{Đáp số:} $60$\\
60 phút.
}
\end{ex}

\begin{ex} %[2H14V73]
Tính: $45\text{ kg} + 35\text{ kg}$ theo ki-lô-gam.
\shortans{80}
\loigiai{
\textbf{Đáp số:} $80$\\
$45 + 35 = 80\text{ kg}$.
}
\end{ex}

\begin{ex} %[2H14V73]
Tính: $60\text{ l} - 25\text{ l}$ theo lít.
\shortans{35}
\loigiai{
\textbf{Đáp số:} $35$\\
$60 - 25 = 35\text{ l}$.
}
\end{ex}

\begin{ex} %[2H14V73]
Lan có 1 000 đồng, mua bút hết 700 đồng. Hỏi Lan còn lại bao nhiêu đồng?
\shortans{300}
\loigiai{
\textbf{Đáp số:} $300$\\
$1000 - 700 = 300$ đồng.
}
\end{ex}

\begin{ex} %[2H14C73]
Nếu ngày 5 tháng 5 là thứ Hai thì ngày 12 tháng 5 cùng năm đó là thứ mấy? (Viết: thứ Hai)
\shortans{thứ Hai}
\loigiai{
\textbf{Đáp số:} thứ Hai\\
Cách nhau 7 ngày nên là thứ Hai.
}
\end{ex}

\newpage

\TaoPhieuKiemTra{Chủ đề 14 -- Ôn tập cuối năm}{74}{Ôn tập yếu tố thống kê, xác suất}{https://loptoanthaylinh.vn/thi-online/toan-2/bai-74}
\textbf{Phần 1. Trắc nghiệm lựa chọn}
\setcounterpageref{ex}{0}

\begin{ex} %[2X14N74]
Trong một túi chỉ có các viên bi màu xanh. Lấy ra 1 viên bi, sự kiện "Lấy được viên bi xanh" là
\choice
{Có thể}
{Không thể}
{\True Chắc chắn}
{Chắc chắn không}
\loigiai{
Chỉ có bi xanh nên chắc chắn lấy được bi xanh.\\
\textbf{Đáp án của câu hỏi là C}
}
\end{ex}

\begin{ex} %[2X14N74]
Trong một túi chỉ có các viên bi màu xanh. Lấy ra 1 viên bi, sự kiện "Lấy được viên bi đỏ" là
\choice
{Chắc chắn}
{Có thể}
{Có thể hoặc không}
{\True Không thể}
\loigiai{
Không có bi đỏ nên không thể lấy được bi đỏ.\\
\textbf{Đáp án của câu hỏi là D}
}
\end{ex}

\begin{ex} %[2X14H74]
Trong túi có cả kẹo dâu và kẹo cam. Lấy ra 1 cái kẹo, sự kiện "Lấy được kẹo dâu" là
\choice
{\True Có thể}
{Chắc chắn}
{Không thể}
{Chắc chắn không}
\loigiai{
Có thể lấy được kẹo dâu.\\
\textbf{Đáp án của câu hỏi là A}
}
\end{ex}

\begin{ex} %[2X14H74]
Kiểm đếm các loại hoa trong lọ: 6 hoa hồng, 4 hoa cúc, 5 hoa đồng tiền. Loại hoa có nhiều nhất là
\choice
{Hoa cúc}
{Hoa đồng tiền}
{Bằng nhau}
{\True Hoa hồng ($6$ bông)}
\loigiai{
Hoa hồng có 6 bông (nhiều nhất).\\
\textbf{Đáp án của câu hỏi là D}
}
\end{ex}

\begin{ex} %[2X14H74]
Tổng số hoa có trong lọ là
\choice
{$14$ bông}
{\True $15$ bông}
{$16$ bông}
{$17$ bông}
\loigiai{
$6 + 4 + 5 = 15$ bông hoa.\\
\textbf{Đáp án của câu hỏi là B}
}
\end{ex}

\begin{ex} %[2X14H74]
Số hoa hồng nhiều hơn số hoa cúc là
\choice
{\True $2$ bông}
{$1$ bông}
{$3$ bông}
{$4$ bông}
\loigiai{
$6 - 4 = 2$ bông.\\
\textbf{Đáp án của câu hỏi là A}
}
\end{ex}

\begin{ex} %[2X14V74]
Biểu đồ tranh ghi số sách của mỗi bạn (mỗi biểu tượng $= 2$ cuốn sách): An có 3 biểu tượng, Bình có 5 biểu tượng. Số sách của Bình là
\choice
{$8$ cuốn}
{$12$ cuốn}
{\True $10$ cuốn}
{$15$ cuốn}
\loigiai{
$2 \times 5 = 10$ cuốn.\\
\textbf{Đáp án của câu hỏi là C}
}
\end{ex}

\begin{ex} %[2X14V74]
Bình có nhiều hơn An số cuốn sách là
\choice
{$2$ cuốn}
{\True $4$ cuốn}
{$6$ cuốn}
{$8$ cuốn}
\loigiai{
An có: $2 \times 3 = 6$ cuốn.\\
$10 - 6 = 4$ cuốn.\\
\textbf{Đáp án của câu hỏi là B}
}
\end{ex}

\begin{ex} %[2X14V74]
Gieo một con xúc xắc 6 mặt. Sự kiện "Số chấm xuất hiện lớn hơn 0 và nhỏ hơn 7" là
\choice
{Có thể}
{Không thể}
{Tùy lần gieo}
{\True Chắc chắn}
\loigiai{
Xúc xắc chỉ có từ 1 đến 6 chấm nên chắc chắn xảy ra.\\
\textbf{Đáp án của câu hỏi là D}
}
\end{ex}

\begin{ex} %[2X14V74]
Thực hành quay chong chóng 20 lần: kim chỉ vào màu đỏ 12 lần, màu xanh 8 lần. Số lần kim chỉ vào màu đỏ nhiều hơn màu xanh là
\choice
{\True $4$ lần}
{$2$ lần}
{$6$ lần}
{$8$ lần}
\loigiai{
$12 - 8 = 4$ lần.\\
\textbf{Đáp án của câu hỏi là A}
}
\end{ex}

\begin{ex} %[2X14V74]
Khảo sát trái cây yêu thích của 35 học sinh: 15 bạn thích cam, 12 bạn thích xoài, còn lại thích táo. Số bạn thích táo là
\choice
{$7$ bạn}
{\True $8$ bạn}
{$9$ bạn}
{$10$ bạn}
\loigiai{
$35 - (15 + 12) = 35 - 27 = 8$ bạn.\\
\textbf{Đáp án của câu hỏi là B}
}
\end{ex}

\begin{ex} %[2X14C74]
Trong hộp có 4 viên bi đỏ và 4 viên bi xanh. Lấy ra ngẫu nhiên 5 viên bi. Khẳng định nào sau đây là CHẮC CHẮN?
\choice
{Lấy được cả 5 viên bi đỏ}
{Lấy được cả 5 viên bi xanh}
{\True Trong 5 viên lấy ra có ít nhất 1 viên bi đỏ và 1 viên bi xanh}
{Lấy được 3 viên bi vàng}
\loigiai{
Vì mỗi màu chỉ có 4 viên nên lấy 5 viên chắc chắn có cả hai màu.\\
\textbf{Đáp án của câu hỏi là C}
}
\end{ex}

\textbf{Phần 2. Trắc nghiệm đúng sai}
\setcounterpageref{ex}{0}

\begin{ex} %[2X14H74]
Trong một hộp có 6 bút chì đen và 2 bút chì đỏ. Xét tính đúng sai khi lấy ra 1 chiếc bút:
\choiceTF
{\True Có thể lấy được bút chì đen.}
{\True Có thể lấy được bút chì đỏ.}
{Chắc chắn lấy được bút chì đen.}
{\True Không thể lấy được bút chì màu xanh (vì không có bút xanh).}
\loigiai{
\begin{itemize}
    \item Có thể lấy bút đen: Đúng.
    \item Có thể lấy bút đỏ: Đúng.
    \item Không chắc chắn đen vì có thể ra đỏ. Ý 3 Sai, Ý 4 Đúng.
\end{itemize}
\begin{itemize}
    \item a) \textbf{Đúng}.
    \item b) \textbf{Đúng}.
    \item c) \textbf{Sai}.
    \item d) \textbf{Đúng}.
\end{itemize}
}
\end{ex}

\begin{ex} %[2X14V74]
Biểu đồ tranh số con gà của 2 bạn (mỗi biểu tượng $= 5$ con gà): Nam có 4 biểu tượng, Hải có 3 biểu tượng. Xét tính đúng sai:
\choiceTF
{\True Nam có 20 con gà ($5 \times 4 = 20$).}
{\True Hải có 15 con gà ($5 \times 3 = 15$).}
{Cả hai bạn có tất cả 40 con gà.}
{\True Cả hai bạn có tất cả 35 con gà ($20 + 15 = 35$).}
\loigiai{
\begin{itemize}
    \item Nam: 20 con: Đúng.
    \item Hải: 15 con: Đúng.
    \item Tổng: $20 + 15 = 35 \neq 40$. Ý 3 Sai, Ý 4 Đúng.
\end{itemize}
\begin{itemize}
    \item a) \textbf{Đúng}.
    \item b) \textbf{Đúng}.
    \item c) \textbf{Sai}.
    \item d) \textbf{Đúng}.
\end{itemize}
}
\end{ex}

\textbf{Phần 3. Trả lời ngắn}
\setcounterpageref{ex}{0}

\begin{ex} %[2X14H74]
Trong hộp có 5 viên bi xanh. Nhắm mắt lấy 1 viên bi, sự kiện "Lấy được bi xanh" là gì? (Viết: chắc chắn, có thể, hoặc không thể)
\shortans{chắc chắn}
\loigiai{
\textbf{Đáp số:} chắc chắn\\
Chắc chắn.
}
\end{ex}

\begin{ex} %[2X14H74]
Trong hộp có 5 viên bi xanh. Nhắm mắt lấy 1 viên bi, sự kiện "Lấy được bi vàng" là gì? (Viết: chắc chắn, có thể, hoặc không thể)
\shortans{không thể}
\loigiai{
\textbf{Đáp số:} không thể\\
Không thể.
}
\end{ex}

\begin{ex} %[2X14V74]
Biểu đồ tranh: mỗi hình bông hoa biểu diễn cho 2 bông hoa thật. Có 6 hình bông hoa, hỏi có bao nhiêu bông hoa thật?
\shortans{12}
\loigiai{
\textbf{Đáp số:} $12$\\
$2 \times 6 = 12$ bông.
}
\end{ex}

\begin{ex} %[2X14V74]
Khảo sát 30 học sinh: 14 bạn thích đá bóng, 10 bạn thích bơi, còn lại thích cờ vua. Hỏi có bao nhiêu bạn thích cờ vua?
\shortans{6}
\loigiai{
\textbf{Đáp số:} $6$\\
$30 - (14 + 10) = 30 - 24 = 6$ bạn.
}
\end{ex}

\begin{ex} %[2X14V74]
Gieo đồng xu 15 lần thu được 9 lần mặt Sấp. Hỏi mặt Ngửa xuất hiện bao nhiêu lần?
\shortans{6}
\loigiai{
\textbf{Đáp số:} $6$\\
$15 - 9 = 6$ lần.
}
\end{ex}

\begin{ex} %[2X14C74]
Có 50 cây xanh cần biểu diễn lên biểu đồ tranh. Biết mỗi biểu tượng cây biểu diễn cho 5 cây thật. Hỏi cần vẽ bao nhiêu biểu tượng cây?
\shortans{10}
\loigiai{
\textbf{Đáp số:} $10$\\
$50 : 5 = 10$ biểu tượng.
}
\end{ex}

\newpage

\TaoPhieuKiemTra{Chủ đề 14 -- Ôn tập cuối năm}{75}{Ôn tập chung cuối năm}{https://loptoanthaylinh.vn/thi-online/toan-2/bai-75}
\textbf{Phần 1. Trắc nghiệm lựa chọn}
\setcounterpageref{ex}{0}

\begin{ex} %[2D14N75]
Số gồm $8$ trăm, $4$ chục và $6$ đơn vị là
\choice
{$864$}
{\True $846$}
{$486$}
{$648$}
\loigiai{
Số đó là $846$.\\
\textbf{Đáp án của câu hỏi là B}
}
\end{ex}

\begin{ex} %[2D14N75]
Số liền sau của số $799$ là
\choice
{$798$}
{$801$}
{\True $800$}
{$810$}
\loigiai{
$799 + 1 = 800$.\\
\textbf{Đáp án của câu hỏi là C}
}
\end{ex}

\begin{ex} %[2D14H75]
Kết quả của phép tính $458 + 237$ là
\choice
{\True $695$}
{$685$}
{$705$}
{$696$}
\loigiai{
$458 + 237 = 695$.\\
\textbf{Đáp án của câu hỏi là A}
}
\end{ex}

\begin{ex} %[2D14H75]
Kết quả của phép tính $835 - 362$ là
\choice
{$463$}
{$483$}
{$474$}
{\True $473$}
\loigiai{
$835 - 362 = 473$.\\
\textbf{Đáp án của câu hỏi là D}
}
\end{ex}

\begin{ex} %[2D14H75]
Kết quả của phép tính $5 \times 8$ là
\choice
{$35$}
{\True $40$}
{$45$}
{$50$}
\loigiai{
$5 \times 8 = 40$.\\
\textbf{Đáp án của câu hỏi là B}
}
\end{ex}

\begin{ex} %[2D14H75]
Kết quả của phép tính $18 : 2$ là
\choice
{$8$}
{$7$}
{\True $9$}
{$6$}
\loigiai{
$18 : 2 = 9$.\\
\textbf{Đáp án của câu hỏi là C}
}
\end{ex}

\begin{ex} %[2D14V75]
Đoạn thẳng thứ nhất dài $35\text{ cm}$, đoạn thứ hai dài $45\text{ cm}$. Cả hai đoạn dài tất cả theo đề-xi-mét là
\choice
{$7\text{ dm}$}
{$9\text{ dm}$}
{$10\text{ dm}$}
{\True $8\text{ dm}$ ($80\text{ cm}$)}
\loigiai{
$35 + 45 = 80\text{ cm} = 8\text{ dm}$.\\
\textbf{Đáp án của câu hỏi là D}
}
\end{ex}

\begin{ex} %[2D14V75]
Mỗi gói kẹo có 5 cái kẹo. Hỏi 9 gói kẹo như thế có tất cả bao nhiêu cái kẹo?
\choice
{$40$ cái}
{\True $45$ cái}
{$50$ cái}
{$35$ cái}
\loigiai{
$5 \times 9 = 45$ cái.\\
\textbf{Đáp án của câu hỏi là B}
}
\end{ex}

\begin{ex} %[2D14V75]
Có 30 bông hoa cắm đều vào 5 lọ hoa. Mỗi lọ có số bông hoa là
\choice
{\True $6$ bông}
{$5$ bông}
{$7$ bông}
{$8$ bông}
\loigiai{
$30 : 5 = 6$ bông hoa.\\
\textbf{Đáp án của câu hỏi là A}
}
\end{ex}

\begin{ex} %[2D14V75]
Một hình tứ giác có 4 cạnh bằng nhau, mỗi cạnh dài $6\text{ cm}$. Chu vi hình tứ giác đó là
\choice
{$20\text{ cm}$}
{$28\text{ cm}$}
{\True $24\text{ cm}$}
{$30\text{ cm}$}
\loigiai{
$6 \times 4 = 24\text{ cm}$.\\
\textbf{Đáp án của câu hỏi là C}
}
\end{ex}

\begin{ex} %[2D14V75]
Tính giá trị biểu thức: $450 + 250 - 320$.
\choice
{$370$}
{$390$}
{$400$}
{\True $380$}
\loigiai{
$700 - 320 = 380$.\\
\textbf{Đáp án của câu hỏi là D}
}
\end{ex}

\begin{ex} %[2D14C75]
Mẹ mua 4 hộp bánh, mỗi hộp có 5 cái bánh. Mẹ đem chia đều số bánh đó cho 2 anh em. Hỏi mỗi người được bao nhiêu cái bánh?
\choice
{\True $10$ cái}
{$8$ cái}
{$12$ cái}
{$15$ cái}
\loigiai{
Tổng số bánh: $5 \times 4 = 20$ cái.\\
Mỗi người được: $20 : 2 = 10$ cái bánh.\\
\textbf{Đáp án của câu hỏi là A}
}
\end{ex}

\textbf{Phần 2. Trắc nghiệm đúng sai}
\setcounterpageref{ex}{0}

\begin{ex} %[2D14H75]
Xét tính đúng sai của các kiến thức Toán lớp 2:
\choiceTF
{\True $1\text{ m} = 100\text{ cm} = 10\text{ dm}$.}
{\True Số lớn nhất có ba chữ số khác nhau là 987.}
{Khối cầu có 6 mặt bằng nhau.}
{\True Khối lập phương có 6 mặt vuông bằng nhau (khối cầu không có mặt phẳng).}
\loigiai{
\begin{itemize}
    \item $1\text{ m} = 100\text{ cm} = 10\text{ dm}$: Đúng.
    \item 987 lớn nhất có 3 chữ số khác nhau: Đúng.
    \item Khối cầu không có mặt phẳng. Ý 3 Sai, Ý 4 Đúng.
\end{itemize}
\begin{itemize}
    \item a) \textbf{Đúng}.
    \item b) \textbf{Đúng}.
    \item c) \textbf{Sai}.
    \item d) \textbf{Đúng}.
\end{itemize}
}
\end{ex}

\begin{ex} %[2D14V75]
Một nông trại có 450 con vịt và 380 con gà. Xét tính đúng sai:
\choiceTF
{\True Tổng số gia cầm là 830 con ($450 + 380 = 830$).}
{\True Số vịt nhiều hơn số gà là 70 con ($450 - 380 = 70$).}
{Nông trại có tất cả 900 con gia cầm.}
{\True Nông trại cần nuôi thêm 70 con nữa để có tất cả 900 con.}
\loigiai{
\begin{itemize}
    \item Tổng: 830 con: Đúng.
    \item Vịt hơn gà 70 con: Đúng.
    \item $830 \neq 900$. Ý 3 Sai, Ý 4 Đúng.
\end{itemize}
\begin{itemize}
    \item a) \textbf{Đúng}.
    \item b) \textbf{Đúng}.
    \item c) \textbf{Sai}.
    \item d) \textbf{Đúng}.
\end{itemize}
}
\end{ex}

\textbf{Phần 3. Trả lời ngắn}
\setcounterpageref{ex}{0}

\begin{ex} %[2D14H75]
Tính: $567 + 285$.
\shortans{852}
\loigiai{
\textbf{Đáp số:} $852$\\
$567 + 285 = 852$.
}
\end{ex}

\begin{ex} %[2D14H75]
Tính: $824 - 356$.
\shortans{468}
\loigiai{
\textbf{Đáp số:} $468$\\
$824 - 356 = 468$.
}
\end{ex}

\begin{ex} %[2D14V75]
Tính: $2 \times 8 + 35$.
\shortans{51}
\loigiai{
\textbf{Đáp số:} $51$\\
$16 + 35 = 51$.
}
\end{ex}

\begin{ex} %[2D14V75]
Tính: $45 : 5 + 18$.
\shortans{27}
\loigiai{
\textbf{Đáp số:} $27$\\
$9 + 18 = 27$.
}
\end{ex}

\begin{ex} %[2D14V75]
Một hình tam giác có độ dài các cạnh lần lượt là $12\text{ cm}, 15\text{ cm}, 18\text{ cm}$. Tính chu vi hình tam giác đó theo xăng-ti-mét.
\shortans{45}
\loigiai{
\textbf{Đáp số:} $45$\\
$12 + 15 + 18 = 45\text{ cm}$.
}
\end{ex}

\begin{ex} %[2D14C75]
Tính nhanh: $375 + 184 + 125 + 216$.
\shortans{900}
\loigiai{
\textbf{Đáp số:} $900$\\
$(375 + 125) + (184 + 216) = 500 + 400 = 900$.
}
\end{ex}

\newpage
\end{document}
